%% Typeset with:  latexmk -pdf paper-v14.tex
%% Needs: amsmath amsthm amssymb array longtable enumitem hyperref geometry
\documentclass[11pt,a4paper]{article}
\IfFileExists{lmodern.sty}{\usepackage{lmodern}}{}   %% scalable T1 fonts
\usepackage[T1]{fontenc}
\usepackage[utf8]{inputenc}
\usepackage[margin=28mm]{geometry}
\usepackage{amsmath,amsthm,amssymb}
\usepackage{array,longtable}
\usepackage{enumitem}
\usepackage[hidelinks]{hyperref}
\usepackage[english]{babel}

%% ---- symbols used by the unicode declarations below ----------------------
\newcommand{\txtsym}[1]{\ifmmode\text{#1}\else#1\fi}
%% every non-ASCII character of the source, declared once:
\DeclareUnicodeCharacter{03D6}{\ensuremath{\varpi}}
\DeclareUnicodeCharacter{03B8}{\ensuremath{\theta}}
\DeclareUnicodeCharacter{221E}{\ensuremath{\infty}}
\DeclareUnicodeCharacter{0393}{\ensuremath{\Gamma}}
\DeclareUnicodeCharacter{03BD}{\ensuremath{\nu}}
\DeclareUnicodeCharacter{0398}{\ensuremath{\Theta}}
\DeclareUnicodeCharacter{03B3}{\ensuremath{\gamma}}
\DeclareUnicodeCharacter{03B5}{\ensuremath{\varepsilon}}
\DeclareUnicodeCharacter{03BA}{\ensuremath{\kappa}}
\DeclareUnicodeCharacter{2057}{\ensuremath{{}^{\prime\prime\prime\prime}}}
\DeclareUnicodeCharacter{2284}{\ensuremath{\not\subset}}
\DeclareUnicodeCharacter{00A7}{\txtsym{\S}}
\DeclareUnicodeCharacter{00AC}{\ensuremath{\neg}}
\DeclareUnicodeCharacter{00B1}{\ensuremath{\pm}}
\DeclareUnicodeCharacter{00B2}{\ensuremath{{}^2}}
\DeclareUnicodeCharacter{00B7}{\ensuremath{\cdot}}
\DeclareUnicodeCharacter{00B9}{\ensuremath{{}^1}}
\DeclareUnicodeCharacter{00D7}{\ensuremath{\times}}
\DeclareUnicodeCharacter{00E1}{\txtsym{\'a}}
\DeclareUnicodeCharacter{00E9}{\txtsym{\'e}}
\DeclareUnicodeCharacter{00ED}{\txtsym{\'\i}}
\DeclareUnicodeCharacter{00F6}{\txtsym{\"o}}
\DeclareUnicodeCharacter{0151}{\txtsym{\H o}}
\DeclareUnicodeCharacter{03A0}{\ensuremath{\Pi}}
\DeclareUnicodeCharacter{03A3}{\ensuremath{\Sigma}}
\DeclareUnicodeCharacter{03A8}{\ensuremath{\Psi}}
\DeclareUnicodeCharacter{03B6}{\ensuremath{\zeta}}
\DeclareUnicodeCharacter{03C0}{\ensuremath{\pi}}
\DeclareUnicodeCharacter{03C3}{\ensuremath{\sigma}}
\DeclareUnicodeCharacter{03C4}{\ensuremath{\tau}}
\DeclareUnicodeCharacter{2013}{\txtsym{\textendash}}
\DeclareUnicodeCharacter{2014}{\txtsym{\textemdash}}
\DeclareUnicodeCharacter{2026}{\ensuremath{\dots}}
\DeclareUnicodeCharacter{2032}{\ensuremath{{}^\prime}}
\DeclareUnicodeCharacter{2070}{\ensuremath{{}^0}}
\DeclareUnicodeCharacter{2075}{\ensuremath{{}^5}}
\DeclareUnicodeCharacter{2076}{\ensuremath{{}^6}}
\DeclareUnicodeCharacter{2079}{\ensuremath{{}^9}}
\DeclareUnicodeCharacter{2080}{\ensuremath{{}_0}}
\DeclareUnicodeCharacter{2081}{\ensuremath{{}_1}}
\DeclareUnicodeCharacter{2082}{\ensuremath{{}_2}}
\DeclareUnicodeCharacter{2083}{\ensuremath{{}_3}}
\DeclareUnicodeCharacter{2084}{\ensuremath{{}_4}}
\DeclareUnicodeCharacter{2085}{\ensuremath{{}_5}}
\DeclareUnicodeCharacter{2086}{\ensuremath{{}_6}}
\DeclareUnicodeCharacter{2087}{\ensuremath{{}_7}}
\DeclareUnicodeCharacter{2088}{\ensuremath{{}_8}}
\DeclareUnicodeCharacter{2089}{\ensuremath{{}_9}}
\DeclareUnicodeCharacter{2074}{\ensuremath{{}^4}}
\DeclareUnicodeCharacter{00B3}{\ensuremath{{}^3}}
\DeclareUnicodeCharacter{2077}{\ensuremath{{}^7}}
\DeclareUnicodeCharacter{2078}{\ensuremath{{}^8}}
\DeclareUnicodeCharacter{2071}{\ensuremath{{}^i}}
\DeclareUnicodeCharacter{2261}{\ensuremath{\equiv}}
\DeclareUnicodeCharacter{21A6}{\ensuremath{\mapsto}}
\DeclareUnicodeCharacter{2205}{\ensuremath{\emptyset}}
\DeclareUnicodeCharacter{2218}{\ensuremath{\circ}}
\DeclareUnicodeCharacter{2225}{\ensuremath{\parallel}}
\DeclareUnicodeCharacter{2229}{\ensuremath{\cap}}
\DeclareUnicodeCharacter{2237}{\ensuremath{::}}
\DeclareUnicodeCharacter{2248}{\ensuremath{\approx}}
\DeclareUnicodeCharacter{2282}{\ensuremath{\subset}}
\DeclareUnicodeCharacter{00BD}{\ensuremath{1/2}}
\DeclareUnicodeCharacter{2044}{\ensuremath{/}}
\DeclareUnicodeCharacter{2115}{\ensuremath{\mathbb{N}}}
\DeclareUnicodeCharacter{211A}{\ensuremath{\mathbb{Q}}}
\DeclareUnicodeCharacter{211D}{\ensuremath{\mathbb{R}}}
\DeclareUnicodeCharacter{2124}{\ensuremath{\mathbb{Z}}}
\DeclareUnicodeCharacter{2192}{\ensuremath{\to}}
\DeclareUnicodeCharacter{2200}{\ensuremath{\forall}}
\DeclareUnicodeCharacter{2203}{\ensuremath{\exists}}
\DeclareUnicodeCharacter{2208}{\ensuremath{\in}}
\DeclareUnicodeCharacter{2209}{\ensuremath{\notin}}
\DeclareUnicodeCharacter{220E}{\ensuremath{\blacksquare}}
\DeclareUnicodeCharacter{2212}{\ensuremath{-}}
\DeclareUnicodeCharacter{221A}{\ensuremath{\surd}}
\DeclareUnicodeCharacter{2227}{\ensuremath{\wedge}}
\DeclareUnicodeCharacter{222A}{\ensuremath{\cup}}
\DeclareUnicodeCharacter{2260}{\ensuremath{\neq}}
\DeclareUnicodeCharacter{2261}{\ensuremath{\equiv}}
\DeclareUnicodeCharacter{2264}{\ensuremath{\le}}
\DeclareUnicodeCharacter{2265}{\ensuremath{\ge}}
\DeclareUnicodeCharacter{226A}{\ensuremath{\ll}}
\DeclareUnicodeCharacter{2286}{\ensuremath{\subseteq}}
\DeclareUnicodeCharacter{230A}{\ensuremath{\lfloor}}
\DeclareUnicodeCharacter{230B}{\ensuremath{\rfloor}}
\DeclareUnicodeCharacter{27E8}{\ensuremath{\langle}}
\DeclareUnicodeCharacter{27E9}{\ensuremath{\rangle}}
\DeclareUnicodeCharacter{27F9}{\ensuremath{\Longrightarrow}}
\DeclareUnicodeCharacter{27FA}{\ensuremath{\Longleftrightarrow}}
\DeclareUnicodeCharacter{2223}{\ensuremath{\mid}}
\DeclareUnicodeCharacter{00B0}{\ensuremath{{}^\circ}}
\DeclareUnicodeCharacter{00C1}{\txtsym{\'A}}
\DeclareUnicodeCharacter{00CA}{\txtsym{\^E}}
\DeclareUnicodeCharacter{00E4}{\txtsym{\"a}}
\DeclareUnicodeCharacter{00EA}{\txtsym{\^e}}
\DeclareUnicodeCharacter{00FC}{\txtsym{\"u}}
\DeclareUnicodeCharacter{0112}{\ensuremath{\bar{E}}}
\DeclareUnicodeCharacter{0169}{\ensuremath{\tilde{u}}}
\DeclareUnicodeCharacter{0394}{\ensuremath{\Delta}}
\DeclareUnicodeCharacter{039B}{\ensuremath{\Lambda}}
\DeclareUnicodeCharacter{03A6}{\ensuremath{\Phi}}
\DeclareUnicodeCharacter{03A9}{\ensuremath{\Omega}}
\DeclareUnicodeCharacter{03B1}{\ensuremath{\alpha}}
\DeclareUnicodeCharacter{03B2}{\ensuremath{\beta}}
\DeclareUnicodeCharacter{03B4}{\ensuremath{\delta}}
\DeclareUnicodeCharacter{03B7}{\ensuremath{\eta}}
\DeclareUnicodeCharacter{03B9}{\ensuremath{\iota}}
\DeclareUnicodeCharacter{03BB}{\ensuremath{\lambda}}
\DeclareUnicodeCharacter{03BC}{\ensuremath{\mu}}
\DeclareUnicodeCharacter{03BE}{\ensuremath{\xi}}
\DeclareUnicodeCharacter{03C1}{\ensuremath{\rho}}
\DeclareUnicodeCharacter{03C6}{\ensuremath{\varphi}}
\DeclareUnicodeCharacter{03C7}{\ensuremath{\chi}}
\DeclareUnicodeCharacter{03C8}{\ensuremath{\psi}}
\DeclareUnicodeCharacter{03C9}{\ensuremath{\omega}}
\DeclareUnicodeCharacter{03D5}{\ensuremath{\phi}}
\DeclareUnicodeCharacter{03F1}{\ensuremath{\varrho}}
\DeclareUnicodeCharacter{0428}{\ensuremath{\mathrm{Sha}}}
\DeclareUnicodeCharacter{1D57}{\ensuremath{{}^{t}}}
\DeclareUnicodeCharacter{2033}{\ensuremath{{}^{\prime\prime}}}
\DeclareUnicodeCharacter{2113}{\ensuremath{\ell}}
\DeclareUnicodeCharacter{2202}{\ensuremath{\partial}}
\DeclareUnicodeCharacter{220F}{\ensuremath{\prod}}
\DeclareUnicodeCharacter{2216}{\ensuremath{\setminus}}
\DeclareUnicodeCharacter{2217}{\ensuremath{*}}
\DeclareUnicodeCharacter{2224}{\ensuremath{\nmid}}
\DeclareUnicodeCharacter{2245}{\ensuremath{\cong}}
\DeclareUnicodeCharacter{2262}{\ensuremath{\not\equiv}}
\DeclareUnicodeCharacter{2295}{\ensuremath{\oplus}}
\DeclareUnicodeCharacter{2297}{\ensuremath{\otimes}}
\DeclareUnicodeCharacter{22A5}{\ensuremath{\perp}}
\DeclareUnicodeCharacter{22EF}{\ensuremath{\cdots}}
\DeclareUnicodeCharacter{2605}{\ensuremath{\star}}
\DeclareUnicodeCharacter{27F8}{\ensuremath{\Longleftarrow}}
\DeclareUnicodeCharacter{2190}{\ensuremath{\leftarrow}}
\DeclareUnicodeCharacter{2191}{\ensuremath{\uparrow}}
\DeclareUnicodeCharacter{2193}{\ensuremath{\downarrow}}
\DeclareUnicodeCharacter{2500}{\ensuremath{-}}
\DeclareUnicodeCharacter{2502}{\ensuremath{|}}
\DeclareUnicodeCharacter{2510}{\txtsym{+}}
\DeclareUnicodeCharacter{2514}{\txtsym{+}}
\DeclareUnicodeCharacter{2518}{\txtsym{+}}
\DeclareUnicodeCharacter{251C}{\txtsym{+}}
\DeclareUnicodeCharacter{2524}{\txtsym{+}}
\DeclareUnicodeCharacter{252C}{\txtsym{+}}
\DeclareUnicodeCharacter{2534}{\txtsym{+}}
\DeclareUnicodeCharacter{253C}{\txtsym{+}}
\DeclareUnicodeCharacter{2034}{\ensuremath{{}^{\prime\prime\prime}}}
\DeclareUnicodeCharacter{1D9C}{\ensuremath{{}^{c}}}
\DeclareUnicodeCharacter{207B}{\ensuremath{{}^{-}}}
\DeclareUnicodeCharacter{21D2}{\ensuremath{\Rightarrow}}
\DeclareUnicodeCharacter{2308}{\ensuremath{\lceil}}
\DeclareUnicodeCharacter{2309}{\ensuremath{\rceil}}
\DeclareUnicodeCharacter{230A}{\ensuremath{\lfloor}}
\DeclareUnicodeCharacter{230B}{\ensuremath{\rfloor}}
\DeclareUnicodeCharacter{1D4F}{\ensuremath{{}^{k}}}
\DeclareUnicodeCharacter{207A}{\ensuremath{{}^{+}}}
\DeclareUnicodeCharacter{207D}{\ensuremath{{}^{(}}}
\DeclareUnicodeCharacter{207E}{\ensuremath{{}^{)}}}
\DeclareUnicodeCharacter{2112}{\ensuremath{\mathcal{L}}}
\DeclareUnicodeCharacter{2118}{\ensuremath{\wp}}
\DeclareUnicodeCharacter{E000}{\ensuremath{\bar{ω}}}
\DeclareUnicodeCharacter{E001}{\ensuremath{\bar{φ}}}
\DeclareUnicodeCharacter{E002}{\ensuremath{\bar{ψ}}}
\DeclareUnicodeCharacter{E003}{\ensuremath{\bar{x}}}
\DeclareUnicodeCharacter{E004}{\ensuremath{\bar{K}}}
\DeclareUnicodeCharacter{E005}{\ensuremath{\bar{Q}}}
\DeclareUnicodeCharacter{E006}{\ensuremath{\hat{π}}}
\DeclareUnicodeCharacter{E007}{\ensuremath{\bar{M}}}
\DeclareUnicodeCharacter{E008}{\ensuremath{\bar{χ}}}
\DeclareUnicodeCharacter{E009}{\ensuremath{\bar{P}}}
\DeclareUnicodeCharacter{E00A}{\ensuremath{\bar{T}}}
\DeclareUnicodeCharacter{E00B}{\ensuremath{\bar{F}}}
\DeclareUnicodeCharacter{E00C}{\ensuremath{\tilde{E}}}
\DeclareUnicodeCharacter{E00D}{\ensuremath{\bar{τ}}}
\DeclareUnicodeCharacter{E00E}{\ensuremath{\tilde{Ω}}}
\DeclareUnicodeCharacter{E00F}{\ensuremath{\bar{π}}}
\DeclareUnicodeCharacter{E010}{\ensuremath{\bar{ϱ}}}
\DeclareUnicodeCharacter{E011}{\ensuremath{\bar{B}}}
\DeclareUnicodeCharacter{E012}{\ensuremath{\tilde{P}}}
\DeclareUnicodeCharacter{E013}{\ensuremath{\hat{Ω}}}

%% ---- named theorem environments (the labels are those of the draft) ------
\newtheoremstyle{named}{\topsep}{\topsep}{\itshape}{0pt}{\bfseries}{.}{ }{\thmnote{#3}}
\theoremstyle{named}
\newtheorem*{namedthm}{}
\newenvironment{mdproof}[1][Proof]{\par\medskip\noindent\textit{#1.}\ }{\par\medskip}
\newlist{leancode}{itemize}{1}
\setlist[leancode]{label={},leftmargin=1.2em,topsep=.5em,itemsep=.15em,parsep=0pt}
\newcommand{\dispf}[1]{\par\smallskip{\centering #1\par}\smallskip}
\setlength{\parskip}{.35em}
\sloppy
\hbadness=10000
\begin{document}

\title{Integrality of the p-adic height pairing at an anomalous prime is equivalent to a point of order p in E(Q\_p)\\[.5em]\large Splitting of the reduction sequence, the vanishing of a Fermat quotient of the p-th division polynomial, and p-integrality of the height matrix on E(Q)/tors are one and the same condition; in rank one v\_p(Reg\_p) ≥ 0 ⟺ split, the non-split value being −1 exactly against the elementary baseline r − 2; modulo p the height matrix has rank exactly one in the non-split case and rank at most two in the split case, its constant being a Fermat quotient of the division polynomial in the first case and a second-order Fermat quotient in the second}
\author{Kiichi \and Shiori \and Rin}
\date{Draft v14, 2026-10-02. Not submitted. See §8 for what is not claimed and §6 for the boundary between the machine-checked part and the paper part.}
\maketitle

\section*{Abstract}

Let E/Q be an elliptic curve, p ≥ 5 a prime of good ordinary reduction, r the rank of E(Q), and Reg\_p the cyclotomic p-adic regulator in the normalisation of [MST] used by [Har], h\_p(P) = 2 log\_p(σ\_p(P)/d(P)) with σ\_p the Mazur–Tate sigma function. Call p \textit{anomalous} when p ∣ \#E(F\_p) (a\_p = 1, that is \#E(F\_p) = p, for p ≥ 7) and E \textit{split at p} when E(Q\_p)[p] ≠ 0.

At an anomalous prime Reg\_p need not be p-integral: counting indices gives v\_p(Reg\_p) ≥ r − 2 when p ∤ ∏\_ℓ c\_ℓ and ≥ r otherwise (Theorem A). That baseline is elementary and it is not ours — the three negative entries of the tables of [MST], the precision the normalisation of [Har] costs at an anomalous prime, a hypothesis of [Ban]. What this note contributes is a local mechanism that governs the deficit, in both directions:

\begin{namedthm}[Integrality theorem for the p-adic regulator at split anomalous primes]
Suppose \#E(F\_p) = p, p ∤ ∏\_ℓ c\_ℓ, and that Λ := E(Q)/tors surjects onto E(F\_p). Then \textbf{every entry of the height matrix M\_Λ on a basis of Λ lies in Z\_p if and only if E(Q\_p)[p] ≠ 0}; in rank one, v\_p(Reg\_p) ≥ 0 ⟺ E(Q\_p)[p] ≠ 0, the non-split value being \textbf{−1 exactly}; and in the split case v\_p(Reg\_p) ≥ max(r − 2, 0) in every rank.
\end{namedthm}

The same mechanism determines the height matrix modulo p, and it does so in every rank. In the non-split case p·M\_Λ is integral and p·M\_Λ mod p is the rank-one matrix c·χχᵗ, χ the reduction character, with c = 2λ₁ (Theorem R, Proposition R″); in the split case M\_Λ itself is integral and M\_Λ mod p vanishes on the kernel of χ, so its rank is at most two, and the coefficient a = h\_p(e₁) mod p equals 2λ\textasciicircum{}\{(2)\}/n² for a second-order Fermat quotient λ\textasciicircum{}\{(2)\} of the same division polynomial (Theorem R\_s, Proposition R″\_s). In both cases v\_p(Reg\_p) = (r − 2) + v\_p(det T) for an integral Schur complement T, which for r = 1 gives the exact value −1 or, on the split side, the criterion v\_p(Reg\_p) = 0 ⟺ λ\textasciicircum{}\{(2)\} ≠ 0. λ is a function of E mod p² and of the class in E(F\_p); λ\textasciicircum{}\{(2)\} is not, so the split value has no law in terms of the reduction alone. In the split case the whole of the height matrix modulo p is the closed form 2λ\_\{T,1\}χ⊗χ + B₀(χ⊗γ + γ⊗χ) in two local constants and the two homomorphisms χ, γ given by the first two layers of the filtration of E(Q\_p) (Theorem K), so that the deficit r − 2 is exactly what those two layers take away; but the second p-adic digit of the height contains the Fermat quotient of the denominator of the point (Theorem N), so the localisation stops after one digit and no third layer of the filtration governs v\_p(Reg\_p) ≥ 2. The same expansion can nevertheless be carried two digits further (Corollaries N⁽³⁾ and N⁽⁴⁾): the third digit is the first at which the normalisation of σ — the p-adic weight-two Eisenstein value E₂ — becomes visible in the height, and the fourth needs no new local quantity at all. Since E₂ modulo p is an invariant of the reduced curve, and since the one new quantity of the third digit lies in the first layer of the formal filtration (Proposition 12), the third digit does have a complete answer: granted one non-degeneracy — that φ\_p′ has valuation exactly one at the p-torsion point, which is now a theorem for every p ≥ 5 (Theorem ND, §3.14) — the j-th layer of the filtration moves exactly the digits of the height from p\textasciicircum{}\{j−1\} on (Theorem G′), and min(v\_p(Reg\_p), 3) is a function of five numbers, three layers of the filtration and two digits of the Fermat quotient of the denominator, and of nothing coarser (Theorem G″). Past the fourth digit the expansion need not be continued by hand. Summed instead of expanded, the ledger is the identity n²h\_p(P) = (1 − p)\textasciicircum{}\{−1\}[p\textasciicircum{}\{−2\}log(1 + p²Θ) + 2log(1 + pF)] + p\textasciicircum{}\{−2\}·ℒ(z(P″)) in two Iwasawa logarithms and one power series ℒ(z) = log(x(P″)·σ\_p(t(P″))²) (\textbf{Theorem N*}), and ℒ is \textbf{even} in the formal logarithm z, with coefficients μ₂ = −c, μ₄ = −b₄/12, μ₆ = −b₂b₄/360 − b₆/30 — invariants of E in which E₂ occurs in the first alone (Proposition 17). Every congruence of the paper is this identity truncated: the term of index j in a logarithm is one digit shallower than p\textasciicircum{}j whenever p ∣ j, and the m-th coefficient of ℒ has valuation 2me − 2 + v\_p(μ\_\{2m\}), so it arrives at the (2me − 1 + v\_p(μ\_\{2m\}))-th digit (Proposition 18). So c is visible at the third digit, the next local invariant b₄ not until the \textbf{seventh}, and μ₆ at the tenth for p = 5 and the eleventh for p ≥ 7 — a depth that depends on p through the denominator of a coefficient. The truncations at the fifth, sixth and seventh digits are Corollaries N⁽⁵⁾, N⁽⁶⁾ and N⁽⁷⁾, checked in 648 boxes, and the identity itself was checked to the tenth digit. Since the argument of ℒ is p·log\_E(t₁) with t₁ a layer of the filtration (Proposition 19), the global input at the k-th digit is (χ, γ, …, γ\_k) together with F mod p\textasciicircum{}\{k−1\}: two kinds of datum, each growing by one per digit, and nothing coarser will do at any depth (\textbf{Theorem G⁽ᵏ⁾}), because a layer moves Θ at the j-th digit and the ledger only at the (j+2)-nd — a gap with no upper bound on j (Proposition 20).

Two further statements concern rational points rather than the parameter space. Along p-power translates the digits of F are not free: F(P + p\textasciicircum{}aZ) ≡ F(P) + p\textasciicircum{}\{a−1\}·(χ(Z)/χ(P))·(q(P) − C(P)ζ(P)/2) (mod p\textasciicircum{}a) for the quadratic form q = n²h\_p, a unit C(P) and a homomorphism ζ (Proposition 21), and on rational points F mod p\textasciicircum{}k is a function of n²h\_p mod p\textasciicircum{}\{k+1\} and of t₁ mod p\textasciicircum{}\{k+2\} (Proposition 22); nevertheless both kinds of necessity in Theorem G⁽ᵏ⁾ are realised by pairs of rational points of one curve (4407a1 at p = 5). And one archimedean datum enters: a rational point in E\_k(Q\_p) has naive height at least 2k·log p (Lemma 23), so if the shortest non-zero point of Λ ∩ E\_1(Q\_p) has height below 4·log p then (χ, γ) is onto and, in rank two, v\_p(Reg\_p) = 0 (Proposition 24) — an upper bound that no local datum gives. The same argument with the canonical height ĥ and the local heights away from p gives the threshold 4·log p + β(E), with β(E) the sum over v ≠ p of the lower bounds of the local heights (Proposition 24ᶜ). Over 3,056 split boxes of rank two the naive threshold accounts for about two thirds of the observed deficit of det(χ, γ) ≡ 0 (207.8 of 305.35 expected cases), and the canonical threshold for all of it in expectation (307.1 of 305.35); over a second, disjoint set of 5,572 split boxes (conductor between 200,000 and 500,000) the figures are 55.8 \% and 99.7 \%, and no box with det(χ, γ) ≡ 0 lies below the canonical threshold in either set, nor in two further sets of 5,782 and 7,861 split boxes of curves outside Cremona's tables [computation, with a numerical archimedean minimum]. Above the threshold det(χ, γ) ≡ 0 occurs with frequency 1/p within the noise on the first two sets; that is an empirical regularity and not a theorem, and on the two further sets it is exceeded just above the threshold, for a reason we do not know (§5.12).

Three things are separated throughout. The \textbf{baseline} r − 2 is a known phenomenon. The equivalence for the \textbf{matrix} is the content of the note, and it holds in every rank. The bound on the \textbf{determinant} improves on the baseline in rank one only, since max(r−2, 0) = r−2 for r ≥ 2.

The proof is local and runs through one element λ ∈ F\_p, the Fermat quotient q\_p(φ\_p(x₀))/2 of the p-th division polynomial at a lift of a nonzero class of E(F\_p). With P = (α/d², β/d³) in lowest terms the height splits into a formal part of valuation b and 2 log\_p(α/β) of valuation c ≥ 1 (Lemma 3); in a model with p ∣ a₁ — always available, since p is odd — one has b ≥ 2 and q\_p(α/β) = −λ (Theorem B), so the height is integral exactly when λ vanishes. Vélu's factorisation of multiplication by p kills λ in the split case (Theorem C, Corollary D); conversely λ = 0 forces splitting, a known criterion (Proposition 6, [LR, Thm 4.1], [DW, Lemma 3.2]) proved here through the Euler identity for the weighted-homogeneous division polynomial and a residue computation at the cusp. CM curves split at every anomalous prime of good reduction (Theorem F). The arithmetic steps are theorems of Lean 4 with Mathlib (no \texttt{sorry}, no \texttt{native\_\allowbreak{}decide}, axioms [propext, Classical.choice, Quot.sound]), the geometric inputs entering as hypotheses. Nothing here bears on the conjecture of Birch and Swinnerton-Dyer; the one statement that assumes its p-adic analogue, Proposition V, is kept apart in §6 and in a companion note.

\textbf{Keywords.} p-adic height, p-adic regulator, anomalous prime, local torsion, division polynomial, Fermat quotient, Vélu isogeny, Tate curve, Lean 4.

\textbf{MSC 2020.} 11G50, 11G07, 11S31, 11Y40, 68V20.

\section*{1. Introduction}

The cyclotomic p-adic height pairing of Mazur–Tate and Schneider attaches to E/Q and a prime p of good ordinary reduction a symmetric bilinear form on E(Q) ⊗ Q\_p, whose determinant Reg\_p on a basis of Λ := E(Q)/tors is the regulator in the p-adic analogue of the conjecture of Birch and Swinnerton-Dyer [MST] [BMS]. Unlike the archimedean regulator it is not known to be non-zero (Schneider's conjecture [SW, Conj. 4.1]), and nothing here bears on that.

This note concerns a coarser invariant, the valuation of Reg\_p. At an \textit{anomalous} prime — p ∣ \#E(F\_p), which for p ≥ 7 means a\_p = 1, the notion of [Maz72] — it can be negative, because the points reducing to O at p form a subgroup of index divisible by p and dividing the determinant by the square of that index costs 2. This is the baseline v\_p(Reg\_p) ≥ r − 2 of Theorem A; it is known (§9), and it is attained.

What this note proves is that the deficit is present exactly when the reduction sequence 0 → E\_1(Q\_p) → E(Q\_p) → E(F\_p) → 0 fails to split, that is exactly when E(Q\_p) contains no point of order p. Splitting makes the height of every point of Λ p-integral — a statement about the pairing, which the baseline gives in no rank — and its failure puts the valuation −1 into the matrix, in every rank, and into the determinant in rank one. The obstruction sits in the second p-adic digit of the height, where it is a Fermat quotient, and the route runs: decomposition of the height (Lemma 3), a criterion expressing that digit through one element λ ∈ F\_p (Theorem B), the vanishing of λ in the split case by Vélu's factorisation of multiplication by p (Theorem C, Corollary D), and its non-vanishing in the non-split case (Proposition 6), which is a known criterion for local torsion.

The materials are classical and we do not call the lemmas new; what we did not find in the sources we read is the equivalence itself, and the presentation of the local criterion through a Fermat quotient. Past the second digit the note becomes a ledger, and §3.11 closes it: the whole expansion is one identity in two Iwasawa logarithms and one even power series, whose truncations are the congruences of §3.12 and whose coefficients say at which digit each invariant of E appears. §3.15 adds one archimedean input — a point deep in the filtration has a large naive height — and what it decides about v\_p(Reg\_p). §9 records where we looked, §6 draws the boundary between what a kernel has checked and what it has not, §7 lists what is open and §8 what is not claimed.

\section*{2. Setting and normalisation}

\textbf{The curve.} E/Q is given by a global minimal Weierstrass equation with coefficients a₁, …, a₆, and p ≥ 5 is a prime of good ordinary reduction; Λ := E(Q)/tors has rank r, and c\_ℓ is the Tamagawa number at ℓ. Write t = −x/y for the parameter of the formal group Ê, E\_n(Q\_p) := \{P : v\_p(t(P)) ≥ n\}, and for P ∈ E(Q)

\begin{quote}
P = (α(P)/d(P)², β(P)/d(P)³),  gcd(α, d) = gcd(β, d) = 1, d ≥ 1.
\end{quote}

v\_p is normalised by v\_p(p) = 1, log\_p is the Iwasawa logarithm with log\_p(p) = 0, and for a p-adic unit u, q\_p(u) := (u\textasciicircum{}\{p−1\} − 1)/p mod p is the Fermat quotient. Then q\_p : Z\_p\textasciicircum{}× → F\_p is a surjective homomorphism and v\_p(log\_p u) = 1 + v\_p(q\_p(u)), so "the second p-adic digit of log\_p u" and "the Fermat quotient of u" are the same datum.

\textbf{The height.} The height is taken in the explicit form of [Har, §2.2], which quotes the definition of [MST]: if P satisfies

\begin{quote}
(A1) P reduces to O at p, and (A2) P reduces to a nonsingular point at every bad prime,
\end{quote}

then h\_p(P) = 2 log\_p(σ\_p(P)/d(P)); for a general point h\_p(Q) := h\_p(n\_\{lcm\}Q)/n\_\{lcm\}² with n\_\{lcm\} := lcm(\#E(F\_p), lcm\_ℓ c\_ℓ). This is 2p times the normalisation of [MST], so the two valuations of Reg\_p differ by r. Against a computer algebra system the sign is worth recording: with this normalisation \textbf{h\_p is minus PARI's canonical p-adic height}, so that n²h\_p(P) = −n²·\texttt{ellpadicheight(E,p,n,P)*[1,−s₂]\textasciitilde{}} exactly, in every digit; valuations and the valuation of Reg\_p are unaffected.

\textit{The letter n.} Only n\_\{lcm\} needs the least common multiple; from §3.9 on, \textbf{n always denotes the Tamagawa product ∏\_ℓ c\_ℓ}, which is what makes nQ satisfy (A2) and is all the proofs use. The other uses of the letter are kept apart by a subscript: n\_\{lcm\} here, n\_Q in §5.1, n\_p := (p−1)/2 in §§3.14 and 4.13. In §2 and Theorem C the index of a division polynomial φ\_n, ψ\_n, ω\_n and the order of a cyclic subgroup are bound variables of those statements. The pairing is h\_p(P, R) = (h\_p(P+R) − h\_p(P) − h\_p(R))/2 and Reg\_p := det(h\_p(P\_i, P\_j)) on a basis of Λ.

\textit{Which σ matters.} σ\_p is the Mazur–Tate sigma function: the unique odd σ(t) = t + c₂t² + ⋯ with \textbf{σ\_p ∈ tZ\_p[[t]]} solving x(t) + c = −(d/ω)((1/σ)(dσ/ω)) with c = (b₂ − E₂(E, ω))/12 [Har, (2)], [MT91]; c₂ = a₁/2. Another constant in place of E₂ shifts h\_p by a multiple of the square of the p-adic elliptic logarithm, and only this choice keeps the coefficients in Z\_p. Every statement below is about σ\_p.

\textit{E₂ modulo p is an invariant of the reduced curve.} For p ≥ 5 and E ordinary at p, Katz's congruence E₂\textasciicircum{}* ≡ E\_\{p+1\}/E\_\{p−1\} (mod p) expresses the p-adic weight-two Eisenstein value through two classical forms, whose values on (E, ω) are isobaric polynomials in E₄(E, ω) = c₄ and E₆(E, ω) = −c₆. Hence

\begin{quote}
\textbf{E₂(E, ω) mod p — and with it c mod p — is a function of (c₄ mod p, c₆ mod p) alone},
\end{quote}

that is, of the reduced pair (Ē, ): the constant of the σ-normalisation is a local quantity at p, and not a global one as its definition through a p-adic modular form makes it look. Two internal checks make the statement independent of the source: E\_\{p−1\} ≡ 1 (mod p) ([Ser73]) and E\_\{p+1\} ≡ E₂ (mod p), the latter from Kummer's congruence B\_\{p+1\}/(p+1) ≡ B₂/2 (mod p) together with σ\_p(n) ≡ σ₁(n) (mod p). The normalisation E₆(E, ω) = −c₆ is forced by the Tate curve — for E\_q : y² + xy = x³ − 5s₃x − (5s₃ + 7s₅)/12 one has c₄ = 1 + 240s₃ = E₄ and c₆ = −1 + 504s₅ = −E₆ — and is calibrated on the Hasse invariant in §5.8. The congruence itself is [known] ([Kat73], p-adic modular forms), and we have not retrieved the original; what §5.8 does is check the consequence against PARI's \texttt{ellpadics2} over 315,953 boxes, which makes the \textit{use} of the congruence here safe without making its proof ours. This local description of c is what allows the quantity W of §3.11 to be identified with a datum of the formal group.

\textit{Which model matters.} h\_p(P) and v\_p(h\_p(P)) are invariants of E/Q, but the two halves b and c of the decomposition of Lemma 3 are not: the substitution (x, y) ↦ (x, y + sx) with s ∈ Z carries a global minimal equation to a global minimal equation, fixes x, α, d, φ\_p, ψ\_p and λ, and moves only β ↦ β + sαd, so b and c trade with each other while min(b, c) stands. Since p is odd, s can be chosen with \textbf{p ∣ a₁}, and that is the model in which Theorem B and Theorem S are proved; in it b ≥ 2e comes for free, and the whole question is c. Statements about b or c alone are statements about a model, and the invariant form of Theorem S is v\_p(h\_p(pP′)) = 1.

\textbf{Anomalous, split, ι.} p is \textit{anomalous} when p ∣ \#E(F\_p), which for p ≥ 7 means \#E(F\_p) = p (§3.1); E is \textit{split at p} when E(Q\_p)[p] ≠ 0; ι says that Λ → E(F\_p) is onto.

\textbf{Division polynomials.} φ\_n, ψ\_n, ω\_n are normalised by x(nQ) = φ\_n/ψ\_n², y(nQ) = ω\_n/ψ\_n³, with φ\_n monic of degree n² in x. For odd n and a rational cyclic D of order n, the Vélu isogeny π\_D : E → E/D has x-part /² with (x) = ∏\_\{\{±T\}⊂D∖\{O\}\}(x − x(T)) monic of degree (n−1)/2 and  monic of degree n. Modulo p, for E ordinary, φ\_p(x) = f(x\textasciicircum{}p) and ψ\_p(x) = g(x\textasciicircum{}p) in F\_p[A, B][X] with f monic of degree p in X, g of degree (p−1)/2 with leading coefficient the Hasse invariant, and f/g² the x-part of the Verschiebung V; when p is anomalous, ker V = E(F\_p).

\textbf{Grades.} \textbf{[Lean]} kernel-checked, no \texttt{sorry}, no \texttt{native\_\allowbreak{}decide}, axioms within [propext, Classical.choice, Quot.sound]; \textbf{[paper]} a complete proof here or cited; \textbf{[computation]} verified over a stated range; \textbf{[known]} in the literature. Where grades mix, the weakest governs.

\section*{3. Main theorems}

Throughout, p ≥ 5 is a prime of good ordinary reduction for E.

\subsection*{3.1 The baseline, and the decomposition of the height}

\begin{namedthm}[Theorem A (baseline) [paper]]
Let G := \{P ∈ E(Q) : (A1) and (A2) hold\} and k := [Λ : image of G]. Then v\_p(Reg\_p) ≥ r − 2 v\_p(k). If moreover p ∤ ∏\_ℓ c\_ℓ, then v\_p(k) ≤ 1, with equality exactly when the image of Λ in E(F\_p) meets E(F\_p)[p] ≠ 0, so that
\dispf{v\_p(Reg\_p) ≥ r − 2 in that case,  v\_p(Reg\_p) ≥ r otherwise.}
\end{namedthm}

The hypothesis is the one the proof uses and is weaker than a\_p = 1: by Hasse, p ∣ \#E(F\_p) forces \#E(F\_p) ∈ \{p, 2p\} for p ≥ 5, and 2p occurs only at p = 5 with a\_p = −4, so for p ≥ 7 it is exactly a\_p = 1, E(F\_p) ≅ Z/p and the condition on Λ is ι, while at p = 5 with \#E(F\_5) = 10 it is strictly weaker than ι. \textbf{Everything after Lemma 4 assumes \#E(F\_p) = p}, and no statistic below mixes the two. The argument is index bookkeeping and nothing more (§4.2), and it is known (§9); we give it because the rest of the paper is measured against it. What it does \textit{not} give is control on the entries of the height matrix, only on the determinant (§3.2).

For P satisfying (A1) and (A2), with e := v\_p(t(P)) ≥ 1:

\begin{namedthm}[Lemma 1 [paper]]
σ\_p(P)/d(P) ∈ Z\_p\textasciicircum{}×.
\textbf{Lemma 2 [paper].} h\_p(P) ∈ 2pZ\_p; more precisely v\_p(h\_p(P)) = 1 + v\_p(q\_p(σ\_p(P)/d(P))).
\textbf{Lemma 3 (decomposition) [paper].} With u(P) := σ\_p(t)/t ∈ 1 + p\textasciicircum{}eZ\_p and n\_ℓ(P) := v\_ℓ(α) − v\_ℓ(β) ∈ Z,
\dispf{h\_p(P) = 2 log\_p(u(P)) + 2 Σ\_\{ℓ≠p\} n\_ℓ(P) log\_p(ℓ) =: 2Λ\_σ + 2A,}
and, with b := v\_p(Λ\_σ) and c := v\_p(A): b ≥ e (and b ≥ 2e when p ∣ a₁), and \textbf{c ≥ 1}.
\end{namedthm}

By §2 the pair (b, c) depends on the model and min(b, c) does not; in the model with p ∣ a₁ one has b ≥ 2e ≥ 2 for free.

\subsection*{3.2 The main theorem}

\begin{namedthm}[Integrality theorem for the p-adic regulator at split anomalous primes\textbf{ (\textit{Split-integrality theorem}) }[paper]]
Suppose \#E(F\_p) = p, p ∤ ∏\_ℓ c\_ℓ and ι. Then, with M\_Λ the height matrix on a basis of Λ,
\dispf{\textbf{(i)} every entry of M\_Λ lies in Z\_p  ⟺  E(Q\_p)[p] ≠ 0;}
\dispf{\textbf{(ii)} if r = 1, then v\_p(Reg\_p) ≥ 0 ⟺ E(Q\_p)[p] ≠ 0, and in the non-split case v\_p(Reg\_p) = −1 exactly;}
\dispf{\textbf{(iii)} if E(Q\_p)[p] ≠ 0, then v\_p(Reg\_p) ≥ max(r − 2, 0) in every rank.}
Without ι the hypothesis of Theorem A is not met and that theorem already gives v\_p(Reg\_p) ≥ r.
\end{namedthm}

Three things are being said, and only some of them are new against Theorem A.

\textit{The pairing.} That the \textbf{entries} are integral, and that they cease to be exactly when E(Q\_p)[p] = 0, does not follow from any determinant bound. The direction ⟸ of (i) is the mechanism of §4.8; the direction ⟹ is Theorem S, which produces an entry of valuation −1 whenever the reduction sequence does not split. This is what the proof gives, and it holds in every rank.

\textit{The determinant.} Since max(r−2, 0) = r−2 for r ≥ 2, the bound on v\_p(Reg\_p) is new only in \textbf{rank one}, where the baseline −1 becomes 0 and, by (ii), is attained exactly in the non-split case: Theorem A says the rank-one regulator may fail to be integral, and (ii) says which curves those are.

\textit{The baseline.} r − 2 in higher rank is the elementary count of §4.2 and is not ours.

\subsection*{3.3 The criterion: one element of F\_p}

Let \#E(F\_p) = p, P′ ∈ E(Q) with reduction of exact order p and satisfying (A2), P := pP′, e := v\_p(d(P)).

\begin{namedthm}[Lemma 4 (splitting and depth) [paper]]
E(Q\_p)[p] ≠ 0 ⟺ e ≥ 2.
\textbf{Lemma 4′ (the same, read locally) [paper].} Let (x) := ∏\_\{\{±T\} ⊂ E(F\_p)∖\{O\}\} (x − x(T)) ∈ F\_p[x], monic of degree (p−1)/2. Then ψ\_p ≡ u·\textasciicircum{}p (mod p) for some u ∈ F\_p\textasciicircum{}×; hence p ∣ ψ\_p′ and
\dispf{p² ∣ ψ\_p(x₀ + pk) − ψ\_p(x₀)  for all x₀, k ∈ Z\_p.}
Consequently, for an integral lift x₀ of a class  ∈ E(F\_p)∖\{O\}: v\_p(ψ\_p(x₀)) ≥ 1, \textbf{min(v\_p(ψ\_p(x₀)), 2) depends on  only} and not on the lift, and E(Q\_p)[p] ≠ 0 ⟺ v\_p(ψ\_p(x₀)) ≥ 2.
\end{namedthm}

The \textit{exact} value of v\_p(ψ\_p(x₀)) is \textbf{not} a function of  alone once it is ≥ 2 (§5.2), which is why only the cut at 2 is used; Lemma 4′ is what makes the measurements of §5.2 legitimate.

\begin{namedthm}[Theorem B (criterion) [paper; arithmetic step in Lean]]
Put λ := q\_p(φ\_p(P′))/2 ∈ F\_p. If e ≥ 2, or if the model is chosen with p ∣ a₁, then q\_p(α(P)/β(P)) = −λ, hence \textbf{c ≥ 2 ⟺ λ = 0}.
\textbf{Theorem B′ (the general model) [paper].} Without either hypothesis q\_p(α(P)/β(P)) = −λ − w₁/2, where w₁ := a₁·ω\_p(P′)·(ψ\_p(P′)/p)/φ\_p(P′)² mod p.
\textbf{Lemma 5 (λ is local and finite) [paper].} φ\_p mod p is a polynomial in x\textasciicircum{}p, so φ\_p(x) mod p² depends only on x mod p. Hence λ = λ(E mod p², ) is a function of the reduction of E modulo p² and of the class  ∈ E(F\_p)∖\{O\} alone: not of the global point, not of the height, not of the lift.
\end{namedthm}

By Theorem B and Lemma 3 the second p-adic digit of the height is λ, in a model chosen with p ∣ a₁. The rest of §3 determines λ.

\subsection*{3.4 The two directions of the criterion}

\begin{namedthm}[Theorem C (Vélu factorisation) [paper]]
Let K be a field of characteristic 0, E/K given by a Weierstrass equation, n ≥ 3 odd, D ⊂ E() a Galois-stable cyclic subgroup of order n, and T ∈ D∖\{O\} with x(T) ∈ K. Then, with  monic of degree n as in §2,
\dispf{φ\_n(x(T)) = (x(T))\textasciicircum{}n.}
\textbf{Corollary D [paper].} If p is anomalous and E(Q\_p)[p] ≠ 0, then λ() = 0 for every  ∈ E(F\_p)∖\{O\}.
\textbf{Corollary E [paper; arithmetic step in Lean].} In the situation of §3.3, if E is split at p then c ≥ 2.
\end{namedthm}

The converse holds, and the equivalence it produces is a known criterion for local torsion.

\begin{namedthm}[Proposition 6 (a known criterion, read as a Fermat quotient) [paper; known]]
Let p ≥ 5 be a prime of good ordinary reduction with \#E(F\_p) = p. Then for every  ∈ E(F\_p)∖\{O\},
\dispf{λ() = 0  ⟺  E(Q\_p)[p] ≠ 0.}
Neither side depends on the lift x₀ or on the model, and both are determined by E mod p².
\end{namedthm}

The equivalence is known. It is the case n = 1 of Gross's tameness criterion [Gro90, p. 514, §§14–15], quoted verbatim as [LR, Thm 4.1]: for E/Q with good ordinary reduction at p, j ≠ 0, 1728 and E[p] irreducible as a Gal(/Q)-module, a decomposition group D\_n at p is diagonalisable if and only if j\_E ≡ j₀ (mod p\textasciicircum{}\{n+1\}) for odd p, j₀ the j-invariant of the canonical lift. The two hypotheses there are not ours — j = 0 occurs among the CM curves of Theorem F — and the step from "D\_1 diagonalisable" to "E(Q\_p)[p] ≠ 0" is [OY, Lem. 4.1], which for good ordinary reduction and odd p gives E(F)[p] ≠ 0 ⟺ E[p] ≅ F\_p(1) ⊕ F\_p ⟺ E[p] tame with no hypothesis on j or on irreducibility, and quotes [Gro90, Prop. 13.7 and §17] for the equivalence with the congruence of j-invariants; and it is the content of [DW, Lemma 3.2 and Cor. 3.5], where the primes with E(Q\_p)[p] ≠ 0 are the \textit{local torsion primes} and the criterion reads: of the lifts of Ē to Z/p², exactly those congruent to the canonical lift have a rational point of order p. What we did not find in the sources we read is the presentation of that condition as the vanishing of a Fermat quotient of φ\_p at a single point (§9); the proof in §4.6 is independent of the sources above and gives the statement for all p ≥ 5 without a bound.

\subsection*{3.5 The non-split case: the exact value}

\begin{namedthm}[Theorem S [paper; arithmetic steps in Lean]]
Suppose \#E(F\_p) = p, p ∤ ∏\_ℓ c\_ℓ, and let P′ ∈ E(Q) have reduction of exact order p and reduce to a nonsingular point at every bad prime; put P := pP′. If E(Q\_p)[p] = 0 then
\dispf{v\_p(h\_p(P)) = 1 exactly,  hence  v\_p(h\_p(P′)) = −1.}
\end{namedthm}

So in the non-split case the height of a point of Λ∖Λ₀ is not p-integral, which is the direction ⟹ of (i) and the exact value in (ii). The invariance is worth stating: the quantity that Theorem S pins down is v\_p(h\_p(pP′)), not b or c separately, and the model with p ∣ a₁ is where the proof runs (§2, §4.7).

\subsection*{3.6 Two further properties of λ}

\begin{namedthm}[Theorem H (scaling) [paper; arithmetic step in Lean]]
Let p be anomalous, Q ∈ E(Z\_p) with  ≠ O, and m prime to p. Then q\_p(φ\_p(x(mQ))) = m²·q\_p(φ\_p(x(Q))), that is λ(m ) = m² λ().
\textbf{Corollary H′ [paper, in Lean].} λ vanishes at one nonzero class of E(F\_p) if and only if it vanishes at all of them; λ is a quadratic, not a linear, function on E(F\_p) ≅ Z/p.
\textbf{Corollary H″ (the same one digit deeper) [paper].} Let p be anomalous \textbf{and split at p}, Q ∈ E(Z\_p) with  ≠ O, and m prime to p. Then, with the second Fermat quotient q\textasciicircum{}\{(2)\} of §3.9 and λ\textasciicircum{}\{(2)\} its half,
\dispf{q\textasciicircum{}\{(2)\}(φ\_p(x(mQ))) = m²·q\textasciicircum{}\{(2)\}(φ\_p(x(Q))),  that is  λ\textasciicircum{}\{(2)\}(x(mQ)) = m²·λ\textasciicircum{}\{(2)\}(x(Q)).}
\textit{Proof.} Both sides are defined, since q\_p(φ\_p) = 2λ = 0 in the split case (Corollary D). In the identity (★) of §4.9, φ\_p(x(mQ))·ψ\_m(Q)\textasciicircum{}\{2p²\} = φ\_p(x(Q))\textasciicircum{}\{m²\}(1 + δ), the correction satisfies v\_p(δ) ≥ 2e, because the monic φ\_m of degree m² is evaluated at X = x(pQ) of valuation −2e; and e ≥ 2 by Lemma 4, E being split. Hence v\_p(δ) ≥ 4 \textgreater{} 3, so q\textasciicircum{}\{(2)\}(1 + δ) = 0; ψ\_m(Q)\textasciicircum{}\{2p²\} is a p-th power, and (v\textasciicircum{}\{p−1\})\textasciicircum{}\{2p²\} ∈ 1 + p³Z\_p for a unit v, so q\textasciicircum{}\{(2)\} kills it as well. Applying q\textasciicircum{}\{(2)\} to (★) exactly as q\_p was applied in Theorem H gives the claim. ∎
\textbf{Theorem L (the canonical subgroup contributes a p-th power) [paper].} Let C := Ê[p], which is Q\_p-rational for good ordinary reduction, π\_C : E → E″ := E/C the Vélu isogeny with x-part φ\_C/ψ\_C², and \_C the dual with x-part φ\_/ψ\_², the leading coefficient of φ\_ being A. For Q ∈ E(Z\_p) with  ≠ O and ξ := φ\_C()/ψ\_C()²,
\dispf{A·φ\_p() = (ψ\_C()²)\textasciicircum{}p · φ\_(ξ),}
so in Q\_p\textasciicircum{}×/(Q\_p\textasciicircum{}×)\textasciicircum{}p the class of φ\_p() is that of φ\_(ξ)/A, and A = c² with \_C\textasciicircum{}*ω\_E = c·ω\_\{E″\}.
\end{namedthm}

\subsection*{3.7 The CM case}

\begin{namedthm}[Theorem F [paper; known, = [DW, Prop. 2.1(3)]]]
Suppose E/Q has complex multiplication by an order in an imaginary quadratic field k, and p ≥ 5 is a prime of good reduction which is ordinary (equivalently, split in k) with a\_p = 1. Then E(Q\_p)[p] ≠ 0.
This is [DW, Prop. 2.1(3)], where it is stated for p \textgreater{} 3 with the converse as well: \textit{if E has complex multiplication, then a prime p \textgreater{} 3 of good reduction is a local torsion prime for E if and only if a\_p(E) = 1}, "local torsion prime" being E(Q\_p)[p] ≠ 0 by definition there. Three differences with the statement above are worth recording: [DW] needs only p \textgreater{} 3; it has the converse, which on our side is Lemma 4 (a\_p ≠ 1 leaves E(F\_p) with no p-torsion, so E(Q\_p)[p] = 0); and it needs no condition on the conductor f of the order — an earlier version of this note assumed p ∤ f, which for E/Q is automatic anyway, the thirteen possible orders having f ∈ \{1, 2, 3\}. The one-line reason given in [DW], that a CM elliptic curve is a canonical lift of each of its reductions (their Lemma 3.2), is the argument of §4.9 in another dress. The proof in §4.9 is kept because it is self-contained, not because the statement is ours.
\textbf{Corollary F′ [paper].} If in addition p ∤ ∏\_ℓ c\_ℓ, then v\_p(Reg\_p) ≥ max(r − 2, 0); in rank one, v\_p(Reg\_p) ≥ 0. So a CM curve never shows the deficit at a prime of good ordinary reduction.
\end{namedthm}

\subsection*{3.8 A rank-theoretic refinement of the baseline}

\begin{namedthm}[Theorem J [paper; determinant step in Lean]]
Assume \#E(F\_p) = p, ι, p ∤ ∏\_ℓ c\_ℓ. Let Λ₀ := ker(Λ → E(F\_p)), M\_\{Λ₀\} the height matrix on a basis of Λ₀ (entries in pZ\_p by Lemma 2), and ρ\_0 := rank\_\{F\_p\}((M\_\{Λ₀\}/p) mod p). Then v\_p(Reg\_p) = v\_p(det M\_\{Λ₀\}) − 2 ≥ 2r − 2 − ρ\_0.
\end{namedthm}

With ρ\_0 ≤ r this recovers Theorem A; the stronger bound max(2r−4, 0) that its shape suggests is \textbf{false} (§5.3). §3.9 replaces the bound by an equality, at the price of a quantity (det T) that is not local: in the notation there, ρ\_0 = 1 + rank(T mod p), and Theorem J is sharp exactly when T mod p is invertible.

\subsection*{3.9 The height matrix modulo p}

Theorem S pins the valuation of the height of one point. The next statements pin the whole matrix modulo p, on both sides of the dichotomy. Throughout, χ : Λ → E(F\_p) ≅ Z/p is the reduction, onto by ι, Λ₀ = ker χ, and a generator ₀ of E(F\_p) is fixed so that χ takes values in F\_p with χ(₀) = 1.

\begin{namedthm}[Theorem R (non-split) [paper; the linear algebra in Lean]]
Suppose \#E(F\_p) = p, p ∤ ∏\_ℓ c\_ℓ, ι, and E(Q\_p)[p] = 0. Then there is c ∈ F\_p\textasciicircum{}× with
\dispf{p·h\_p(P, R) ∈ Z\_p  and  p·h\_p(P, R) ≡ c·χ(P)·χ(R) (mod p)  for all P, R ∈ Λ.}
Hence v\_p(h\_p(P,R)) = −1 ⟺ χ(P)χ(R) ≠ 0, and v\_p(h\_p(P,R)) ≥ 1 when χ(P) = χ(R) = 0. In particular p·M\_Λ mod p is the matrix c·χχᵗ, of rank exactly one; and c is well defined in F\_p\textasciicircum{}×/(F\_p\textasciicircum{}×)² once the generator is forgotten, since replacing ₀ by m ₀ multiplies c by m\textasciicircum{}\{−2\}.
\textbf{Proposition R″ [paper].} In the situation of Theorem R, c = 2λ₁ where λ₁ := λ(₀).
\textbf{Corollary R′ [paper; the determinant bookkeeping in Lean].} In the situation of Theorem R take a basis of Λ with χ(e₁) = 1 and e₂, …, e\_r ∈ Λ₀. Then A := p·M\_Λ satisfies A₁₁ ∈ Z\_p\textasciicircum{}×, A\_\{1j\} ∈ pZ\_p (j ≥ 2) and A\_\{ij\} ∈ p²Z\_p (i, j ≥ 2), so A = D A′ D with D = diag(1, p, …, p) and A′ integral; with T the Schur complement of the unit A′₁₁,
\dispf{v\_p(Reg\_p) = (r − 2) + v\_p(det T),  so v\_p(Reg\_p) = r − 2 ⟺ det T ∈ Z\_p\textasciicircum{}×.}
For r = 1 the matrix T is empty and v\_p(Reg\_p) = −1 exactly, which is part (ii) of the main theorem again, this time without Theorem A.
\end{namedthm}

So in the non-split case the answer to "which entry carries the −1" is: exactly those (i, j) with χ(P\_i)χ(P\_j) ≠ 0, and the whole matrix modulo p is determined by the single element λ₁ ∈ F\_p, a Fermat quotient of φ\_p at one point — no rational point enters. What is \textit{not} determined by local data is det T, and that is where the exact value of v\_p(Reg\_p) still escapes (§7).

In the split case p·M\_Λ ≡ 0 and Theorem R says nothing; what carries information is M\_Λ itself, integral by part (i) of the main theorem.

\begin{namedthm}[Theorem R\_s (split) [paper; the linear algebra in Lean]]
Suppose \#E(F\_p) = p, p ∤ ∏\_ℓ c\_ℓ, ι, and E(Q\_p)[p] ≠ 0. Put  := M\_Λ mod p, a symmetric bilinear form on V := Λ/pΛ ≅ F\_p\textasciicircum{}r, and  the induced linear form, nonzero by ι. Then  vanishes on the hyperplane ker  = Λ₀/pΛ; consequently, for any v ∈ V with (v) = 1,
\dispf{(x, y) = (x, v)·(y) + (x)·(y, v) − (v, v)·(x)·(y),}
so  is expressed by the two linear forms  and (·, v) and \textbf{ρ₁ := rank  ≤ 2}. In a basis with χ(e₁) ≠ 0 and e₂, …, e\_r ∈ Λ₀ one has  = [[a, bᵗ], [b, 0]] with a = h\_p(e₁) mod p and b\_j = h\_p(e₁, e\_j) mod p, and
\dispf{ρ₁ = 2 ⟺ b ≠ 0;  ρ₁ = 1 ⟺ b = 0 and a ≠ 0;  ρ₁ = 0 ⟺ b = 0 and a = 0.}
\end{namedthm}

\begin{namedthm}[Corollary R′\_s [paper; the p\textasciicircum{}\{r−2\} step in Lean]]
In the situation of Theorem R\_s,
\dispf{v\_p(Reg\_p) ≥ r − ρ₁,}
and more precisely: if ρ₁ = 2 then v\_p(Reg\_p) = (r − 2) + v\_p(det T\_s) with T\_s an integral (r−2)×(r−2) Schur complement; if ρ₁ = 1 then v\_p(Reg\_p) ≥ r − 1; if ρ₁ = 0 then v\_p(Reg\_p) ≥ r. For r = 2 this says v\_p(Reg\_p) = 0 ⟺ ρ₁ = 2.
\end{namedthm}

Since ρ₁ ≤ 2 always, Corollary R′\_s contains part (iii) of the main theorem — v\_p(Reg\_p) ≥ max(r − 2, 0) — and adds when the bound is attained. The two sides now read alike: non-split gives rank exactly 1 for p·M\_Λ mod p, split gives rank at most 2 for M\_Λ mod p, and in both the deficit of v\_p(Reg\_p) below r − 2 is zero while the excess is the valuation of a Schur complement.

The constant a is again a Fermat quotient, one digit deeper. In the split case λ = 0 by Corollary D, so φ\_p(P′)\textasciicircum{}\{p−1\} ∈ 1 + p²Z\_p and the \textbf{second Fermat quotient}

\begin{quote}
q\textasciicircum{}\{(2)\}(u) := (u\textasciicircum{}\{p−1\} − 1)/p² mod p,  defined for u ∈ Z\_p\textasciicircum{}× with q\_p(u) = 0,
\end{quote}

is available; it is a homomorphism on that subgroup, since (1+p²α)(1+p²β) = 1 + p²(α + β + p²αβ) (the Lean statement \texttt{fermatQuot2\_\allowbreak{}add}).

\begin{namedthm}[Proposition R″\_s [paper]]
In the situation of Theorem R\_s let e₁ ∈ Λ with χ(e₁) ≠ 0, n := ∏\_ℓ c\_ℓ, P′ := n e₁, and λ\textasciicircum{}\{(2)\} := q\textasciicircum{}\{(2)\}(φ\_p(P′))/2 ∈ F\_p. Then
\dispf{a = h\_p(e₁) mod p = 2λ\textasciicircum{}\{(2)\}/n².}
In particular for r = 1: \textbf{v\_p(Reg\_p) = 0 ⟺ λ\textasciicircum{}\{(2)\} ≠ 0.}
\end{namedthm}

The analogy with Proposition R″ is exact — c = 2λ₁ there, a = 2λ\textasciicircum{}\{(2)\}/n² here — with one difference that decides §7(3). λ is a function of E mod p² and of the class  alone (Lemma 5, because p ∣ φ\_p′), whereas φ\_p′ is \textit{not} divisible by p², so λ\textasciicircum{}\{(2)\} moves when the lift x₀ is replaced by x₀ + pk: it is a function of E mod p³ and of the point P′ modulo p², not of . There is therefore no law for the split value of v\_p(Reg\_p) in terms of the reduction alone, and §7(3) is answered in that negative form.

One consequence of Proposition R″ can be stated for a whole isogeny class. Let φ : E → E″ be an isogeny over Q of degree m prime to p. Then a\_p, hence anomalousness, is shared; the reduction  : E(F\_p) → E″(F\_p) is an isomorphism, its kernel having order dividing m; and E(Q\_p)[p] ≅ E″(Q\_p)[p], so splitting is shared as well.

\begin{namedthm}[Corollary R‴ [paper, conditional; computation]]
Assume in addition the functoriality h\_p\textasciicircum{}\{E″\}(φ(P)) = m·h\_p\textasciicircum{}\{E\}(P) of the cyclotomic p-adic height. Then, in the non-split case, λ\_\{E″\}(()) = m·λ\_E() for every  ∈ E(F\_p)∖\{O\}; equivalently, since a change of generator multiplies λ₁ by a square,
\dispf{(λ₁(E)/p)·(λ₁(E″)/p) = (m/p)   in \{±1\}.}
So the quadratic residue class of λ₁ is constant on the isogeny class when every degree is a square modulo p, and splits it into two otherwise.
\end{namedthm}

The functoriality assumed there is the expected one for the canonical cyclotomic height and we did not verify it in a source, which is why the grade is conditional; the identity itself is [computation] over 2,098 pairs (§5.6). A direct local proof of λ\_\{E″\}∘ = m·λ\_E — one that does not pass through a global point and the height — we do not have.

Two further consequences of Proposition R″\_s are quoted throughout §§3.10–3.12 and §4.11, and we state them here. The first is the behaviour of λ\textasciicircum{}\{(2)\} under a change of lift — Proposition 9 of §3.10 read modulo p — and it is stated separately because what the later sections use is its \textit{slope}.

\begin{namedthm}[Proposition R⁗ (the slope) [paper; the arithmetic in Lean]]
Let p be anomalous and split,  ∈ E(F\_p)∖\{O\}, x₀ ∈ Z\_p a lift of , and G := φ\_p′/p, integral by Lemma 5. Then for every k ∈ Z\_p
\dispf{λ\textasciicircum{}\{(2)\}(x₀ + pk) = λ\textasciicircum{}\{(2)\}(x₀) + k·s(),  s() := −G()/(2φ\_p()) ∈ F\_p,}
and \textbf{s() ≠ 0}; so k ↦ λ\textasciicircum{}\{(2)\}(x₀ + pk) is a bijection of F\_p and λ\textasciicircum{}\{(2)\}, unlike λ, does depend on the lift.
\textit{Proof.} The congruence is Proposition 9 (§3.10, proved in §4.11) read modulo p: there Θ(x₀ + pk) ≡ Θ(x₀) + (p−1)kG(x₀)/u − pk²H(x₀)/(2u) (mod p²) with u := φ\_p(x₀) and Θ ≡ 2λ\textasciicircum{}\{(2)\} (mod p), whence 2λ\textasciicircum{}\{(2)\}(x₀ + pk) ≡ 2λ\textasciicircum{}\{(2)\}(x₀) − kG()/φ\_p(); the denominator is a unit by Corollary D. The non-vanishing of s is the non-degeneracy (ND) of §3.13 — φ\_p() being a unit, s() ≠ 0 ⟺ G() ≠ 0 ⟺ v\_p(φ\_p′(x₀)) = 1 — and (ND) holds for every p ≥ 5 by Theorem ND (§3.14). ∎
\end{namedthm}

\begin{namedthm}[Corollary R⁗′ (localisation of the off-diagonal block) [paper; the identity in Lean]]
In the situation of Theorem R\_s take a basis of Λ with χ(e₁) = 1 and e₂, …, e\_r ∈ Λ₀, put n := ∏\_ℓ c\_ℓ and, for j ≥ 2,
\dispf{κ\_j := (x(n(e₁ + e\_j)) − x(n e₁))/p mod p ∈ F\_p.}
Then b\_j = h\_p(e₁, e\_j) mod p satisfies \textbf{b\_j = s(₀)·κ\_j/n²} with ₀ the reduction of x(n e₁); hence \textbf{ρ₁ = 2 ⟺ κ\_j ≠ 0 for some j}.
\textit{Proof.} Both x(n(e₁ + e\_j)) and x(n e₁) are lifts of the same class ₀, since χ(e\_j) = 0, and they differ by pκ\_j; so κ\_j is well defined. Polarising, 2b\_j ≡ a(e₁ + e\_j) − a(e₁) − a(e\_j) with a(P) := h\_p(P) mod p. Here a(e\_j) = 0 by Lemma 2, and by Proposition R″\_s a(e₁ + e\_j) − a(e₁) = 2(λ\textasciicircum{}\{(2)\}(x(n(e₁+e\_j))) − λ\textasciicircum{}\{(2)\}(x(n e₁)))/n², which by Proposition R⁗ with k = κ\_j is 2s(₀)κ\_j/n². The equivalence is the last display of Theorem R\_s together with s ≠ 0. ∎
\end{namedthm}

Everything in §3.9 assumes \#E(F\_p) = p. One line of Hasse's bound says how much room there is beyond that, and the statements of §3.9 fill it.

\begin{namedthm}[Proposition R⁺ (the theorems of §3.9 read on the p-part of the reduction) [paper; the counting in Lean]]
Let p ≥ 5 have good ordinary reduction, p ∣ N := \#E(F\_p), p ∤ ∏\_ℓ c\_ℓ, and put m := N/p. Let χ : Λ → E(F\_p)[p] ≅ Z/p be the p-part of the reduction and assume χ onto. Then Theorem S, Theorem R, Proposition R″ and Corollary R′ hold as stated.
Moreover \textbf{p ∥ N always}, and the only case not already covered by \#E(F\_p) = p is
\dispf{\textbf{N = 2p, which by Hasse forces p = 5} (Lean \texttt{card\_\allowbreak{}cases}, \texttt{two\_\allowbreak{}p\_\allowbreak{}forces\_\allowbreak{}five}, \texttt{sq\_\allowbreak{}not\_\allowbreak{}dvd}).}
\end{namedthm}

\begin{mdproof}[Proof]
Three changes. (i) If χ(P) = 0 then the order of the reduction of P divides m, so mnP reduces to O and satisfies (A2) as well, Lemma 2 puts h\_p(mnP) ∈ 2pZ\_p, and h\_p(P) = h\_p(mnP)/(mn)² with p ∤ mn gives h\_p(P) ∈ 2pZ\_p (Lean \texttt{dvd\_\allowbreak{}of\_\allowbreak{}unit\_\allowbreak{}sq\_\allowbreak{}mul}). Lemma 2 needs both (A1) and (A2), which is why the Tamagawa factor n is carried here as everywhere else. (ii) If χ(P) ≠ 0 then P′ := mnP has reduction of order exactly p, so Theorem S applies to P′ unchanged and v\_p(h\_p(P)) = −1 exactly. (iii) The linear algebra (R0) is untouched: pΛ ⊆ ker χ and ker  is a hyperplane, as before. For the numerical claim, p ∣ N and p ≥ 5 with N ≤ p + 1 + 2√p \textless{} p² give p² ∤ N, so N ∈ \{p, 2p\}; and N = 2p forces (p − 1)² ≤ 4p, hence p = 5. ∎
\end{mdproof}

So the hypothesis \#E(F\_p) = p of the main theorem excludes exactly one family — p = 5 with a\_5 = −4 — and on that family the statements of §3.9 survive in the form above. §5.3 records the rank-four curve where this is what is needed.

\subsection*{3.10 The height matrix modulo p in closed form, and the second digit}

Here E is split at p, \#E(F\_p) = p, p ∤ ∏\_ℓ c\_ℓ, ι holds, and n := ∏\_ℓ c\_ℓ. Two steps of the filtration of E(Q\_p) carry the whole of ; the \textit{second} p-adic digit of the height then leaves the filtration altogether.

\textbf{Two homomorphisms.} By Lemma 4 the group D := E(Q\_p)[p] has order p and D ∩ E\_1(Q\_p) = 0, while D∖\{O\} → E(F\_p)∖\{O\} is a bijection (proof of Corollary D); hence D → E(F\_p) is an isomorphism and

\begin{quote}
E(Q\_p) = D ⊕ E\_1(Q\_p).
\end{quote}

Fix T₀ ∈ D reducing to the chosen generator G of E(F\_p), let χ be the discrete logarithm of the reduction with respect to G, and for Q ∈ E(Z\_p) with  ≠ O write Q = χ(Q)T₀ + R with R ∈ E\_1(Q\_p) — well defined, since pT₀ = O — and

\begin{quote}
γ(Q) := t(R)/p mod p ∈ F\_p,  t = −x/y the formal parameter.
\end{quote}

Then χ and γ are homomorphisms, γ(T₀) = 0, γ|\_\{E\_1(Q\_p)\} = t/p mod p, and ker χ ∩ ker γ = E\_2(Q\_p).

\begin{namedthm}[Lemma 7 (the first digit of the x-coordinate) [paper; the arithmetic step in Lean]]
Let P ∈ E(Z\_p) with  ≠ O and let R ∈ E\_1(Q\_p) with t := t(R). Then R ↦ x(P + R) lies in Z\_p[[t]] and
\dispf{x(P + R) = x(P) + c₁(P)·t + O(t²),  c₁(P) := 2y(P) + a₁x(P) + a₃.}
Since \#E(F\_p) = p is odd, Ē(F\_p)[2] = \{O\}, so \textbf{c₁(P) ∈ Z\_p\textasciicircum{}×}.
\end{namedthm}

With κ\_j as in Corollary R⁗′, Lemma 7 gives κ\_j = c₁·γ(n e\_j), so κ\_j ≠ 0 ⟺ v\_p(t(n e\_j)) = 1, and Corollary R⁗′ turns into

\begin{quote}
\textbf{ρ₁ = 2 ⟺ γ|\_\{Λ₀\} ≠ 0 ⟺ Λ₀ ⊄ E\_2(Q\_p)},
\end{quote}

a condition on the image of Λ in the \textit{second} layer of the filtration, of exactly the same nature as the hypothesis ι in the first.

\begin{namedthm}[Proposition 8 (two purely local identities) [paper]]
For every class  ∈ E(F\_p)∖\{O\}, with s as in Proposition R⁗, c₁ as in Lemma 7, and λ\_T() := λ\textasciicircum{}\{(2)\}(x(T)) for the unique T ∈ D∖\{O\} above ,
\dispf{(L1) s()·c₁() = B₀·χ(),  (L2) λ\_T() = λ\_\{T,1\}·χ()²,}
where B₀ := s(G)·c₁(G) and λ\_\{T,1\} := λ\_T(G) depend only on E/Q\_p and on the choice of G. No global point and no height enter.
\end{namedthm}

\begin{namedthm}[Theorem K (the height matrix modulo p) [paper; the bilinear algebra in Lean]]
Under the hypotheses of Theorem R\_s, for all P, R ∈ Λ,
\dispf{h\_p(P) ≡ 2λ\_\{T,1\}·χ(P)² + 2B₀·χ(P)γ(P) (mod p),}
\dispf{(P, R) = 2λ\_\{T,1\}·χ(P)χ(R) + B₀·(χ(P)γ(R) + γ(P)χ(R)).}
The constants λ\_\{T,1\} and B₀ depend only on E/Q\_p; the only global input is the pair of homomorphisms (χ, γ) : Λ → F\_p × F\_p, and \textbf{the Tamagawa product n cancels}.
\end{namedthm}

Theorem K contains Theorem R\_s (the form is built from two linear functionals, so ρ₁ ≤ 2), Corollary R′\_s, Proposition R″\_s and Corollary R⁗′ in one line. Its consequences are read off:

\begin{namedthm}[Corollary K′ [paper]]
(i) ρ₁ = 2 ⟺ (χ, γ) : Λ → F\_p² is onto; (ii) if γ = cχ on Λ then  = (2λ\_\{T,1\} + 2B₀c)·χ⊗χ, of rank ≤ 1, and of rank 0 exactly when c = −λ\_\{T,1\}/B₀; (iii) in rank one, \textbf{v\_p(Reg\_p) ≥ 1 ⟺ γ(e₁)/χ(e₁) = −λ\_\{T,1\}/B₀}, one residue class out of p; (iv) with Λ₁ := Λ ∩ E\_2(Q\_p) = ker χ ∩ ker γ of rank r − 2 when (χ, γ) is onto, and Ĉ the Gram matrix of p\textasciicircum{}\{−1\}h\_p on Λ₁, the Schur complement of Corollary R′\_s satisfies T\_s ≡ p·Ĉ (mod p²), so
\dispf{\textbf{v\_p(Reg\_p) = r − 2 ⟺ det Ĉ ≢ 0 (mod p)}.}
\end{namedthm}

So the deficit r − 2 is exactly what the two local layers χ and γ take away, and the residual global content of v\_p(Reg\_p) in rank ≥ 3 is a non-degeneracy of p\textasciicircum{}\{−1\}h\_p one layer deeper — a statement of Schneider's kind, moved but not removed. The exact formula v\_p(Reg\_p) = (r − 2) + v\_p(det Ĉ) needs det Ĉ to be a unit and fails otherwise (§5.7).

\textbf{The second digit.} Theorem K settles h\_p modulo p. The natural continuation would be a third homomorphism γ₂ : Λ₁ → E\_2/E\_3, and the natural guess is that v\_p(Reg\_p) ≥ 2 is a condition on the image of Λ in E(Q\_p)/E\_3. It is not. Write, for a unit u with q\_p(u) = 0,

\begin{quote}
Q₂(u) := (u\textasciicircum{}\{p−1\} − 1)/p² ∈ Z\_p,  Θ(x₀) := Q₂(φ\_p(x₀)),
\end{quote}

so that q\textasciicircum{}\{(2)\}(u) = Q₂(u) mod p and Θ(x₀) ≡ 2λ\textasciicircum{}\{(2)\}(x₀) (mod p). Θ is defined at every lift x₀ of a class of E(F\_p)∖\{O\} in the split case, by Corollary D and Lemma 5.

\begin{namedthm}[Theorem N (the second digit) [paper]]
In the situation of Theorem R\_s let P ∈ Λ with χ(P) ≠ 0, put P′ := nP and write x(P′) = a/d′² in lowest terms (so p ∤ d′). Then
\dispf{\textbf{n²·h\_p(P) ≡ (1 + p)·Θ(x(P′)) + 2p·q\_p(d′)  (mod p²)},}
where q\_p(d′) = (d′\textasciicircum{}\{p−1\} − 1)/p mod p is the Fermat quotient of the rational integer d′.
\end{namedthm}

Modulo p this is Proposition R″\_s, since Θ ≡ 2λ\textasciicircum{}\{(2)\}. Modulo p² a second term appears, and it is not a function of the image of P in E(Q\_p) at all: d′ is the denominator of the point, a global integer, and two points with the same image in E(Q\_p)/E\_k for every k may have different q\_p(d′). Hence:

\begin{namedthm}[Remark (no third homomorphism)]
The localisation of Theorem K does not continue past the first digit. There is no γ₂ in the answer: v\_p(Reg\_p) ≥ 2 is not a condition on the position of Λ in the filtration of E(Q\_p), and the induction "layer k ↦ a homomorphism to E\_k/E\_\{k+1\}" stops at k = 1. What replaces it is the pair (a local p-adic number, the Fermat quotient of a global denominator); see Corollary N′.
\end{namedthm}

The local half of the second digit is nevertheless completely described, and by a single p-adic number.

\begin{namedthm}[Proposition 9 (the second digit under lifting) [paper; the arithmetic step in Lean]]
With G := φ\_p′/p (integral by Lemma 5), H := dG/dx and u := φ\_p(x₀), for every k ∈ Z\_p
\dispf{Θ(x₀ + pk) ≡ Θ(x₀) + (p − 1)·k·G(x₀)/u − p·k²·H(x₀)/(2u)  (mod p²),}
where G(x₀) and u are needed modulo p² and H(x₀) modulo p. Modulo p this is Proposition R⁗, with slope s = −G()/(2u).
\end{namedthm}

\begin{namedthm}[Proposition 10 (the bad lift) [paper; the arithmetic step in Lean]]
Fix a class  ∈ E(F\_p)∖\{O\} and T ∈ D∖\{O\} above it. Then δ ↦ Θ(x(T) + pδ) is a bijective isometry of Z\_p; in particular it has a unique zero δ\_*(E, p, ) ∈ Z\_p, a local invariant, and for every lift x₀ = x(T) + pδ
\dispf{\textbf{v\_p(Θ(x₀)) = v\_p(δ − δ\_*)}.}
\end{namedthm}

Proposition 10 is the statement that all layers of the local half are governed by the distance to one p-adic number: no new local datum appears at any depth. Combining it with Theorem N gives the criterion at the second layer.

\begin{namedthm}[Corollary N′ (rank one, second layer) [paper]]
Let r = 1 and v\_p(Reg\_p) ≥ 1, that is p ∣ Θ(x(P′)) by Proposition R″\_s. Then
\dispf{\textbf{v\_p(Reg\_p) ≥ 2 ⟺ Θ(x(P′))/p + 2·q\_p(d′) ≡ 0 (mod p)},}
and v\_p(Θ(x(P′)))/p is computed from δ\_* and δ by Proposition 10. So the second layer is again one residue class out of p, but the class is displaced by a quantity — 2q\_p(d′) — that is invisible in E(Q\_p).
\end{namedthm}

Finally the bookkeeping of Lemma 3 needs one correction before Theorem N can be proved, and it is the reason the model does not appear in Theorem N.

\begin{namedthm}[Lemma 3′ (the formal part, sharpened) [paper]]
In the notation of Lemma 3, b = v\_p(Λ\_σ) ≥ min(e + v\_p(a₁), 2e). In particular the bound b ≥ 2e stated in Lemma 3 for a model with p ∣ a₁ is correct for e = 1 and too strong for e ≥ 2, where the right bound is e + v\_p(a₁). It is enough for Proposition R″\_s, which needs b ≥ 3; at the next digit the term a₁·t(pP′) of 2Λ\_σ has to be carried, and it cancels against the term a₁φωψ of the Weierstrass equation (§4.11, Lean \texttt{a1\_\allowbreak{}correction\_\allowbreak{}cancels}).
\end{namedthm}

\subsection*{3.11 The ledger in closed form}

Notation as in §3.10: P ∈ Λ with χ(P) ≠ 0, P′ := nP, P″ := pP′, e := v\_p(ψ\_p(x(P′))) = v\_p(d(P″)) ≥ 2, x(P′) = a/d′² in lowest terms. Write

\begin{quote}
Θ := Q₂(φ\_p(x(P′))),  F := (d′\textasciicircum{}\{p−1\} − 1)/p ∈ Z,  c := (b₂ − E₂(E, ω))/12,  W := c/(p⁴·x(P″)),
\end{quote}

so that q\_p(d′) = F mod p, and 1/x(P″) = ψ\_p(x(P′))²/φ\_p(x(P′)) has valuation 2e ≥ 4, so that W ∈ Z\_p. Let z := log\_E(t) be the formal logarithm, normalised by dz = ω, so that v\_p(z(Q)) = v\_p(t(Q)) for Q ∈ E₁(Q\_p). Put X := x(P″), Y := y(P″), t := t(P″), z := z(P″) and

\begin{quote}
R := a₂/X + a₄/X² + a₆/X³ − a₁Y/X² − a₃Y/X³,
\end{quote}

so that ω\_p(P′)² = φ\_p(x(P′))³·(1 + R) is the Weierstrass equation of E cleared of denominators at P″.

The exact identity of the Remark opening §4 can be summed instead of expanded. Summed, it has no digits in it: what is left is two Iwasawa logarithms of explicit units and one power series in the formal parameter.

\begin{namedthm}[Theorem N* (the ledger in closed form) [paper]]
Under the hypotheses of Theorem N, as an identity in Q\_p,
\dispf{\textbf{n²·h\_p(P) = (1/(1 − p))·[ p\textasciicircum{}\{−2\}·log(1 + p²Θ) + 2·log(1 + pF) ] + p\textasciicircum{}\{−2\}·ℒ(z)},}
where
\dispf{\textbf{ℒ(z) := log( x(P″)·σ\_p(t(P″))² ) = 2·log(σ\_p(t)/t) − log(1 + R)},}
both logarithms being Iwasawa's. Moreover v\_p(ℒ(z)) = 2e whenever p ∤ c.
\end{namedthm}

\begin{mdproof}[Proof]
The Remark opening §4 gives p²·n²h\_p(P) = 2log(σ\_p(t)/t) + 2log(φ/ω) − 2p²·log(d′) with φ := φ\_p(x(P′)), ω := ω\_p(P′). From ω² = φ³(1 + R) one has 2log(φ/ω) = log(φ²/ω²) = −log φ − log(1 + R). For a unit u, log u = (p − 1)\textasciicircum{}\{−1\}·log(u\textasciicircum{}\{p−1\}), so with φ\textasciicircum{}\{p−1\} = 1 + p²Θ and d′\textasciicircum{}\{p−1\} = 1 + pF,
\end{mdproof}

\begin{quote}
−log φ − 2p²·log d′ = −(p − 1)\textasciicircum{}\{−1\}·[ log(1 + p²Θ) + 2p²·log(1 + pF) ].
\end{quote}

Dividing by p² and using (1 − p)\textasciicircum{}\{−1\} = −(p − 1)\textasciicircum{}\{−1\} gives the display, with ℒ(z) = 2log(σ\_p(t)/t) − log(1 + R). The two expressions for ℒ agree because

\begin{quote}
t² = X²/Y² = (φ/ψ²)²·(ψ³/ω)² = φ²ψ²/ω² = ψ²/(φ·(1 + R)) = 1/(X·(1 + R)),
\end{quote}

where ψ := ψ\_p(x(P′)), using t = −x/y, x = φ/ψ², y = ω/ψ³ at P″ and ω² = φ³(1 + R). Hence X·σ\_p(t)² = (σ\_p(t)/t)²·Xt² = (σ\_p(t)/t)²/(1 + R), and taking logarithms gives the second form. The valuation is Lemma 1 together with Proposition 17. ∎

The series ℒ converges at z = z(P″). By Proposition 17 its m-th coefficient has v\_p(μ\_\{2m\}) ≥ −v\_p((2m + 2)!) ≥ −(2m + 2)/(p − 1), while v\_p(z\textasciicircum{}\{2m\}) = 2me with e ≥ 2 and p ≥ 5, so the m-th term has valuation ≥ 2me − (2m + 2)/(p − 1) → ∞: reading ℒ term by term modulo p\textasciicircum{}k, as Proposition 18 does, is legitimate. As a series in the formal parameter t the convergence is immediate instead, both 2log(σ\_p(t)/t) and log(1 + R) being logarithms of 1 + (something of positive valuation).

Nothing in Theorem N* refers to a digit. The two logarithms are functions of the two integers Θ and F alone; the third term is a function of the position of P″ in the formal group, and the next two propositions say what kind of function.

\begin{namedthm}[Proposition 17 (ℒ is an even series with invariant coefficients) [paper; the coefficient identities in Lean]]
As a power series in z,
\dispf{\textbf{ℒ(z) = μ₂z² + μ₄z⁴ + μ₆z⁶ + ⋯},  \textbf{no odd power occurring},}
and the coefficients are invariants of E over Q:
\dispf{\textbf{μ₂ = −c},  \textbf{μ₄ = −b₄/12},  \textbf{μ₆ = −b₂b₄/360 − b₆/30},}
with \textbf{μ\_\{2m\} ∈ Z[1/(2m + 2)!][b₂, b₄, b₆, E₂]} — so v\_p(μ\_\{2m\}) ≥ −v\_p((2m + 2)!), which for p ≥ 5 is ≥ −(2m + 2)/(p − 1) — and \textbf{E₂ occurring in μ₂ alone}. The denominators are \textbf{not} bounded independently of m: μ₆ has a 5, μ₈ a 7, and μ₁₄ an 11 and a 13. That is the mechanism by which the digit at which a coefficient becomes visible can depend on p (Proposition 18).
\end{namedthm}

\begin{mdproof}[Proof]
Mazur and Tate's characterisation of σ\_p is the differential equation (d/dz)²·log σ\_p(z) = −(x(z) + c), and x(z) + b₂/12 = ℘(z) with ℘(z) = z\textasciicircum{}\{−2\} + Σ\_\{k≥1\}π\_kz\textasciicircum{}\{2k\} the Weierstrass series of the lattice of E, π₁ = c₄/240, π₂ = c₆/6048. Write a := −b₂/12. Then
\end{mdproof}

\begin{quote}
(d/dz)²·log(σ\_p(z)/z) = z\textasciicircum{}\{−2\} − ℘(z) − a − c = −(a + c) − Σ\_\{k≥1\}π\_kz\textasciicircum{}\{2k\},
\end{quote}

whose right-hand side is even; σ\_p is odd and σ\_p(z)/z ∈ 1 + zZ\_p[[z]], so log(σ\_p(z)/z) is even with no constant term, and integrating twice gives

\begin{quote}
2·log(σ\_p(z)/z) = −(a + c)z² − 2·Σ\_\{k≥1\}π\_kz\textasciicircum{}\{2k+2\}/((2k + 1)(2k + 2)).
\end{quote}

Also x(z)z² = z²℘(z) + az² = 1 + az² + π₁z⁴ + π₂z⁶ + ⋯ is a series in z², so log(x(z)z²) is even, and ℒ = log(x(z)z²) + 2log(σ\_p(z)/z) is even: that is the first claim. For the coefficients, expand log(1 + az² + π₁z⁴ + π₂z⁶ + ⋯):

\begin{itemize}[leftmargin=1.4em,itemsep=.2em,topsep=.3em]
\item z²: a from the first logarithm and −(a + c) from the second, so \textbf{μ₂ = −c}; explicitly −b₂/12 + E₂/12 = −(b₂ − E₂)/12, since −(a + c) = E₂/12 (Lean \texttt{mu2\_\allowbreak{}eq});
\item z⁴: π₁ − a²/2 from the first, −π₁/6 from the second, so \textbf{μ₄ = π₁ − a²/2 − π₁/6 = (c₄ − b₂²)/288 = −b₄/12} by c₄ = b₂² − 24b₄ — and \textbf{E₂ has cancelled} (Lean \texttt{mu4\_\allowbreak{}eq});
\item z⁶: π₂ − aπ₁ + a³/3 from the first, −π₂/15 from the second; with c₄ = b₂² − 24b₄ and c₆ = −b₂³ + 36b₂b₄ − 216b₆ the three terms in b₂³ cancel exactly (Lean \texttt{mu6\_\allowbreak{}no\_\allowbreak{}b2\_\allowbreak{}cube}: −1/6480 + 1/2880 − 1/5184 = 0), leaving \textbf{μ₆ = −b₂b₄/360 − b₆/30} (Lean \texttt{mu6\_\allowbreak{}eq}).
\end{itemize}

Only b₂, b₄, b₆ and, through c, E₂ enter. For the denominators, π\_k is a polynomial in g₂ = c₄/12 and g₃ = c₆/216 with denominator dividing (2k + 2)!, and the second integration divides the coefficient of z\textasciicircum{}\{2k+2\} by (2k + 1)(2k + 2); the logarithm of 1 + az² + π₁z⁴ + ⋯ contributes no worse denominator, whence μ\_\{2m\} ∈ Z[1/(2m + 2)!][b₂, b₄, b₆, E₂]. ∎

The coefficients μ\_\{2m\} carry even numerical subscripts throughout; they are not the unit series M(t) of §4.12 nor the homomorphism μ of §7.

\begin{namedthm}[Proposition 18 (which terms are visible at the k-th digit) [paper; the valuation arithmetic in Lean]]
In the situation of Theorem N*, reading the right-hand side modulo p\textasciicircum{}k leaves exactly the terms of valuation ≤ k − 1, and the valuations are
| series | j-th (m-th) term | valuation (equality when v\_p(Θ) = v\_p(F) = 0) | |---|---|---| | p\textasciicircum{}\{−2\}·log(1 + p²Θ) | (−1)\textasciicircum{}\{j+1\}p\textasciicircum{}\{2j−2\}Θ\textasciicircum{}j/j | \textbf{2j − 2 − v\_p(j)} | | 2·log(1 + pF) | 2(−1)\textasciicircum{}\{j+1\}p\textasciicircum{}jF\textasciicircum{}j/j | \textbf{j − v\_p(j)} | | p\textasciicircum{}\{−2\}·ℒ(z) | μ\_\{2m\}z\textasciicircum{}\{2m\}/p² | \textbf{2me − 2 + v\_p(μ\_\{2m\})} |
In general the first two entries are \textbf{lower bounds}, the excess being j·v\_p(Θ) and j·v\_p(F); the non-generic case occurs in this paper, since Corollaries N″ and N‴ treat boxes with Θ = pΘ₁. A term of valuation v is first visible at the \textbf{(v + 1)-th} digit, that is modulo p\textasciicircum{}\{v+1\}; in this counting Θ occupies the first digit and F the second, which is Proposition R″\_s and Theorem N. In the generic case e = 2 the m-th coefficient of ℒ is first visible at the \textbf{(4m − 1 + v\_p(μ\_\{2m\}))-th} digit: μ₂ = −c at the \textbf{third}, μ₄ = −b₄/12 at the \textbf{seventh}, μ₆ at the \textbf{tenth} for p = 5 and at the \textbf{eleventh} for p ≥ 7.
\end{namedthm}

\begin{mdproof}[Proof]
The three lines are the term-by-term valuations of the Iwasawa logarithm, with v\_p(Θ) ≥ 0, v\_p(F) ≥ 0 and v\_p(z) = e; the first two are equalities exactly when Θ and F are units, and lower bounds otherwise. The arithmetic of the subscript is Lean \texttt{padicVal\_\allowbreak{}eq\_\allowbreak{}zero\_\allowbreak{}of\_\allowbreak{}lt}, \texttt{log\_\allowbreak{}term\_\allowbreak{}depth\_\allowbreak{}small}, \texttt{log\_\allowbreak{}term\_\allowbreak{}depth\_\allowbreak{}five}, \texttt{log\_\allowbreak{}term\_\allowbreak{}depth\_\allowbreak{}ne\_\allowbreak{}five}, \texttt{log\_\allowbreak{}term\_\allowbreak{}factor}, \texttt{theta\_\allowbreak{}depth\_\allowbreak{}j5\_\allowbreak{}p5}, \texttt{theta\_\allowbreak{}depth\_\allowbreak{}j5\_\allowbreak{}p7}, \texttt{log\_\allowbreak{}term\_\allowbreak{}j10\_\allowbreak{}five}, \texttt{log\_\allowbreak{}term\_\allowbreak{}j10\_\allowbreak{}seven}; the three digits of the last sentence are \texttt{lambda\_\allowbreak{}depth\_\allowbreak{}mono}, \texttt{lambda\_\allowbreak{}depth\_\allowbreak{}m2\_\allowbreak{}e2}, \texttt{val\_\allowbreak{}thirty}, \texttt{val\_\allowbreak{}three\_\allowbreak{}sixty}, \texttt{val\_\allowbreak{}thirty\_\allowbreak{}seven}, \texttt{mu6\_\allowbreak{}visible\_\allowbreak{}at\_\allowbreak{}ten}. The factor v\_p(j) is the whole content of the first two lines: the j-th term of a logarithm is one digit shallower than p\textasciicircum{}j whenever p ∣ j, which is what puts the term with j = p inside the p\textasciicircum{}\{j−1\} digit. The third line is why μ₆, whose denominator is divisible by p for p = 5 alone, is visible one digit earlier at p = 5 than at p ≥ 7. ∎
\end{mdproof}

Theorem N* also fixes the coordinate in which the deep digits are to be read. The quantity W of §3.11 uses 1/x(P″) where ℒ uses z², and the two are not equal.

\begin{namedthm}[Lemma 21 (the two coordinates) [paper]]
As power series in the formal logarithm,
\dispf{\textbf{z² − 1/x(z) = −(b₂/12)·z⁴ + O(z⁶)},  \textbf{no odd power occurring}.}
Hence v\_p(z(P″)² − 1/x(P″)) ≥ 4e, and p\textasciicircum{}\{−2\}(z(P″)² − 1/x(P″)) has valuation ≥ 4e − 2 ≥ 6 for e ≥ 2: \textbf{1/x(P″) serves in place of z² through the sixth digit and not at the seventh}.
\end{namedthm}

\begin{mdproof}[Proof]
The proof of Proposition 17 gives x(z)z² = 1 + az² + π₁z⁴ + ⋯ with a = −b₂/12, a series in z². Inverting, 1/x(z) = z²(1 + az² + π₁z⁴ + ⋯)\textasciicircum{}\{−1\} = z² − az⁴ + O(z⁶), again with only even powers, so z² − 1/x(z) = az⁴ + O(z⁶). The valuation follows from v\_p(z) = e, v\_p(b₂) ≥ 0 and p ≥ 5. ∎
\end{mdproof}

The leading coefficient is b₂ and not a₁: the terms of R linear in a₁ cancel against a₁t to leading order, so no explanation through a₁ alone can reach the fourth power of z. §5.10 measures the valuation in all 497 boxes with b₂ ≠ 0, where it is 4e + v\_p(b₂) exactly.

\subsection*{3.12 The digits of the height}

Every congruence for n²h\_p(P) in this paper is Theorem N* read modulo a power of p, the terms being selected by Proposition 18. Theorem N is the case k = 2, and the cases k = 3, 4 are the two theorems of the fourth digit; from k = 5 on the truncation is mechanical, and what is stated is what was measured.

\begin{namedthm}[Corollary N⁽³⁾ (the third digit) [paper; the arithmetic in Lean]]
Under the hypotheses of Theorem N,
\dispf{\textbf{n²·h\_p(P) ≡ (1 + p + p²)Θ + 2p(1 + p)F − p²F² − (p²/2)Θ² − c/(p²·x(P″))  (mod p³)}.}
\end{namedthm}

\begin{namedthm}[Corollary N⁽⁴⁾ (the fourth digit) [paper; the arithmetic in Lean]]
In the same situation,
\dispf{\textbf{n²·h\_p(P) ≡ (1 + p + p² + p³)Θ + 2p(1 + p + p²)F − p²(1 + p)F² + (2/3)p³F³ − (p²/2)(1 + p)Θ² − c/(p²·x(P″))  (mod p⁴)}.}
\textbf{No new local quantity enters}: by Proposition 18 the next coefficient of ℒ is μ₄, and its digit is the seventh.
\end{namedthm}

Every term on the right of the two congruences is invariant under (x, y) ↦ (x, y + sx), since Θ, F, c and x(P″) are: the odd division polynomials ψ\_n and φ\_n are polynomials in b₂, b₄, b₆, b₈, the b\_i are invariant under that substitution, x does not move, and E₂(E, ω) does not either. So the congruences are statements about invariants of E; the model with p ∣ a₁ is where the proof runs (§4.12), not a hypothesis of the statement.

\begin{namedthm}[Corollary N⁽⁵⁾ (the fifth digit) [computation: 648 of 648 boxes]]
In the same situation,
\dispf{\textbf{n²·h\_p(P) ≡ (1 + p + p² + p³ + p⁴)Θ − (p²/2)(1 + p + p²)Θ² + (p⁴/3)Θ³}}
\dispf{\textbf{+ 2p(1 + p + p² + p³)F − p²(1 + p + p²)F² + (2/3)p³(1 + p)F³ − (1/2)p⁴F⁴ + (2/5)p⁵F⁵}}
\dispf{\textbf{− c/(p²·x(P″))  (mod p⁵)}.}
The term (2/5)p⁵F⁵ is 2p⁴F⁵ ≡ 2p⁴F at p = 5, by Fermat, and ≡ 0 (mod p⁵) for p ≥ 7: it is the first appearance of the factor v\_p(j) of Proposition 18. \textbf{No new local quantity enters.}
\end{namedthm}

\begin{namedthm}[Corollary N⁽⁶⁾ (the sixth digit) [computation: 648 of 648 boxes]]
In the same situation,
\dispf{\textbf{n²·h\_p(P) ≡ (1 + p + p² + p³ + p⁴ + p⁵)Θ − (p²/2)(1 + p + p² + p³)Θ² + (p⁴/3)(1 + p)Θ³}}
\dispf{\textbf{+ 2p(1 + p + p² + p³ + p⁴)F − p²(1 + p + p² + p³)F² + (2/3)p³(1 + p + p²)F³}}
\dispf{\textbf{− (1/2)p⁴(1 + p)F⁴ + (2/5)p⁵(1 + p)F⁵ − c/(p²·x(P″))  (mod p⁶)}.}
\textbf{No new local quantity enters}, and 1/x(P″) may still be used in place of z²: by Lemma 21 the error z² − 1/x(P″) has valuation ≥ 4e ≥ 2e + 4, so it is invisible modulo p⁶.
The geometric factor of the F⁵ term is (1 + p) and not 1, because at p = 5 the term (2/5)p⁵F⁵ has valuation 4 and not 5: the factor v\_p(j) of Proposition 18 shifts the geometric series as well. For p ≥ 7 the second summand has valuation 6 and is invisible modulo p⁶, so the one display serves both primes.
\end{namedthm}

\begin{namedthm}[Corollary N⁽⁷⁾ (the seventh digit) [computation: 648 of 648 boxes]]
In the same situation,
\dispf{\textbf{n²·h\_p(P) ≡ (1 + p + p² + p³ + p⁴ + p⁵ + p⁶)Θ − (p²/2)(1 + p + p² + p³ + p⁴)Θ² + (p⁴/3)(1 + p + p²)Θ³ − (p⁶/4)Θ⁴}}
\dispf{\textbf{+ 2p(1 + p + p² + p³ + p⁴ + p⁵)F − p²(1 + p + p² + p³ + p⁴)F² + (2/3)p³(1 + p + p² + p³)F³}}
\dispf{\textbf{− (1/2)p⁴(1 + p + p²)F⁴ + (2/5)p⁵(1 + p + p²)F⁵ − (1/3)p⁶F⁶ + (2/7)p⁷F⁷}}
\dispf{\textbf{+ p\textasciicircum{}\{−2\}·( −c·z(P″)² − (b₄/12)·z(P″)⁴ )  (mod p⁷)}.}
At this depth Θ⁴ and F⁶ enter for every p; the term (2/7)p⁷F⁷ is effective at p = 7 alone, where its valuation is 6, and the third summand of the geometric factor of F⁵ is effective at p = 5 alone — both are invisible modulo p⁷ at the other prime, so the one display serves both. \textbf{b₄ is the first new local invariant after c, and the seventh digit is where it enters}: dropping the term μ₄z⁴ leaves 232 of the 648 boxes, and keeping 1/x(P″) in place of z² leaves 310 of 648 (§5.10). The congruence is sharp — the two sides differ at valuation exactly 7 — in 499 of the 648 boxes.
\end{namedthm}

The pattern of the coefficients is Proposition 18 and nothing else, so no further corollary is worth displaying. What is worth recording is how far the closed form was tested and what each ingredient of it is needed for. Over the 648 boxes of §5.10 — conductor ≤ 12000, p ∈ \{5, 7\}, a\_p = 1, \#E(F\_p) = p, p ∤ ∏\_ℓ c\_ℓ, ι, split, e ≥ 2:

At k = 3, 4 and 5 the truncation of Theorem N* is the corollary displayed above and holds in 648 of the 648 boxes; since ℒ contributes the single term μ₂z² there, the truncations of ℒ do not differ from it. From k = 6 on:

\begin{center}\begin{small}
\begin{tabular}{>{\raggedright\arraybackslash}p{0.071\linewidth}>{\raggedright\arraybackslash}p{0.26\linewidth}>{\raggedright\arraybackslash}p{0.22\linewidth}>{\raggedright\arraybackslash}p{0.213\linewidth}>{\raggedright\arraybackslash}p{0.166\linewidth}}
\hline
k & Theorem N* truncated by Proposition 18 & ℒ truncated at z⁴ & ℒ truncated at z² & ℒ without its odd powers \\
\hline
6 & \textbf{648/648} & 648/648 & 648/648 & 648/648 \\
7 & \textbf{648/648} & 648/648 & 232/648 & 648/648 \\
8 & \textbf{648/648} & 648/648 & 146/648 & 648/648 \\
9 & \textbf{648/648} & 648/648 & 125/648 & 648/648 \\
10 & \textbf{648/648} & 392/648 (p = 7: \textbf{316/316}; p = 5: 76/332) & 114/648 & 648/648 \\
\hline
\end{tabular}
\end{small}\end{center}

The one entry that two independent runs do not agree on is the third column at k = 10: the run of §5.10 gives 114 and an independent recomputation at p-adic precision 40 and series order 26 gives 113. The difference is one box; it has not been resolved, and nothing in the paper turns on it.

Three readings of the table are the content of Proposition 18. First, the last column: \textbf{ℒ needs no odd power at any depth reached}, which is the first claim of Proposition 17 under test. Second, the third column: truncating ℒ at z² is exact through k = 6 and fails at k = 7, which is μ₄ arriving at its predicted digit. Third, the second column at k = 10: truncating at z⁴ still holds for \textbf{all} 316 boxes with p = 7 and fails for 256 of the 332 with p = 5, which is v\_p(μ₆) = −1 for p = 5 and 0 for p ≥ 7 — a local invariant whose digit depends on p, the same mechanism as the factor v\_p(j) on the side of the logarithms.

\begin{namedthm}[Remark (where the normalisation of σ becomes visible, and where it stops)]
The third digit is the first at which the choice of c — that is, of E₂ — is visible in the height: by Proposition 18 the term μ₂z² has valuation 2e − 2 = 2 for e = 2, hence the third digit. The remark of §2, that another constant in place of E₂ shifts h\_p by a multiple of the square of the p-adic elliptic logarithm, is this statement with the digit made explicit; and the square of the elliptic logarithm is exactly z², so the remark of §2 is the term μ₂z² itself. By Proposition 17 the third digit is also the \textbf{last} at which the normalisation is visible: \textbf{E₂ occurs in μ₂ alone}, and every deeper coefficient μ₄, μ₆, … is a polynomial in b₂, b₄, b₆ with rational coefficients — the ring is Z[1/(2m + 2)!] and not Z[1/6], which is what makes the digit of μ\_\{2m\} depend on p (Proposition 17). Unlike Θ, the term c/(p²x(P″)) is invariant under x ↦ u²x as well (c ↦ u²c and x(P″) ↦ u²x(P″)), so it does not share the model dependence recorded in the Remark of §4.12.
\end{namedthm}

\subsection*{3.13 The minimal global data}

Descending to the deeper layers gives two more criteria of the same shape as Corollary N′.

\begin{namedthm}[Corollary N″ (the third layer) [paper; the bookkeeping in Lean]]
Let r = 1 and v\_p(Reg\_p) ≥ 2, and write Θ = pΘ₁, S := (Θ₁ + 2F)/p. Then
\dispf{\textbf{v\_p(Reg\_p) ≥ 3 ⟺ S ≡ F² + W (mod p)}.}
\end{namedthm}

\begin{namedthm}[Corollary N‴ (the fourth layer) [paper; the bookkeeping in Lean]]
Let r = 1 and v\_p(Reg\_p) ≥ 3, and write S − F² − W = pU. Then
\dispf{\textbf{v\_p(Reg\_p) ≥ 4 ⟺ 6U + 6(S − F²) + 4F³ ≡ 0 (mod p)},}
and together with the three earlier criteria min(v\_p(Reg\_p), 4) is determined (Lean \texttt{quintichotomy}).
\end{namedthm}

\textbf{W is a datum of the formal group.} The quantity W is the only new ingredient of the third digit, and it is not a new \textit{global} one.

\begin{namedthm}[Proposition 12 (W in the first layer) [paper; the bookkeeping in Lean]]
Let T := χ(P)·T₀ and t₁ := t(nP − T) ∈ pZ\_p, with digits t₁/p = γ + pγ₂ + p²γ₃ + ⋯. Then
\dispf{\textbf{t(P″) = [p](t₁) ≡ p²γ (mod p³)},  hence  \textbf{W ≡ c·γ² (mod p)}  and  \textbf{e = 2 ⟺ γ ≢ 0 (mod p)}.}
So Corollary N″ reads v\_p(Reg\_p) ≥ 3 ⟺ S ≡ F² + c·γ² (mod p), with c local by the Remark of §2 and γ the first layer of the filtration: W adds no global datum to (χ, γ, F).
\end{namedthm}

\begin{namedthm}[Proposition 11 (the x-coordinate to three digits) [paper]]
With x ∘ τ\_T(t) = x(T) + c₁t + κt² + νt³ + ⋯ ∈ Z\_p[[t]] (Lemma 7) and Γ := γ + pγ₂ + p²γ₃, the quantity δ := (x(nP) − x(T))/p satisfies
\dispf{δ ≡ c₁·Γ + p·κ·Γ² + p²·ν·Γ³  (mod p³);}
in particular δ mod p³ depends only on the image of nP in E(Q\_p)/E₄(Q\_p).
\end{namedthm}

The last piece is the rate at which the layers of the filtration reach the digits of Θ. It needs one non-degeneracy, which is the non-vanishing of the slope of Proposition R⁗.

\begin{namedthm}[Theorem G′ (how deep a layer reaches) [paper; the valuation arithmetic in Lean]]
Let T ∈ D∖\{O\}, X(t) := x(τ\_T(R(t))) ∈ Z\_p[[t]] (Lemma 7) and Θ(t) := Q₂(φ\_p(X(t))). Assume
\dispf{(ND)  v\_p(φ\_p′(x(T))) = 1,  equivalently G(x(T)) ≢ 0 (mod p),  equivalently s() ≠ 0.}
Then for every t₁ ∈ pZ\_p and every j ≥ 1
\dispf{\textbf{v\_p(Θ(t₁ + p\textasciicircum{}j) − Θ(t₁)) = j − 1}.}
In particular Θ(t₁) mod p\textasciicircum{}\{j−1\} is a function of t₁ mod p\textasciicircum{}j, and no coarser datum will do: the j-th layer of the filtration moves exactly the digits of Θ from p\textasciicircum{}\{j−1\} on.
\end{namedthm}

(ND) is a condition on (Ē, x()) alone, since φ\_p′ = p·G with G ∈ Z[a₁, …, a₆][x] (Lemma 5). It holds whenever Res\_x(G, ψ\_p mod p) ≠ 0 in F\_p, because for anomalous p the x-coordinate of every point of E(F\_p)∖\{O\} is a root of ψ\_p mod p (ker V = E(F\_p), §2). \textbf{That resultant is non-zero for every p ≥ 5} (Theorem ND, §3.14, proved in §4.13), so (ND) is a hypothesis no longer; §5.8 records the exhaustive verification for p ≤ 41 that preceded the proof. The corresponding statement for ψ\_p needs no hypothesis either: v\_p(ψ\_p′(x(T))) = 1 exactly, always (§4.12).

\begin{namedthm}[Theorem G″ (the minimal global data at the third digit) [paper; the sharpness in Lean]]
Let p ≥ 5 be good ordinary, anomalous and split, p ∤ ∏\_ℓ c\_ℓ, ι, r = 1. Fix E, p, T₀ and the local invariants x(T), c₁, κ, ν, δ\_*, c. Then min(v\_p(Reg\_p), 3) is a function of
\dispf{\textbf{(χ(P), γ(P), γ₂(P), γ₃(P), F mod p²)}}
— three layers of the formal filtration of E(Q\_p) and two digits of the Fermat quotient of the denominator — and of nothing coarser. The refinement is: min(v\_p, 1) uses (χ, γ); min(v\_p, 2) adds (γ₂, F mod p); min(v\_p, 3) adds (γ₃, F mod p²).
\end{namedthm}

The two halves of Theorem G″ come from the two halves of Theorem G′. \textit{Sufficiency}: by the inequality half, Θ mod p³ is determined by t₁ mod p⁴, that is by (γ, γ₂, γ₃), and by Corollary N⁽³⁾ and Corollary N″ the three criteria use Θ only modulo p³ and F only modulo p². \textit{Necessity of γ₃}: by the sharp half with j = 3, two points with the same (χ, γ, γ₂) and different γ₃ have Θ differing at valuation exactly 2, hence S mod p differs, hence at most one value of γ₃ satisfies the criterion of Corollary N″. \textit{Necessity of F}: the pair of curves 3888p1 and 3267d1 at p = 7 are isomorphic over Q\_p, have the same Θ mod p² and differ in v\_p(Reg\_p) by the term q\_p(d′) alone (§5.6, §5.8); on one curve, the two points of 4407a1 of §5.13 do the same (Proposition 21 (ii); read pointwise, v\_p(n²h\_p(P)) in place of v\_p(Reg\_p), since 4407a1 has rank two).

\begin{namedthm}[Proposition 19 (the argument of ℒ is a layer of the filtration) [paper; the valuation step in Lean]]
Let T := χ(P)·T₀ and t₁ := t(nP − T) ∈ pZ\_p as in Proposition 12. Then P″ = p·(nP − T), and since log\_E is a homomorphism of formal groups,
\dispf{\textbf{z(P″) = p·log\_E(t₁)},  hence  \textbf{e = 1 + v\_p(t₁)}.}
So ℒ(z(P″)) is a function of t₁ — that is, of the layers (γ, γ₂, γ₃, …) — and of the invariants μ₂, μ₄, μ₆, … of Proposition 17: \textbf{the third term of Theorem N* brings in no global datum beyond the filtration} (Lean \texttt{layer\_\allowbreak{}depth}; verified in 40 of 40 boxes, §5.10).
\end{namedthm}

Proposition 19 makes Proposition 12 the case of the third digit: W ≡ cγ² (mod p) is the term μ₂z² read modulo the first digit it occupies. It also says which coordinate to use deeper. W = c/(p⁴x(P″)) is a quantity \textbf{of the digits up to the sixth}: by Lemma 21 the difference z² − 1/x(P″) has valuation ≥ 4e, so p\textasciicircum{}\{−2\}(z² − 1/x(P″)) has valuation ≥ 4e − 2 ≥ 6, and from the seventh digit on the formal logarithm z(P″) = p·log\_E(t₁), not x(P″), is the right coordinate.

The two ingredients of Theorem N* that a layer of the filtration moves — Θ and ℒ — respond to it at different rates, and the whole of what follows rests on the gap between them.

\begin{namedthm}[Proposition 20 (the two rates) [paper; the valuation arithmetic in Lean]]
In the situation of Theorem G′ and Proposition 19, let e = 1 + v\_p(t₁) ≥ 2 and t₁′ := t₁ + p\textasciicircum{}j with j ≥ 1, and write z′ := z(t₁′) = p·log\_E(t₁′). Then
(i) \textbf{v\_p(Θ(t₁′) − Θ(t₁)) = j − 1} (Theorem G′); (ii) \textbf{v\_p(z′ − z) = 1 + j}; (iii) \textbf{v\_p(ℒ(z′) − ℒ(z)) ≥ min\_\{m ≥ 1\} [ v\_p(μ\_\{2m\}) + (1 + j) + (2m − 1)·min(e, 1 + j) ]}, with equality when 1 + j ≠ e and the minimum is attained by one m alone. In particular, if
\dispf{\textbf{v\_p(c) \textless{} 2·min(e, 1 + j) + v\_p(μ₄)}  — which holds whenever p ∤ c, since min(e, 1 + j) ≥ 2 —}
then the term m = 1 leads and the bound reads \textbf{v\_p(c) + min(e, 1 + j) + (1 + j)}, with equality for 1 + j ≠ e; whereas when c = 0 — which forces j(E) = 0, and then μ₄ = 0 as well in the models met here — the leading term is m = 3 and the value is \textbf{(1 + j) + 5·min(e, 1 + j) + v\_p(μ₆)} for 1 + j ≠ e; (iv) hence \textbf{v\_p( p\textasciicircum{}\{−2\}(ℒ(z′) − ℒ(z)) ) ≥ j + 1 = (j − 1) + 2}, using only e ≥ 2, j ≥ 1, p ≥ 5 and v\_p(c) ≥ 0 — c = 0 deletes the m = 1 term and makes the bound deeper, so \textbf{(iv) and (v) need no hypothesis on c}; (v) hence the difference of the two sides of Theorem N*, with F held fixed, has valuation \textbf{exactly j − 1}.
So the ledger lags the digits of Θ by at least two digits, \textbf{at every depth}: (iv) puts no upper bound on j.
\end{namedthm}

\begin{mdproof}[Proof]
(ii) In log\_E(t) = Σ\_\{n≥1\}b\_nt\textasciicircum{}n/n with b\_n ∈ Z and b₁ = 1, the k-th term of the Taylor expansion at t₁ with increment h = p\textasciicircum{}j is h\textasciicircum{}k·L\textasciicircum{}\{(k)\}(t₁)/k!, and the coefficients of L\textasciicircum{}\{(k)\}/k! are b\_n·C(n−1, k−1)/k, since C(n,k)/n = C(n−1,k−1)/k; so its denominator divides k and the k-th term has valuation ≥ kj − v\_p(k). For k ≥ 2 and p ≥ 5, p\textasciicircum{}\{v\_p(k)\} ≤ k gives v\_p(k) \textless{} k − 1 ≤ (k−1)j, so every term past the first is deeper than the first, which is h·L′(t₁) with L′ ∈ 1 + tZ\_p[[t]] a unit, of valuation exactly j. Multiplying by p gives 1 + j.
\end{mdproof}

(iii) Δ(z²) = (z′ + z)(z′ − z) with z′ + z = 2z + Δz, and v\_p(2z) = e since 2 is a unit, v\_p(Δz) = 1 + j by (ii); the two differ unless 1 + j = e, and then v\_p(z′ + z) = min(e, 1 + j). For every m ≥ 1, Δ(z\textasciicircum{}\{2m\}) = Δz·Σ\_i z′\textasciicircum{}i z\textasciicircum{}\{2m−1−i\} has valuation \textbf{(1 + j) + (2m − 1)·min(e, 1 + j) whenever p ∤ 2m}, and ≥ that in general: the term z′\textasciicircum{}i z\textasciicircum{}\{2m−1−i\} has valuation i(1 + j) + (2m − 1 − i)e, which is ≥ (2m − 1)·min(e, 1 + j) with equality for exactly one i when 1 + j ≠ e — the exponent is min(e, 1 + j) and \textbf{not e}, since for 1 + j \textless{} e it is z′ and not z that is the shallower factor — while at 1 + j = e all 2m terms are of that valuation and their sum is 2m·z\textasciicircum{}\{2m−1\} to leading order. Summing over m and using min(e, 1 + j) ≥ 2 gives the bound of (iii), the terms with m ≥ 2 lying at ≥ (1 + j) + 2(2m − 1); with v\_p(μ\_\{2m\}) ≥ −v\_p((2m+2)!) ≥ −(2m+2)/(p−1) ≥ −(m+1)/2 (Proposition 17) the inequality 4m − 4 ≥ (m+1)/2 for m ≥ 2 leaves those terms at ≥ j + 3 after the factor p\textasciicircum{}\{−2\}. The m = 1 term, when it is present, gives v\_p(c) + min(e, 1+j) + (1+j) − 2 ≥ j + 1, which is (iv); when c = 0 that term is absent and the bound is the deeper one from m ≥ 2 (Lean \texttt{ledger\_\allowbreak{}lags\_\allowbreak{}any\_\allowbreak{}m}, uniform in m and using neither c ≠ 0 nor v\_p(μ\_\{2m\}) ≥ 0). The displayed condition for m = 1 to lead is Lean \texttt{m1\_\allowbreak{}leads\_\allowbreak{}iff}, and \texttt{m1\_\allowbreak{}leads\_\allowbreak{}of\_\allowbreak{}unit\_\allowbreak{}c} is the case p ∤ c. The equality in (iii) needs a unique minimising m, since two summands of equal valuation may cancel: that is what happens at the boundary row 1 + j = e, where Δ(z\textasciicircum{}\{2m\}) = z\textasciicircum{}\{2m\}((1 + w)\textasciicircum{}\{2m\} − 1) with w := Δz/z a unit, so the m-th leading term drops \textbf{exactly when (1 + w)\textasciicircum{}\{2m\} = 1 in F\_p} — for m = 1 the criterion 2u + 1 ≡ 0 of §5.11, and for (p − 1) ∣ 2m every u with u ≢ −1 (Lean \texttt{boundary\_\allowbreak{}cancel\_\allowbreak{}m1}, \texttt{boundary\_\allowbreak{}cancel\_\allowbreak{}pow}, \texttt{boundary\_\allowbreak{}cancel\_\allowbreak{}of\_\allowbreak{}card}).

(v) (1−p)\textasciicircum{}\{−1\}p\textasciicircum{}\{−2\}[log(1 + p²Θ′) − log(1 + p²Θ)] = (1−p)\textasciicircum{}\{−1\}·ΔΘ·[1 − p²(Θ′ + Θ)/2 + ⋯], where the bracket and (1−p)\textasciicircum{}\{−1\} are units, so the Θ side has valuation v\_p(ΔΘ) = j − 1 by (i), while the ℒ side is ≥ j + 1 by (iv); j − 1 \textless{} j + 1. The subscript arithmetic is Lean \texttt{theta\_\allowbreak{}rate\_\allowbreak{}digit}, \texttt{ledger\_\allowbreak{}rate\_\allowbreak{}digit}, \texttt{gap\_\allowbreak{}is\_\allowbreak{}e}, \texttt{ledger\_\allowbreak{}lags}, \texttt{ledger\_\allowbreak{}lags\_\allowbreak{}general}, \texttt{uniform\_\allowbreak{}in\_\allowbreak{}k}. ∎

\textbf{Translates by p-th powers.} Proposition 20 compares two values of the parameter t₁ with F held fixed. On the rational points of one curve F is not a free coordinate: the next proposition computes how F moves along P + p\textasciicircum{}aΛ, and the one after it says that on rational points F is a function of the height and of t₁.

\begin{namedthm}[Proposition 21 (translates by p-th powers) [paper; the arithmetic and the linear algebra in Lean]]
In the situation of Theorem N, let Λ ⊆ E(Q) be free, q := n²h\_p|\_Λ with polarisation q(P, Z) := (q(P + Z) − q(P) − q(Z))/2, χ : Λ → Z/p the reduction character, and
\dispf{ζ(Z) := t(p·nZ)/p² mod p,}
a homomorphism Λ → Z/p (since p·nZ ∈ E\_2(Q\_p)) with ζ(P) ≠ 0 exactly when e(P) = 2. For χ(P) ≠ 0 put
\dispf{C(P) := −(φ\_p′(x(P′))/p)·(2y(P′) + a₁x(P′) + a₃)/φ\_p(x(P′)) mod p,}
a unit: v\_p(φ\_p′(x(P′))) = 1 by (ND) (Theorem ND), since x(P′) ≡ x(T) (mod p); φ\_p(x(P′)) is a unit since x(P″) = φ\_p(x(P′))/ψ\_p(x(P′))² with v\_p(x(P″)) = −2e = −2v\_p(ψ\_p(x(P′))); and 2y + a₁x + a₃ does not vanish at P′ modulo p, since E(F\_p) has odd order. Then for Z ∈ Λ and a ≥ 1
\dispf{\textbf{F(P + p\textasciicircum{}aZ) ≡ F(P) + p\textasciicircum{}\{a−1\}·(χ(Z)/χ(P))·(q(P) − C(P)ζ(P)/2)  (mod p\textasciicircum{}a)}.}
In particular: (i) if χ(Z) = 0 then F(P + p\textasciicircum{}aZ) ≡ F(P) (mod p\textasciicircum{}a); (ii) if χ(Z) ≠ 0 and p ∣ q(P), the (a−1)-st digit of F moves if and only if e(P) = 2 — exactly at valuation a − 1 when e(P) = 2, and at valuation ≥ a when e(P) ≥ 3; (iii) if k ≥ 2 and S ∈ Λ ∩ ker χ has ζ(S) ≠ 0, then R := p\textasciicircum{}\{k−1\}S has depth exactly k, the points P and P + R agree in (χ, γ, …, γ\_\{k−1\}) and in F mod p\textasciicircum{}\{k−1\} and differ in γ\_k, and v\_p(q(P + R) − q(P)) = k − 1 exactly (Proposition 20 (v)); so whenever v\_p(q(P)) ≥ k, min(v\_p(q), k) separates the two points. (For S ∈ ker χ the conditions ζ(S) ≠ 0 and γ(S) ≠ 0 are equivalent, each saying that S has depth exactly one; Theorem G⁽ᵏ⁾ and §7(9) use the second.)
\end{namedthm}

\begin{mdproof}[Proof]
Four steps. (1) \textit{The linear term of Θ.} If nR ∈ E\_j(Q\_p) with j ≥ 1, then Θ(P + R) − Θ(P) ≡ C(P)·t(nR)/p (mod p\textasciicircum{}j). Indeed Δ := t₁(P + R) − t₁(P) ≡ t(nR) (mod p\textasciicircum{}\{j+1\}) by the formal group law; X(t) := x(T ⊕ ι(t)) has X′(t₁) ≡ (2y + a₁x + a₃)(P′) (mod p) from the invariant differential, so h := X(t₁ + Δ) − X(t₁) ≡ X′(t₁)Δ (mod p\textasciicircum{}\{2j\}); and since φ\_p mod p is a polynomial in x\textasciicircum{}p, the Taylor coefficients φ\_p\textasciicircum{}\{(k)\}/k! with 2 ≤ k \textless{} p are divisible by p, so φ\_p(x₀ + h) − φ\_p(x₀) ≡ φ\_p′(x₀)h (mod p\textasciicircum{}\{2j+1\}). Taking Q₂ gives [(Φ + δ)\textasciicircum{}\{p−1\} − Φ\textasciicircum{}\{p−1\}]/p² ≡ (p − 1)Φ\textasciicircum{}\{p−2\}δ/p² (mod p\textasciicircum{}\{2j\}) with (p − 1)Φ\textasciicircum{}\{p−2\} ≡ −Φ\textasciicircum{}\{−1\}, and v\_p(φ\_p′) = 1 gives the claim. (2) \textit{The kernel of χ.} For S ∈ ker χ, q(S) ≡ 0 and 2q(P, S) ≡ C(P)ζ(S) (mod p). Modulo p Theorem N* reads q(Q) ≡ Θ(Q), the F term lying in pZ\_p and the ℒ term at valuation ≥ 2e − 2 ≥ 2; by step (1) with j = 1 and t(nS)/p ≡ t(p·nS)/p² (mod p), q(P ± S) − q(P) ≡ ±C(P)ζ(S); the second difference gives 2q(S) ≡ 0 (Lean \texttt{second\_\allowbreak{}difference}), and 2q(P, S) = q(P + S) − q(P) − q(S). (3) \textit{A linear form.} L\_P(Z) := 2q(P, Z) − C(P)ζ(Z) mod p is additive and vanishes on ker χ by step (2), so L\_P(Z)·χ(P) = L\_P(P)·χ(Z) with L\_P(P) = 2q(P) − C(P)ζ(P) (Lean \texttt{ratio\_\allowbreak{}of\_\allowbreak{}ker\_\allowbreak{}le}, which needs only the inclusion of kernels). (4) \textit{The translate.} For R = p\textasciicircum{}aZ the quadratic identity q(P + R) − q(P) = p\textasciicircum{}a·2q(P, Z) + p\textasciicircum{}\{2a\}q(Z) (Lean \texttt{translate\_\allowbreak{}quadratic}) gives D := q(P + R) − q(P) ≡ p\textasciicircum{}a·2q(P, Z) (mod p\textasciicircum{}\{a+1\}) (Lean \texttt{shift\_\allowbreak{}quadratic}). Since nR ∈ E\_\{a+1\}(Q\_p), step (1) and t(p\textasciicircum{}\{a−1\}W) ≡ p\textasciicircum{}\{a−1\}t(W) put the Θ side of D at p\textasciicircum{}aC(P)ζ(Z) (mod p\textasciicircum{}\{a+1\}), and Proposition 20 (iv) puts the ℒ side at valuation ≥ a + 2; so the F side of Theorem N* is p\textasciicircum{}a·L\_P(Z) (mod p\textasciicircum{}\{a+1\}) (Lean \texttt{ledger\_\allowbreak{}digit}), and 2(1 − p)\textasciicircum{}\{−1\}[log(1 + pF′) − log(1 + pF)] ≡ 2p(F′ − F). Step (3) turns L\_P(Z) into (χ(Z)/χ(P))·L\_P(P). Part (ii) is Lean \texttt{kappa\_\allowbreak{}ne\_\allowbreak{}zero} (with q ≡ 0, C ≠ 0, ζ ≠ 0 and p ≠ 2 the coefficient is non-zero), and the separation in (iii) is Lean \texttt{separates}. ∎
\end{mdproof}

Step (2) is a second proof, from the linear term of Θ alone, of the statement of Theorem K that M\_Λ mod p vanishes on ker χ. The point of Proposition 21 is where q(P) enters: \textbf{only modulo p}. Along P + p\textasciicircum{}aΛ the digit of F at p\textasciicircum{}\{a−1\} is a function of the local quantities χ, q(P) mod p, C(P) and ζ, because the (a+1)-st digit of the translate of a quadratic form falls back to the form modulo p.

\begin{namedthm}[Proposition 22 (Theorem N* inverted) [paper; the polynomial identity in Lean]]
Since log is an isometry 1 + pZ\_p → pZ\_p, the term (2/(1 − p))·log(1 + pF) determines F modulo each power of p, and Theorem N* can be solved for F:
\dispf{\textbf{F mod p\textasciicircum{}k is a function of n²h\_p(P) mod p\textasciicircum{}\{k+1\}, of t₁ mod p\textasciicircum{}\{k+2\}, and of the local invariants},}
Θ mod p\textasciicircum{}\{k+1\} being a function of t₁ mod p\textasciicircum{}\{k+2\} (Theorem G′) and p\textasciicircum{}\{−2\}ℒ a function of t₁ and μ₂, μ₄, … (Proposition 19). For k = 2, with
\dispf{A := [n²h\_p(P) − (1 + p + p²)Θ + (p²/2)Θ² + c/(p²·x(P″))]/p,}
\dispf{\textbf{F ≡ A/2 − pA/2 + pA²/8  (mod p²)}.}
\end{namedthm}

\begin{mdproof}[Proof]
Corollary N⁽³⁾ reads 2p(1 + p)F − p²F² ≡ pA (mod p³), that is 2(1 + p)F − pF² ≡ A (mod p²). If 8F = 4A − 4pA + pA² + 8p²m then 16(1 + p)(8F) − p(8F)² − 64A ∈ p²Z (Lean \texttt{inverse\_\allowbreak{}n3}), and p ∤ 64. The solution is unique modulo p²: F ↦ 2(1 + p)F − pF² is injective modulo p², the difference of two values being (F₁ − F₂)(2(1 + p) − p(F₁ + F₂)) with the second factor a unit. ∎
\end{mdproof}

So on the rational points of one curve \textbf{the two kinds of datum of Theorem G⁽ᵏ⁾ are not independent coordinates}: F mod p\textasciicircum{}\{k−1\} is read off the height and the layers. Theorem G⁽ᵏ⁾ is a statement about the parameter space of Theorem N*, and Proposition 21 is what that parameter space looks like along p-power translates. Both are checked in the form displayed, in §5.13.

With Theorem G′, Proposition 19 and Proposition 20 the description at the third digit extends to every digit, in a shape that says how the global input grows.

\begin{namedthm}[Theorem G⁽ᵏ⁾ (the minimal global data at the k-th digit) [paper; the valuation arithmetic in Lean]]
Let p ≥ 5 be good ordinary, anomalous and split, p ∤ ∏\_ℓ c\_ℓ, ι, r = 1, (ND). Fix E, p, T₀ and the local invariants x(T), c₁, κ, ν, δ\_*, and μ₂, μ₄, …, μ\_\{2m₀\}, where m₀ is the largest m with 2me − 2 + v\_p(μ\_\{2m\}) ≤ k − 1. Then min(v\_p(Reg\_p), k) is a function of
\dispf{\textbf{(χ, γ, γ₂, …, γ\_k)  and  F mod p\textasciicircum{}\{k−1\}}.}
— and of nothing coarser. The global input is therefore of \textbf{two kinds at every depth} — the position of Λ in the filtration of E(Q\_p), and the digits of the Fermat quotient of the denominator — and each grows by one per digit, so the list is of 2k numbers (Lean \texttt{data\_\allowbreak{}count}, \texttt{data\_\allowbreak{}count\_\allowbreak{}step}, \texttt{data\_\allowbreak{}count\_\allowbreak{}linear}). \textbf{Theorem G″ is the case k = 3.}
\textit{The sense of the minimality.} It is minimality as a function on the parameter space that Theorem N* defines — t₁ mod p\textasciicircum{}\{k+1\} and F mod p\textasciicircum{}\{k−1\}, with the local invariants fixed. \textbf{Both kinds of necessity are nevertheless realised by rational points of one curve.} For the layer: if S ∈ Λ ∩ ker χ has γ(S) ≠ 0, then R := p\textasciicircum{}\{k−1\}S has depth exactly k and F(P + R) ≡ F(P) (mod p\textasciicircum{}\{k−1\}) by Proposition 21 (i), so whenever v\_p(n²h\_p(P)) ≥ k the two points agree in (χ, γ, …, γ\_\{k−1\}) and in F mod p\textasciicircum{}\{k−1\}, differ in γ\_k, and are separated by min(v\_p, k) (Proposition 21 (iii)); at 4407a1 and 5956a1 with p = 5 and k = 3 this happens in 32 pairs of 32, and at k = 4 in 2 of 2 (§5.13). These curves have rank two, so the statement is read pointwise: v\_p(n²h\_p(P)) in place of v\_p(Reg\_p), Theorem N* being an identity for each P with χ(P) ≠ 0. For the digit of F: besides the pair of curves 3888p1, 3267d1 at k = 3, the two points P = −5g₁ + 6g₂ and P + 25(−g₁ + g₂) of 4407a1 at p = 5 share t₁ mod p⁴, differ in the second digit of F alone, and have v₅(n²h\_p) = 7 and 2 (Proposition 21 (ii)). \textbf{What rational points do not give is independence of the two kinds of datum along p-power translates}: on P + p\textasciicircum{}aΛ the digit of F at p\textasciicircum{}\{a−1\} is determined by χ, q(P) mod p, C(P) and ζ (Proposition 21), and on rational points F mod p\textasciicircum{}\{k−1\} is a function of the height and of t₁ (Proposition 22).
\end{namedthm}

\textbf{The list is empty at a CM cell.} At a cell with c = 0 — necessarily j(E) = 0, and then μ₄ = 0 as well — the smallest m with 2me − 2 + v\_p(μ\_\{2m\}) ≤ k − 1 is m = 3, so for e = 2 the list μ₂, …, μ\_\{2m₀\} is \textbf{empty for every k ≤ 10}: v\_p(p\textasciicircum{}\{−2\}ℒ(z)) = 6e − 2 = 10, the ledger contributes nothing before the eleventh digit, and min(v\_p(Reg\_p), k) is a function of (χ, γ, …, γ\_k) and F mod p\textasciicircum{}\{k−1\} alone. Corollary N″ reads there v\_p(Reg\_p) ≥ 3 ⟺ S ≡ F² (mod p), since W ≡ cγ² = 0; the verification over the twenty such cells is in §5.11, and the witness for the necessity of F in Theorem G″ — the pair 3888p1, 3267d1 — is one of them (Lean \texttt{cm\_\allowbreak{}ledger\_\allowbreak{}invisible\_\allowbreak{}before\_\allowbreak{}eleventh}, \texttt{cm\_\allowbreak{}no\_\allowbreak{}local\_\allowbreak{}invariant\_\allowbreak{}before\_\allowbreak{}eleventh}). The lag there is ten digits and not two.

\begin{mdproof}[Proof]
\textit{Sufficiency}, uniform in k: by Theorem N*, v\_p(Reg\_p) ≥ k is a condition on Θ mod p\textasciicircum{}k, on F mod p\textasciicircum{}\{k−1\} and on ℒ(z(P″)) mod p\textasciicircum{}\{k+2\}; Θ mod p\textasciicircum{}k is a function of t₁ mod p\textasciicircum{}\{k+1\} by the inequality half of Theorem G′, and ℒ(z(P″)) is a function of t₁ and of μ₂, …, μ\_\{2m₀\} by Propositions 17, 18 and 19, so the list suffices.
\end{mdproof}

\textit{Necessity of γ\_k.} Take two data agreeing in (χ, γ, …, γ\_\{k−1\}) and in F mod p\textasciicircum{}\{k−1\} and differing in γ\_k. By Proposition 20 (v) with j = k the two values of n²h\_p differ at valuation exactly k − 1, so they are not both ≡ 0 (mod p\textasciicircum{}k) and min(v\_p(Reg\_p), k) separates them (Lean \texttt{exact\_\allowbreak{}pred\_\allowbreak{}not\_\allowbreak{}at}, \texttt{no\_\allowbreak{}cancel}, \texttt{at\_\allowbreak{}most\_\allowbreak{}one\_\allowbreak{}vanishes}, \texttt{layer\_\allowbreak{}is\_\allowbreak{}needed}). \textit{Necessity of the (k−1)-st digit of F.} Replacing F by F + p\textasciicircum{}\{k−2\} changes 2log(1 + pF) by a quantity of valuation exactly k − 1, since 2 and 1 + pF are units (Lean \texttt{fermat\_\allowbreak{}digit}, \texttt{fermat\_\allowbreak{}digits\_\allowbreak{}needed}), and by Proposition 20 (iv) the ledger cannot cancel it. \textit{Nothing deeper contributes.} By Theorem G′ the layer γ\_\{k+1\} moves Θ only at valuation k, and the k-th digit of F moves the second logarithm only at valuation k; neither is visible modulo p\textasciicircum{}k.

The two rates are what makes this uniform in k: Proposition 20 (iv) assumes only e ≥ 2, j ≥ 1, p ≥ 5, with no upper bound on j, so there is no depth at which the local term catches up with the layers. The larger e is, the further the ledger lags — its response to γ\_j sits at the (e + j)-th digit — and the minimality holds with more room to spare. ∎

This closes, at the third digit, the question left open by the Remark of §3.10: there is no third homomorphism into E₂/E₃ that carries the answer, but there \textit{is} a finite list, and the whole of the non-filtration content is the pair of digits F mod p². What happens deeper is Theorem N*: the list grows by one layer and one digit of F per digit of the height, and the local side stays finite — by Proposition 17 the new local invariants are μ₄ = −b₄/12 at the seventh digit, μ₆ at the tenth or eleventh, and so on, all of them classical invariants of E over Q.

\subsection*{3.14 The non-degeneracy (ND) is a theorem}

Theorem G′ and the sharp half of Theorem G″ used (ND), the non-vanishing of the slope, and v5 of this note verified it by exhaustion for p ≤ 41. It is a theorem for every p ≥ 5. The route is a statement about division polynomials alone, so this subsection works over \_p with the short model

\begin{quote}
Ē : y² = u(x) = x³ + a₄x + a₆,  Δ₀ = 4a₄³ + 27a₆²,  weights w(x) = 2, w(a₄) = 4, w(a₆) = 6,
\end{quote}

and p ≥ 5, so that E₄ = −48a₄ and E₆ = 864a₆ identify F\_p[a₄, a₆] with the ring of modular forms mod p. Since [p] = V∘F in characteristic p, ψ\_p mod p is a polynomial in x\textasciicircum{}p (§2): write

\begin{quote}
ψ\_p mod p = Ψ(x\textasciicircum{}p),  deg\_X Ψ = n\_p := (p−1)/2  (the letter n is the Tamagawa product ∏\_ℓ c\_ℓ elsewhere),  lc\_X Ψ = A (the Hasse invariant),  h := Ψ/A, φ\_p mod p = Φ(x\textasciicircum{}p),  Φ monic of degree p in X,  x∘V = Φ/(A²h²),  w\textasciicircum{}\{(p)\}(X) := X³ + a₄\textasciicircum{}pX + a₆\textasciicircum{}p,
\end{quote}

so that h is the kernel polynomial of the Verschiebung V : Ē\textasciicircum{}\{(p)\} → Ē and Φ/(A²h²) is its x-coordinate. Let K := ψ\_p′/p mod p and S₁ := (φ\_p − 4uK²)/ψ\_p, so that (ND) ⟺ S₁(x()) ≠ 0 (§4.11). Let

\begin{quote}
D := −6a₆·∂/∂a₄ + (4/3)a₄²·∂/∂a₆,
\end{quote}

the derivation that raises the weight by 2 and kills Δ₀; in the normalisation above D is \textit{minus} the Serre derivative, and on the q-expansion of a form of weight ≡ 0 (mod p) it acts as −q d/dq.

\begin{namedthm}[Lemma 13 (two identities for an odd separable isogeny) [paper; the arithmetic in Lean]]
Let λ : E′ → E be separable of odd degree m with kernel polynomial h monic of degree (m−1)/2, and write x∘λ = P/h², y∘λ = y′Ω/h³, λ\textasciicircum{}*ω\_E = c·ω\_\{E′\}. Then P′h − 2Ph′ = cΩ and u\_\{E′\}Ω² = P³ + a₄Ph⁴ + a₆h⁶, and at every simple root ξ of h with P(ξ) ≠ 0
\dispf{\textbf{c²P(ξ) = 4u\_\{E′\}(ξ)h′(ξ)²}  and  \textbf{c²P′(ξ) = 2[u\_\{E′\}·(h′)²]′(ξ)}.}
For λ = [p] this gives p²φ\_p(x\_T) = 4u(x\_T)ψ\_p′(x\_T)² and the identity (I-2) of §4.12. For λ = V, where c² = A² because the numerator of x∘λ has leading coefficient c\textasciicircum{}\{−2\}, it gives Φ(ξ) = 4w\textasciicircum{}\{(p)\}(ξ)h′(ξ)² and Φ′(ξ) = 2[w\textasciicircum{}\{(p)\}(h′)²]′(ξ) at every root ξ of h.
\end{namedthm}

\begin{namedthm}[Proposition 14 (the second p-adic digit of ψ\_p is read off the first) [paper]]
For p ≥ 5,
\dispf{\textbf{A·(ψ\_p′/p) = u\textasciicircum{}\{(p−1)/2\}·(dΨ/dX)(x\textasciicircum{}p)  in F\_p[a₄, a₆][x]},}
so u\textasciicircum{}\{(p−1)/2\} divides K and 4uK² = 4w\textasciicircum{}\{(p)\}(x\textasciicircum{}p)·h′(x\textasciicircum{}p)² is a polynomial in x\textasciicircum{}p.
\end{namedthm}

\begin{namedthm}[Proposition 15 (S₁ is a derivative of ψ\_p) [paper]]
For p ≥ 5,
\dispf{\textbf{φ\_p ≡ 4u·(ψ\_p′/p)² + D(ψ\_p²)  (mod p)},  equivalently  \textbf{S₁ = 2·D ψ\_p},}
equivalently A²Φ = 4w\textasciicircum{}\{(p)\}(Ψ′)² + A²·D(Ψ²). The weight is rigid: w(DΨ) = p² + 1 = w(S₁ in X), and there is no monomial of weight 2, so no correction term of the shape c·Ψ can be added.
\end{namedthm}

The two propositions are what (ND) rests on. With them the resultant Res\_x(G, ψ\_p mod p) factors — not the resultant Res\_X(ũ, g) of §4.6(3), which is a different one — and can be counted at the cusp by the method of §4.6(3), which gives the theorem:

\begin{namedthm}[Theorem ND (the non-degeneracy holds always) [paper]]
For every prime p ≥ 5 there is c\_S ∈ F\_p\textasciicircum{}× with
\dispf{\textbf{Res\_x(S₁, ψ\_p mod p) = c\_S·Δ₀\textasciicircum{}\{m\_S\}  in F\_p[a₄, a₆]},  m\_S = p²(p²−1)/24,}
and hence (ND) holds for \textbf{every} elliptic curve over \_p and every  ∈ Ē(F\_p)∖\{O\}: at an ordinary point the resultant is non-zero, so S₁ and ψ\_p mod p have no common root, and at a supersingular point ψ\_p mod p is a non-zero constant and (ND) is vacuous. Consequently \textbf{Theorem G′ and Theorem G″ are unconditional}, and so are Proposition R⁗ and Proposition 10.
\end{namedthm}

\begin{namedthm}[Proposition 16 (the three constants in closed form) [paper]]
With n\_p = (p−1)/2, M = p(p²−1)/24, m\_K = (p²−1)(p²−3)/24 and m\_p = m\_K + m\_S, in F\_p:
\dispf{\textbf{Res\_X(DΨ, Ψ) = γ\_p·Δ₀\textasciicircum{}\{M\}},  \textbf{γ\_p = (−1)\textasciicircum{}\{n\_p(n\_p+1)/2\}·n\_p!·(−16)\textasciicircum{}\{M\}},  \textbf{c\_S = 2\textasciicircum{}n\_p·γ\_p},}
\dispf{\textbf{Res\_x(K, ψ\_p mod p) = c\_K·Δ₀\textasciicircum{}\{m\_K\}},  \textbf{c\_K = ε\_p·2\textasciicircum{}\{m\_K\}}  with  \textbf{ε\_p = −1 ⟺ p ≡ 3 (mod 8)},}
and c\_p = (−1)\textasciicircum{}\{p(p−1)/2\}c\_Kc\_S for Res\_x(G, ψ\_p mod p) = c\_p·Δ₀\textasciicircum{}\{m\_p\}. The sign ε\_p is the second supplementary law: it is (2|p) times, when p ≡ 1 (mod 4), a second copy of (2|p).
\end{namedthm}

The closed form for c\_K, and with it the Legendre symbol, is the shape of the observation recorded in v5 for p ≤ 53; it is now a consequence of Proposition 14. Note what Theorem ND does \textit{not} say: it is a statement about the reduced curve, and none of the hypotheses of §3.2 or §3.10 are weakened by it.

\subsection*{3.15 Depth and height: an archimedean input}

Every statement so far is local at p: the data are the layers of the filtration of E(Q\_p) and the digits of Fermat quotients, and §7(3) records that the local data do not decide whether the Schur complement of Corollary R′\_s is a unit. One global datum of a different kind does decide it in part, and it is archimedean: how deep a rational point can lie in the filtration is bounded by its naive height. For Q ∈ E(Q), Q ≠ O, write h(x(Q)) := log max(|α(Q)|, d(Q)²) for the naive height of x(Q) = α/d².

\begin{namedthm}[Lemma 23 (depth and height) [paper; the inequality in Lean]]
Let p be a prime of good reduction, the equation minimal at p, and Q ∈ E\_k(Q\_p) with k ≥ 1. Then p\textasciicircum{}k ∣ d(Q), and hence
\dispf{\textbf{h(x(Q)) ≥ 2k·log p},}
with no error term. Equivalently, k ≤ h(x(Q))/(2 log p): the depth of a rational point in the formal filtration is bounded by its naive height.
\end{namedthm}

\begin{mdproof}[Proof]
For Q ∈ E\_1(Q\_p) with v\_p(x(Q)) = −2m one has v\_p(y(Q)) = −3m and v\_p(t(Q)) = m, so Q ∈ E\_k(Q\_p) exactly when p\textasciicircum{}k ∣ d(Q); and d² ≤ max(|α|, d²). ∎ (Lean \texttt{log\_\allowbreak{}denom\_\allowbreak{}ge}, \texttt{height\_\allowbreak{}ge\_\allowbreak{}of\_\allowbreak{}depth}, \texttt{depth\_\allowbreak{}lt\_\allowbreak{}of\_\allowbreak{}height\_\allowbreak{}lt}.)
\end{mdproof}

The inequality has an exact form, which says where the slack goes.

\begin{namedthm}[Lemma 23′ (the slack) [paper; Lean]]
Write d(Q) = p\textasciicircum{}D·c′ with p ∤ c′. Then
\dispf{\textbf{h(x(Q)) = 2 log c′ + 2D·log p + max(0, log|α(Q)| − 2 log d(Q))}.}
Hence, with s := h(x(Q)) − 2D·log p, one has c′ ≤ e\textasciicircum{}\{s/2\}; equality h(x(Q)) = 2D·log p holds if and only if c′ = 1 and |α(Q)| ≤ d(Q)²; and \textbf{s \textless{} 2 log 2 forces c′ = 1}.
\end{namedthm}

\begin{mdproof}[Proof]
h = max(log|α|, 2 log d) = 2 log d + max(0, log|α| − 2 log d), and 2 log d = 2D log p + 2 log c′. ∎ (Lean \texttt{cofactor\_\allowbreak{}log\_\allowbreak{}le}, \texttt{cofactor\_\allowbreak{}eq\_\allowbreak{}one\_\allowbreak{}of\_\allowbreak{}slack\_\allowbreak{}lt}, \texttt{max\_\allowbreak{}eq\_\allowbreak{}iff}.)
\end{mdproof}

The slack s splits into the part of the denominator prime to p, 2 log c′, and an archimedean excess max(0, log|x(Q)|); §5.12 measures that the second saturates first. Both lemmas are elementary, and we regard them as known (§8).

Combined with Corollary K′(i), Lemma 23 gives an upper bound for v\_p(Reg\_p) that no local datum gives.

\begin{namedthm}[Proposition 24 (an upper bound from the shortest vector) [paper; the linear algebra in Lean]]
Let p ≥ 5 be good ordinary, \#E(F\_p) = p, p ∤ ∏\_ℓ c\_ℓ, E(Q) torsion-free, ι, and E split at p; write Λ = E(Q) and
\dispf{h\_min := min\{ h(x(Q)) : 0 ≠ Q ∈ Λ ∩ E\_1(Q\_p) \}.}
\dispf{\textbf{If h\_min \textless{} 4·log p, then (χ, γ) : Λ → F\_p² is onto}, hence ρ₁ = 2 (Corollary K′(i)).}
In rank two this gives \textbf{v\_p(Reg\_p) = 0} (Corollary R′\_s); in rank r ≥ 3 it gives v\_p(Reg\_p) = (r − 2) + v\_p(det T\_s) with T\_s ≡ p·Ĉ (mod p²) (Corollary K′(iv)).
\end{namedthm}

\begin{mdproof}[Proof]
χ is onto by ι. If (χ, γ) is not onto, its image is a line projecting onto the first factor, that is the graph of γ = cχ, so γ vanishes on ker χ = Λ ∩ E\_1(Q\_p), that is Λ ∩ E\_1(Q\_p) ⊆ E\_2(Q\_p). Then every non-zero Q ∈ Λ ∩ E\_1(Q\_p) has h(x(Q)) ≥ 4 log p by Lemma 23 with k = 2, against h\_min \textless{} 4 log p. ∎ (Lean \texttt{surjective\_\allowbreak{}pair\_\allowbreak{}iff}, \texttt{surjective\_\allowbreak{}of\_\allowbreak{}shallow\_\allowbreak{}kernel\_\allowbreak{}element}.) Whether h\_min is taken over all of Λ ∩ E\_1(Q\_p) or over (Λ ∩ E\_1(Q\_p)) ∖ pΛ does not change the truth of h\_min \textless{} 4 log p: in the split case pΛ ⊆ E\_2(Q\_p), so every non-zero element of pΛ has h(x(Q)) ≥ 4 log p by Lemma 23.
\end{mdproof}

Three remarks place it. (a) \textit{Torsion.} If \#E(F\_p) = p then the torsion of E(Q) prime to p injects into E(F\_p), so p ∤ \#E(Q)\_tors already makes E(Q) torsion-free. (b) \textit{The split hypothesis cannot be dropped.} In the non-split case E(Q\_p) ≅ Z\_p and pP has depth exactly one for every P with χ(P) ≠ 0, so Λ ∩ E\_1 ⊆ E\_2 never happens and the proposition has nothing to say; in the split case pΛ ⊆ E\_2 automatically, and for r = 2 the quotient (Λ ∩ E\_1)/pΛ has order p, so \textbf{(χ, γ) fails to be onto if and only if p² ∣ d(Q) for one — equivalently every — Q ∈ (Λ ∩ E\_1) ∖ pΛ}, a test with no p-adic computation in it. (c) \textit{What it is and is not.} h\_min is computed from the archimedean height pairing and an integer search; the condition is sufficient and not necessary, and it bounds v\_p(Reg\_p) from above by an archimedean datum where §7(3) says that no local datum does. Cut by the naive height of the proposition, the sweep of §5.12 shows that it accounts for \textbf{about 68 \%} of an observed deficit in the frequency of det(χ, γ) ≡ 0 (207.8 of 305.35; 55.8 \% on the second set of boxes). The canonical form below, Proposition 24ᶜ, accounts for \textbf{all of it in expectation} (100.6 \%, and 99.7 \% on the second set): where the deficit ends is explained by a theorem. That above the threshold the frequency is 1/p is an empirical regularity, not a consequence of either proposition.

\textit{The canonical height.} Let the equation be a global minimal model and ĥ the canonical height in the normalisation of PARI (ĥ ≈ h(x)), so that ĥ = Σ\_v λ′\_v with λ′\_v := 2(λ\_v + (1/12) log|Δ|\_v) (the terms in Δ cancel by the product formula). At the good prime p, λ′\_p(Q) = 2D(Q)·log p with D(Q) = v\_p(d(Q)). Write R(Q) := ĥ(Q) − 2D(Q)·log p = Σ\_\{v≠p\} λ′\_v(Q) = 2 log c′ + λ′\_∞(Q) + Σ\_\{ℓ bad\} corr\_ℓ(Q)·log ℓ, where corr\_ℓ(Q) ≤ 0 is the component correction (zero when Q has non-singular reduction at ℓ). This decomposition and the table of corrections are the known theory of local heights [known; [SilH]]; we use them and do not prove them. Put

\begin{quote}
\dispf{β(E) := m\_∞ + Σ\_\{ℓ bad\} m\_ℓ·log ℓ,  m\_∞ := min\_\{E(R)∖O\} λ′\_∞,  m\_ℓ := min over the components of E(Q\_ℓ)/E\_0(Q\_ℓ) of corr\_ℓ,}
\end{quote}

with m\_ℓ = 0 when c\_ℓ = 1, −⌊N/2⌋⌈N/2⌉/N for split I\_N, −N/4 for non-split I\_N with N even, −1/2, −2/3, −1, −4/3, −3/2 for III, IV, I₀*, IV*, III*, and for I\_n* the smaller of −1 and −(n+4)/4.

\begin{namedthm}[Proposition 24ᶜ (the canonical threshold) [paper; the arithmetic and the logic in Lean]]
Under the hypotheses of Proposition 24 and with the equation globally minimal, let ĥ\_min := min\{ ĥ(Q) : 0 ≠ Q ∈ Λ ∩ E\_1(Q\_p) \}.
\dispf{\textbf{If ĥ\_min \textless{} 4·log p + β(E), then (χ, γ) is onto}, hence ρ₁ = 2, and in rank two v\_p(Reg\_p) = 0.}
\end{namedthm}

\begin{mdproof}[Proof]
R(Q) ≥ β(E) termwise: 2 log c′ ≥ 0, λ′\_∞ ≥ m\_∞, corr\_ℓ ≥ m\_ℓ (Lean \texttt{beta\_\allowbreak{}le\_\allowbreak{}remainder\_\allowbreak{}of\_\allowbreak{}local\_\allowbreak{}bounds}). If (χ, γ) is not onto then, as in Proposition 24, every non-zero Q ∈ Λ ∩ E\_1 has D(Q) ≥ 2, so ĥ(Q) = 2D log p + R(Q) ≥ 4 log p + β(E) (Lean \texttt{canonical\_\allowbreak{}ge\_\allowbreak{}of\_\allowbreak{}depth}), against the hypothesis. ∎ (Lean \texttt{surjective\_\allowbreak{}of\_\allowbreak{}canonical\_\allowbreak{}lt\_\allowbreak{}four\_\allowbreak{}log\_\allowbreak{}add\_\allowbreak{}beta}.) As for Proposition 24, taking ĥ\_min over all of Λ ∩ E\_1 or outside pΛ does not change the truth of the hypothesis.
\end{mdproof}

\textit{Remarks.} (d) The two propositions are two sufficient conditions, neither implying the other: β(E) can be of either sign (its median over the boxes of §5.12 is −0.07·log p), and 23 + 43 boxes of §5.12 lie below the naive threshold and above the canonical one. (e) The lower bound R ≥ β(E) holds for every reduction type with a bounded constant; it fails only if a hypothesis is dropped (p a bad prime, a non-minimal model, the component corrections omitted). (f) The local height formulas, the component corrections and the value of m\_∞ are not in Lean; they enter the Lean statements as hypotheses. m\_∞ is computed numerically, so the box-by-box application in §5.12 is [computation]. (g) The proof uses only that every non-zero Q ∈ Λ ∩ E\_1 has R(Q) ≥ β(E). Since corr\_ℓ depends only on the component of Q at ℓ [known], β(E) may be replaced by β\_Λ(E) := m\_∞ + min\{ Σ\_ℓ corr\_ℓ(κ)·log ℓ : κ in the image of Λ ∩ E\_1 in ∏\_ℓ Φ\_ℓ \} ≥ β(E), a finite computation; β\_Λ \textgreater{} β exactly when that image misses the minimal component at some ℓ. The Lean statements take a bound at each ℓ separately (m\_ℓ ≤ corr\_ℓ); β\_Λ needs instead the joint bound Σ\_ℓ corr\_ℓ·log ℓ ≥ β\_Λ − m\_∞, a change of hypothesis that we have not made in Lean [paper].

What the threshold leaves open is the frequency of det(χ, γ) ≡ 0 above it. The next lemma says that the frequency 1/p is a local count, so that any departure from it is a correlation between the local class of the box and the global height.

\begin{namedthm}[Lemma 25 (the image modulo E\_2) [paper]]
Let p ≥ 5 be good, \#E(F\_p) = p, ι, E split at p, and φ : Λ → E(Q\_p)/E\_2(Q\_p) the reduction. Then
(i) E(Q\_p)/E\_2(Q\_p) ≅ (Z/p)²;
(ii) \textbf{det(χ, γ) ≡ 0 if and only if φ(Λ) has order p}, and then φ(Λ) is one of the p lines of (Z/p)² other than E\_1/E\_2; otherwise φ is onto;
(iii) if Λ′ ⊆ Λ is a subgroup of index prime to p, then φ(Λ′) = φ(Λ). In particular the verdict det(χ, γ) ≡ 0 depends only on the images of the generators of any such Λ′ in E(Q\_p)/E\_2(Q\_p) ≅ E(Z/p²Z) [known], so, for generators with p-integral coordinates, only on the equation and those coordinates modulo p².
\end{namedthm}

\begin{mdproof}[Proof]
(i) E(Z\_p) is an extension of (F\_p) ≅ Z/p by E\_1(Q\_p), and E\_1 ≅ Z\_p with E\_2 = pE\_1 through the formal logarithm (p ≥ 3) [known]. So E(Q\_p) is Z\_p or Z\_p ⊕ Z/p; in the first case E\_1 = pE(Q\_p) and the quotient by E\_2 is cyclic of order p², in the second the quotient is (Z/p)², and the second is the split case. (ii) φ(Λ ∩ E\_1) = φ(Λ) ∩ E\_1/E\_2, and since χ is onto, \#φ(Λ) = p·\#φ(Λ ∩ E\_1). As in the proof of Proposition 24, det(χ, γ) ≡ 0 exactly when Λ ∩ E\_1 ⊆ E\_2, that is when φ(Λ ∩ E\_1) = 0. (iii) φ(Λ) is a p-group and Λ/Λ′ has order prime to p. ∎
\end{mdproof}

For the family of set C of §5.12 (two points with x-coordinates x₁ and x₁ + 1), counting the classes (x₁, y₁, y₂) modulo p² uniformly, the split classes with a\_p = 1 have det(χ, γ) ≡ 0 in exactly 40 of 200 at p = 5 and 112 of 784 at p = 7, that is 1/p; and in the 6,818 boxes of set C where p does not divide the index of the two points in Λ, the verdict computed from the classes modulo p² agrees with the measured one in every box [computation]. Lemma 25 is elementary and should be regarded as known.

\textbf{When the closure is small.} If Λ ∩ E\_1(Q\_p) ⊆ E\_3(Q\_p), then the map Λ → E(Q\_p)/E\_3(Q\_p) factors through Λ/(Λ ∩ E\_1) ≅ Z/p — no hypothesis on γ is needed for this — so its image has index p². If moreover γ(P₁) = 0 for one P₁ with χ(P₁) ≠ 0, then \textbf{every layer up to the third is a multiple of χ}: γ ≡ 0 and γ₂ = c₂χ on Λ (Lean \texttt{factor\_\allowbreak{}through\_\allowbreak{}chi}). That γ₂ is a homomorphism, which the Lean statement takes as a hypothesis, holds because γ ≡ 0 puts t₁ in p²Z\_p, where the formal group law is addition modulo p⁴. By Theorem K,  = 2λ\_\{T,1\}χ⊗χ; if λ\_\{T,1\} ≠ 0 then v\_p(h\_p(P)) = 0 at every P with χ(P) ≠ 0 (Lean \texttt{two\_\allowbreak{}ne\_\allowbreak{}zero\_\allowbreak{}of\_\allowbreak{}ne\_\allowbreak{}two}, \texttt{diag\_\allowbreak{}ne\_\allowbreak{}zero}), and in rank two ρ₁ = 1 and v\_p(Reg\_p) ≥ 1. [paper.] The two boxes of this kind that §5.12 finds, 328563a1 and 477950bg1 at p = 7, have ρ₁ = 1 and v\_p(Reg\_p) = 1 [computation]. Such a box cannot realise the necessity of γ₃ in Theorem G″ on its rational points, since no point with χ(P) ≠ 0 has v\_p(h\_p(P)) ≥ 1 there; the realisation of §3.13 uses Proposition 21 instead.

\section*{4. Proofs}

\begin{namedthm}[Remark (the ledger is not an induction)]
The congruences of Theorem N and of Corollaries N⁽³⁾–N⁽⁷⁾ are not statements proved one after another: they are the p-adic expansion, digit by digit, of the single exact identity
\dispf{\textbf{p²·n²h\_p(P) = h\_p(P″) = 2 log\_p(σ\_p(t)/t) + 2 log\_p(φ/ω) − 2p²·log\_p(d′)},  t = t(P″),}
in which φ = φ\_p(x(P′)), ω = ω\_p(P′). Theorem N* is this identity summed rather than expanded, and Proposition 18 is the table below in closed form. What is visible at digit k of n²h\_p — that is, modulo p\textasciicircum{}k, equivalently modulo p\textasciicircum{}\{k+2\} in h\_p(P″) — are exactly the terms of valuation ≤ k + 1, and the table of valuations settles which those are (e = v\_p(t) ≥ 2, model with p ∣ a₁):
| term | source | valuation | first digit k | |---|---|---|---| | p²Θ | log\_p(φ/ω) | 2 | 0 | | p³F | denominator | 3 | 1 | | p⁴Θ², p⁶Θ³ | log\_p(φ/ω), Q₂ | 4, 6 | 2, 4 | | (−1)\textasciicircum{}\{j+1\}p\textasciicircum{}\{2j\}Θ\textasciicircum{}j/j, j ≥ 2 | log\_p(φ/ω) | \textbf{2j − v\_p(j)} | 2j − 2 − v\_p(j) | | p⁴F², p⁵F³ | denominator | 4, 5 | 2, 3 | | 2(−1)\textasciicircum{}\{j+1\}p\textasciicircum{}\{2+j\}F\textasciicircum{}j/j, j ≥ 2 | denominator | \textbf{2 + j − v\_p(j)} | j − v\_p(j) | | a₁t + a₁ωψ/φ² | σ, Weierstrass | \textbf{cancel} (remainder ≥ 2v\_p(a₁) + 2e ≥ 6) | — | | the t² terms, combined: μ₂z² = −c·z² | σ (t²), η (a₂), the a₁ remainders | 2e | 2 (for e = 2) | | c₄t³, a₃ωψ³/φ³ | σ (t³), Weierstrass | 3e, and \textbf{no new invariant} (Proposition 17) | — | | μ\_\{2m\}z\textasciicircum{}\{2m\}, m ≥ 2 | ℒ (Proposition 17) | \textbf{2me + v\_p(μ\_\{2m\})} | 2me − 2 + v\_p(μ\_\{2m\}) | | a₄ψ⁴/φ², a₆ψ⁶/φ³ | Weierstrass (inside log(1 + R), hence inside ℒ) | 4e, 6e | 6, 10 |
Two features of the table are corrected and generalised by Proposition 18, and the corrections matter from the fifth digit on. First, \textbf{the j-th term of a logarithm is one digit shallower than p\textasciicircum{}j whenever p ∣ j}: the factor v\_p(j) is what brings (2/5)p⁵F⁵ into the fifth digit at p = 5 (Corollary N⁽⁵⁾) and what makes the depth of a term depend on p. Second, \textbf{the individual coefficients of σ\_p are not the right unit of account}: by Proposition 17 the combination that survives is the even series ℒ, whose m-th coefficient μ\_\{2m\} is an invariant of E. So the coefficients of σ\_p of odd index bring in nothing — the term c₄t³ of the row above cancels against a₃ωψ³/φ³ and a reparametrisation of μ₂z² — and the first local invariant after c is μ₄ = −b₄/12, at the \textbf{seventh} digit, not the fifth. The denominator descends at rate p\textasciicircum{}\{k−v\_p(k)\} and ℒ at rate p\textasciicircum{}\{2me−2+v\_p(μ\_\{2m\})\}: the two non-local ingredients of Lemma 3 arrive at different speeds, which is one answer — a structural rather than an intrinsic one — to the question of §7 about q\_p(d′).
\end{namedthm}

\subsection*{4.1 Lemmas 1–3}

\textit{Lemma 1.} σ\_p(t)/t ∈ 1 + tZ\_p[[t]] and v\_p(t) = e ≥ 1, so v\_p(σ\_p(t)) = e; from v\_p(x(P)) = −2e and gcd(α, d) = 1 we get v\_p(d) = e and v\_p(α) = 0, whence v\_p(σ\_p(P)/d(P)) = 0. ∎

\textit{Lemma 2.} For w ∈ Z\_p\textasciicircum{}× one has w\textasciicircum{}\{p−1\} ∈ 1 + pZ\_p and log(1 + px) ∈ pZ\_p, so v\_p(log\_p w) ≥ 1, refined as in §2; apply this to w = σ\_p(P)/d(P). ∎

\textit{Lemma 3.} From t = −x/y and x = α/d², y = β/d³ we get t/d = −α/β, hence σ\_p(P)/d(P) = u(P)·(−α(P)/β(P)) with u = σ\_p(t)/t ∈ 1 + p\textasciicircum{}eZ\_p. Taking log\_p, using log\_p(−1) = 0 and factoring α/β into primes ℓ ≠ p (both are prime to p by Lemma 1), gives the decomposition with n\_ℓ ∈ Z; each ℓ is a p-adic unit, so v\_p(log\_p ℓ) ≥ 1 by Lemma 2, and c ≥ 1 as the n\_ℓ are rational integers. Finally u ∈ 1 + p\textasciicircum{}eZ\_p gives b ≥ e, and b ≥ 2e when p ∣ a₁, the coefficient c₂ = a₁/2 of σ\_p then being divisible by p. ∎

\subsection*{4.2 Theorem A}

(1) If P, R, P+R ∈ G, Lemma 2 puts all three heights in 2pZ\_p and p is odd, so h\_p(P,R) ∈ pZ\_p; the matrix M\_G on a basis of the image of G has entries in pZ\_p, and every term of the Leibniz expansion of det M\_G has valuation ≥ r. (2) A change of basis between G and Λ has determinant ±k, so det M\_G = k² det M\_Λ and v\_p(Reg\_p) ≥ r − 2v\_p(k). (3) Λ/(image of G) embeds into E(F\_p) × ∏\_ℓ Φ\_ℓ(F\_ℓ), so for p ∤ ∏\_ℓ c\_ℓ the p-part comes from E(F\_p), which by Hasse has order \textless{} p², its p-part cyclic of order at most p. Hence v\_p(k) ≤ 1, with equality exactly when p ∣ \#E(F\_p) and the reduction of Λ meets E(F\_p)[p]. ∎

\subsection*{4.3 Lemmas 4 and 4′}

(⟹) Let τ ∈ E(Q\_p) have order p. For p ≥ 5 the formal group Ê(pZ\_p) is torsion-free, so  ≠ O, and \#E(F\_p) = p gives E(F\_p) = ⟨⟩; hence R := P′ − jτ ∈ E\_1(Q\_p) for some j and pP′ = pR. For a formal group over Z\_p, [p](t) = pt + (higher terms with Z\_p coefficients), so v\_p([p](t)) ≥ min(1 + v\_p(t), 2v\_p(t)) ≥ 2 whenever v\_p(t) ≥ 1. Thus e = v\_p(t(pP′)) ≥ 2.

(⟸, contrapositive) If E(Q\_p)[p] = 0 then E(Q\_p) is torsion-free, prime-to-p torsion injecting into E(F\_p) of order p, and being a compact p-adic Lie group of dimension 1 it is ≅ Z\_p with E\_1(Q\_p) = pE(Q\_p) of index p. Torsion-freeness lets one \textit{extend} the formal logarithm by ℓ(P) := log\_E(pP)/p, an injective homomorphism with ℓ(E(Q\_p)) = Z\_p and ℓ(E\_1(Q\_p)) = pZ\_p, since log\_E : E\_1(Q\_p) ≅ pZ\_p for p ≥ 5. So ℓ(P′) ∈ Z\_p\textasciicircum{}× when ′ ≠ O, v\_p(log\_E(pP′)) = 1, and as v\_p(log\_E(t)) = v\_p(t) on E\_1, e = 1. ∎

\textit{Lemma 4′.} E being ordinary at p, E[p] over \_p is an extension of the étale Z/p by μ\_p, with p distinct geometric points each of multiplicity p, and Frobenius acts on E[p]\textasciicircum{}\{ét\} through the unit root of X² − a\_pX + p ≡ X² − X (mod p), hence trivially; so E[p]\textasciicircum{}\{ét\} ⊂ E(F\_p), with equality as \#E(F\_p) = p. The divisor of ψ\_p mod p on the x-line is therefore p·Σ\_\{±T\}(x(T)) over the (p−1)/2 classes of E(F\_p)∖\{O\}, which matches deg ψ\_p − (p−1)/2 since ψ\_p has leading coefficient p; hence ψ\_p ≡ u·\textasciicircum{}p (mod p), p ∣ ψ\_p′, and p² ∣ ψ\_p(x₀+pk) − ψ\_p(x₀) by the step of Lemma 5, so min(v\_p(ψ\_p(x₀)), 2) is a function of . For the last claim lift  to Q ∈ E(Z\_p): as φ\_p and ψ\_p² are coprime with resultant a power of Δ ∈ Z\_p\textasciicircum{}× [known] and v\_p(x(pQ)) \textless{} 0, we get v\_p(ψ\_p(x(Q))) = v\_p(t(pQ)), and Lemma 4 gives E(Q\_p)[p] ≠ 0 ⟺ v\_p(ψ\_p(x₀)) ≥ 2. ∎

\subsection*{4.4 Theorem B, Theorem B′, Lemma 5}

\textit{Step 1 (denominators do not matter).} Write x(P′) = a/d′², y(P′) = b/d′³ with p ∤ d′, possible as ′ ≠ O. Homogenising, Φ := d′\textasciicircum{}\{2p²\}φ\_p(P′), Ψ := d′\textasciicircum{}\{p²−1\}ψ\_p(P′),  := d′\textasciicircum{}\{3p²\}ω\_p(P′) are integers, and

\begin{quote}
x(P) = Φ/(d′Ψ)², y(P) = /(d′Ψ)³, so d(P) = d′|Ψ|, α = ±Φ, β = ±, α/β = (φ\_p(P′)/ω\_p(P′))·d′\textasciicircum{}\{−p²\}.
\end{quote}

That d(P) = d′|Ψ| with no cancellation is where (A2) is used: it is what the factor ∏\_ℓ c\_ℓ in the definition of P′ secures, and it fails without it (§5.5). As q\_p is a homomorphism and q\_p(d′\textasciicircum{}\{p²\}) = 0,

\begin{quote}
q\_p(α/β) = q\_p(φ\_p(P′)) − q\_p(ω\_p(P′)).  (∗)
\end{quote}

\textit{Step 2 (the Weierstrass equation).} Multiply it at pP′ by ψ\_p(P′)\textasciicircum{}6 and write φ = φ\_p(P′), ω = ω\_p(P′), ψ = ψ\_p(P′):

\begin{quote}
ω² + a₁φωψ + a₃ωψ³ = φ³ + a₂φ²ψ² + a₄φψ⁴ + a₆ψ⁶.
\end{quote}

Since v\_p(ψ) = v\_p(d(P)) = e (as p ∤ d′), every term of ω² − φ³ other than a₁φωψ has valuation ≥ 2e, and that one has valuation v\_p(a₁) + e. Hence ω² ≡ φ³ (mod p²) under either hypothesis of Theorem B: e ≥ 2, or p ∣ a₁ with e ≥ 1.

\textit{Step 3.} Write φ\textasciicircum{}\{p−1\} = 1 + pA′, ω\textasciicircum{}\{p−1\} = 1 + pB′. Raising ω² ≡ φ³ (mod p²) to the power p−1 gives 2B′ ≡ 3A′ (mod p), so with λ := A′/2 mod p = q\_p(φ)/2 we get q\_p(φ) = 2λ, q\_p(ω) = 3λ, and (∗) gives q\_p(α/β) = −λ, while c ≥ 2 ⟺ q\_p(α/β) = 0 (§2). In the form 2(A′ − B′) = −A′ the implication λ ≠ 0 ⟹ q\_p(α/β) ≠ 0 needs no division. Under neither hypothesis the term a₁φωψ survives modulo p²; dividing by φ³ and using q\_p(1+pw) = −w gives the correction w₁ of Theorem B′. ∎

\textit{Lemma 5.} In characteristic p, [p] = V ∘ F and E is defined over F\_p, so x ∘ [p] is a rational function of x\textasciicircum{}p; hence φ\_p mod p is a polynomial in x\textasciicircum{}p and φ\_p′ ≡ 0 (mod p). For a polynomial f with p ∣ f′(x) one has p² ∣ f(x + pk) − f(x) (the Lean statement \texttt{eval\_\allowbreak{}sub\_\allowbreak{}eval\_\allowbreak{}dvd\_\allowbreak{}sq}), so φ\_p(x) mod p² depends only on x mod p. ∎

\subsection*{4.5 Theorem C, Corollary D, Corollary E}

\textit{Theorem C.} D being Galois-stable of order n, π\_D and \_D are defined over K and [n] = \_D ∘ π\_D. Fix a Weierstrass model of E/D and write x ∘ \_D = φ\_/ψ\_² with φ\_(X) = Σ\_\{i≤n\} c\_i X\textasciicircum{}i, c\_n = A ≠ 0 and deg ψ\_² = n − 1. Clearing denominators in x ∘ [n] = (φ\_/ψ\_²)∘(/²),

\begin{quote}
N(x) := Σ\_\{i≤n\} c\_i (x)\textasciicircum{}i (x)\textasciicircum{}\{2(n−i)\},  Den(x) := \textasciicircum{}\{2n\}·ψ\_(/²)²,  φ\_n/ψ\_n² = N/Den.
\end{quote}

Now deg(\textasciicircum{}i \textasciicircum{}\{2(n−i)\}) = n² − n + i, so i = n alone attains the top degree: deg N = n² with leading coefficient A ( being monic) and deg Den = n² − 1. Since x ∘ [n] : P¹ → P¹ has degree n², and the degree of a reduced representation is the larger of the two, cancelling a common factor of degree d would give n² − d = n²; so d = 0, N/Den is reduced, and as φ\_n and ψ\_n² are coprime with φ\_n monic, N = A·φ\_n. At x = x(T) one has (x(T)) = 0, so every term with i \textless{} n vanishes and N(x(T)) = A·(x(T))\textasciicircum{}n. ∎

\textit{Corollary D.} Let τ ∈ E(Q\_p) have order p; as in Lemma 4,  ≠ O. D := ⟨τ⟩ is Q\_p-rational with D∖\{O\} ⊂ E(Q\_p), so x(T) ∈ Z\_p and Theorem C gives φ\_p(x(T)) = (x(T))\textasciicircum{}p. Both sides are p-adic units, φ\_p = x(pQ)ψ\_p² with v\_p(x(pQ)) = −2e and v\_p(ψ\_p) = e, so q\_p(φ\_p(x(T))) = p·q\_p((x(T))) = 0 in F\_p. The reduction D∖\{O\} → E(F\_p)∖\{O\} is injective (D ∩ Ê = 0) between sets of cardinality p − 1, hence bijective, so x(T) mod p runs over all nonzero classes; by Lemma 5, λ ≡ 0. ∎

\textit{Corollary E.} Lemma 4 gives e ≥ 2, Corollary D gives λ = 0, Theorem B gives q\_p(α/β) = 0, that is c ≥ 2. ∎

\subsection*{4.6 Proposition 6, the direction λ = 0 ⟹ split}

The statement is [known] (§3.4); the argument below is independent of the sources cited there and carries no bound on p. It is given in outline. The arithmetic steps are the Lean theorems of \texttt{PadicHeightValuation\{5,6,7\}} (§6); the analytic steps use the Tate parametrisation as recorded in [ATAEC] and [Dok].

\textit{(1) An Euler identity.} φ\_p is weighted homogeneous of weight 2p² for weights (4, 6, 2) on (A, B, x); giving X = x\textasciicircum{}p the weight 2p and reducing 4A∂\_Aφ\_p + 6B∂\_Bφ\_p + 2x∂\_xφ\_p = 2p²φ\_p modulo p, both the right-hand side and ∂\_x f(x\textasciicircum{}p) vanish, leaving 4A·∂\_A f + 6B·∂\_B f = 0 in F\_p[A, B][X]. Hence A ∣ ∂\_B f; put ũ := (∂\_B f)/A.

\textit{(2) One slope.} For E\_s : y² = x³ + (a₄ + sp)x + a₆ a Taylor expansion of φ\_p in its coefficients gives λ\_s = λ\_0 − s·B₄/2, and likewise a slope B in the a₆ direction. λ is invariant under (x, y) ↦ (u²x, u³y), φ\_p transforming by the p-th power u\textasciicircum{}\{2p²\}; with u = 1 + pt that invariance reads 4a₄B₄ + 6a₆B ≡ 0, which is (1) again, so both slopes are multiples of one scalar ρ() = −ũ()/(4f()), the denominator being non-zero by Vélu's formulas in the shape f(X₀) = 4Y₀²g′(X₀)²/c\_V². By Serre–Tate theory the family E\_s exhausts the deformations of Ē over Z/p² and exactly one member splits, where λ = 0 by Corollary D. An affine function with non-zero slope vanishes once, so ρ ≠ 0 gives the assertion.

\textit{(3) A dichotomy.} Res\_X(ũ, g) has weight (p−1)(p−2)(p+3) = 12 m\_ũ, where m\_ũ := (p−1)(p−2)(p+3)/12 (so m\_ũ = 400, 561, 1001 at p = 17, 19, 23). The constants of this resultant are written c\_ũ, m\_ũ throughout; the letters c\_p, m\_p are reserved for the other resultant, Res\_x(G, ψ\_p mod p) of Proposition 16, whose exponent is m\_p = (p²−1)(2p²−3)/24 — 6,900 at p = 17, not 400. Take the transversal family E\_t : y² = x³ − 3x + (2 + t) over \_p((t)) at the cusp, of split multiplicative reduction with v(q) = v(Δ) = 1 [ATAEC, Thm V.5.3], [Dok, Thm A.1]; the Tate parametrisation [ATAEC, Thm V.3.1] gives ker V = ⟨q⟩ mod q\textasciicircum{}p, so the roots of g have v(β\_j − 1) = j (1 ≤ j ≤ n := (p−1)/2), those of f have valuations \{0, 1, 1, …, n, n\}, and the Newton polygon of f gives v(f(β\_j)) = j(p−j). Here ũ = −(1/3)(d/dt)f, and differentiation lowers a valuation by at most 1, so ord\_Δ Res\_X(ũ, g) ≥ Σ\_j (j(p−j) − 1) = m\_ũ; weighted homogeneity then forces Res\_X(ũ, g) = c\_ũ·Δ\textasciicircum{}\{m\_ũ\} with c\_ũ ∈ F\_p constant. For each p, either ρ ≠ 0 for every ordinary curve over \_p and every point of ker V∖\{O\}, or ρ = 0 for all of them.

\textit{(4) Which side.} E\textasciicircum{}\{(p)\} is the member of the same family at parameter t\textasciicircum{}p, hence defined over K\textasciicircum{}p, so d/dt reaches only the argument of the Tate parameter and v(dβ\_j/dt) = j − 1 exactly. Of the p−2j+1 roots of f at distance j from β\_j, the two flanking ones contribute 1/(1−u) + 1/(1−u\textasciicircum{}\{−1\}) = 1 between them, so f′(β\_j)/f(β\_j) has valuation exactly −j with leading coefficient p − 2j ≡ −2j. The chain rule gives 12ρ\_j = (d/dt)log f(β\_j) − f′(β\_j)(dβ\_j/dt)/f(β\_j) = (−j² + 2j²)/t + O(1): the residues −j² and −2j² differ by j² ≠ 0, so v(ρ\_j) = −1, ord\_Δ Res\_X(ũ, g) = m\_ũ exactly, and c\_ũ ≠ 0. ∎

The naive estimate fails — the second term has valuation −1, not ≥ 0 — and what closes the argument is that the two residues differ.

\subsection*{4.7 Theorem S}

Choose the model with p ∣ a₁ (§2); e = 1 since E is not split (Lemma 4). Then b ≥ 2e = 2 by Lemma 3, and by Theorem B, applicable in this model for e ≥ 1, q\_p(α/β) = −λ with λ ≠ 0 by Proposition 6; so c = 1. From h\_p(P)/2 = Λ\_σ + A with b ≥ 2 \textgreater{} 1 = c we get v\_p(h\_p(P)/2) = 1, and p is odd, so v\_p(h\_p(P)) = 1. Finally h\_p is a quadratic form, so h\_p(P) = p²h\_p(P′) and v\_p(h\_p(P′)) = −1. ∎

\subsection*{4.8 The main theorem}

Put n := ∏\_ℓ c\_ℓ, so p ∤ n. As h\_p is a quadratic form, v\_p(h\_p(P)) = v\_p(h\_p(nP)); and nP satisfies (A2) for every P ∈ E(Q).

\textit{(⟸ of (i)), case P ∈ Λ∖Λ₀.} Then  has order exactly p, \#E(F\_p) being p, and so does the reduction of P′ := nP as p ∤ n; P″ := pP′ satisfies (A1) and (A2). As E is split, Lemma 4 gives e = v\_p(d(P″)) ≥ 2, hence b ≥ e ≥ 2 by Lemma 3 and c ≥ 2 by Corollary E, so v\_p(h\_p(P″)) ≥ 2; and h\_p(P″) = p²h\_p(P′) gives v\_p(h\_p(P)) = v\_p(h\_p(P′)) ≥ 0.

\textit{(⟸ of (i)), case P ∈ Λ₀.} Then nP satisfies (A1) and (A2), so h\_p(nP) ∈ 2pZ\_p by Lemma 2 and v\_p(h\_p(P)) ≥ 1.

So v\_p(h\_p(P)) ≥ 0 for every P ∈ Λ. Polarising, h\_p(P,R) ∈ Z\_p since p is odd; every entry of the height matrix on any basis of Λ lies in Z\_p, every term of the Leibniz expansion has valuation ≥ 0, and v\_p(Reg\_p) ≥ 0. With Theorem A this is (iii).

\textit{(⟹ of (i)).} Suppose E(Q\_p)[p] = 0. By ι there is P ∈ Λ whose reduction has order p; with P′ := nP, Theorem S gives v\_p(h\_p(P′)) = −1, hence v\_p(h\_p(P)) = −1. If every entry of M\_Λ were in Z\_p then so would be the value of the quadratic form at every point of Λ (the Lean statement \texttt{quadratic\_\allowbreak{}value\_\allowbreak{}mem\_\allowbreak{}of\_\allowbreak{}entries\_\allowbreak{}mem}), a contradiction.

\textit{(ii).} Let r = 1 and let Q generate Λ; by ι, Q ∉ Λ₀ and Reg\_p = h\_p(Q). If E is split then v\_p(Reg\_p) ≥ 0 by the first case above. If not, Theorem S gives v\_p(Reg\_p) = −1 exactly. ∎

The choice of lattice is what carries the proof: what is needed is v\_p(h\_p) ≥ 0 on Λ, not v\_p(h\_p) ≥ 2 on Λ₀, and for a point of Λ∖Λ₀ it is enough that its reduction have order p.

\subsection*{4.9 Theorems H, L, F, J}

\textit{Theorem H.} The composition rule ψ\_\{ab\}(R) = ψ\_b(R)\textasciicircum{}\{a²\}ψ\_a(bR) at (a,b) = (p,m) and (m,p), with φ\_k(x(R)) = x(kR)ψ\_k(R)² and φ\_m monic of degree m² over Z[a₁,…,a₆] evaluated at X := x(pQ) of valuation −2e ≤ −2, gives the exact identity

\begin{quote}
φ\_p(x(mQ))·ψ\_m(Q)\textasciicircum{}\{2p²\} = φ\_p(x(Q))\textasciicircum{}\{m²\}·(1 + δ),  δ ∈ p²Z\_p.  (★)
\end{quote}

Only even powers of ψ\_m(Q) occur, so (★) is well posed for every m prime to p. Its three division-polynomial values are p-adic units, ψ\_m(Q)\textasciicircum{}\{2p²\} is a p-th power and 1 + δ ∈ 1 + p²Z\_p, so q\_p applied to (★) gives the claim; Corollary H′ follows since E(F\_p) ≅ Z/p and m² ≠ 0 in F\_p. ∎

\textit{Theorem L.} Repeat the computation of Theorem C with D = C = Ê[p], Galois-stable for good ordinary reduction: N(x) = ψ\_C(x)\textasciicircum{}\{2p\}·φ\_(φ\_C(x)/ψ\_C(x)²), the degree count is unchanged, so N = A·φ\_p and evaluating at  gives the identity, the factor (ψ\_C()²)\textasciicircum{}p being a p-th power. The leading coefficient of the x-part of \_C is c\textasciicircum{}\{−2\}, since t\_E(\_C R) = c·t″(R) + O(t″²), whence A = c². ∎

\textit{Theorem F.} p splits in k, so k ⊂ Q\_p and the CM order acts on E\_\{/Q\_p\} by Q\_p-endomorphisms; by Deuring's lifting theorem End(E\_\{/Q\_p\}) → End(Ē\_\{/F\_p\}) is an isomorphism for ordinary reduction, so Frobenius ϕ lifts to π with π +  = 1 and π = p. Put ϱ := 1 − π = , of degree p; on the special fibre  = 1 − ϕ is separable, as (1−ϕ)\textasciicircum{}*ω = ω ≠ 0, so ker ϱ is étale of order p over Z\_p. On it π = 1, so Frobenius acts trivially, and its points, being étale, are defined over Q\_p\textasciicircum{}\{ur\}, where Galois acts through Frobenius. Hence ker ϱ ⊂ E(Q\_p). ∎

\textit{Theorem J.} [Λ : Λ₀] = p gives det M\_\{Λ₀\} = p²·det M\_Λ; the entries of M\_\{Λ₀\} lie in pZ\_p by Lemma 2, so M\_\{Λ₀\} = p·N, and if N mod p has rank ρ\_0, Smith normal form gives v\_p(det N) ≥ r − ρ\_0. ∎

\subsection*{4.10 Theorems R and R\_s, Propositions R″ and R″\_s, Corollaries R′ and R′\_s}

\textit{The linear algebra.} Let K be a field of characteristic ≠ 2, B a symmetric bilinear form on a K-vector space V, χ a linear form and v ∈ V with χ(v) = 1; put W := ker χ.

\begin{namedthm}[(R0) [Lean \texttt{bilin\_\allowbreak{}rank\_\allowbreak{}one\_\allowbreak{}of\_\allowbreak{}zero\_\allowbreak{}iff}]]
If B(x,x) = 0 ⟺ χ(x) = 0, then c := B(v,v) ≠ 0 and B(x,y) = c·χ(x)·χ(y).
**(R0s) [Lean \texttt{bilin\_\allowbreak{}eq\_\allowbreak{}of\_\allowbreak{}vanishing\_\allowbreak{}on\_\allowbreak{}hyperplane}].** If χ(x) = 0 ⟹ B(x,x) = 0, then B(x,y) = B(x,v)·χ(y) + χ(x)·B(y,v) − B(v,v)·χ(x)·χ(y).
\end{namedthm}

Both begin with polarisation inside W: for w, w′ ∈ W the vector w + w′ lies in W, so expanding B(w+w′, w+w′) = 0 leaves 2B(w,w′) = 0 and B|\_\{W×W\} = 0 since char K ≠ 2. (R0) has more: W ⊥ v, for if a := B(w,v) ≠ 0 then s := −B(v,v)/(2a) makes x := s w + v satisfy χ(x) = 1 while B(x,x) = s²B(w,w) + 2 a s + B(v,v) = 0, against the direction ⟹; and c ≠ 0 is the contrapositive of ⟸ at v. Decomposing x = (x − χ(x)v) + χ(x)v and likewise y then gives B = c·χ⊗χ. In (R0s) only B|\_\{W×W\} = 0 is available, and the same decomposition gives the displayed identity. Its right-hand side involves only χ and the linear form B(·, v), so ker χ ∩ ker B(·,v) is contained in the radical of B and rank B ≤ 2 — recorded in Lean as \texttt{bilin\_\allowbreak{}zero\_\allowbreak{}of\_\allowbreak{}mem\_\allowbreak{}radical} and, in determinant form, as the vanishing of every 3×3 Gram determinant (\texttt{gram\_\allowbreak{}det\_\allowbreak{}three\_\allowbreak{}eq\_\allowbreak{}zero}).

\textit{Theorem R.} Put Q(P) := p·h\_p(P) and B(P,R) := p·h\_p(P,R) = (Q(P+R) − Q(P) − Q(R))/2, legitimate as p is odd. For P ∈ Λ₀, Lemma 2 gives v\_p(Q(P)) ≥ 2; for P ∈ Λ∖Λ₀, Theorem S gives v\_p(Q(P)) = 0 exactly, its reduction having order p in E(F\_p). So Q is Z\_p-valued, B is a Z\_p-valued symmetric bi-additive form, and since pΛ ⊆ Λ₀ both descend to V := Λ/pΛ ≅ F\_p\textasciicircum{}r, where  ≠ 0 by ι has the hyperplane Λ₀/pΛ as kernel. The two cases say exactly that the zero set of the quadratic form (x) = (x,x) is that hyperplane — here the \textit{exactness} in Theorem S is what is used, an inequality v\_p ≥ −1 leaving the reduction undecided — and (R0) applies with p ≥ 5 odd. The valuation statements follow: c·χ(P)χ(R) ≠ 0 means p ∤ B(P,R), that is v\_p(h\_p(P,R)) = −1; c·χ(P)χ(R) = 0 with one of the two in Λ∖Λ₀ means v\_p ≥ 0; and both in Λ₀ gives v\_p ≥ 1 by Lemma 2 polarised. ∎

\textit{Proposition R″.} For a unit u, u\textasciicircum{}\{p−1\} = 1 + p q with q ≡ q\_p(u), so (p−1)log\_p u = log(1 + pq) ≡ pq (mod p²) and, as p − 1 ≡ −1, log\_p u ≡ −p·q\_p(u) (mod p²). Take P′ ∈ Λ∖Λ₀ multiplied by n := ∏\_ℓ c\_ℓ so that (A2) holds, and P″ := pP′, in a model with p ∣ a₁. Lemma 3 gives h\_p(P″) = 2Λ\_σ + 2A with v\_p(Λ\_σ) ≥ 2e ≥ 2 and A = log\_p(α/β), and Theorem B gives q\_p(α/β) = −λ, so A ≡ pλ (mod p²) and h\_p(P″) ≡ 2pλ (mod p²). From h\_p(P″) = p²h\_p(P′) we get p·h\_p(P′) ≡ 2λ (mod p). Theorem H gives λ(m ) = m²λ(), that is λ = χ²λ₁ on E(F\_p), and c is the value at χ = 1. ∎

\textit{Corollary R′.} A basis with χ(e₁) = 1 and e₂, …, e\_r ∈ Λ₀ exists because Λ₀ ⊂ Λ has index p with cyclic quotient. Theorem R gives A₁₁ = p h\_p(e₁) ≡ c ≠ 0, a unit, and A\_\{1j\} ≡ 0 for j ≥ 2; for i, j ≥ 2 Lemma 2 polarised gives h\_p(e\_i, e\_j) ∈ pZ\_p, so A\_\{ij\} ∈ p²Z\_p. Hence A = D A′ D with A′ integral, det A = p\textasciicircum{}\{2(r−1)\} det A′ (Lean \texttt{det\_\allowbreak{}conj\_\allowbreak{}diagonal\_\allowbreak{}one\_\allowbreak{}p}) and det A = p\textasciicircum{}r Reg\_p, so Reg\_p = p\textasciicircum{}\{r−2\} det A′ (Lean \texttt{reg\_\allowbreak{}eq\_\allowbreak{}of\_\allowbreak{}det\_\allowbreak{}eq}). As A′₁₁ is a unit, det A′ = A′₁₁·det T. ∎

\textit{Theorem R\_s.} By part (i) of the main theorem M\_Λ is integral. For P, R ∈ Λ₀ Lemma 2 puts h\_p(P), h\_p(R), h\_p(P+R) in 2pZ\_p, so h\_p(P,R) ∈ pZ\_p as p is odd; that is,  vanishes on the hyperplane ker  = Λ₀/pΛ of V = Λ/pΛ, and (R0s) applies. In a basis with e₂, …, e\_r ∈ Λ₀ the block (i, j ≥ 2) of  is therefore zero, so  = [[a, bᵗ],[b, 0]]; the leading 2×2 minor on (e₁, e\_j) is −b\_j², so the rank is 2 as soon as some b\_j ≠ 0, and otherwise the matrix is a·E₁₁ of rank 1 or 0. ∎

\textit{Corollary R′\_s.} Write M := M\_Λ in that basis, so M\_\{ij\} ∈ pZ\_p for i, j ≥ 2. Suppose ρ₁ = 2. A change of basis inside Λ₀ ⊗ Z\_p — which multiplies det M by the square of a unit and so does not move v\_p — brings b to (b₁, 0, …, 0) with b₁ ∈ Z\_p\textasciicircum{}×. The leading 2×2 block S = [[M₁₁, b₁],[b₁, pC]] has det S ≡ −b₁² ≢ 0, hence S ∈ GL₂(Z\_p); its Schur complement in M is p·T\_s with T\_s integral of size r − 2, because the off-block entries lie in pZ\_p and S\textasciicircum{}\{−1\} is integral. So det M = det S·p\textasciicircum{}\{r−2\}·det T\_s and v\_p(Reg\_p) = (r − 2) + v\_p(det T\_s). If ρ₁ ≤ 1 then p ∣ M\_\{1j\} for j ≥ 2 as well, so rows 2, …, r are divisible by p, det M = p\textasciicircum{}\{r−1\}·det N with N integral, and v\_p(Reg\_p) ≥ r − 1; if moreover a ≡ 0 then N mod p has vanishing determinant and v\_p(Reg\_p) ≥ r. In all three cases v\_p(Reg\_p) ≥ r − ρ₁, and the uniform bound p\textasciicircum{}\{r−2\} ∣ det M holds directly from the Leibniz expansion: in the term of a permutation σ the factors avoiding the block divisible by p are the one in row 1 and the one in column 1, at most two of the r factors (the Lean statement \texttt{pow\_\allowbreak{}dvd\_\allowbreak{}det\_\allowbreak{}of\_\allowbreak{}offdiag\_\allowbreak{}block}). For r = 2 the matrix T\_s is empty and v\_p(Reg\_p) = 0 ⟺ ρ₁ = 2, the case ρ₁ ≤ 1 giving ≥ 1 (the Lean statement \texttt{det\_\allowbreak{}two\_\allowbreak{}dvd\_\allowbreak{}iff}). ∎

\textit{Proposition R″\_s.} Work in a model with p ∣ a₁ (§2); this changes neither h\_p nor φ\_p, only the pair (b, c). Put P′ := n e₁ with n = ∏\_ℓ c\_ℓ, P″ := pP′ and e := v\_p(d(P″)), so e ≥ 2 by Lemma 4 since E is split. In the Weierstrass equation multiplied by ψ\_p(P′)\textasciicircum{}6 (Step 2 of §4.4), every term of ω² − φ³ other than a₁φωψ has valuation ≥ 2e ≥ 4, while a₁φωψ has valuation ≥ v\_p(a₁) + e ≥ 3; hence

\begin{quote}
ω² ≡ φ³ (mod p³).
\end{quote}

Corollary D gives q\_p(φ) = 2λ = 0 and q\_p(ω) = 3λ = 0, so φ\textasciicircum{}\{p−1\} = 1 + p²A₂ and ω\textasciicircum{}\{p−1\} = 1 + p²B₂; raising the congruence to the power p − 1 gives (1 + p²B₂)² ≡ (1 + p²A₂)³ (mod p³), that is 2B₂ ≡ 3A₂ (mod p), the same relation one digit deeper (the Lean statement \texttt{dvd\_\allowbreak{}three\_\allowbreak{}mul\_\allowbreak{}sub\_\allowbreak{}two\_\allowbreak{}mul\_\allowbreak{}two}). With λ\textasciicircum{}\{(2)\} := A₂/2 this reads q\textasciicircum{}\{(2)\}(φ) = 2λ\textasciicircum{}\{(2)\} and q\textasciicircum{}\{(2)\}(ω) = 3λ\textasciicircum{}\{(2)\}. The factor d′\textasciicircum{}\{p²\} of Step 1 of §4.4 has q\_p = 0 and q\textasciicircum{}\{(2)\} = 0, since (1 + p q)\textasciicircum{}\{p²\} ≡ 1 (mod p³) for p ≥ 3; and q\textasciicircum{}\{(2)\} is additive on its domain, so

\begin{quote}
q\textasciicircum{}\{(2)\}(α/β) = q\textasciicircum{}\{(2)\}(φ) − q\textasciicircum{}\{(2)\}(ω) = −λ\textasciicircum{}\{(2)\}.
\end{quote}

For a unit u with q\_p(u) = 0, log\_p u = log(1 + p²q\textasciicircum{}\{(2)\}(u) + O(p³))/(p−1) ≡ −p²q\textasciicircum{}\{(2)\}(u) (mod p³). Hence A = log\_p(α/β) ≡ p²λ\textasciicircum{}\{(2)\} (mod p³), while the formal part has v\_p(Λ\_σ) ≥ 2e ≥ 4, so h\_p(P″) ≡ 2p²λ\textasciicircum{}\{(2)\} (mod p³). Finally h\_p(P″) = p²n²h\_p(e₁), whence n²·h\_p(e₁) ≡ 2λ\textasciicircum{}\{(2)\} (mod p) and, n being prime to p, a = 2λ\textasciicircum{}\{(2)\}/n² (the Lean statement \texttt{split\_\allowbreak{}const\_\allowbreak{}eq}). For r = 1 the class χ(e₁) is nonzero and Reg\_p = h\_p(e₁), so v\_p(Reg\_p) = 0 ⟺ a ≠ 0 ⟺ λ\textasciicircum{}\{(2)\} ≠ 0. ∎

\subsection*{4.11 Lemma 7, Proposition 8, Theorems K and N, Propositions 9 and 10}

\textit{Lemma 7.} Translation τ\_P is an automorphism of the smooth model over Z\_p and τ\_P(O) = P has non-zero reduction, so x is regular at P and x ∘ τ\_P ∈ Ô\_\{E,O\} = Z\_p[[t]]. The linear coefficient is read from the invariant differential: ω = dx/(2y + a₁x + a₃) is translation-invariant and ω = (1 + a₁t + ⋯)dt, so d(x ∘ τ\_P) = τ\_P\textasciicircum{}*(dx) = (2y + a₁x + a₃)(τ\_P(·))·ω and the coefficient of t at t = 0 is 2y(P) + a₁x(P) + a₃. Finally 2y + a₁x + a₃ is divisible by p only at  ∈ Ē[2], and \#Ē(F\_p) = p is odd, so only at  = O. For κ\_j apply this with P = n e₁ and R = n e\_j ∈ E\_1(Q\_p), t = t(n e\_j) of valuation ≥ 1; the first digit of the difference is c₁·(t/p) mod p = c₁·γ(n e\_j). ∎

\textit{Proposition 8.} Corollary H″, applied to the multiplication [m] with p ∤ m, gives F(mQ) = m²F(Q) for F(Q) := λ\textasciicircum{}\{(2)\}(x(Q)) — Theorem H one digit deeper. By Proposition R⁗ and Lemma 7, for Q ∈ E(Z\_p) with  ≠ O,

\begin{quote}
F(Q) = λ\_T(χ(Q)) + s(χ(Q))·c₁(χ(Q))·γ(Q),
\end{quote}

since x(Q) = x(χ(Q)T₀) + p·c₁γ(Q) + O(p²) and λ\textasciicircum{}\{(2)\} depends on the lift only through its first digit. On the coset χ = 1 the value γ(Q) runs over all of F\_p, because t : E\_1(Q\_p) → pZ\_p is a bijection. Under Q ↦ mQ one has χ ↦ mχ and γ ↦ mγ, so comparing the two sides of F(mQ) = m²F(Q) as affine functions of γ(Q) gives λ\_T(m) = m²λ\_T(1) for the constant terms and m·s(m)c₁(m) = m²·s(1)c₁(1) for the linear ones. ∎

\textit{Theorem K.} Proposition R″\_s gives h\_p(P) ≡ 2λ\textasciicircum{}\{(2)\}(x(nP))/n² = 2F(nP)/n² (mod p). Insert the display of the previous proof with χ(nP) = nχ(P), γ(nP) = nγ(P) and then (L1), (L2):

\begin{quote}
F(nP) = λ\_\{T,1\}n²χ(P)² + B₀·nχ(P)·nγ(P) = n²(λ\_\{T,1\}χ(P)² + B₀χ(P)γ(P)),
\end{quote}

and n² cancels. Polarising gives the bilinear form; the Lean statements \texttt{master}, \texttt{master\_\allowbreak{}symm}, \texttt{master\_\allowbreak{}add\_\allowbreak{}left}, \texttt{master\_\allowbreak{}smul\_\allowbreak{}left} record its shape and bilinearity. ∎

\textit{Corollary K′.} (i)–(iii) are linear algebra over F\_p on the form χ⊗χ and χ⊗γ + γ⊗χ (Lean \texttt{master\_\allowbreak{}mem\_\allowbreak{}radical}, \texttt{master\_\allowbreak{}eq\_\allowbreak{}of\_\allowbreak{}gam\_\allowbreak{}eq\_\allowbreak{}smul}, \texttt{master\_\allowbreak{}offdiag\_\allowbreak{}ne\_\allowbreak{}iff}, \texttt{master\_\allowbreak{}diag\_\allowbreak{}eq\_\allowbreak{}zero\_\allowbreak{}iff}, \texttt{master\_\allowbreak{}bad\_\allowbreak{}class\_\allowbreak{}unique}). For (iv) take a basis e₁ with χ(e₁) = 1, e₂ ∈ ker χ with γ(e₂) ≠ 0, and a basis of Λ₁; then M\_Λ has the shape [[a, b, 0 …], [b, pC₂₂, p(∗)], [0, p(∗), pĈ]] with b ∈ Z\_p\textasciicircum{}× by Theorem K, the Schur complement of the leading 2×2 block is p·Ĉ modulo p², and Corollary R′\_s applies (Lean \texttt{det\_\allowbreak{}three\_\allowbreak{}layer}, \texttt{det\_\allowbreak{}sq\_\allowbreak{}dvd\_\allowbreak{}iff}). ∎

\textit{Proposition 9.} By Lemma 5, φ\_p′ = p·G with G ∈ Z[a₁, …, a₆][x], and φ\_p″ = p·H. Taylor at x₀ with h = pk gives

\begin{quote}
φ\_p(x₀ + pk) ≡ φ\_p(x₀) + p²·k·G(x₀) + p³·k²·H(x₀)/2 (mod p⁴),
\end{quote}

the cubic term being p⁴k³·(H′/6)(x₀) with p ∤ 6 (Lean \texttt{eval\_\allowbreak{}shift\_\allowbreak{}second\_\allowbreak{}digit}, in the integral form 2(P(r+pk) − P(r)) − 2p²kG(r) − p³k²H(r) ≡ 0 mod p⁴). Write ε for the correction, so p² ∣ ε, and u := φ\_p(x₀), u\textasciicircum{}\{p−1\} = 1 + p²Θ(x₀). Then (u + ε)\textasciicircum{}\{p−1\} ≡ u\textasciicircum{}\{p−1\}(1 + (p−1)ε/u) (mod p⁴), because ε² is divisible by p⁴ (Lean \texttt{fermatQuot\_\allowbreak{}second\_\allowbreak{}move}), whence

\begin{quote}
(u + ε)\textasciicircum{}\{p−1\} − 1 ≡ p²Θ(x₀) + (p − 1)ε/u (mod p⁴).
\end{quote}

Dividing by p² and using (p−1)p ≡ −p (mod p²) gives the displayed formula. ∎

\textit{Proposition 10.} By Proposition 9 the function δ ↦ Θ(x(T) + pδ) is, modulo p, the affine map Θ\_T + 2sδ with s ≠ 0 by §4.6 — the quadratic term carries a factor p. An analytic function Z\_p → Z\_p whose reduction modulo p is a bijection of F\_p and whose derivative is a unit at every point is a bijective isometry of Z\_p, by Hensel's lemma applied at each residue; in particular 0 is attained exactly once. The equality of valuations is the factorisation f(δ) − f(δ\_\textit{) = (δ − δ\_})·(b + p·q) with b = 2s a unit, which is Lean \texttt{root\_\allowbreak{}factor} and \texttt{dvd\_\allowbreak{}iff\_\allowbreak{}dvd\_\allowbreak{}sub\_\allowbreak{}root}; the uniqueness of the bad class at each layer without assuming a root is \texttt{unique\_\allowbreak{}bad\_\allowbreak{}class}, and its existence is \texttt{exists\_\allowbreak{}bad\_\allowbreak{}class}, one step of Hensel written by hand. ∎

\textit{Lemma 3′.} u(P) = σ\_p(t)/t = 1 + c₂t + c₃t² + ⋯ with c₂ = a₁/2 and v\_p(t) = e, so v\_p(log\_p u) ≥ min(v\_p(c₂) + e, 2e). For e = 1 and p ∣ a₁ this is 2 = 2e; for e ≥ 2 the first term is the smaller one. ∎

\textit{Theorem N.} Work in a global minimal model with p ∣ a₁ (§2), and put P′ := nP, P″ := pP′, e := v\_p(d(P″)) ≥ 2 (Lemma 4). Four terms have to be carried modulo p⁴, where before only their leading behaviour was needed.

(1) \textit{The formal part.} By Lemma 3′, 2Λ\_σ ≡ a₁·t(P″) (mod p⁴), of valuation v\_p(a₁) + e ≥ 3; the next term is 2c₃t² of valuation ≥ 2e ≥ 4.

(2) \textit{The Weierstrass equation.} Multiplying at P″ by ψ\_p(P′)\textasciicircum{}6 and writing φ = φ\_p(P′), ω = ω\_p(P′), ψ = ψ\_p(P′) with v\_p(ψ) = e and φ, ω units,

\begin{quote}
ω² + a₁φωψ + a₃ωψ³ = φ³ + a₂φ²ψ² + a₄φψ⁴ + a₆ψ⁶,
\end{quote}

the terms a₃ωψ³, a₄φψ⁴, a₆ψ⁶ having valuation ≥ 3e ≥ 6 and a₂φ²ψ² valuation ≥ 2e ≥ 4. Hence ω² ≡ φ³ − a₁φωψ (mod p⁴). Raising to the power p − 1 and writing φ\textasciicircum{}\{p−1\} = 1 + p²A, ω\textasciicircum{}\{p−1\} = 1 + p²B with A = Θ(x(P′)), one gets 2B ≡ 3A + a₁ωψ/(p²φ²) (mod p²), so

\begin{quote}
Q₂(φ/ω) = A − B ≡ −Θ/2 − a₁ωψ/(2p²φ²) (mod p²).
\end{quote}

(3) \textit{The denominator.} By Step 1 of §4.4, α/β = (φ\_p(P′)/ω\_p(P′))·d′\textasciicircum{}\{−p²\}, and (1 + p q)\textasciicircum{}\{p²\} ≡ 1 + p³q (mod p⁴) for p ≥ 5, so Q₂(d′\textasciicircum{}\{p²\}) ≡ p·q\_p(d′) (mod p²); equivalently −p²log\_p(d′) ≡ p³q\_p(d′) (mod p⁴). \textbf{This is the term that Proposition R″\_s could drop and Theorem N cannot.}

(4) \textit{The logarithm.} For a unit u with q\_p(u) = 0, (p−1)log\_p u = log(1 + p²Q₂(u)) ≡ p²Q₂(u) (mod p⁴), and 1/(p−1) ≡ −(1 + p) (mod p²), so log\_p u ≡ −p²(1 + p)Q₂(u) (mod p⁴).

Putting (2), (3), (4) into A = log\_p(α/β) and adding (1),

\begin{quote}
h\_p(P″) = 2Λ\_σ + 2A ≡ a₁t(P″) + (1 + p)(p²Θ + a₁ωψ/φ²) + 2p³q\_p(d′) (mod p⁴).
\end{quote}

The two a₁ terms cancel: t(P″) = −x(P″)/y(P″) = −φψ/ω, so a₁t(P″) + a₁ωψ/φ² = a₁ψ·(ω/φ² − φ/ω) = a₁ψ·(ω² − φ³)/(φ²ω), which by (2) is −a₁²ψ²/φ + O(p⁴) of valuation ≥ 2v\_p(a₁) + 2e ≥ 6; and the surviving p·a₁ωψ/φ² of (1 + p) has valuation ≥ 1 + 1 + e ≥ 4. Hence

\begin{quote}
h\_p(P″) ≡ (1 + p)p²Θ(x(P′)) + 2p³q\_p(d′) (mod p⁴),
\end{quote}

and h\_p(P″) = p²n²h\_p(P) gives the statement. The cancellation is the Lean statement \texttt{a1\_\allowbreak{}correction\_\allowbreak{}cancels}; the two-digit criteria that Theorem N feeds are \texttt{second\_\allowbreak{}digit\_\allowbreak{}zero\_\allowbreak{}iff} and \texttt{second\_\allowbreak{}digit\_\allowbreak{}criterion}. ∎

\textit{Corollary N′.} Write Θ = pΘ₁. By Theorem N, n²h\_p(P) ≡ p(Θ₁ + 2q\_p(d′)) (mod p²), and p ∤ n. ∎

\subsection*{4.12 Corollaries N⁽³⁾ and N⁽⁴⁾, Theorem G′, Propositions 11, 12 and 19, Corollaries N″ and N‴}

Throughout, the model has p ∣ a₁ (§2), P′ = nP, P″ = pP′, e = v\_p(t(P″)) ≥ 2, φ = φ\_p(x(P′)), ψ = ψ\_p(x(P′)), ω = ω\_p(P′), and the identity of the Remark opening §4 is expanded one digit at a time. Only the terms newly reached are discussed; the rest is §4.11.

\textit{The coefficient of t² in σ\_p.} Write σ\_p(t)/t = 1 + c₂t + c₃t² + ⋯ and g₁ := 2c₃ − c₂². From the defining differential equation x(t) + c = −(d/ω)((1/σ)(dσ/ω)) with D := w\textasciicircum{}\{−1\}d/dt, w = 1 + a₁t + (a₁² + a₂)t² + ⋯ and g := σ′/σ = t\textasciicircum{}\{−1\} + g₀ + g₁t + ⋯ (so g₀ = c₂ = a₁/2), one has (g/w)′ = −w·(x + c). Now w·(x + c) = t\textasciicircum{}\{−2\} + 0·t\textasciicircum{}\{−1\} + c + (a₃ + a₁c)t + ⋯ — the vanishing of the t\textasciicircum{}\{−1\} coefficient is what makes σ\_p odd — so g/w = t\textasciicircum{}\{−1\} + C − ct − (a₃ + a₁c)t²/2 + ⋯, while 1/w = 1 − a₁t − a₂t² − ⋯ gives for the coefficient of t¹ in g/w the value g₁ − a₁g₀ − a₂. Comparing,

\begin{quote}
\textbf{g₁ = a₁g₀ + a₂ − c = a₁²/2 + a₂ − c},  c = (b₂ − E₂(E, ω))/12,
\end{quote}

and the t⁰ coefficient g₀ − a₁ = C = −a₁/2 is consistent with c₂ = a₁/2.

\textit{Corollary N⁽³⁾.} Four stages, each one digit deeper than §4.11.

(1) \textit{The formal part.} 2Λ\_σ = a₁t + g₁t² + O(t³) with v\_p(a₁t) ≥ 1 + e ≥ 3, v\_p(g₁t²) ≥ 2e ≥ 4 and v\_p(t³) ≥ 3e ≥ 6, so modulo p⁵ the t² term survives and the t³ term does not.

(2) \textit{The Weierstrass equation.} Modulo p⁵ the surviving corrections are a₁φωψ and a₂φ²ψ², so η := (ω² − φ³)/φ³ = −a₁ωψ/φ² + a₂ψ²/φ + O(p⁶) with v\_p(η) ≥ 3. Writing φ\textasciicircum{}\{p−1\} = 1 + p²Θ and ω\textasciicircum{}\{p−1\} = 1 + p²B and expanding (1 + p²B)² = (1 + p²Θ)³(1 + η)\textasciicircum{}\{p−1\} modulo p⁶ (η² has valuation ≥ 6) gives 2B + p²B² ≡ 3Θ + 3p²Θ² + (p−1)η/p² + 3(p−1)Θη (mod p⁴), hence B ≡ 3Θ/2 + (3/8)p²Θ² + (p−1)η/(2p²) (mod p³) (Lean \texttt{sq\_\allowbreak{}cube\_\allowbreak{}step}, in the integral form 8B ≡ 12Θ + 3p²Θ² + 4E).

(2′) \textit{Q₂ is not additive on quotients.} From (1 + p²B)(1 + p²Q) = 1 + p²Θ with Q := Q₂(φ/ω),

\begin{quote}
\textbf{Q = Θ − B + p²B(B − Θ) + O(p⁴)}
\end{quote}

(Lean \texttt{q2\_\allowbreak{}quotient}, as the exact identity Q − (Θ − B + p²B(B − Θ)) = p⁴B²Q), so with B ≡ 3Θ/2 (mod p) one gets Q ≡ −Θ/2 + (3/8)p²Θ² − (p−1)η/(2p²) (mod p³). \textbf{This stage is what an induction on digits omits}: the second digit only needed Q ≡ Θ − B.

(3) \textit{The denominator.} −2p²log\_p(d′) ≡ 2p³(1 + p)F − p⁴F² (mod p⁵), from 1/(p−1) ≡ −(1 + p + p²) (mod p³) and the quadratic term of log(1 + pF) (Lean \texttt{denominator\_\allowbreak{}third\_\allowbreak{}digit}, \texttt{neg\_\allowbreak{}inv\_\allowbreak{}pred\_\allowbreak{}cube}).

(4) \textit{The logarithm.} 2log\_p(φ/ω) ≡ −2p²(1 + p + p²)Q + p⁴Q² (mod p⁵). Inserting Q from (2′): the Θ term is p²(1 + p + p²)Θ; the Θ² terms are −(3/4)p⁴Θ² from Q and +(1/4)p⁴Θ² from p⁴Q², total −(1/2)p⁴Θ²; the η term is (1 + p + p²)(p−1)η = (p³ − 1)η ≡ −η (mod p⁵).

Finally −η = a₁ωψ/φ² − a₂ψ²/φ, and a₁t(P″) + a₁ωψ/φ² = a₁ψ(ω² − φ³)/(φ²ω) = −a₁²ψ²/φ + O(p⁷) has valuation ≥ 2v\_p(a₁) + 2e ≥ 6, so the two a₁ terms cancel modulo p⁵ as they did modulo p⁴. What is left is g₁t² − a₂ψ²/φ, and t² = φ²ψ²/ω² ≡ ψ²/φ (mod p⁷), so it equals (g₁ − a₂)ψ²/φ = (a₁²/2 − c)ψ²/φ ≡ −c·ψ²/φ (mod p⁶) when p ∣ a₁. Since ψ²/φ = 1/x(P″), dividing h\_p(P″) by p² gives Corollary N⁽³⁾. The intermediate constant a₁²/2 − c belongs to this computation in this model and \textbf{not to the statement}: the model-free identity behind the same step is 2c₃ − a₂ − (3/4)a₁² = −c in the coefficients c₃, c₄ of σ\_p(t)/t (Lean \texttt{sigma\_\allowbreak{}third\_\allowbreak{}coeff}, which writes c₃ as g₁; its companion \texttt{sigma\_\allowbreak{}fourth\_\allowbreak{}coeff} is a₃ − 2c₄ = −a₁(a₁²/2 + a₂ + E₂/8)), the three further contributions being −a₂/X, (a₁²/2)(Y/X²)² and the a₁ remainder −a₁²/X. Both identities are Proposition 17 at m = 1 read in the parameter t instead of z. ∎

\textit{Corollary N⁽⁴⁾.} The same four stages one digit deeper. What is new: (a) 1/(p−1) ≡ −(1 + p + p² + p³) (mod p⁴) (Lean \texttt{inv\_\allowbreak{}pred\_\allowbreak{}quartic}); (b) the cubic term of the logarithm of the denominator, −2p²log\_p(d′) ≡ 2p³(1 + p + p²)F − p⁴(1 + p)F² + (2/3)p⁵F³ (mod p⁶) (Lean \texttt{denominator\_\allowbreak{}fourth\_\allowbreak{}digit}, multiplied by 6 to stay in Z); (c) the two contributions of Θη — one from the term (1/4)(p−1)Θη of Q and one from p⁴(1 + p)Q² — cancel, the remainder having valuation ≥ 7 (Lean \texttt{theta\_\allowbreak{}eta\_\allowbreak{}cancels}); (d) everything of valuation 3e ≥ 6 dies: the t³ coefficient of σ\_p, the term a₃ωψ³ of the Weierstrass equation, and p⁶Θ³ (Lean \texttt{cubic\_\allowbreak{}term\_\allowbreak{}vanishes}). No mixed term ΘF appears. ∎

\textit{Corollaries N″ and N‴.} With Θ = pΘ₁ and Θ₁ + 2F = pS, Corollary N⁽³⁾ collapses to n²h\_p ≡ p²(S − F² − W) (mod p³), which is Corollary N″ (Lean \texttt{cube\_\allowbreak{}dvd\_\allowbreak{}sub\_\allowbreak{}crit}, \texttt{cube\_\allowbreak{}dvd\_\allowbreak{}iff\_\allowbreak{}crit}, \texttt{quadrichotomy}). With S − F² − W = pU, Corollary N⁽⁴⁾ collapses to 6·n²h\_p ≡ p³(6U + 6(S − F²) + 4F³) (mod p⁴) and p ∤ 6, which is Corollary N‴ (Lean \texttt{fourth\_\allowbreak{}digit\_\allowbreak{}reduction}, \texttt{quartic\_\allowbreak{}dvd\_\allowbreak{}iff\_\allowbreak{}crit}, \texttt{quintichotomy}). ∎

\textit{Proposition 12.} Good ordinary reduction means the formal group has height one, so [p](t) ≡ (unit)·t\textasciicircum{}p (mod p) and the coefficients c\_i of [p](t) = pt + Σ\_\{i≥2\}c\_it\textasciicircum{}i are divisible by p for 2 ≤ i \textless{} p. With t₁ = pγ + p²γ₂ + ⋯ this gives v\_p(c\_it₁\textasciicircum{}i) ≥ 1 + 2 = 3 for 2 ≤ i \textless{} p and v\_p(c\_pt₁\textasciicircum{}p) ≥ p ≥ 5, hence t(P″) = [p](t₁) ≡ p²γ (mod p³) (Lean \texttt{low\_\allowbreak{}terms\_\allowbreak{}vanish}, \texttt{p\_\allowbreak{}mult\_\allowbreak{}second\_\allowbreak{}digit}); in particular e = v\_p(t(P″)) equals 2 exactly when γ ≢ 0. Then 1/x(P″) = t(P″)²(1 + a₁t(P″) + ⋯) with v\_p(a₁t³) ≥ 6 (Lean \texttt{inv\_\allowbreak{}x\_\allowbreak{}correction}), so W = c·t(P″)²/p⁴ ≡ c·γ² (mod p) (Lean \texttt{w\_\allowbreak{}local}, \texttt{w\_\allowbreak{}ne\_\allowbreak{}zero\_\allowbreak{}iff}). ∎

\textit{Proposition 11.} Lemma 7 read to order three: x ∘ τ\_T ∈ Z\_p[[t]], and δ = (x(nP) − x(T))/p = (x∘τ\_T(t₁) − x(T))/p with t₁ = pΓ + O(p⁴). Substituting and dividing by p gives the display, the coefficients being in Z\_p. ∎

\textit{Theorem G′.} Four steps.

(1) \textit{The x-coordinate moves by exactly p\textasciicircum{}j.} By Lemma 7, X := x ∘ τ\_T ∈ Z\_p[[t]] and X′(0) = c₁(T) = 2y(T) + a₁x(T) + a₃ ∈ Z\_p\textasciicircum{}×; since X′ ∈ Z\_p[[t]] and t₁ ∈ pZ\_p, X′(t₁) ≡ X′(0) (mod p) is again a unit. Writing X(t₁ + h) − X(t₁) = Σ\_\{k≥1\}b\_kh\textasciicircum{}k with b\_k ∈ Z\_p the coefficients of the shifted series and b₁ = X′(t₁), the term k = 1 has valuation j and the terms k ≥ 2 valuation ≥ 2j \textgreater{} j. Hence v\_p(X(t₁ + p\textasciicircum{}j) − X(t₁)) = j.

(2) \textit{All divided derivatives of φ\_p below the p-th are divisible by p.} By §2, φ\_p(x) ≡ f(x\textasciicircum{}p) (mod p) with f ∈ F\_p[a₁, …, a₆][X]; hence, as polynomials in two variables over F\_p,

\begin{quote}
φ\_p(x + y) ≡ f((x + y)\textasciicircum{}p) = f(x\textasciicircum{}p + y\textasciicircum{}p) ≡ φ\_p(x) + (terms in y\textasciicircum{}p only),
\end{quote}

so the coefficient of y\textasciicircum{}k in φ\_p(x + y) — which is φ\_p\textasciicircum{}\{(k)\}/k!, a polynomial with coefficients in Z[a₁, …, a₆] — is divisible by p for every 1 ≤ k ≤ p − 1. For k = 1 this is Lemma 5, φ\_p′ = p·G.

(3) \textit{The value of φ\_p moves by exactly p\textasciicircum{}\{j+1\}.} Put x₀ := X(t₁) and Δ := X(t₁ + p\textasciicircum{}j) − X(t₁), so v\_p(Δ) = j ≥ 1 by (1). By (2),

\begin{quote}
φ\_p(x₀ + Δ) − φ\_p(x₀) = p·G(x₀)·Δ + p·Δ²·B + Δ\textasciicircum{}p·C  with B, C ∈ Z\_p,
\end{quote}

B collecting the divided derivatives of orders 2 ≤ k ≤ p − 1 divided by p and C those of order ≥ p. The three valuations are j + 1, ≥ 2j + 1 ≥ j + 2 and ≥ pj ≥ 5j ≥ j + 4, and G(x₀) ≡ G(x(T)) (mod p) is a unit by (ND). Hence the difference has valuation exactly j + 1 (Lean \texttt{shift\_\allowbreak{}depth\_\allowbreak{}exact}; the inequality half, which needs no non-degeneracy, is \texttt{shift\_\allowbreak{}depth\_\allowbreak{}dvd}).

(4) \textit{Raising to the power p − 1 and dividing by p².} Φ := φ\_p(x₀) is a unit (Corollary D), and Φ' = Φ + D with v\_p(D) = j + 1 ≥ 2 by (3). From (Φ + D)\textasciicircum{}\{p−1\} − Φ\textasciicircum{}\{p−1\} = (p − 1)Φ\textasciicircum{}\{p−2\}D + D²C′ (Lean \texttt{pow\_\allowbreak{}sub\_\allowbreak{}pow\_\allowbreak{}expand}) and v\_p(D²) ≥ 2(j + 1) ≥ j + 3, and since p ∤ (p − 1)Φ\textasciicircum{}\{p−2\}, the left-hand side has valuation exactly j + 1 (Lean \texttt{unit\_\allowbreak{}pow\_\allowbreak{}diff\_\allowbreak{}exact}). Dividing by p², which is what Q₂ does, gives v\_p(Θ(t₁ + p\textasciicircum{}j) − Θ(t₁)) = j − 1 (Lean \texttt{quotient\_\allowbreak{}by\_\allowbreak{}p\_\allowbreak{}sq}, \texttt{theta\_\allowbreak{}sensitivity\_\allowbreak{}exact}; the determination half is \texttt{theta\_\allowbreak{}digit\_\allowbreak{}determined}, the case j = 3 is \texttt{three\_\allowbreak{}layers\_\allowbreak{}determine} and its sharp companion \texttt{layer\_\allowbreak{}three\_\allowbreak{}moves\_\allowbreak{}digit\_\allowbreak{}two}). ∎

\textit{On (ND), and why the same statement for ψ\_p is free.} Consider Ψ := ψ\_p ∘ X ∈ Z\_p[[t]]. Since [p](T + R(t)) = R([p](t)), the identity x(mQ) = φ\_m/ψ\_m² at m = p reads

\begin{quote}
φ\_p(X(t)) = x(R([p](t)))·Ψ(t)² = [p](t)\textasciicircum{}\{−2\}·u([p](t))·Ψ(t)²,  u ∈ Z\_p[[s]], u(0) = 1,
\end{quote}

and with [p](t) = pt·M(t), M(0) = 1, and Ψ(t) = t·Ψ₁(t) this becomes φ\_p(X(t)) = p\textasciicircum{}\{−2\}M(t)\textasciicircum{}\{−2\}u([p](t))Ψ₁(t)². At t = 0, φ\_p(x(T)) = p\textasciicircum{}\{−2\}Ψ₁(0)², and φ\_p(x(T)) is a unit (Corollary D), so v\_p(Ψ₁(0)) = 1. Since Ψ₁(0) = Ψ′(0) = ψ\_p′(x(T))·c₁ with c₁ a unit,

\begin{quote}
\textbf{v\_p(ψ\_p′(x(T))) = 1 exactly} — no hypothesis needed.
\end{quote}

For φ\_p the same route gives nothing: the invariant differential [p]\textasciicircum{}*ω = pω says, after clearing denominators, G·ψ\_p − 2φ\_p·K = B with K := ψ\_p′/p and B a rational function regular at x(T), and at x = x(T), where ψ\_p vanishes, this determines K(x(T)) and not G(x(T)); differentiating once more determines K′(x(T)) and again not G(x(T)), because G enters only in the product p·G·K. So (ND) is a statement about φ\_p modulo p², one order beyond what the invariant differential sees. That extra order is exactly what Proposition 14 supplies — the p² digit of ψ\_p is a function of the p digit — and §4.13 turns it into a proof of (ND) for every p (Theorem ND). It is the same non-degeneracy that Propositions R⁗ and 10 use, so those are unconditional as well.

\begin{namedthm}[Remark (Theorem N is a statement about a p-minimal model)]
Scaling the model by x′ = u²x, y′ = u³y with p ∤ u multiplies φ\_p(x) by u\textasciicircum{}\{2p²\} by the weights of the division polynomials, so
\dispf{\textbf{Θ ↦ Θ + 2p·q\_p(u)  (mod p²)},}
while d′ is unchanged if gcd(u, d′) = 1. As n²h\_p(P) is intrinsic, the right-hand sides of Theorem N and Corollaries N⁽³⁾, N⁽⁴⁾ depend on the model, and those theorems are statements about a model minimal at p. This is a partial answer to the question of §7 about q\_p(d′): looking for an intrinsic meaning of q\_p(d′) by itself is the wrong request, since it is paired with the normalisation of the model — with the normalisation of σ. The third-digit term c/(p²x(P″)) is, by contrast, invariant under scaling as well.
\end{namedthm}

\textit{Theorem N* and Propositions 17 and 18.} Proved where they are stated (§3.11). Three remarks on how they stand to the four stages above. (i) Stages (2), (2′) and (4) — the Weierstrass equation, the failure of additivity of Q₂ on quotients, and the expansion of the logarithm — are, in the closed form, the single step log(φ/ω) = −(1/2)(log φ + log(1 + R)); the carries that an induction on digits has to track are the carries of the Iwasawa logarithm, and they are not tracked twice. (ii) Stage (1) and the cancellation of a₁ are, in the closed form, the statement that ℒ is even: the odd powers of z are exactly the terms that the model with p ∣ a₁ was chosen to suppress, and Proposition 17 suppresses them in every model. (iii) What the closed form does not contain is any claim about which terms a \textit{later} digit can produce; that is Proposition 18, whose proof is an inspection of the three series and not of the expansion, and which is why the caveat of §8 on the valuation table is weaker in v8 than the table itself would suggest.

\textit{Proposition 19.} T := χ(P)T₀ reduces to the same point of E(F\_p) as P′ = nP, by the definition of χ, so P′ − T ∈ E₁(Q\_p) and t₁ := t(P′ − T) ∈ pZ\_p. Since pT = O, P″ = pP′ = p(P′ − T). The formal logarithm is a homomorphism E₁(Q\_p) → Q\_p with log\_E([p](t)) = p·log\_E(t), whence z(P″) = p·log\_E(t₁). For t₁ ∈ pZ\_p one has v\_p(log\_E(t₁)) = v\_p(t₁), so e = v\_p(t(P″)) = v\_p(z(P″)) = 1 + v\_p(t₁) (Lean \texttt{layer\_\allowbreak{}depth}). ∎

\subsection*{4.13 Lemma 13, Propositions 14, 15, 16 and Theorem ND}

\textit{Lemma 13.} Write (∗) P′h − 2Ph′ = cΩ and (∗∗) u\_\{E′\}Ω² = P³ + a₄Ph⁴ + a₆h⁶; both come from x∘λ = P/h², y∘λ = y′Ω/h³ and λ\textasciicircum{}*ω\_E = c·ω\_\{E′\} by clearing denominators in ω\_E = dx/2y. At a root ξ of h, (∗) gives Ω(ξ) = −2P(ξ)h′(ξ)/c, and (∗∗) then gives c²P(ξ) = 4u\_\{E′\}(ξ)h′(ξ)². Differentiating (∗) and (∗∗) once and evaluating at ξ gives −P′(ξ)h′(ξ) − 2P(ξ)h″(ξ) = cΩ′(ξ) and u′Ω² + 2uΩΩ′ = 3P²P′ at ξ; eliminating Ω′ and dividing by P(ξ)² gives the second identity. No formal group and no Tate curve are needed. ∎

\textit{Proposition 14.} Both sides are reductions of elements of Z[a₄, a₆][x], so it suffices to prove the identity at every ordinary point, the ordinary locus being dense and open. Lift Ē to a curve E over a complete discrete valuation ring R with residue field F, v(p) = 1, ordinary. Let C := Ê[p] ⊂ E[p] be the connected component — the canonical subgroup, of order p. Its non-zero points have v(x) = −2/(p−1) while the remaining p² − p points of E[p] have v(x) ≥ 0, so ψ\_p = λ·M·N with M := p·(kernel polynomial of C) of degree n\_p, N of degree pn and λ ∈ R\textasciicircum{}×; the valuations of the roots put M ∈ R[x], and the leading coefficient of N is a unit, so the whole drop in degree from (p²−1)/2 to p(p−1)/2 in ψ\_p mod p comes from M. Hence

\begin{quote}
\textbf{M ≡ A₀ (mod p) for a constant A₀ ∈ R\textasciicircum{}×, so M′ ≡ 0 (mod p)}.
\end{quote}

Let π : E → E\textasciicircum{}\textit{ := E/C, with reduction F, and  : E\textasciicircum{}} → E with reduction V, so ∘π = [p] and c\_πc\_\{\} = p; choose the Weierstrass model of E\textasciicircum{}\textit{ whose reduction is Ē\textasciicircum{}\{(p)\}. Writing x∘π = /M² and y∘π = y/M³ with ,  ∈ R[x] we have (∗) for π, Z := x∘π ≡ x\textasciicircum{}p, /M³ ≡ u\textasciicircum{}\{n\_p\}, and c\_π ≡ 0 (mod p), say c\_π = pϖ. The roots of N are the x-coordinates of the points π(Q) with Q ∈ ker , so up to a unit N = M\textasciicircum{}\{2n\_p\}·h\textasciicircum{}}(Z) with h\textasciicircum{}\textit{ the kernel polynomial of , h\textasciicircum{}} ≡ h (mod p). Differentiating ψ\_p = λMN and using 2n\_p + 1 = p,

\begin{quote}
ψ\_p′ = λ(2n\_p+1)M\textasciicircum{}\{2n\_p\}M′h\textasciicircum{}\textit{(Z) + λM\textasciicircum{}\{2n\_p+1\}(h\textasciicircum{}})′(Z)·c\_π/M³,
\end{quote}

so that \textbf{K = ψ\_p′/p = λM\textasciicircum{}\{2n\_p\}M′h\textasciicircum{}\textit{(Z) + λϖM\textasciicircum{}\{2n\_p−2\}(h\textasciicircum{}})′(Z)}: the factor p of c\_π and the factor p of 2n\_p + 1 = p cancel the denominator. Reducing, the first term dies by M′ ≡ 0 and the second gives K ≡ κ·u\textasciicircum{}\{n\_p\}h′(x\textasciicircum{}p) with κ ∈ F\textasciicircum{}× independent of x. Comparing leading coefficients, lc\_x K = (p²−1)/2 ≡ −1/2 and lc\_x(u\textasciicircum{}nh′(x\textasciicircum{}p)) = n\_p ≡ −1/2, both in degree (p²−3)/2, so κ = 1. ∎ (In particular λϖA₀\textasciicircum{}p = 1, whence c\_\{\} ≡ A: the Hasse invariant is the action of V on the invariant differential, as it must be.)

\textit{Proposition 15, the reduction to one congruence.} By Proposition 14, 4uK² = 4w\textasciicircum{}\{(p)\}(x\textasciicircum{}p)h′(x\textasciicircum{}p)², so Proposition 15 is the identity A²Φ = 4w\textasciicircum{}\{(p)\}(h′)²A² /A² + A²D(Ψ²) in X, that is

\begin{quote}
Φ − 4w\textasciicircum{}\{(p)\}(h′)² = D(A²h²) = 2A(DA)h² + 2A²h·Dh.
\end{quote}

By Lemma 13 the left side vanishes at every root of h, and its degree is ≤ p − 1 because the leading coefficients cancel (4n\_p² = (p−1)² ≡ 1), so it equals h·C₁ with deg\_X C₁ ≤ n\_p. Since h is monic of degree n\_p, the displayed identity is equivalent to a congruence modulo h together with the coefficient of X\textasciicircum{}n\_p:

\begin{quote}
\textbf{(R′) [w\textasciicircum{}\{(p)\}(h′)²]′ + A²h′·Dh ≡ 0 (mod h)}  and  \textbf{(L) [X\textasciicircum{}\{p−1\}](Φ − 4w\textasciicircum{}\{(p)\}(h′)²) = 2A·DA},
\end{quote}

the first being the residue at each root — using C₁(ξ) = −2[w\textasciicircum{}\{(p)\}(h′)²]′(ξ)/h′(ξ) from Lemma 13 — and the second the top coefficient. In terms of the roots, (R′) says A²·Dξ\_i = (w\textasciicircum{}\{(p)\})′(ξ\_i) + 4w\textasciicircum{}\{(p)\}(ξ\_i)Σ\_\{j≠i\}(ξ\_i − ξ\_j)\textasciicircum{}\{−1\}: it is the Kodaira–Spencer derivative of the étale p-division points.

\textit{Proposition 15, the Tate family.} Each coefficient of the identity is weight-homogeneous in (a₄, a₆), and the kernel of the q-expansion map F\_p[E₄, E₆] → F\_p[[q]] is generated by A − 1, which is not homogeneous; so q-expansions are injective in each weight ([Ser73], [SwD73]) and it is enough to prove Proposition 15 for the Tate family. There A ≡ 1, Ē\textasciicircum{}\{(p)\} = E\_\{q\textasciicircum{}p\}, and the isogeny E\_\{q\textasciicircum{}p\} → E\_q induced by q\textasciicircum{}\{pZ\} ⊂ q\textasciicircum{}Z is V, its kernel being \{q\textasciicircum{}j\}. In the short model, with Q := q\textasciicircum{}p, X\_Q(W) = Σ\_\{m∈Z\}Q\textasciicircum{}mW(1−Q\textasciicircum{}mW)\textasciicircum{}\{−2\} − 2s₁(Q) + 1/12 and Y\_Q = y\_Q + x\_Q/2, the invariant differential is dX/2Y = dW/W, so W∂\_W X = 2Y and W∂\_W Y = u\_Q′(X); the roots are ξ\_i = X\_Q(q\textasciicircum{}i) and η\_i := Y\_Q(q\textasciicircum{}i), i = 1, …, n\_p. Three cancellations prove (R′).

(a) Since Q = q\textasciicircum{}p, the chain rule gives q dξ\_i/dq = pQ∂\_QX + i·W∂\_WX, and pQ = 0 in characteristic p:

\begin{quote}
\textbf{q dξ\_i/dq = 2i·η\_i},  hence  Dξ\_i = −q dξ\_i/dq = −2iη\_i.
\end{quote}

(b) Let θ\_Q(V) := (1−V)∏\_\{m≥1\}(1−Q\textasciicircum{}mV)(1−Q\textasciicircum{}m/V), so that x\_Q(W₁) − x\_Q(W₂) = −c₀²W₂·θ\_Q(W₁W₂)θ\_Q(W₁/W₂) · θ\_Q(W₁)\textasciicircum{}\{−2\}θ\_Q(W₂)\textasciicircum{}\{−2\} with c₀ = ∏(1−Q\textasciicircum{}m), the two sides having the same divisor and the same principal part at W₁ = 1. Substituting into H(W) := h(X\_Q(W)) = ∏\_i(X\_Q(W) − ξ\_i) and using 2n\_p + 1 = p,

\begin{quote}
H(W) = C·G(W)·θ\_Q(W)\textasciicircum{}\{−p\},  G(W) := ∏\_\{j=−n\_p\}\textasciicircum{}\{n\_p\}θ\_Q(Wq\textasciicircum{}j),
\end{quote}

and G is a theta function for q, by a rearrangement of the exponents alone: the factors (1 − q\textasciicircum{}aW) occurring in G are exactly those with a ≥ −n\_p, each once, because every integer ≥ −n\_p is uniquely pm + j with m ≥ 0 and
\begin{center}\begin{small}
\begin{tabular}{>{\raggedright\arraybackslash}p{0.056\linewidth}>{\raggedright\arraybackslash}p{0.916\linewidth}}
\hline
j & ≤ n\_p; the factors (1 − q\textasciicircum{}b/W) are exactly those with b ≥ n\_p + 1. Comparing with θ\_q(W), \\
\hline
\hline
\end{tabular}
\end{small}\end{center}

\begin{quote}
G(W) = (−1)\textasciicircum{}nq\textasciicircum{}\{−n\_p(n\_p+1)/2\}·W\textasciicircum{}n\_p·θ\_q(W),  hence  \textbf{H(W) = C′·W\textasciicircum{}n\_p·θ\_q(W)·θ\_Q(W)\textasciicircum{}\{−p\}}.
\end{quote}

(c) θ\_Q(W)\textasciicircum{}p is a p-th power, so its logarithmic derivative vanishes in characteristic p. With D := W d/dW,

\begin{quote}
\textbf{L := DH/H = n\_p + D log θ\_q(W)},  D log θ\_q(W) = −W/(1−W) + Σ\_\{m≥1\}[−q\textasciicircum{}mW/(1−q\textasciicircum{}mW) + q\textasciicircum{}m/(W − q\textasciicircum{}m)].
\end{quote}

At W = q\textasciicircum{}i(1 + s) the term m = i is exactly 1/s, and the value at s = 0 of the rest telescopes: writing q\textasciicircum{}m/(q\textasciicircum{}i−q\textasciicircum{}m) = −(1−q\textasciicircum{}\{i−m\})\textasciicircum{}\{−1\} for m \textless{} i and = q\textasciicircum{}\{m−i\}(1−q\textasciicircum{}\{m−i\})\textasciicircum{}\{−1\} for m \textgreater{} i, the sum equals Σ\_\{k=1\}\textasciicircum{}\{i−1\}(q\textasciicircum{}k−1)/(1−q\textasciicircum{}k) = −(i−1). So L = 1/s + κ + O(s) with \textbf{κ = n\_p + 1 − i}, and since D = (1+s)d/ds and H = h₁s + h₂s² + ⋯ one has κ = 1 + h₂/h₁ and D²H/DH|\_\{s=0\} = 1 + 2h₂/h₁ = 2κ − 1 = p − 2i ≡ −2i. Finally DH = 2Y·h′(X) and D²H = 4Y²h″(X) + 2h′(X)u\_Q′(X), so at W = q\textasciicircum{}i

\begin{quote}
η\_i·D²H/DH = (w\textasciicircum{}\{(p)\})′(ξ\_i) + 4w\textasciicircum{}\{(p)\}(ξ\_i)Σ\_\{j≠i\}(ξ\_i−ξ\_j)\textasciicircum{}\{−1\} = −2iη\_i = Dξ\_i = A²Dξ\_i,
\end{quote}

which is (R′). (Replacing i by p − i changes η\_i to −η\_i and i to −i and leaves both sides unchanged.)

For (L), expand the identity divided by A²h²: Φ/h² = 4w\textasciicircum{}\{(p)\}(h′/h)² + 2A·DA + 2A²Dh/h. As X → ∞ the coefficient of X matches because 4n\_p² ≡ 1, and the constant term is (L) in the form φ₁ + 6σ = 2A·DA with φ₁ := [X\textasciicircum{}\{p−1\}]Φ and σ := Σ\_iξ\_i. In the Tate family Φ/h² is the x-coordinate of V, so it equals X\_q(W) at X = X\_Q(W), and lim\_\{W→1\}(X\_q − X\_Q) = 0 because the m = 0 terms cancel and Σ\_\{m≠0\}c\textasciicircum{}m(1−c\textasciicircum{}m)\textasciicircum{}\{−2\} = 2s₁(c); hence φ₁ = −2σ. Summing over the whole kernel, Σ\_\{j=1\}\textasciicircum{}\{p−1\}x\_Q(q\textasciicircum{}j) = 2s₁(q) − 2p·s₁(Q) ≡ 2s₁(q), so σ = n\_p/12 + s₁(q) = −1/24 + s₁(q) = −E₂/24; and 2A·DA ≡ −E₂/6 because DA = −(θA − ((p−1)/12)E₂A) ≡ −E₂/12 on q-expansions. So φ₁ + 6σ = 4σ = −E₂/6 = 2A·DA. This proves (L), hence Proposition 15. ∎

\textit{Theorem ND.} By Proposition 15, S₁ mod p = 2(DΨ)(x\textasciicircum{}p), and Res\_x(f(x\textasciicircum{}p), g(x\textasciicircum{}p)) = Res\_X(f, g)\textasciicircum{}p, so Res\_x(S₁, ψ\_p) = (2\textasciicircum{}n\_p·Res\_X(DΨ, Ψ))\textasciicircum{}p. Writing Ψ = A∏(X−ξ\_i) and applying D coefficientwise, DΨ(ξ\_i) = −A·(Dξ\_i)·∏\_\{j≠i\}(ξ\_i−ξ\_j), so with the sign of the resultant convention

\begin{quote}
Res\_X(DΨ, Ψ) = (−1)\textasciicircum{}\{n\_p(n\_p−1)/2\}·A\textasciicircum{}\{2n\_p\}·∏\_i(Dξ\_i)·∏\_\{i\textless{}j\}(ξ\_i−ξ\_j)².
\end{quote}

In the Tate family ord\_q A = 0, ord\_q(Dξ\_i) = i and ord\_q(ξ\_i − ξ\_j) = min(i, j), so ord\_q of the resultant is Σ\_i i + 2Σ\_\{i\textless{}j\}min(i, j) = n\_p(n\_p+1)/2 + n\_p(n\_p²−1)/3 = M = p(p²−1)/24, which is exactly 1/12 of its weight n\_p(p²+1) + n\_p(p²−1) − 2pn² = 12M. By §4.6(3) the only irreducible weight-homogeneous factor vanishing at the cusp is Δ₀, so Res\_X(DΨ, Ψ) = γ\_pΔ₀\textasciicircum{}\{M\} with γ\_p ∈ F\_p\textasciicircum{}× and c\_S = (2\textasciicircum{}n\_pγ\_p)\textasciicircum{}p = 2\textasciicircum{}n\_pγ\_p ≠ 0. ∎

\textit{Proposition 16.} The leading coefficients in the last computation are ∏\_i(Dξ\_i) → (−1)\textasciicircum{}n\_p n\_p! (by (a) above, Dξ\_i = −2iη\_i = −iq\textasciicircum{}i(1 + ⋯)), ∏\_\{i\textless{}j\}(ξ\_i−ξ\_j)² → 1 and Δ₀ → −q/16, which gives γ\_p and c\_S. For c\_K, Proposition 14 factors the resultant:

\begin{quote}
Res\_x(K, ψ\_p) = Res\_x(u, ψ\_p)\textasciicircum{}n\_p·(A\textasciicircum{}\{n\_p−1\}Res\_X(h′, h))\textasciicircum{}p,  Res\_x(u, ψ\_p) = (−1)\textasciicircum{}nA³∏\_iw\textasciicircum{}\{(p)\}(ξ\_i) = (−1)\textasciicircum{}nA³∏\_iη\_i²,  Res\_X(h′, h) = ∏\_ih′(ξ\_i),
\end{quote}

and in the Tate family η\_i = q\textasciicircum{}i/2 + ⋯ and ∏\_ih′(ξ\_i) = (−1)\textasciicircum{}\{n\_p(n\_p−1)/2\}q\textasciicircum{}\{n\_p(n\_p²−1)/3\}(1 + ⋯). Hence the q-order is n\_p·n\_p(n\_p+1) + p·n\_p(n\_p²−1)/3 = m\_K — again 1/12 of the weight (p²−1)(p²−3)/2 — and

\begin{quote}
c\_K = (−1)\textasciicircum{}\{n\_p + n\_p(n\_p−1)/2 + m\_K\}·2\textasciicircum{}\{4m\_K − 2n\_p²\}.
\end{quote}

Now m\_K is even: with j := (p²−1)/24 ∈ Z for p ≥ 5 one has m\_K = (p²−1)(p²−3)/24 = j(24j−2) = 2j(12j−1). (No divisibility of j is needed, and none holds — j = 1 at p = 5.) So the sign is (−1)\textasciicircum{}\{n\_p(n\_p+1)/2\} = (−1)\textasciicircum{}\{(p²−1)/8\} = (2|p). For the exponent, 2n\_p² = (p−1)·n\_p ≡ 0 (mod p−1), while 3m\_K = (p−1)·(p+1)(p²−3)/8 and (p+1)(p²−3)/8 = ab/2 with a = (p+1)/2 and b = (p²−3)/2 odd: for p ≡ 3 (mod 4) the factor a is even and 3m\_K ≡ 0 (mod p−1), so 2\textasciicircum{}\{3m\_K\} = 1; for p ≡ 1 (mod 4) both are odd and 3m\_K ≡ (p−1)/2 (mod p−1), so 2\textasciicircum{}\{3m\_K\} = (2|p). Since 2\textasciicircum{}\{4m\_K−2n\_p²\} = 2\textasciicircum{}\{m\_K\}·2\textasciicircum{}\{3m\_K−2n\_p²\}, this gives c\_K = ε\_p2\textasciicircum{}\{m\_K\} with ε\_p = (2|p) for p ≡ 3 (mod 4) and ε\_p = (2|p)² = 1 for p ≡ 1 (mod 4), that is ε\_p = −1 exactly for p ≡ 3 (mod 8). The formula for c\_p is §4.11. ∎

\section*{5. Computations}

All computations use PARI/GP 2.17.3 with \texttt{elldata} (Cremona's tables) on one core, with positive and negative controls; none is a proof, and no number is quoted outside its stated range.

\textbf{The convention, made explicit.} \texttt{ellpadicheight(E,p,n,P)} returns the coordinates [f, g] of the height on the basis (ω, η = xω) of H¹\_dR ⊗ Q\_p; the height of §2 is

\begin{quote}
h\_p(P) = −( f − s₂·g ),  s₂ := \texttt{ellpadics2(E,p,n)} = b₂/12 − E₂(E, ω)/12,
\end{quote}

since s₂ is precisely the constant c of [Har, (2)] that singles out σ\_p: \textbf{the first component alone is a different height}, the one belonging to c = 0, whose σ is not in tZ\_p[[t]]. As a control this reproduces two entries of the tables of [MST], v = −1 for 37A at p = 53 (§4.2) and v = +1 for 5077A at p = 5 (§4.4), their −2 shifted by r. Every count below was recomputed with h\_p, the two conventions differing by a correction of valuation ≥ 0 that never moves a negative valuation.

\subsection*{5.1 The decomposition}

Rank one, conductor ≤ 900, a\_p = 1, p ≤ 43, p ∤ ∏c\_ℓ, ι, precision 14 — \textbf{474 boxes} — with P = n\_QQ, n\_Q := lcm(ord(), ∏\_ℓ c\_ℓ), b := v\_p(h\_p(P)/2 − log\_p(α/β)), c := v\_p(log\_p(α/β)), e := v\_p(d(P)). Lemma 3 (c ≥ 1; b ≥ e; b ≥ 2e in the 223 boxes with p ∣ a₁) and Corollary E (in the 45 split boxes) held without exception, and v\_p(h\_p(P)) = min(b, c) in 470 boxes, the other 4 having b = c with cancelling leading terms. Among the 429 non-split boxes c ≥ 2 occurred 26 times, all of them in Cremona's model with a₁ odd (§5.5). The choice of σ moves b in 98 boxes and v\_p(h\_p) in none.

\subsection*{5.2 λ = 0 ⟺ split}

λ is measured with no rational point in sight: by Lemma 5 it is a function of E mod p² and , and by Lemma 4′ splitting is v\_p(ψ\_p(x₀)) ≥ 2 for an integral lift. Over all curves of conductor N ≤ 1200 (any rank), 5 ≤ p ≤ 43, a\_p = 1 — \textbf{1,445 boxes}, controls deg φ\_p = p² and leading coefficient 1 — these held with \textbf{no exception}: the scaling λ(m ) = m² λ() for m = 2, …, p−1 (Theorem H); ψ\_p ≡ u·\textasciicircum{}p (mod p) (Lemma 4′); the invariance of min(v\_p(ψ\_p(x₀)), 2) and of λ under the lifts x₀ + kp, k ≤ 3; and \textbf{λ = 0 ⟺ split} (158 split boxes, 158 with λ = 0, the same ones), which is Proposition 6. The identity (★) was checked as an exact rational identity in 403 boxes and Theorem B′ in 310, both without exception, and splitting read independently, by \texttt{polrootspadic} on ψ\_p, agreed wherever the PARI stack allowed it (1,234 boxes). One limit: the \textit{exact} value of v\_p(ψ\_p(x₀)) is not invariant under change of lift, failing in 85 boxes — which is why only the cut at 2 is used, and is all that Lemma 4′ proves.

Proposition 6 was also tested where it is stated, over F\_p rather than over the tables: running (A, B) over F\_p² for the anomalous curves with 5 ≤ p ≤ 157 — 4,824 curves, 260,258 classes  — the slope ρ of §4.6 was non-zero without exception, the family E\_s contained exactly one split member, and the zero of λ\_s sat at that member. The identity Res\_X(ũ, g) = c\_ũ·Δ\textasciicircum{}\{m\_ũ\} was verified exactly for p = 17, 19, 23 (m\_ũ = 400, 561, 1001) at more points than the degree bound, and the cusp estimates — the valuations of the roots of f and g, ord\_t Res\_X(ũ, g) = m\_ũ, and the ratio of leading coefficients predicted by 12ρ\_j = j²/t — held for p ≤ 29 without exception.

\subsection*{5.3 The height matrix, the regulator, and a bound that is too strong}

Conductor ≤ 1500, p ≤ 19, a\_p = 1, p ∤ ∏c\_ℓ, ι, ranks 1 and 2 — \textbf{659 boxes}: the minimum valuation of the entries of M\_Λ is 0, 1 or 2 in the 71 split boxes and −1 in every one of the 588 non-split boxes, which is part (i) of the main theorem in both directions.

Rank 1, conductor ≤ 2000, p ≤ 23, a\_p = 1 — \textbf{1,107 boxes}: in the 8 where ι fails, v\_p(Reg\_p) ∈ \{1,2\} ≥ r as Theorem A requires; of the other 1,099, every one of the 976 non-split boxes has v\_p(Reg\_p) = −1 exactly and every one of the 123 split boxes has v\_p(Reg\_p) ∈ \{0,1,2\}. Part (ii) accounts for the first; the distribution inside \{0,1,2\} is not accounted for by anything here.

One box beyond the reach of these sweeps is \textbf{a box of ours}. The rank-four curve of [MST, §4.5] — y² + xy = x³ − x² − 79x + 289, conductor 234446, at p = 5 — we recomputed in full: a\_5 = −4, so \textbf{\#E(F\_5) = 10 = 2p and E(F\_5) ≅ Z/10}, which puts it outside the hypothesis \#E(F\_p) = p of the main theorem and inside Proposition R⁺ (§3.9), the case N = 2p that Hasse allows for p = 5 alone. In our normalisation v\_5(Reg\_5) = 2 = r − 2, so \textbf{the baseline is attained at rank four}, which no sweep here reaches (the sweeps below stop at rank 3). What we measured: ∏\_ℓ c\_ℓ = 4, so p ∤ ∏\_ℓ c\_ℓ; the discrete logarithms of a basis of Λ in E(F\_5)[5] are (3, 5, 6, 8) with gcd 1, so \textbf{ι holds and χ is onto}; v\_5(ψ\_5(x₀)) = 1 for all four points of order 5 and for all five lifts of each, so the curve is \textbf{non-split}; the valuations of the entries of M\_Λ are

\begin{quote}
[[−1, 0, −1, −1], [0, 1, 0, 0], [−1, 0, −1, −1], [−1, 0, −1, −1]],
\end{quote}

the minimum being −1 as part (i) of the main theorem requires on the non-split side (here through Theorem S as extended by Proposition R⁺); with χ = (4, 0, 3, 4) the congruence p·M\_Λ ≡ c·χχᵗ (mod p) holds in all sixteen entries with c = 3 = 2λ₁, λ₁ = 4, which is Theorem R and Proposition R″ \textbf{verified at rank four}; and det T is a unit, so Corollary R′ gives the value 2 rather than a bound. Sweeping the same curve over p ≤ 120 produced no other anomalous prime, so this curve does not supply a rank-four box with \#E(F\_p) = p.

Rank ≥ 2, conductor ≤ 9000, p ≤ 19 — \textbf{261 boxes}: 23 split boxes of rank 2 with v\_p(Reg\_p) ∈ \{0,1\}, so the bound 0 is attained. The stronger bound max(2r−4, 0) suggested by the shape of Theorem J is \textbf{refuted} by a split box of rank 3, where v\_p(Reg\_p) = 1 = max(r−2,0) \textless{} 2 and where in two boxes the height form does not vanish on the kernel of the reduction character, as that bound would require. Those boxes are not in this range: this range contains no anomalous box of rank ≥ 3 at all, and widening it to conductor 20000 and p ≤ 23 produces two, 16811a1 at p = 23 and 18562c1 at p = 19, both non-split. The split boxes of rank ≥ 3 first occur at conductor 105643, at p = 19, and they are the 48 boxes of sweep F in §5.6.

\subsection*{5.4 CM curves}

Over all curves of conductor ≤ 1500 (8,163 curves, 138 with CM) and p ≤ 43, all \textbf{14} anomalous CM boxes are split, against 178 of 1,935 in the non-CM boxes; the primes occurring are p = 7, 19, 37 for curves with j = 0, which is what a\_p = 1 with π = p forces. This is the content of Theorem F, not independent evidence for it.

\subsection*{5.5 The model, and the exact value in the non-split case}

Conductor ≤ 3000, p ≤ 23, rank ≥ 1, a\_p = 1, p ∤ ∏c\_ℓ, ι, with P′ = (∏c\_ℓ)·Q for a generator Q of non-zero reduction and P = pP′ computed exactly — \textbf{1,765 boxes}, in both Cremona's model and the model with p ∣ a₁. In the model with p ∣ a₁ the identity q\_p(α/β) = −λ of Theorem B held in \textbf{1,765 of 1,765}, and \textbf{no non-split box had c ≥ 2}; the positive control, 185 split boxes with λ = 0 and c ≥ 2, held throughout. In Cremona's model the same identity failed in 719 boxes, all of them with e = 1 and a₁ odd, and 90 non-split boxes had c ≥ 2, all of them with a₁ odd; the shift of §2 removes all 90. This is the numerical form of the remark that b and c are quantities of a model while min(b, c) is not.

The relation d(P) = d′·|Ψ| — no cancellation, Step 1 of §4.4 — held in every box, and fails when the factor ∏\_ℓ c\_ℓ is dropped from P′: (A2) is used at exactly that point. On the side of the height matrix, conductor ≤ 900 and p ≤ 13 with \texttt{ellpadicheightmatrix} gave \textbf{239 boxes}: min\_\{i,j\} v\_p(M\_Λ[i,j]) = −1 in all 206 non-split boxes and ≥ 0 in all 33 split ones.

\subsection*{5.6 The height matrix modulo p, on both sides}

Same convention h\_p = −(f − s₂g), precision 10, \texttt{elldata}, splitting read locally as v\_p(ψ\_p(x₀)) ≥ 2 (Lemma 4′), χ\_i = elllog of the reduction of the i-th generator, boxes with all χ\_i = 0 (ι failing) discarded.

\textit{Non-split, Theorem R and Proposition R″.}

\begin{center}\begin{small}
\begin{tabular}{>{\raggedright\arraybackslash}p{0.063\linewidth}>{\raggedright\arraybackslash}p{0.215\linewidth}>{\raggedright\arraybackslash}p{0.063\linewidth}>{\raggedright\arraybackslash}p{0.114\linewidth}>{\raggedright\arraybackslash}p{0.063\linewidth}>{\raggedright\arraybackslash}p{0.127\linewidth}>{\raggedright\arraybackslash}p{0.14\linewidth}>{\raggedright\arraybackslash}p{0.102\linewidth}}
\hline
sweep & range & boxes & non-split & split & congruence fails & rank ≠ 1 (non-split) & p·M not integral \\
\hline
A & rank ≥ 2, conductor ≤ 12000, p ≤ 19 & 460 & 414 & 46 & \textbf{0} & \textbf{0} & \textbf{0} \\
B & rank ≥ 3, conductor ≤ 20000, p ≤ 43 & 4 & 4 & 0 & \textbf{0} & \textbf{0} & \textbf{0} \\
D & rank ≥ 1, conductor ≤ 12000, p ≤ 19 & 8,811 & 7,759 & 1,052 & \textbf{0} & \textbf{0} & \textbf{0} \\
\hline
\end{tabular}
\end{small}\end{center}

The negative control is essential and it holds: over the 46 split boxes of sweep A the rank of p·M\_Λ mod p is \textbf{0 in all 46}, never 1, so "rank exactly one" separates the two cases rather than being an identity. In the 4 boxes where one generator falls into Λ₀ the predicted A\_\{ij\} ∈ p²Z\_p held. Over the 7,759 non-split boxes of sweep D the distribution of v\_p(Reg\_p) − (r − 2) is 7,723 / 28 / 8 / 0 for the values 0 / 1 / 2 / ≥ 3, with \textbf{no negative value}: Corollary R′ predicts that the excess is v\_p(det T) ≥ 0 and that it is usually 0. Proposition R″ was checked as the exact ratio c/λ₁ = 2 in \textbf{7,759 of 7,759}, λ₁ being computed from φ\_p(x₀) mod p² alone with no rational point in sight.

\textit{Split, Theorem R\_s and Propositions R″\_s.}

\begin{center}\begin{small}
\begin{tabular}{>{\raggedright\arraybackslash}p{0.071\linewidth}>{\raggedright\arraybackslash}p{0.21\linewidth}>{\raggedright\arraybackslash}p{0.071\linewidth}>{\raggedright\arraybackslash}p{0.114\linewidth}>{\raggedright\arraybackslash}p{0.099\linewidth}>{\raggedright\arraybackslash}p{0.053\linewidth}>{\raggedright\arraybackslash}p{0.128\linewidth}>{\raggedright\arraybackslash}p{0.142\linewidth}}
\hline
sweep & range & split boxes & M not integral &  nonzero on Λ₀ & ρ₁ \textgreater{} 2 & ρ₁ ≠ predicted from (a, b) & v\_p(Reg\_p) \textless{} r − ρ₁ \\
\hline
E & rank ≥ 1, conductor ≤ 12000, p ≤ 19 & 1,052 & \textbf{0} & \textbf{0} & \textbf{0} & \textbf{0} & \textbf{0} \\
F & rank ≥ 3, conductor ≤ 400000, p ≤ 19 & 48 & \textbf{0} & \textbf{0} & \textbf{0} & \textbf{0} & \textbf{0} \\
\hline
\end{tabular}
\end{small}\end{center}

Sweep E is the same 1,052 boxes as the split column of sweep D, now measured without the factor p; its rank distribution is ρ₁ = 0 in 135 boxes, 1 in 872, 2 in 45, and 3 or more in none. Split boxes of rank 2 are 46: 45 with ρ₁ = 2 and v\_p(Reg\_p) = 0, one with ρ₁ = 1 and v\_p(Reg\_p) = 1, matching the equivalence v\_p(Reg\_p) = 0 ⟺ ρ₁ = 2 of Corollary R′\_s. Sweep F exists because ρ₁ ≤ 2 is vacuous for r ≤ 2; all \textbf{48} split boxes of rank 3 have ρ₁ = 2 exactly, with v\_p(Reg\_p) = 1 = r − 2 in 43 of them and 2 in the other 5. The distribution of v\_p(Reg\_p) over the 1,006 split boxes of rank one is 871 / 114 / 17 / 4 for the values 0 / 1 / 2 / 3.

The negative control here is a synthetic one, since ρ₁ ≤ 2 would be empty if random matrices satisfied it: over 4,000 symmetric matrices of sizes 3 to 5 with entries drawn uniformly modulo p³, p ∈ \{5, …, 19\}, the rank modulo p was 3 or more in 3,844 and exactly 2 in 151, and the form vanished on the coordinate hyperplane in \textbf{1}. So neither "rank ≤ 2" nor "vanishing on the hyperplane" is a property of symmetric matrices in general, and 48 of 48 at rank 3 is not an artefact of the sample.

Proposition R″\_s was checked as a = 2λ\textasciicircum{}\{(2)\}/n² in \textbf{1,052 of 1,052} boxes of sweep E, with the positive control λ = 0 holding in all 1,052 (as Corollary D requires in the split case) and with λ\textasciicircum{}\{(2)\} computed from φ\_p(x(n e₁)) mod p³ and the Tamagawa product alone. For rank one the equivalence a ≠ 0 ⟺ v\_p(Reg\_p) = 0 held in \textbf{1,006 of 1,006}. Finally the lift dependence: replacing x₀ by x₀ + kp for k = 1, 2, 3 moved λ\textasciicircum{}\{(2)\} in \textbf{1,052 of 1,052} boxes while, as the control, it moved λ in \textbf{0} — which is Lemma 5 on one side and its failure one digit deeper on the other.

\textit{The isogeny formula.} Over the curves of conductor ≤ 12000 with 5 ≤ p ≤ 19, anomalous, non-split, and every isogenous partner of degree prime to p — \textbf{2,098 pairs} — the identity (λ₁(E)/p)·(λ₁(E″)/p) = (m/p) of Corollary R‴ held \textbf{without exception}, with λ₁ computed independently on the two minimal models from φ\_p(x₀) mod p² alone. The two obvious alternatives fail: the product of the symbols equals 1 in 696 pairs (exactly those with m a square modulo p, where the two readings agree) and equals (−m/p) in 660. Splitting agreed on the two curves in all 2,098 pairs, and λ₁ was nonzero on both in all of them, as Proposition 6 requires in the non-split case.

\textit{One earlier observation withdrawn.} Sweep D also measured whether v\_p(s₂·g) \textgreater{} v\_p(f), the inequality under which the first component of \texttt{ellpadicheightmatrix} alone would have the same valuation as h\_p. It \textbf{fails in 15 of the 8,811 boxes} (869a1, 1318b1, 2678k1, 5484b1, 6576f1, 8162i1, 8962b1, 9109a1, 9559d1, 9861c1, 11557a1 at p = 5; 10176c1, 11267a1 at p = 7; 4384b1, 10950m1 at p = 17). The first component is a different height (§5) and not merely a different normalisation of the same one: with it in place of h\_p the ratio of Proposition R″ comes out as −2 rather than 2 in every one of 1,423 boxes of conductor ≤ 3000. Statements about valuations taken with the first component alone are therefore correct up to those 15 boxes and wrong in sign for c.

\subsection*{5.7 The closed form modulo p, and the second digit}

All sweeps here are over curves of good ordinary reduction with \#E(F\_p) = p, 5 ≤ p ≤ 19, p ∤ ∏\_ℓ c\_ℓ, ι, split, and with M\_Λ integral; the ranges differ and are stated each time.

\textit{Lemma 7.} Over two disjoint sweeps — rank exactly 2 to conductor 60000 (617 boxes) and rank ≥ 3 to conductor 400000 (48 boxes) — the identity κ\_j = c₁·γ(n e\_j) held in \textbf{713 of 713} pairs, and κ\_j ≠ 0 ⟺ v\_p(t(n e\_j)) = 1 in all 713. The depth distribution was 689 pairs at v\_p(t) = 1, 23 at 2 and 1 at 3, so the negative side is not vacuous. Controls: the wrong constant x₀ in place of c₁ holds in 125 of 713 by coincidence, and forgetting the factor n breaks 579 of the 600 pairs with n ≢ 1 (mod p).

\textit{Proposition 8.} Purely local, with no global point and no height: over the \textbf{372} split anomalous boxes of conductor ≤ 3000, with all p − 1 classes in each box (\textbf{2,876} classes), (L1) held in 2,876 of 2,876 and (L2) in 2,876 of 2,876. Controls: "s/χ constant" holds in 689 of 2,876 and "s·/χ constant" in 575, of which 372 each are the base class itself, so 317 and 203 are coincidences; positive control, the non-split constancy of λ/χ² (Theorem H), held in \textbf{2,613 of 2,613} classes.

\textit{Theorem K.} Over rank ≥ 1 and conductor ≤ 12000 — 1,051 boxes, 1,143 entries — and over rank ≥ 3 and conductor ≤ 400000 — 48 boxes, \textbf{all six entries} of each Gram matrix, 288 entries — the closed form held with \textbf{no exception}. Controls: replacing γ by a constant, and dropping the B₀ term, each hold in 0 of the 48 rank-3 boxes (and in 111 of the 1,051 boxes of the first sweep, where the rank is usually 1 and the two forms agree). Corollary K′(iv) was checked on the 48 rank-3 boxes: the equivalence v\_p(Reg\_p) = r − 2 ⟺ det Ĉ ≢ 0 held in \textbf{48 of 48}, while the exact formula v\_p(Reg\_p) = (r − 2) + v\_p(det Ĉ) fails in one of them (310329c1 at p = 5: v\_p(det M\_Λ) = 2 against 3).

\textit{Theorem N.} Over conductor ≤ 15720 — \textbf{1,386 boxes and 1,457 entries} with χ ≠ 0, of which 1,313 boxes of rank one, contributing one entry each, and 73 of rank two, contributing 144 — the congruence n²h\_p(P) ≡ (1 + p)Θ(x(nP)) + 2p·q\_p(d′) (mod p²) held with \textbf{no exception}. The distribution of e = v\_p(d(pnP)) was 1,294 / 141 / 20 / 2 for e = 2, 3, 4, 5. The controls are informative:

\begin{center}\begin{small}
\begin{tabular}{>{\raggedright\arraybackslash}p{0.565\linewidth}>{\raggedright\arraybackslash}p{0.407\linewidth}}
\hline
control & holds in \\
\hline
drop the term 2p·q\_p(d′) & 377 of 1,457 — exactly the entries with q\_p(d′) = 0 \\
drop the factor (1 + p), i.e. the term p·Θ & 199 of 1,457 \\
q\_p(n) in place of q\_p(d′) & 310 of 1,457 \\
q\_p(d′ + p) in place of q\_p(d′) & \textbf{0 of 1,457} \\
\hline
\end{tabular}
\end{small}\end{center}

\textit{The extension that fails.} The congruence \textbf{without} the denominator term, n²h\_p(P) ≡ Θ(x(nP)) (mod p²), holds in \textbf{15 of 113} boxes of conductor ≤ 2000 — the frequency of a coincidence, about 1/p — and the equality of valuations v\_p(n²h\_p(P)) = v\_p(Θ(x(nP))) fails in 3 of the 10 boxes there with v\_p ≥ 1 (603e1 and 869a1 at p = 5, 1520h1 at p = 7). So the second digit of the height is not the second digit of the Fermat quotient, and the local model of Propositions 9 and 10 governs Θ and not h\_p.

\textit{Corollary N′.} Among the rank-one boxes of conductor ≤ 15720 there are \textbf{184} with v\_p(Reg\_p) ≥ 1, of which \textbf{30} have v\_p(Reg\_p) ≥ 2. The criterion v\_p(Reg\_p) ≥ 2 ⟺ Θ/p + 2q\_p(d′) ≡ 0 held in \textbf{184 of 184}; the control that ignores the q\_p(d′) term held in 149, so that term decides \textbf{35 of the 184} boxes. The frequencies of v\_p ≥ 1 by prime were 87, 61, 9, 13, 6, 8 at p = 5, 7, 11, 13, 17, 19.

\textit{Propositions 9 and 10.} Over the same 372 local boxes, with every class and 2p + 1 lifts per class — \textbf{67,444} pairs — the second-digit expansion of Proposition 9 held in 67,444 of 67,444; the control that drops the k² term differs from the truth in 55,040 of them, so the term is not decoration. The bad lift δ\_\textit{ was computed digit by digit and verified to p³ in \textbf{372 of 372} boxes, and v\_p(Θ(x₀)) = v\_p(δ − δ\_}) held in \textbf{113 of 113} rank-one boxes of conductor ≤ 2000.

\textit{λ\_T one digit deeper.} Modulo p² the identity Θ\_T(m) = m²Θ\_T(1) with the standard lift of m holds in 282 of 2,504 class pairs, and with the Teichmüller lift m\textasciicircum{}p in 580 of 2,504. Both frequencies are those of coincidences; no law is claimed (§7(7)).

\subsection*{5.8 The third and fourth digits, the locality of W, and the sensitivity}

All sweeps here are over curves of good ordinary reduction with \#E(F\_p) = p, 5 ≤ p ≤ 19, p ∤ ∏\_ℓ c\_ℓ, ι, and χ(P) ≠ 0; ranges are stated each time. Everything was calibrated before it was run, and the calibrations are reported first because two of them are the arithmetic backbone of §4.12.

\textit{Calibration.} (i) φ\_p(x)/ψ\_p(x)² = x(pP) with deg φ\_p = p² and integral coefficients: \textbf{36 of 36} (twelve curves of conductor ≤ 600 × p = 5, 7, 11). (ii) p ∣ φ\_p′ and p ∣ ψ\_p′ coefficientwise, the two halves of the shape of §2: \textbf{20 of 20} each (five curves × p = 5, 7, 11, 13); negative control, p ∣ f\_\{p−1\}′ for the even-index division polynomial, \textbf{0 of 20}. (iii) the two statements of §4.12 about the first coefficient — X′(0) = 2y(T) + a₁x(T) + a₃ and its being a unit: \textbf{129 of 129} boxes of conductor ≤ 3000, and (iv) v\_p(ψ\_p′(x(T))) = 1 exactly: \textbf{129 of 129} in the same boxes. (v) the σ-normalisation: \texttt{ellpadics2(ellinit([17,42]),13,4)} reproduces the one worked example of PARI's documentation. (vi) E₆(E, ω) = −c₆ against the Hasse invariant, \textbf{250 of 250} (conductor ≤ 12000, p ≤ 19); with E₆ = +c₆ only 150 of 250, and the three smallest cases were checked by hand (p = 5: c₄ ≡ 2A and the Hasse invariant 2A; p = 7: −c₆ ≡ 3B and 3B; p = 11: c₄(−c₆) ≡ 9AB and 20AB ≡ 9AB).

\textit{E₂ modulo p is local (§2).} Over conductor ≤ 12000 and p ≤ 19, the congruence E\_\{p+1\}(c₄, −c₆)/E\_\{p−1\}(c₄, −c₆) ≡ b₂ − 12·\texttt{ellpadics2}(E, p) (mod p) held in \textbf{315,953 of 315,953} boxes. The classical forms were computed from their q-expansions and re-expressed in E₄, E₆ by solving for the isobaric coefficients, with the denominators checked prime to p in every case. Controls: E₆ ↦ +c₆ gives 54,109; dropping b₂ gives 65,189; E\_\{p+1\} alone, without dividing by E\_\{p−1\}, gives 67,566; junk gives 34,020 — the rate of a coincidence at these p. Negative control: over the \textbf{58,785} supersingular boxes in the same range, E\_\{p−1\}(c₄, −c₆) ≡ 0 in all of them and \texttt{ellpadics2} returned a p-adic number in \textbf{none}, so the two sides agree also about where the quantity is undefined. Finally, over the 315,121 pairs of boxes with the same (p, c₄ mod p, c₆ mod p), the two values of E₂ mod p agreed \textbf{without exception}, which is the statement of the Remark of §2 read directly.

\textit{Corollaries N⁽³⁾ and N⁽⁴⁾.} Over conductor ≤ 20000 and rank 1 — \textbf{1,678 boxes}, with e = v\_p(t(P″)) distributed 1,482 / 164 / 30 / 2 over e = 2, 3, 4, 5 — both congruences held with \textbf{no exception}, Corollary N⁽³⁾ modulo p³ and Corollary N⁽⁴⁾ modulo p⁴. The controls are what make the coefficients load-bearing:

\begin{center}\begin{small}
\begin{tabular}{>{\raggedright\arraybackslash}p{0.61\linewidth}>{\raggedright\arraybackslash}p{0.174\linewidth}>{\raggedright\arraybackslash}p{0.174\linewidth}}
\hline
control (same 1,678 boxes) & N⁽³⁾ holds in & N⁽⁴⁾ holds in \\
\hline
drop the term c/(p²x(P″)) & 378 & 284 \\
change its sign to + & 378 & — \\
coefficient of Θ² equal to +1/4 (the naive guess) & 229 & — \\
c replaced by b₂/12, i.e. E₂ = 0 & 377 & — \\
c replaced by a number built from the conductor & 409 & — \\
drop the Θ² term & 246 & — \\
drop the F³ term & — & 410 \\
F² without the factor (1 + p) & — & 410 \\
Θ² without the factor (1 + p) & — & 246 \\
Corollary N⁽³⁾ read modulo p⁴ & — & 280 \\
add p³Θ³ & — & 246 \\
add p³ΘF & — & 582 \\
add p³ × junk & — & 292 \\
\hline
\end{tabular}
\end{small}\end{center}

The numbers in the high 200s and 300s are the counts of boxes where the touched term vanishes for a reason (e ≥ 3 makes c/(p²x(P″)) ≡ 0 modulo p, which is 196 boxes, and q\_p(d′) = 0 accounts for most of the rest); the numbers near 410 are two controls that both remove one term of the p³ digit and therefore fail on the same boxes. \textbf{Every control fails.}

\textit{Corollaries N″ and N‴.} Among the same boxes there are \textbf{41} with v\_p(Reg\_p) ≥ 2, of which \textbf{9} have v\_p ≥ 3 and \textbf{1} has v\_p ≥ 4 (19779a1 at p = 5). Corollary N″ decided all \textbf{41 of 41} correctly in both directions; the control that drops W decided 30, so W changes the verdict in \textbf{11} of the 41. Corollary N‴ was correct in \textbf{9 of 9}, the controls dropping U or F³ each in 8 of 9 — with a single positive box that is the absence of a violation and not a test of the criterion (§7(11)).

\textit{Proposition 12 and the locality of W.} Over conductor ≤ 20000 and rank 1, with γ computed independently through the group law in E(Q\_p) — \textbf{1,676 boxes} — the congruence W ≡ c·γ² (mod p) held in \textbf{1,676 of 1,676}, and e = 2 ⟺ γ ≢ 0 in 1,676 of 1,676. Controls: c·γ in place of c·γ² holds in 579, W ≡ 0 in 378, γ² without c in 401. Proposition 11 held in \textbf{1,676 of 1,676}; dropping the ν (t³) term gives 351 and dropping γ₃ gives 237.

\textit{Theorem G′, the three steps separately.} Over \textbf{400} boxes of conductor ≤ 5400 and rank 1, each of the three stages of §4.12 was measured on its own, for j = 1, …, 5:

\begin{center}\begin{small}
\begin{tabular}{>{\raggedright\arraybackslash}p{0.588\linewidth}>{\raggedright\arraybackslash}p{0.066\linewidth}>{\raggedright\arraybackslash}p{0.066\linewidth}>{\raggedright\arraybackslash}p{0.066\linewidth}>{\raggedright\arraybackslash}p{0.066\linewidth}>{\raggedright\arraybackslash}p{0.066\linewidth}}
\hline
measured quantity & j = 1 & 2 & 3 & 4 & 5 \\
\hline
v\_p(X(t₁ + p\textasciicircum{}j) − X(t₁)) — step (1), predicted j & \textbf{1} & \textbf{2} & \textbf{3} & \textbf{4} & \textbf{5} \\
v\_p(φ\_p(x₀ + p\textasciicircum{}j) − φ\_p(x₀)) — step (3), predicted j + 1 & \textbf{2} & \textbf{3} & \textbf{4} & \textbf{5} & \textbf{6} \\
v\_p(Θ(t₁ + p\textasciicircum{}j) − Θ(t₁)) — the theorem, predicted j − 1 & \textbf{0} & \textbf{1} & \textbf{2} & \textbf{3} & \textbf{4} \\
\hline
\end{tabular}
\end{small}\end{center}

each entry holding in \textbf{400 of 400} boxes, with no box off the predicted value. In the same boxes the two hypotheses of step (1) were verified directly: c₁ = 2y(T) + a₁x(T) + a₃ is a unit in 400 of 400, and v\_p(ψ\_p′(x(T))·c₁) = 1 in 400 of 400.

\textit{Negative controls for the sensitivity.} The exponent j + 1 of step (3) comes from Lemma 5 and from nothing else. Replacing φ\_p by a polynomial that is \textbf{not} a polynomial in x\textasciicircum{}p modulo p, the slope drops to j: for x² + x the valuation equals j in 346 of the 400 boxes and is deeper in the remaining 54 (there the derivative 2x₀ + 1 happens to be divisible by p), and for the even-index division polynomial f\_\{p−1\} it equals j in 272 and is deeper in 128 — never j + 1 systematically. So the extra digit is the divisibility of φ\_p′ by p and not an artefact of the perturbation p\textasciicircum{}j.

\textit{The negative control for the split hypothesis.} Over the full sweep to conductor 20000 — \textbf{14,168 boxes} — the equivalence of Lemma 4, split ⟺ e ≥ 2, held in \textbf{14,168 of 14,168}: the 1,678 split boxes have e ∈ \{2, 3, 4, 5\} (1,482 / 164 / 30 / 2) and the \textbf{12,490} non-split boxes have e = 1 in \textbf{every} case. On a sample of 300 boxes of each kind, ψ\_p had an integral p-adic root in \textbf{300 of 300} split boxes and in \textbf{0 of 300} non-split ones. This is the sharp form of the negative control for Theorems G′ and G″: in the non-split case there is no p-torsion point T in E(Q\_p), hence no formal parameter t₁ = t(nP − T) and no layers γ, γ₂, γ₃ at all. The sensitivity statement itself is local and survives — φ\_p′ is divisible by p whatever the reduction — but the chain that turns a digit of Θ into a layer of E(Q\_p), and with it Theorem G″, exists only on the split side.

\textit{The non-degeneracy (ND).} Two computations. (a) Over the 1,678 split boxes, G() = φ\_p′()/p ≢ 0 (mod p) in \textbf{1,678 of 1,678}. (b) Exhaustively: for each p with 5 ≤ p ≤ 19, for \textbf{every} pair (a₄, a₆) ∈ F\_p² with 4a₄³ + 27a₆² ≠ 0 — that is for every short Weierstrass model over F\_p, hence for every isomorphism class, each class occurring several times, 20 + 42 + 110 + 156 + 272 + 342 = \textbf{942 models} — the resultant condition holds: gcd(G mod p, ψ\_p mod p) = 1 in \textbf{942 of 942}, with no common root, ordinary and supersingular alike (\textbf{52} of the 942 are anomalous and \textbf{110} supersingular). For the 52 anomalous curves the conclusion was also read off directly: G(x()) ≠ 0 at \textbf{all 732} non-zero points  of E(F\_p). Since Lemma 5 is an identity over Z, this settles (ND) for every curve over Q with 5 ≤ p ≤ 19; the exhaustion was later carried to p ≤ 41 (§5.9), and then replaced by the proof of Theorem ND, which covers every p ≥ 5. That "942 of 942" is not vacuous was checked in two ways. First, the shape of the divisor: ψ\_p mod p has exactly (p−1)/2 roots in F\_p, each of multiplicity exactly p, in every anomalous box — which is the statement ψ\_p ≡ (const)·g(x\textasciicircum{}p) of §2 read off, with g separable and split over F\_p because ker V = E(F\_p). Second, a random polynomial of the same degree as G, deg\_x G = p² − (p+3)/2, shares a root with ψ\_p mod p in \textbf{813 of 2,080} trials (39.1\%), against the predicted 1 − (1 − 1/p)\textasciicircum{}\{(p−1)/2\} of 36.0\% to 38.5\%; and the derivatives of the neighbouring division polynomials f\_\{p+m\}, m = 1, …, 4, do share a root with ψ\_p mod p in many boxes (for instance 6, 6, 0 and 12 of the 12 anomalous curves over F\_13). So the coprimality of G with ψ\_p mod p is a property of G and not of degrees.

\textit{The 576 pairs, and what separates them.} Over conductor ≤ 20000, rank 1, p ≤ 19 there are \textbf{576} pairs of curves isomorphic over Q\_p with different v\_p(Reg\_p); the pairs (v\_lo, v\_hi) are (0,1) 475 times, (0,2) 9, (0,3) 74, (1,2) 2, (1,3) 16, and no pair has v\_lo ≥ 2. Hence \textbf{all 576 are separated by the ledger of Theorem N}: the difference sits in Θ mod p in 558 pairs, in Θ/p only in 4, in \textbf{q\_p(d′) only in 2}, in both in 12, and in neither in 0. The two pairs separated by the denominator alone — 3888p1 against 3267d1, and 19044a1 against 3267d3, both at p = 7 — are the shortest exhibition of the content of Theorem N, and the necessity of F in Theorem G″. Both pairs have j = 0, hence c = 0: the first local term of the ledger is absent there and, by Proposition 18, the whole of p\textasciicircum{}\{−2\}ℒ is invisible before the eleventh digit, so the separation by q\_p(d′) at the second digit is as clean as it can be. That is not a CM phenomenon, though. \textbf{432 of the 576 pairs have j = 0} — the boxes with j = 0 fall into two classes of 38, at p = 7 and at p = 19, and the number of pairs inside a class grows quadratically — and the rate of pairs separated at the second digit alone is 16 of 432 against 2 of 144, which is no difference at all (the expectation apportioned by the number of pairs is 13.5 against 16 observed). Theorem N has no c in it: c enters first at the third digit, through W ≡ cγ², so the way the pairs separate could not have depended on the vanishing of c. Structurally, among curves with the same j-invariant anomalousness forces the twisting factor to be a p-adic square, so every same-j pair surviving the filters is Q\_p-isomorphic: \textbf{6,295 of 6,295} in this sweep (4,889 with j ≠ 0 and 1,406 with j = 0).

\textit{The scaling law.} Θ ↦ Θ + 2p·q\_p(u) under x ↦ u²x, y ↦ u³y with p ∤ u (the Remark of §4.12): \textbf{226 of 226}.

\subsection*{5.9 The identities behind Theorem ND, and the constants}

Everything in this subsection is a check of §3.14 and §4.13, in PARI/GP, with the controls stated. None of it is a proof; the proofs are in §4.13, and the numbers are what preceded them and what tested them afterwards.

\textit{The exhaustion that preceded the proof.} (ND) was verified for \textbf{every} elliptic curve over F\_p with 5 ≤ p ≤ 41 — 942 curves for p ≤ 19 (§5.8) and \textbf{5,220} more for p = 23, 29, 31, 37, 41, so \textbf{6,162} in all — by gcd(G mod p, ψ\_p mod p) = 1, and, on the anomalous curves, directly at all non-zero points of E(F\_p): \textbf{732} points for p ≤ 19 (§5.8) and \textbf{4,480} for p = 23, …, 41, so \textbf{5,212} in all. For 5 ≤ p ≤ 19 the stronger statement Res\_x(G, ψ\_p mod p) = c\_p·Δ₀\textasciicircum{}\{m\_p\} was proved as a polynomial identity (symbolically for p = 5, 7; for p = 11, …, 19 by interpolation in the weight-homogeneous shape Res = a₄\textasciicircum{}\{3m\}P(a₆²/a₄³), deg P ≤ m, over a Galois ring), with c\_p = 2, 2, 1, 6, 13, 11.

\textit{The two identities.} Proposition 14 holds as an identity in F\_p[a₄, a₆][x] for p = 5, 7, 11, 13, 17, 19 (7.0 s at p = 19) and its corollaries u\textasciicircum{}\{(p−1)/2\} ∣ K and K/u\textasciicircum{}n\_p ∈ F\_p[x\textasciicircum{}p] for p ≤ 19. Proposition 15 holds as an identity in F\_p[a₄, a₆][X] — in the form A²Φ = 4w\textasciicircum{}\{(p)\}(Ψ′)² + A²D(Ψ²), whose degree in X is p — for \textbf{p = 5, …, 23} (574 s at p = 23, 286 MB of PARI stack) and at 56 sampled points for p ≤ 43, the derivation D being taken by dual numbers over F\_p[y]/(y²). Controls: replacing D by ∂/∂a₄, or w\textasciicircum{}\{(p)\} by u, breaks the identity at p = 5; replacing ψ\_p by ψ\_\{p−2\} breaks Lemma 13 itself (0 of 96).

\textit{The Tate family.} With Ψ built from ψ\_p over F\_p[[q]] and h = Ψ/A: the identification ψ\_p mod p = A∏\_i(x\textasciicircum{}p − ξ\_i) with ξ\_i = X\_\{q\textasciicircum{}p\}(q\textasciicircum{}i), and A ≡ 1, hold to q\textasciicircum{}\{30\} for p = 5, 7; the congruence (R′) in the form [w\textasciicircum{}\{(p)\}(h′)²]′ − A²h′·(q dh/dq) ≡ 0 (mod h) holds \textbf{identically to q\textasciicircum{}\{30\} for p = 5, 7, 11, 13}, while the same expression with the opposite sign is non-zero from q² (p = 5) and q³ (p = 7) on — this is the control that fixes the sign of D on the Tate family. The three steps of §4.13(c) were measured separately: q dξ\_i/dq = 2iη\_i and (R′) in root form to q\textasciicircum{}\{22\} for p = 5, 7, 11 and every i; the product rearrangement G(W) = (−1)\textasciicircum{}nq\textasciicircum{}\{−n\_p(n\_p+1)/2\}W\textasciicircum{}n\_pθ\_q(W) to a relative error 3.4·10\textasciicircum{}\{−71\} at q = 1/10, W = 0.37; the constant term κ = n\_p + 1 − i, the telescoped sum −(i−1) and the ratio D²H/DH = p − 2i to 20 digits. For (L): φ₁ + 2σ = 0 and φ₁ + 6σ = 2A·DA to q\textasciicircum{}\{22\} for p = 5, 7.

\textit{The constants.} At ordinary points of F\_p², for one box per prime, the closed forms of Proposition 16 were compared with the measured resultants for \textbf{every p with 5 ≤ p ≤ 151} (34 primes): c\_K agreed in \textbf{34 of 34} and c\_S in \textbf{33 of 34}, the single exception being a sample point at p = 19 at which lc(S₁) = 2·DA vanishes, so that deg\_x S₁ drops from 171 to 152 and the resultant does not specialise; at the other four sample points of p = 19 the value is the predicted 7. The values c\_K = 4, 4, 10, 10, 1, 12, 1, 28, 1, 11, 1, 42 for p = 5, …, 43 were measured before the closed form was proved, and c\_K(47) = 1, c\_K(53) = 52 were predicted from the closed form before being measured. The factorisation of the resultant used in §4.13 — Res\_x(K, ψ\_p) = Res\_x(u, ψ\_p)\textasciicircum{}n\_p·(A\textasciicircum{}\{n\_p−1\}Res\_X(h′, h))\textasciicircum{}p, with the two factors as displayed there — holds at sample points for p = 5, 7, 11, 13, and γ\_p = 2, 3, 10, 2 there agrees with the closed form.

\textit{The constant c\_p in one line.} Assembling the two counts of Proposition 16 gives, in F\_p,

\begin{quote}
\textbf{c\_p = ε\_p·(−1)\textasciicircum{}\{n\_p + n\_p(n\_p+1)/2 + M\}·n\_p!·2\textasciicircum{}\{m\_K + n\_p + 4M\}},
\end{quote}

with n\_p = (p − 1)/2, M = p(p² − 1)/24, m\_K = (p² − 1)(p² − 3)/24, ε\_p = −1 exactly when p ≡ 3 (mod 8), and every exponent read modulo p − 1; the regrouping of the exponents is Lean \texttt{cp\_\allowbreak{}one\_\allowbreak{}line}. It agrees with the measured c\_p for all twelve primes 5 ≤ p ≤ 43: the six values 2, 2, 1, 6, 13, 11 for p ≤ 19 recorded above, and \textbf{22, 12, 30, 14, 9, 1} for p = 23, 29, 31, 37, 41, 43, which were printed from the closed form before being measured. One convention has to be fixed for this to hold: \textbf{Δ₀ = 4a₄³ + 27a₆² = −Δ/16}, and not Δ itself — dividing by Δ instead makes eight of the twelve values disagree.

\subsection*{5.10 Theorem N* to the tenth digit}

Everything here is a check of §§3.11–3.12, in PARI/GP, over one fixed set of boxes. The boxes are: conductor ≤ 12000, p ∈ \{5, 7\}, a\_p = 1 (so \#E(F\_p) = p), p ∤ ∏\_ℓ c\_ℓ, ι, split, \textbf{rank ≥ 1} with P the first generator in Cremona's tables, and e ≥ 2 — \textbf{648 boxes}, 332 with p = 5 and 316 with p = 7, the distribution of e being [551, 82, 13, 2] over e = 2, 3, 4, ≥ 5. The point P″ is computed p-adically: x(P″) = φ\_p(x₀)/ψ\_p(x₀)² at x₀ = x(P′) + O(p\textasciicircum{}\{30\}), y(P″) from the Weierstrass equation (the sign of the square root is immaterial, since z(−Q) = −z(Q) and only z² and z⁴ occur), t = −x/y, and z = log\_E(t) by the formal logarithm. For the deep digits ℒ is built as a series: \texttt{ellformallog} and \texttt{serreverse} give x(z), the differential equation of Proposition 17 is integrated twice with \texttt{intformal}, c is PARI's \texttt{ellpadics2}, and ℒ(z) = log(x(z)z²) + 2log(σ\_p(z)/z) is evaluated at z(P″).

\textit{Calibration, before anything new.} The same generating rule — take Theorem N*, keep the terms of valuation ≤ k − 1 by Proposition 18, discard the rest — reproduces \textbf{verbatim} the right-hand sides of Corollary N⁽³⁾ and Corollary N⁽⁴⁾ as they stand in §3.12, and at k = 5 it produces the term (2/5)p⁵F⁵ by itself, with no term put in by hand. Each of the three holds in \textbf{648 of 648} boxes. Beyond k = 5 the generated truncation and the \textbf{displayed} congruence are tested separately, the displayed one exactly as it is printed in §3.12: the corrected Corollary N⁽⁶⁾ holds in \textbf{648 of 648} and the full Corollary N⁽⁷⁾ in \textbf{648 of 648}, the latter sharply (valuation exactly 7) in 499 of them. Printing the geometric factor of the F⁵ term as 1 rather than (1 + p) fails in \textbf{195} boxes, all with p = 5 and 5 ∤ F, which is how the factor was found: a truncation that the machine generates correctly can still be mis-transcribed, so every displayed congruence is now checked in the form in which it is displayed. The counts 648, 332, 316 and the distribution of e agree with the independent sweep of §5.3's filter. Restricting the same filter to rank exactly 1 gives 628 boxes (321 with p = 5, 307 with p = 7, the distribution of e being [534, 80, 13, 1]); the congruences of §3.12 assume no rank, since Theorem N* is a statement about one point with χ(P) ≠ 0, so the wider filter is the one used.

\textit{The coefficients of ℒ, by an independent route.} μ₂ = −c, μ₄ = −b₄/12, μ₆ = −b₂b₄/360 − b₆/30 and the vanishing of every odd coefficient were read off the formal series built as above — which uses the differential equation but not the hand computation of §3.11 — for the four curves [0,0,0,−1,1], [1,−1,0,−79,289], [1,2,3,4,5], [0,1,1,−2,0] with c = 7/3, 11, −1/5, 0: \textbf{4 of 4}, to the order z\textasciicircum{}\{12\}. The curve [1,2,3,4,5] has a₁ ≠ 0 and a₃ ≠ 0, which is what makes the vanishing of the odd coefficients a test and not a consequence of the model.

\textit{The table of §3.12.} The four columns were measured at k = 6, 7, 8, 9, 10, and the counts are those displayed there. The split by p at k = 10 is the sharp form of the third line of Proposition 18: truncating ℒ at z⁴ holds in \textbf{316 of the 316} boxes with p = 7 and fails in 256 of the 332 with p = 5, because v\_5(μ₆) = −1 and v\_7(μ₆) = 0.

\textit{Controls at k = 6.} Eight modifications of the right-hand side, on the same 648 boxes: dropping the p⁵ tail of the Θ series, 111; dropping the p³ tail of F², 204; dropping the p² tail of F³, 204; dropping the (1 + p) of F⁴, 204; dropping the (1 + p) of Θ³, 111; dropping ℒ altogether, 41 — and in the 15 boxes with e ≥ 4, where Proposition 18 says ℒ is invisible modulo p⁶, \textbf{15 of 15}; adding the term (a₃ − 2c₄)/(p²y(P″)), which the expansion by hand had suggested as the next local ingredient, 398; adding μ₄z⁴, \textbf{648}, since its valuation is 6 and it cannot change the congruence modulo p⁶ — the last two being a positive and a negative control on the same quantity.

\textit{Controls on the logarithms, split by p.} At k = 8, dropping the term with j = 5 of the Θ series — valuation 8 − v\_p(5), hence 7 at p = 5 and 8 at p = 7 — leaves \textbf{316 of 316} at p = 7 and 64 of 332 at p = 5. At k = 10, dropping the term with j = 10 of the F series — valuation 10 − v\_p(10), hence 9 at p = 5 and 10 at p = 7 — leaves \textbf{316 of 316} at p = 7 and 117 of 332 at p = 5. In both, one of the two primes is untouched and the other collapses, which is the factor v\_p(j) measured rather than asserted.

\textit{z² against 1/x(P″).} The valuation of z² − 1/x(P″), minus 2e, over the 648 boxes: 4 in 393 boxes, 5 in 38, 6 in 52, 7 in 75, 8 in 39, 9 in 2, 10 in 2, and \textbf{≥ 11 in 47} (11 in 26, 12 in 9, 14 in 1, 15 in 1, 16 in 2, 17 in 6, 18 in 1, 20 in 1) — 648 in all. The lower bound 4 is Lemma 21: z² − 1/x(z) = −(b₂/12)z⁴ + O(z⁶), so the valuation is 4e + v\_p(b₂) as soon as the tail is deeper, and never less than 4e. In the \textbf{497} boxes with b₂ ≠ 0 the measured valuation is 4e + v\_p(b₂) \textbf{exactly}, in all 497, and z² − 1/x(P″) + (b₂/12)z⁴ is deeper in every one of them; since the excess over 2e is then 2e + v\_p(b₂), the 393 boxes of excess 4 are exactly those with e = 2 and p ∤ b₂ (b₂ is a p-adic unit in 453 of the 497, and b₂ = 0 in the remaining 151 boxes, where the leading term is absent and the valuation is deeper still). Hence p\textasciicircum{}\{−2\}(z² − 1/x(P″)) has valuation ≥ 4e − 2 ≥ 6: \textbf{1/x(P″) may be used through the sixth digit and not at the seventh}, where it leaves 310 of 648. This is the computation behind the statement in §3.13 that W is a quantity of the digits up to the sixth. The explanation through a₁ that an earlier version gave is not the right one — it reaches only 2e + 2 — and the terms of R linear in a₁ cancel against a₁t to leading order, so a₁ cannot be the source of the fourth power.

\textit{Proposition 19.} In the first 40 boxes of conductor ≤ 2500, T ∈ E(Q\_p)[p] was built by Hensel from the p-adic roots of ψ\_p, t₁ = t(P′ − T) taken for each of the four choices of sign, and z(P″) compared with p·log\_E(t₁): \textbf{40 of 40}, and e = 1 + v\_p(t₁) in \textbf{40 of 40}. The sign of T changes z(P″) by a sign and z², z⁴ not at all.

\subsection*{5.11 The two rates, over the cells of §5.10}

Proposition 20 was measured on the filter of §5.10 — the same boxes, called cells here — split by conductor and by e into three sweeps. \textbf{Sweep a is the part of sweep b with conductor ≤ 3000, run separately and at a lower working precision; the distinct cells are those of sweeps b and c, which are disjoint (e = 2 against e ≥ 3).} In each cell T ∈ E(Q\_p)[p] was built by Hensel from the p-adic roots of ψ\_p, t₁ = t(nP − T) by the group law in E(Q\_p), X(t) = x(τ\_T(R(t))) as a series, Θ(t) = Q₂(φ\_p(X(t))), z(t) = p·log\_E(t), and ℒ as the series obtained by integrating the differential equation of Proposition 17 twice — which uses the equation but not the hand computation of §3.11.

\begin{center}\begin{small}
\begin{tabular}{>{\raggedright\arraybackslash}p{0.116\linewidth}>{\raggedright\arraybackslash}p{0.209\linewidth}>{\raggedright\arraybackslash}p{0.07\linewidth}>{\raggedright\arraybackslash}p{0.289\linewidth}>{\raggedright\arraybackslash}p{0.116\linewidth}>{\raggedright\arraybackslash}p{0.116\linewidth}}
\hline
sweep & conductor & e & cells & cells with c = 0 & cells with c ≠ 0 \\
\hline
a & ≤ 3000 & = 2 & \textbf{103} (p = 5: 47; p = 7: 56) & 8 & \textbf{95} \\
b & ≤ 6000 & = 2 & \textbf{259} (126; 133) & 14 & \textbf{245} \\
c & ≤ 12000 & ≥ 3 & \textbf{97} (e = 3: 82; e = 4: 13; e = 5: 2) & 6 & \textbf{91} \\
\hline
\end{tabular}
\end{small}\end{center}

The cell counts are the counts of the filter of §5.10 restricted to the stated range and to nothing else: they are read off the list of 648 boxes mechanically. Sweeps b and c together are \textbf{356 cells}, of which \textbf{336} have c ≠ 0. The filter fixes P as the first generator at any rank, and the six cells of rank two in this range — three in sweep b, three in sweep c — behave exactly like the rest, so nothing below depends on a restriction to rank one.

\textit{The cells where c vanishes.} Among the 356 cells the constant c of Proposition 17 — PARI's \texttt{ellpadics2} — vanishes in exactly \textbf{20}, and those 20 are exactly the cells with \textbf{j = 0}: 20 of 20 and 336 of 336, the positive and the negative control on the same criterion. This is the vanishing of the p-adic weight-two Eisenstein series at a CM point, here for complex multiplication by Q(√−3). \textbf{All 20 lie at p = 7, and the reason is the trace and not the splitting of p.} For j = 0 the curve has complex multiplication by Q(√−3), so at a prime of good ordinary reduction a\_p² − 4p = −3f² for some f ∈ Z; anomalousness a\_p = 1 therefore forces \textbf{4p = 1 + 3f²}, hence p ≡ 1 (mod 3), and among 5 ≤ p ≤ 19 only p = 7 (f = 3) and p = 19 (f = 5) qualify — 4·13 − 1 = 51 = 3·17 is not of that shape, so p = 13 is excluded although it splits in the field, while p = 5, 11, 17 are ≡ 2 (mod 3) and supersingular (Lean \texttt{cm\_\allowbreak{}prime\_\allowbreak{}mod\_\allowbreak{}three}, \texttt{cm\_\allowbreak{}no\_\allowbreak{}anomalous\_\allowbreak{}of\_\allowbreak{}mod\_\allowbreak{}two}, \texttt{cm\_\allowbreak{}seven\_\allowbreak{}nineteen}, \texttt{cm\_\allowbreak{}not\_\allowbreak{}thirteen}). Over conductor ≤ 20000 the sweep of §5.8 has \textbf{76 boxes with j = 0, 38 at p = 7 and 38 at p = 19}, and \textbf{all 76 are split and all 76 have c = 0} (76 of 76 and 1,602 of 1,602); the 20 cells here are the part of those 76 that the filter of §5.10 keeps. Since the corrected Proposition 20 (iii) predicts a value at c = 0 as well, these cells are no longer set aside in the ℒ rows: they are the rows where the leading term is m = 3.

\begin{center}\begin{small}
\begin{tabular}{>{\raggedright\arraybackslash}p{0.432\linewidth}>{\raggedright\arraybackslash}p{0.241\linewidth}>{\raggedright\arraybackslash}p{0.285\linewidth}}
\hline
measured & Proposition 20 says & result \\
\hline
v\_p(Θ(t₁ + p\textasciicircum{}j) − Θ(t₁)), j = 1..7 & exactly j − 1 & \textbf{356 of 356}, every row \\
v\_p(z(t₁ + p\textasciicircum{}j) − z(t₁)) & exactly 1 + j & \textbf{356 of 356} \\
e = 1 + v\_p(t₁) (Proposition 19) & — & \textbf{356 of 356} \\
v\_p(ℒ(z′) − ℒ(z)), 1 + j ≠ e & min\_m [ v\_p(μ\_\{2m\}) + (1 + j) + (2m − 1)·min(e, 1+j) ] & \textbf{2,136 of 2,136 rows}, all 356 cells, j = 1..7 \\
— the same rows against the old wording v\_p(c) + min(e, 1+j) + (1 + j) &  & \textbf{2,016 of 2,016} where c ≠ 0; \textbf{false in all 120 rows} where c = 0 \\
v\_p(ℒ(z′) − ℒ(z)), 1 + j = e & at least the same, with the cancellation of the boundary row & \textbf{356 of 356 rows} at the boundary value — see below \\
v\_p of the difference of the two sides of Theorem N*, F fixed & exactly j − 1 & \textbf{356 of 356} \\
v\_p(2log(1 + p(F + p\textasciicircum{}i)) − 2log(1 + pF)), i = 0..3 & exactly i + 1 & \textbf{356 of 356} \\
\hline
\end{tabular}
\end{small}\end{center}

\textit{Controls.} The \textbf{positive} control is the row j = 3, which reproduces the sharpness used by Theorem G″ — Θ moving at valuation exactly 2 — with the same instruments. The \textbf{negative} controls are the rows j = 5, 6, 7, where nothing moves modulo p⁴: the fifth and deeper layers do not enter the verdict at k = 4, and at j = 4 the count is 0 of 356, as it must be. A second \textbf{independent route} measures ℒ by the single term −cz² instead of the whole series and gets the same valuation in all \textbf{336} cells where c ≠ 0.

\textit{(iii) row by row, and the unconditional form.} Over the 356 cells and j = 1, …, 7 — \textbf{2,492 rows} — the measured v\_p(ℒ(z′) − ℒ(z)) was compared with the displayed right-hand side of Proposition 20 (iii) itself, the minimum over m of v\_p(μ\_\{2m\}) + (1 + j) + (2m − 1)·min(e, 1 + j), the coefficients μ\_\{2m\} read off the series for m ≤ 8 and the prediction formed from (e, j, v\_p(μ\_\{2m\})) by a separate program. Equality in \textbf{2,136 of the 2,136 rows} with 1 + j ≠ e, and the minimum is attained by one m alone in all 2,492 of them: m = 1 in the 2,352 rows with c ≠ 0 and m = 3 in the 140 rows with c = 0. The \textbf{unconditional form} — v\_p(c) + min(e, 1+j) + (1 + j) with no hypothesis on c — is false in exactly the \textbf{140 rows with c = 0}, where its right-hand side is infinite while the measured valuation is finite and equals (1 + j) + 5·min(e, 1 + j): 13, 13, 14, …, 18 at the eleven cells with e = 2 that cancel at the boundary row, 12, 13, 14, …, 18 at the three that do not, and 12, 19 or 21, 19, 20, …, 23 at the six with e = 3 — v\_p(μ₆) = 0 in all twenty. Its hypothesis v\_p(c) \textless{} 2·min(e, 1+j) + v\_p(μ₄) holds in all 2,352 rows with c ≠ 0, the narrowest margin being \textbf{one digit} (2115g1 and 2502e1, where v\_p(c) = 3 against the threshold 4).

\textit{A box with c ≠ 0 where the hypothesis on c is needed.} The positive control for the hypothesis needs a box with c deep, so the search was widened to conductor ≤ 40000 over 5 ≤ p ≤ 19: \textbf{3,214 split boxes} (1,009 at p = 5, 976 at p = 7, 301, 484, 151, 293), of which \textbf{61} have v\_p(c) ≥ 3 — v\_p(c) = 3 in 8 boxes, 4 in 49 and 6 in 4. In 60 of the 61 the hypothesis still holds, because v\_p(μ₄) rises with v\_p(c) in this sample (all 49 boxes with v\_p(c) = 4 have v\_p(μ₄) = 2). The exception is \textbf{27729a1 at p = 7}, with e = 3, v₇(c) = 6 and p ∤ b₄: at j = 1 the threshold is 2·min(3, 2) + 0 = 4, the m = 2 term leads, and the measured valuation is \textbf{8 = v\_p(μ₄) + (1 + j) + 3·min(e, 1 + j)}, while the unconditional form gives ≥ 6 + 2 + 2 = 10. \textbf{The unconditional lower bound is false there by two digits, and Proposition 20 (iii) gives 8 exactly.} Over all 111 boxes measured outside the cells of §5.10 — the 61 deep ones and the 50 boxes with j = 0 at p = 19 up to conductor 40000 — Proposition 20 (iii) holds in \textbf{666 of 666} rows with 1 + j ≠ e and the boundary formula in \textbf{111 of 111}, while the unconditional form fails in 351 rows: 350 at c = 0 and the one at 27729a1. The six rows where two values of m tie (27729a1 at j = 2, …, 7) showed no cancellation, so equality held there as well — the statement claims it only for a unique minimiser.

\textit{The boundary row.} When 1 + j = e the two factors of Δ(z²) = (2z + Δz)Δz have equal valuation and the leading terms can cancel. Writing t₁ = p\textasciicircum{}\{e−1\}u with u a unit, z = p\textasciicircum{}e·u + O(p\textasciicircum{}\{e+1\}) and Δz = p\textasciicircum{}\{1+j\} + O(p\textasciicircum{}\{2+j\}), so at 1 + j = e the sum is 2z + Δz ≡ p\textasciicircum{}e(2u + 1), and the cancellation occurs \textbf{exactly when u ≡ −1/2 (mod p)}. Over the 336 cells with c ≠ 0 that criterion is exact: the ℒ row is deeper than v\_p(c) + min(e, 1+j) + (1 + j) in \textbf{67 cells, and those 67 are precisely the cells with u ≡ −1/2 (mod p)} — 67 of 67 and 269 of 269, and the same at two working precisions. Split by e the count is 44 of 245 at e = 2, 22 of 76 at e = 3, 0 of 13 at e = 4 and 1 of 2 at e = 5. Since u is a unit the heuristic frequency is \textbf{1/(p − 1) and not 1/p}: it predicts 51.3, 16.2, 2.8 and 0.4 respectively, and 70.8 over the 336 cells against 67 observed. Every deviation from the general bound is on the slow side — the ledger lagging more, not less — so the conclusion of Proposition 20 is strengthened.

\textit{The boundary row for a general m, and how deep the cancellation goes.} Write w := Δz/z, a unit at the boundary row, with w ≡ u\textasciicircum{}\{−1\} (mod p). Then Δ(z\textasciicircum{}\{2m\}) = z\textasciicircum{}\{2m\}((1 + w)\textasciicircum{}\{2m\} − 1), so the leading term of the m-th summand of ℒ cancels \textbf{exactly when (1 + w)\textasciicircum{}\{2m\} = 1 in F\_p}, that is when (u + 1)\textasciicircum{}\{2m\} = u\textasciicircum{}\{2m\}. For m = 1 this is 2u + 1 ≡ 0, the criterion above; when (p − 1) ∣ 2m — for instance p = 7 and 2m = 6, which is the leading term at a cell with c = 0 — \textbf{every} u with u ≢ −1 cancels. Over the 20 cells with c = 0, all at p = 7 with leading term m = 3, the cancellation happens in exactly the 17 cells with u ≢ −6 and fails in the three with u = 6 (1521a1, 1521a2, 3600bd1): \textbf{20 of 20}, the positive and the negative control on one criterion. The depth of the cancellation is v\_p((1 + w)\textasciicircum{}\{2m\} − 1) − v\_p(w) as well, which for m = 1 is v\_p(2z + Δz) − e: over the 336 cells with c ≠ 0 the excess of the measured ℒ row over the general bound equals that quantity in \textbf{336 of 336}, with the excess 0 in 269 cells, 1 in 54, 2 in 12 and 3 in 1. So the boundary row is not an exception with an unknown size: \textbf{which} cells lag more is one residue condition, and \textbf{by how much} is the valuation of one explicit unit. With w in hand the value at the boundary row is min\_m [ v\_p(μ\_\{2m\}) + 2m·e + v\_p((1 + w)\textasciicircum{}\{2m\} − 1) ], and that predicted the measured valuation in \textbf{356 of 356} boundary rows, the 20 cells with c = 0 included (Lean \texttt{boundary\_\allowbreak{}cancel\_\allowbreak{}m1}, \texttt{boundary\_\allowbreak{}cancel\_\allowbreak{}pow}, \texttt{boundary\_\allowbreak{}cancel\_\allowbreak{}of\_\allowbreak{}card}, \texttt{boundary\_\allowbreak{}p7\_\allowbreak{}m3}, \texttt{boundary\_\allowbreak{}p7\_\allowbreak{}exception}).

\textit{The criterion of the boundary row at p = 19, as a negative control.} At p = 7 the criterion (1 + w)\textasciicircum{}\{2m\} = 1 is satisfied by every u ≢ −1 when 2m = 6, because 6 = p − 1; at p = 19 it is not, and the same formula must then make the cancellation the exception. The sweep of §5.8 extended to conductor 40000 has \textbf{50 boxes with j = 0 at p = 19} (38 of them within conductor 20000, which reproduces the count there), all split, all with c = 0 and all with b₄ = 0, so the leading term is m = 3 in every one. The congruence (u + 1)\textasciicircum{}6 ≡ u\textasciicircum{}6 (mod 19) has exactly five solutions, u ≡ 2, 7, 9, 11, 16, and the boundary row is deeper than the general value in exactly the \textbf{12 of the 50 boxes whose u is one of those five} — \textbf{50 of 50}, against 17 of 20 at p = 7. The heuristic frequency 5/18 = 27.8\% sits against 24\% observed. One criterion, two primes, opposite verdicts, and right both times.

\textit{The residue u is a datum of the generator, not of (E, p).} Replacing P by −P replaces T by −T and t₁ by −t₁ + a₁t₁² + ⋯, so u = t₁/p\textasciicircum{}\{e−1\} is replaced by −u modulo p: measured in \textbf{356 of 356} cells, with no exception. Hence (E, p) determines only the class \{±u\}, and the sign is the convention by which the generator is chosen in \texttt{elldata}. The criterion of the boundary row is unaffected: it is stated for an ordered pair t₁, t₁ + p\textasciicircum{}j, and if both increments ±p\textasciicircum{}j are allowed it reads u ≡ ±1/2, which is sign-free, so the heuristic frequency 1/(p − 1) stands. It does mean that a count of raw residues is a count of a convention: at p = 7 the 178 cells give χ² = 12.85 on 5 degrees of freedom for the raw u — an apparent excess at u = 5 — while the classes \{±u\} give 58, 67, 53 against 59.3, χ² = 1.70 on 2 degrees of freedom, and at p = 5 the classes give 87, 91 against 89.0. The whole of the raw excess is the split 21 against 46 \textbf{inside} the class \{2, 5\}, and it survives counting isogenous curves once (39 against 24.7 at p = 7, with the classes then 53, 51, 44): it is a property of how generators are tabulated and there is nothing p-adic left in it.

\subsection*{5.12 Depth and height over the split boxes of rank two}

\textit{The boxes.} \texttt{elldata}, conductor ≤ 200000, p ∈ \{5, 7, 11, 13, 17, 19\}, rank exactly 2, a\_p = 1, \#E(F\_p) = p, p ∤ ∏\_ℓ c\_ℓ, p ∤ \#E(Q)\_tors, ι — \textbf{25,374 curves and 27,355 boxes}, of which \textbf{3,056 are split} (988, 830, 321, 448, 138, 331 by p) and 24,299 non-split. E(Q) is torsion-free in all 27,355, as remark (a) of §3.15 requires. Nothing p-adic is computed: by Lemma 23 a rational point Q lies in E\_k(Q\_p) ∖ E\_\{k+1\}(Q\_p) exactly when v\_p(d(Q)) = k, splitting is read as D₁ := v\_p(d(pP₁)) ≥ 2 for a generator P₁ with χ(P₁) ≠ 0, and det(χ, γ) ≡ 0 as D\_Q := v\_p(d(Q)) ≥ 2 for Q ∈ (Λ ∩ E\_1) ∖ pΛ (remark (b) of §3.15). The sweep takes 85 seconds.

\textit{Calibration.} On the six boxes whose ρ₁ had been measured from the p-adic height matrix, the test agrees in 6 of 6 (det(χ, γ) ≡ 0 at 8577a1 and 16011e1 at p = 5 and 12208i1 at p = 13; ≢ 0 at 2343c1, 3333b1, 4407a1). Three elements of (Λ ∩ E\_1) ∖ pΛ — the first, third and fifth candidates in order of norm — give the same verdict D\_Q ≥ 2 in \textbf{3,056 of 3,056} split boxes; the value of D\_Q itself can differ (at 16011e1, p = 5, another choice of three elements has depths 2, 2 and 3), and it is the verdict, not the value, that remark (b) makes a property of the box. The negative control, since the argument of remark (b) needs splitting: the three elements disagree in 4,291 of the 24,299 non-split boxes, a number that depends on the choice of elements (7,279 for another choice).

\textit{The frequency, and its cause.} If γ were a random multiple of χ on the two generators, det(χ, γ) ≡ 0 would occur with probability 1/p, and Σ1/p = 405.35 over the split boxes. It occurs in \textbf{100}.

\begin{center}\begin{small}
\begin{tabular}{>{\raggedright\arraybackslash}p{0.417\linewidth}>{\raggedright\arraybackslash}p{0.176\linewidth}>{\raggedright\arraybackslash}p{0.241\linewidth}>{\raggedright\arraybackslash}p{0.11\linewidth}}
\hline
measured & observed & naive expectation & ratio \\
\hline
det(χ, γ) ≡ 0 (rank 2, split, 3,056 boxes) & \textbf{100} & 405.35 & \textbf{0.247} \\
γ(P) ≡ 0 (rank 1, split, 13,926 boxes) & 1,805 & 1,847.2 & 0.977 \\
γ(P₁) ≡ 0 (rank 2, split) & 420 & 405.4 & 1.036 \\
χ(P) = 0 (rank 1, all 133,079 boxes) & 2,975 & 14,559.7 & 0.204 \\
\hline
\end{tabular}
\end{small}\end{center}

The statistics that fall short are exactly those that ask for a power of p in the denominator of a rational point; γ(P) ≡ 0 asks it of pP and not of P, and shows no deficit. Stratified by u := h\_min/log p, with h\_min the smallest canonical height ĥ on (Λ ∩ E\_1) ∖ pΛ found in the box [−2p, 2p]² of coefficients, in the normalisation of PARI — which is not the naive height of Proposition 24, nor the threshold of Proposition 24ᶜ, which adds β(E): over the 3,056 boxes the naive minimum minus the canonical minimum ranges from −1.72 to +7.61 times log p, and is negative in about one box in six. The last column cuts the same boxes by the naive minimum h(x(Q)) instead (boxes, det(χ, γ) ≡ 0, Σ1/p, ratio):

\begin{center}\begin{small}
\begin{tabular}{>{\raggedright\arraybackslash}p{0.122\linewidth}>{\raggedright\arraybackslash}p{0.109\linewidth}>{\raggedright\arraybackslash}p{0.131\linewidth}>{\raggedright\arraybackslash}p{0.109\linewidth}>{\raggedright\arraybackslash}p{0.109\linewidth}>{\raggedright\arraybackslash}p{0.337\linewidth}}
\hline
u & split boxes & det(χ, γ) ≡ 0 & Σ1/p & ratio & cut by the naive height \\
\hline
0 ≤ u \textless{} 2 & 506 & \textbf{0} & 74.3 & \textbf{0.000} & 0 boxes (Lemma 23, k = 1) \\
2 ≤ u \textless{} 4 & 1,732 & 15 & 236.5 & 0.063 & 1,528 · \textbf{0} · 207.8 · \textbf{0.000} \\
4 ≤ u \textless{} 6 & 510 & 59 & 62.4 & \textbf{0.946} & 1,031 · 52 · 140.0 · \textbf{0.371} \\
6 ≤ u \textless{} 8 & 178 & 18 & 20.1 & 0.895 & 307 · 36 · 37.1 · 0.971 \\
8 ≤ u \textless{} 12 & 100 & 5 & 9.4 & 0.531 & 144 · 7 · 16.1 · 0.433 \\
12 ≤ u \textless{} 20 & 27 & 3 & 2.2 & 1.363 & 43 · 5 · 3.9 · 1.281 \\
u ≥ 20 & 3 & 0 & 0.4 & 0.000 & 3 · 0 · 0.4 · 0.000 \\
\hline
\end{tabular}
\end{small}\end{center}

Cut by the canonical height, the ratio returns to 1 at u ≈ 4. Counting only the boxes with u ≥ 4 gives the expectation 94.5 against 100 observed (ratio 1.058); moving the threshold to 2, 3, 4.5 or 6 gives 0.30, 0.57, 1.36 and 3.11, so on that scale the agreement is not a fit. \textbf{Proposition 24 is a statement about the naive height, and cut by the naive height the picture is different.} Below 4 log p there is no box with det(χ, γ) ≡ 0 among 1,528, against 207.8 expected; the 15 boxes with det(χ, γ) ≡ 0 in the canonical stratum 2 ≤ u \textless{} 4 all have naive height between 4.92 and 7.65 times log p. Between 4 and 6 log p the ratio is 0.37, and above the naive threshold 100 are observed against 197.5 (ratio 0.51). So Proposition 24 accounts for 207.8 of the deficit of 405.35 − 100 = 305.35, about 68 \%. \textbf{Proposition 24ᶜ accounts for the rest.} The 15 boxes all have β(E) \textless{} 0 and (ĥ\_min − β)/log p between 4.000 and 5.154: they lie above the canonical threshold 4 log p + β(E), as Proposition 24ᶜ requires, and the return of the ratio to 1 near u ≈ 4 is the median of β(E)/log p, −0.07, seen on the scale β = 0. The naive column was computed twice, by separate programs: once from the shortest candidates in ĥ with a second pass, in the boxes first found at u ≥ 4, over every element of ĥ ≤ 4 log p + 12, and once by enumerating every element of Λ with ĥ ≤ 6 log p + 14, recognising ker χ by p ∣ d(Q) with no discrete logarithm, and doubling the bound in the 33 boxes where none was found. Each naive minimum is therefore an upper bound, exact whenever the minimising element lies within the bound; the two computations agree in every entry, the boundaries at 4 and 6 are fixed within the observed range of ĥ − h, and the strata above 6 are approximate. On the side of χ: over the 66,013 boxes of rank one with h(x(P)) \textless{} 2 log p there is no box with χ(P) = 0 (Lemma 23 with k = 1), and above u = 6 the ratio is 0.98 to 1.12.

The next table cuts the same boxes, and the 5,572 split boxes of the extension below, by the threshold of each proposition (boxes below the threshold, det(χ, γ) ≡ 0 among them, Σ1/p over them, share of the deficit) [computation, with m\_∞ a numerical minimum]. We call it table 5.12ᶜ.

\begin{center}\begin{small}
\begin{tabular}{>{\raggedright\arraybackslash}p{0.222\linewidth}>{\raggedright\arraybackslash}p{0.253\linewidth}>{\raggedright\arraybackslash}p{0.249\linewidth}>{\raggedright\arraybackslash}p{0.22\linewidth}}
\hline
cut & set A: conductor ≤ 200000, 3,056 boxes, 100 with det ≡ 0, deficit 305.35 & set B: conductor (200000, 500000], 5,572 boxes, 283, deficit 458.15 & A + B: 8,628 boxes, 383, deficit 763.50 \\
\hline
naive, h\_min \textless{} 4 log p (Prop. 24) & 1,528 · \textbf{0} · 207.8 · 68.05 \% & 1,874 · \textbf{0} · 255.6 · 55.78 \% & 3,402 · \textbf{0} · 463.4 · 60.7 \% \\
canonical, ĥ\_min \textless{} 4 log p + β(E) (Prop. 24ᶜ) & 2,192 · \textbf{0} · 307.1 · \textbf{100.6 \%} & 3,251 · \textbf{0} · 456.8 · \textbf{99.70 \%} & 5,443 · \textbf{0} · 763.9 · \textbf{100.05 \%} \\
canonical with β\_Λ (remark (g)) & 2,221 · \textbf{0} · 311.7 · 102.07 \% & 3,320 · \textbf{0} · 467.8 · 102.11 \% & 5,541 · \textbf{0} · 779.5 · 102.09 \% \\
either & 2,215 · 0 · 310.1 · 101.6 \% & 3,294 · 0 · 462.8 · 101.0 \% & 5,509 · 0 · 772.9 · 101.2 \% \\
naive only & 23 · 0 · 3.0 · 1.0 \% & 43 · 0 · 6.0 · 1.3 \% & 66 · 0 · 9.0 · 1.2 \% \\
control: ĥ\_min \textless{} 4 log p (β := 0) & 2,239 · \textbf{15} · 311.0 & 3,366 · \textbf{28} · 471.6 & 5,605 · \textbf{43} · 782.6 \\
above the canonical threshold: observed / Σ1/p & 100 / 98.2 = 1.018 & 283 / 284.4 = 0.995 & 383 / 382.6 = 1.001 \\
\hline
\end{tabular}
\end{small}\end{center}

The same canonical cut on the two sets C and D described below, which lie outside Cremona's tables (the naive cut was not computed there):

\begin{center}\begin{small}
\begin{tabular}{>{\raggedright\arraybackslash}p{0.349\linewidth}>{\raggedright\arraybackslash}p{0.27\linewidth}>{\raggedright\arraybackslash}p{0.338\linewidth}}
\hline
cut & set C: 5,782 boxes, 812 with det ≡ 0, Σ1/p 830.3 & set D: 7,861 boxes, 788, Σ1/p 1,087.4 \\
\hline
canonical, ĥ\_min \textless{} 4 log p + β(E) (Prop. 24ᶜ) & 695 · \textbf{0} · 97.8 & 2,039 · \textbf{0} · 317.3 · 106.0 \% of the deficit 299.4 \\
canonical with β\_Λ (remark (g)) & 715 · \textbf{0} · 101.0 & 2,063 · \textbf{0} · 321.3 \\
above the canonical threshold: observed / Σ1/p & 812 / 732.5 = 1.109 & 788 / 770.1 = 1.023 \\
\hline
\end{tabular}
\end{small}\end{center}

(On set C the deficit, 18.3, is too small for a share to mean anything: there the frequency above the threshold is above 1/p, at p = 5 alone, 527 against 445.2.)

Controls: with β(E) replaced by 0 the canonical cut lets 15 and 28 boxes with det(χ, γ) ≡ 0 below the threshold; permuting β(E) among the boxes (1,000 times) lets 8 to 24 and 26 to 50, never 0; a constant β (the median, 0, or −0.5·log p) gives 15, 15 and 4 on set A and 26, 28 and 17 on set B. Only the curve's own β(E) gives 0. On every enumerated non-zero element of Λ in all 8,628 boxes R(Q) ≥ β(E) holds; dropping the component corrections breaks it in 2,833 and 5,176 boxes. The minimum m\_∞ is taken on a grid with golden-section refinement and checked by a second route through the x-coordinate (all of set A; 412 boxes of set B, all those within 0.05 log p of the threshold): the second route never goes below the first. Since a sampled minimum is at least the true one, the computed β(E) is at least the true one, so "no box with det(χ, γ) ≡ 0 below the threshold" cannot be produced by the grid missing a narrow valley; what the error could inflate is the share. The smallest margins are 7.5·10⁻⁵ log p on the side of det(χ, γ) ≡ 0 (434832b1, p = 5) and 1.6·10⁻⁴ log p on the explained side (106947h1). At I\_n* with c\_ℓ = 2 (205 and 524 boxes) the rational component was not determined and the smaller correction was used, which can only lower β(E). By p, the ratio above the canonical threshold over A + B is 0.91, 1.02, 0.96, 1.10, 1.46, 1.16 for p = 5, 7, 11, 13, 17, 19 (139/153.0, 117/115.3, 32/33.4, 49/44.6, 18/12.3, 28/24.1).

\textit{Just above the threshold: an effect whose cause is not known.} Stratified by (ĥ\_min − β)/log p over A + B: [4, 4.5) 118/97.2 = 1.21, [4.5, 5) 51/63.0 = 0.81, [5, 6) 81/84.3 = 0.96, [6, 8) 76/80.3 = 0.95, [8, 12) 37/43.8 = 0.85, ≥ 12 20/14.1 = 1.41. The excess just above the threshold is 1.46 on set A and 1.10 on set B (Poisson z = +2.6 and +0.8). To decide whether it is noise we fixed, before computing, two further sets and a rule (Poisson z \textgreater{} 2 on the set alone: an effect; z \textless{} 1: noise; for the last set also the premise that above the threshold the frequency is within z = ±2 of 1/p at each p). Set C: rank-two curves y² = x³ + Ax + B through two small points whose x-coordinates differ by 1. Set D: rank-two curves y² = x³ + Ax + B with random |A| ≤ 1000, |B| ≤ 10⁴ on which two independent points of naive height at most 10⁴ were found. On C and D the rank is fixed by ellrank with equal bounds, and the generators are saturated at all primes ≤ 200, which we assume to be enough; neither set meets A, B or the other. In the stratum [4, 4.5) (observed det(χ, γ) ≡ 0 against Σ1/p, Poisson z):

\begin{center}\begin{small}
\begin{tabular}{>{\raggedright\arraybackslash}p{0.054\linewidth}>{\raggedright\arraybackslash}p{0.214\linewidth}>{\raggedright\arraybackslash}p{0.074\linewidth}>{\raggedright\arraybackslash}p{0.234\linewidth}>{\raggedright\arraybackslash}p{0.207\linewidth}>{\raggedright\arraybackslash}p{0.133\linewidth}}
\hline
set & conductor & split boxes & [4, 4.5), threshold with β(E) & with β\_Λ(E) & above the threshold \\
\hline
A & ≤ 200000 (Cremona) & 3,056 & 45 / 30.8 = 1.46 (+2.6) & 46 / 30.1 = 1.53 (+2.9) & 1.02 \\
B & (200000, 500000] (Cremona) & 5,572 & 73 / 66.4 = 1.10 (+0.8) & 80 / 65.4 = 1.22 (+1.8) & 0.99 \\
A + B & ≤ 500000 & 8,628 & 118 / 97.2 = 1.21 (+2.1) & 126 / 95.5 = 1.32 (+3.1) & 1.00 \\
C & 512,144 to 5.8·10¹², median 1.1·10¹⁰ & 5,782 & 99 / 50.7 = 1.95 (+6.8) & 100 / 50.7 = 1.97 (+6.9) & 1.11 \\
D & 508,528 to 1.0·10¹¹, median 3.5·10⁹ & 7,861 & 94 / 51.8 = 1.82 (+5.9) & 96 / 50.2 = 1.91 (+6.5) & 1.02 \\
\hline
\end{tabular}
\end{small}\end{center}

On D the premise holds (|z| ≤ 1.3 at each p). On C the frequency above the threshold is 1.18 times 1/p at p = 5, which is not seen on D (0.98); on C the excess in [4, 4.5) remains when the positions above the threshold are permuted within each p (z = +6.6). By the rule the excess is an effect on C and on D. It is not traced to the archimedean place, nor to the component corrections once the curves are matched. By Lemma 25 the frequency 1/p is exact for the local classes, so the effect is a correlation between the class of a box modulo E\_2 and its heights.

Whether the size of the effect — 1.21 on A + B against 1.95 and 1.82 on C and D — follows the conductor, β(E)/log p or the way the set was built was tested by a second rule fixed before counting: strata by the tertiles of the conductor within D, the halves of the conductor within A + B, and the tertiles of β(E)/log p over all four sets; a score test for a log-linear trend in a Poisson model within D and within the β-strata (summed over the sets), a conditional binomial test for the two halves of A + B, and a Mantel–Haenszel-type stratified conditional binomial test for C + D against A + B within the β-strata. In [4, 4.5) the ratios are 1.60, 1.93, 1.95 over the conductor tertiles of D (trend z = +0.8) and 1.22, 1.21 over the halves of A + B (z = −0.0); the β-trend gives z = +1.5 and C + D against A + B within the β-strata z = +0.8 (with β\_Λ: +0.7, +0.2, +0.9, +1.2). None reaches |z| = 2, so by that rule the size of the effect is attributed neither to the conductor nor to the construction of the sets. The rule has little power, and we do not read the outcome as an absence: on synthetic data in which the ratio moves log-linearly in log N from 1.21 at the median conductor of A + B to 1.88 at that of C + D, it detects the conductor in 14 \% of runs, since within D the conductor moves by a factor of about 6.6 between the tertile boundaries and between the medians of A + B and C + D by about 1.8·10⁴; and the β-strata hardly overlap (82 \% of the boxes of A + B lie in the lowest, 98 \% of those of C + D in the two upper ones). Matched on β(E)/log p ∈ [0.4, 1.3] instead — all 1,378 such boxes from Cremona's tables, which end at conductor 500000 (Σ1/p = 20.3 in [4, 4.5)), against curves built as for D until Σ1/p reached 50 there (10,727 boxes) — the ratios are 30/20.3 = 1.48 and 76/50.1 = 1.52, and a further rule fixed before counting (a Mantel–Haenszel-type test within three sub-strata of β(E)/log p) gives z = −0.1, rate ratio 0.97 with an approximate 95 \% interval from 0.58 to 1.62 (with β\_Λ, z = +0.1), which on synthetic data detects a ratio of 1.2 between the two constructions in at most 9 \% of runs and so decides nothing; no box of either set has det(χ, γ) ≡ 0 below the threshold. Within the curves built as for D alone, which can be multiplied at will, a third rule fixed before counting divided the boxes with β(E)/log p ∈ [0.4, 2.0] into four layers (boundaries 1.0, 1.3, 1.6) and tested for a trend with a Cochran–Armitage-type score test stratified by p and by the range of (A, B), after adding 10,355 boxes from curves with 10⁴ \textless{} |B| ≤ 2·10⁴ that had not been looked at: the ratios in [4, 4.5) are 1.31, 1.58, 1.61 and 1.84 (Σ1/p = 26.7, 43.8, 27.9, 46.3), the trend gives z = +1.2 (with β\_Λ, +1.0) and the ratio of the top layer to the bottom one is estimated at 1.26, with a profile-likelihood 95 \% interval from 0.86 to 1.88. On synthetic data a ratio moving linearly from 1.4 to 1.9 across β(E)/log p ∈ [0.4, 2.0] is detected in 10 \% of runs at this size, and 80 \% would need Σ1/p ≈ 1,300 in the band, so this rule decides nothing either: the excess is not shown to follow β(E)/log p, nor to be independent of it. On the new curves alone the excess recurs (93 against Σ1/p = 59.0, a ratio of 1.58) and the trend is z = +0.1; none of the 6,156 boxes built as for D (2,328 of them new) has det(χ, γ) ≡ 0 below the threshold. A fourth rule fixed before counting compared two predictions without free parameters on 11,008 further boxes built as for D from curves with 10⁴ \textless{} |B| ≤ 2·10⁴ whose (A, B) had not been drawn before. The first is a constant multiple of Σ1/p, with the constant 1.62 taken from the earlier boxes built as for D. The second splits the band into boxes with det(χ, γ) ≡ 0 whose shortest non-zero point Q of Λ ∩ E\_1(Q\_p) has p-adic depth 2 and an excess s \textless{} 0.5 outside p, and boxes with det(χ, γ) ≢ 0 whose Q has depth 1 and s ∈ [2, 2.5) (on all the boxes measured, det(χ, γ) ≢ 0 has always come with depth 1), and predicts the first kind from the boxes outside the band, assuming that s \textless{} 0.5 is equally frequent in both classes and that depth 2 has frequency 1 − 1/p. The band holds 111 boxes with det(χ, γ) ≡ 0 against 91.3 and 116.9 predicted. The Poisson log-likelihood ratio is −14.6, in favour of the constant multiple; the split prediction fails its own resampling (p \textless{} 10⁻⁴) and its fit by layers of β(E)/log p (p = 0.03), while the constant multiple fits by layers (p = 0.19) but falls short in total (z = +2.1; the ratio on these curves is 1.97). So the split does not account for the excess. The ratio in [4, 4.5) differs between sets of curves (1.37 to 1.97), but over the six sets — A + B, C, and four sets built as for D, two with |B| ≤ 10⁴ and two with 10⁴ \textless{} |B| ≤ 2·10⁴ — it is consistent with a single constant (pooled ratio 1.73, Poisson χ² = 7.4 on 5 degrees of freedom, p = 0.19) and shows no trend in |B| (p = 0.41) [computation, outside the rules]; none of the 2,449 new boxes below the threshold has det(χ, γ) ≡ 0. A fifth rule fixed before counting tested the frequency 1 − 1/p of depth 2 among the boxes with det(χ, γ) ≡ 0 at p = 5, which had come out higher on every earlier set: on 8,181 further boxes at p ∈ \{5, 7\} from curves whose (A, B) had not been drawn before, 391 of the 454 boxes with p = 5 and det(χ, γ) ≡ 0 have depth 2 (0.861 against 0.800; one-sided exact binomial test, p = 4.3·10⁻⁴), and at p = 7 461 of 511 (0.902 against 0.857). The frequency 1 − 1/p rests on two assumptions about the map θ : Λ ∩ E\_2(Q\_p) → E\_2(Q\_p)/E\_3(Q\_p) ≅ F\_p, through which depth 2 is decided: that θ is uniformly distributed, and that it is independent of the line of Q modulo p. On these boxes the first is not contradicted (θ = 0 in 17 of 454 against 1/25, and no dependence on (A, B) mod 25 is seen, χ² p = 0.49), while the second fails: the kernel of θ is the line of Q in 46 of the 437 boxes with θ ≠ 0, against 1/6. A point of depth 3 carries at least 2 log p more in ĥ than one of depth 2, so choosing Q by ĥ pushes the kernel of θ off the shortest line; subtracting that amount raises the count from 46 to 58, against 72.8 expected [computation, outside the rules]. None of the 2,330 new boxes below the threshold has det(χ, γ) ≡ 0. A sixth rule fixed before counting chose instead the element Q′ of (Λ ∩ E\_1(Q\_p)) ∖ \{0\} that minimises ĥ − 2(d − 2) log p, with d the p-adic depth — the canonical height with the contribution 2d log p of p to the denominator removed, up to a constant — and asked whether depth at least 3 then has frequency 1/p among the boxes with det(χ, γ) ≡ 0 at p = 5. The depth was read from the p-adic valuation of the formal logarithm on a basis of Λ ∩ E\_1 and checked against the denominators, and the minimum was taken over all elements of depth below 40. On 8,048 further boxes at p ∈ \{5, 7\} from curves whose (A, B) had not been drawn before, 80 of the 433 boxes with p = 5 and det(χ, γ) ≡ 0 have Q′ of depth at least 3 (0.185 against 0.200; one-sided exact binomial tests: P(K ≤ 80) = 0.23 under 1/p, and P(K ≥ 80) = 0.005 under the frequency 0.139 of the earlier boxes with Q chosen by ĥ), and the verdict of the rule is that the frequency is consistent with 1/p; at p = 7 it is 71 of 535 (0.133 against 0.143). On these boxes Q chosen by ĥ itself has depth at least 3 in 69 of 433 (0.159), so the rule separates the two choices by 11 boxes. Over all six sets built as for D the frequency with Q′ is 347 of 1,973 (0.176, P(K ≤ 347) = 0.004 under 1/p) against 287 with Q: removing the depth from the height recovers about half of the deviation, and the rest is not explained [computation, outside the rules; the earlier sets include the one on which the explanation was found]. The threshold and the band do not depend on this choice. When det(χ, γ) ≡ 0 every non-zero element of Λ ∩ E\_1 lies in E\_2, so the argument of Proposition 24ᶜ applies to every such element and none of them, Q′ included, lies below 4 log p + β(E) [paper, under the hypotheses of Proposition 24ᶜ]; and over the six sets Q′ differs from Q in 97 of the 5,351 boxes with det(χ, γ) ≡ 0 and β(E)/log p ∈ [0.4, 2.0], and in one of the 48,168 boxes with det(χ, γ) ≢ 0, three boxes leave [4, 4.5) and none enters, and the ratio there is 1.81 with either choice (593/327.3 and 590/326.7) [computation]. So the excess just above the threshold is not an artefact of choosing the shortest element by ĥ. A seventh rule fixed before counting looked at the shape of Λ ∩ E\_1 instead. With b₁, b₂ a Gauss-reduced basis of Λ ∩ E\_1(Q\_p) for the quadratic form ĥ (b₁ is the element Q above) and r = ĥ(b₂)/ĥ(b₁), the boxes with det(χ, γ) ≡ 0 at p = 5 were divided into five classes of r with boundaries 1.1, 1.3, 1.7 and 2.5, and in each class the frequency of θ(b₁) = 0, that is of depth at least 3 for b₁, was compared with 1/p by an exact two-sided binomial test, with Holm's correction over the five classes. The hypothesis written down beforehand, from the shape of the remaining deviation on the earlier sets, was that only the classes with r near 1 would deviate. On 23,927 further boxes at p ∈ \{5, 7\} from curves whose (A, B) had not been drawn before, the 1,334 boxes with p = 5 and det(χ, γ) ≡ 0 give 0.161, 0.193, 0.203, 0.168 and 0.050 from the first class to the last, and only the last (r ≥ 2.5; 26 of 525, p = 5·10⁻²³) is rejected: by the rule the deviation is not confined to r near 1, and it sits at large r. The frequency of θ(b₂) = 0 is consistent with 1/p in every class (247 of 1,334 in all), and among the boxes with θ ≠ 0 the kernel of θ is the line of b₂ as often as 1/(p + 1) predicts in every class (86 against 86.3 in the last), so the deviation is not an exchange of the two shortest lines. At p = 7 the pattern is the same (only the last class deviates, 68 of 707). A secondary rule, added before these boxes were measured, reads this through Proposition 24ᶜ: under its hypotheses every non-zero element v of Λ ∩ E\_1 has ĥ(v) ≥ β(E) + 2d(v) log p, so depth at least 3 for b₁ requires X(b₁) := (ĥ(b₁) − β(E))/log p ≥ 6 [paper, under the hypotheses of Proposition 24ᶜ], and a null model in which θ is independent of the lattice is not compatible with that inequality. On these boxes no box with X(b₁) \textless{} 6 has θ(b₁) = 0 (0 of 495 at p = 5, 0 of 332 at p = 7), and on the boxes with X(b₁) ≥ 6 the frequency is 174 of 839 (0.207, two-sided p = 0.60; at p = 7, 178 of 1,225), with no class rejected; the deficit against 1/p, 92.8 boxes, is of the size of the mass 495/5 = 99.0 that the inequality excludes. Since the determinant of ĥ on Λ ∩ E\_1 is p²·Reg(Λ), a large r goes with a small ĥ(b₁), which is why the deficit sits in the last class (361 of its 525 boxes have X(b₁) \textless{} 6) [computation; a rule fixed before counting and a secondary rule fixed before measuring; this reading accounts for the counts and is not a theorem]. As an observation outside the rules, the element Q′ of the sixth rule has depth at least 3 in 207 of these 1,334 boxes (0.155), so the agreement with 1/p found by the sixth rule did not recur; on the boxes with X(b₁) ≥ 6 it is 205 of 839 (0.244) [computation]. An eighth rule fixed before counting asked whether the boxes with det(χ, γ) ≡ 0 are pressed against their lower bound more than the others. Let s := X(b₁) − 2d(b₁), with d the p-adic depth; under the hypotheses of Proposition 24ᶜ, s ≥ 0 [paper, under the hypotheses of Proposition 24ᶜ], and the bound sits at X(b₁) = 4 (depth 2) when det(χ, γ) ≡ 0 and at X(b₁) = 2 (depth 1) when det(χ, γ) ≢ 0. The rule compared ρ := \#\{s ∈ [0, 0.5)\}/\#\{s ∈ [0.5, 1)\} between the two classes by Fisher's exact test (two-sided, level 0.05), with p = 5 and p = 7 together. On 13,920 further boxes at p ∈ \{5, 7\} from curves y² = x³ + Ax + B with |A| ≤ 1000 and |B| ≤ 10⁴ whose (A, B) had not been drawn before (β(E)/log p ∈ [0.4, 2.0]: 1,709 boxes with det(χ, γ) ≡ 0 and 12,133 with det(χ, γ) ≢ 0), ρ is 2.57 (296 against 115) and 2.44 (2,017 against 826), p = 0.68, and the verdict of the rule is that there is no difference: both classes crowd just above their own bound to the same degree. The same holds at p = 5 and at p = 7 separately (2.70 against 2.77, 2.32 against 2.06) and after dividing s by the width of the window up to Hermite's bound √(4/3·det) for ĥ(b₁) (2.02 against 2.40, p = 0.07), and a split of the boxes with det(χ, γ) ≢ 0 by the parity of the draw is not rejected either; over the whole range s \textless{} 4 the two distributions do differ (χ² p = 0.012, after the fact). Outside the rules: every box with s \textless{} 0.5 has a point b₁ whose x-coordinate has a power of p as denominator, which the hypotheses of Proposition 24ᶜ force, since a further prime in the denominator adds at least 2 log 2/log p \textgreater{} 0.5 to s [paper, under the hypotheses of Proposition 24ᶜ]; near 0 the fraction of the boxes with s \textless{} 0.5 that have s \textless{} a is 0.005, 0.013, 0.044 and 0.13 for a = 10⁻⁵, 10⁻⁴, 10⁻³ and 10⁻² (det(χ, γ) ≢ 0), close to √(2a), as when λ′\_∞ has a non-degenerate minimum on E(R) near which the point is smoothly distributed; the smallest s is 5·10⁻⁸, so the threshold is approached to that precision; in [4, 4.5) the ratio on these boxes is 2.13 (245 against Σ1/p = 115.2); and none of the 4,325 boxes below the threshold has det(χ, γ) ≡ 0 [computation, outside the rules; the square-root reading is not proved]. So on these boxes the crowding just above the threshold is not a property of the boxes with det(χ, γ) ≡ 0: it reads as a crowding common to both classes just above the bound of each element, with the bound for det(χ, γ) ≡ 0 sitting at X(b₁) = 4, where the boxes with det(χ, γ) ≢ 0 are already two units above theirs; this accounts for where the excess sits and not for its size [computation; a rule fixed before counting; this reading is not a theorem]. A ninth rule fixed before counting predicted the coefficient of this square-root shape curve by curve from the archimedean place. Write s = arch + nonp + bad, where arch = (λ′\_∞(b₁) − m\_∞)/log p, nonp = 2 log d′/log p with d′ the part of the denominator of x(b₁) prime to p, and bad is the excess of the component corrections at b₁ over their minima; under the hypotheses of Proposition 24ᶜ all three are ≥ 0 [paper, under the hypotheses of Proposition 24ᶜ]. Parametrise each real component of E(R) by z = tω₁ (or tω₁ + ω₂/2 on the second component, when there is one), t ∈ R/Z. If λ′\_∞ has non-degenerate minima t\_j at the level m\_∞ on the component of b₁, with second derivatives c\_j in t, and t(b₁) is uniformly distributed on that component, then P(s ≤ a) ≈ C(E)√a for small a, with C(E) = Σ\_j mult\_j · 2√(2 log p/c\_j), where mult\_j = 1 if t\_j ∈ \{0, 1/2\} and 2 otherwise [paper, by the Morse lemma in one variable; the uniformity of t(b₁) is an assumption]; a point on a component whose minimum lies above m\_∞, or with nonp + bad \textgreater{} a, has P(s ≤ a) = 0. With q = e\textasciicircum{}\{2πiτ\} and τ = ω₂/ω₁ normalised to Im τ \textgreater{} 0 and Re τ ∈ \{0, 1/2\}, the second derivative at the 2-division point t = 1/2 of the identity component is 2π²(1 + 8 Σ\_\{n≥1\} q\textasciicircum{}n/(1 + q\textasciicircum{}n)²), and at t = 1/2 on the second component it is 16π² Σ\_\{k≥0\} q\textasciicircum{}\{k+1/2\}/(1 + q\textasciicircum{}\{k+1/2\})² [paper, from the q-expansion of the local height, whose normalisation we have not checked against the source; the two expressions agree with numerical second differences on 198 components to a relative 3·10⁻⁶]. The minimum at t = 1/2 of the identity component degenerates where the first expression vanishes (q ≈ −0.107); near there the shape is a fourth root rather than a square root. The square-root law also accounts for the √(2a) above: if P(s ≤ a) = C√a up to a = 0.5, then P(s ≤ a ∣ s \textless{} 0.5) = √(2a) whatever C is, so the information about the curve lies in C and not in the 2. The rule took the boxes with C(E) \textgreater{} 0 and no flag of degeneracy and compared three things: k, the ratio of the observed to the predicted least-squares coefficient of P(s ≤ a) on √a over a ∈ [10⁻⁴, 10⁻²]; g, the bias of C(E) among the boxes with s ≤ 10⁻² relative to the bias the prediction implies (1 if the prediction holds, 0 if s does not depend on C(E)); and a two-sided Mann–Whitney test of C(E) between the boxes with s ≤ 10⁻² and the others. The verdict that the prediction holds required the test to reject at level 0.05 in the predicted direction and the 95 \% bootstrap intervals of k and g both to contain 1. On 10,346 further boxes at p ∈ \{5, 7\} from curves with |A| ≤ 1000 and |B| ≤ 10⁴ whose (A, B) had not been drawn before (10,280 with β(E)/log p ∈ [0.4, 2.0], of which 1,523 have C(E) \textgreater{} 0 and no flag, and 211 of these have s ≤ 10⁻²), k = 1.10 (0.96 to 1.25), g = 0.76 (0.42 to 1.10) and the test gives p = 6·10⁻⁶, and the verdict of the rule is that the prediction holds. Permuting C(E) among these boxes removes the dependence (p = 0.80; over 1,000 permutations the test rejects in 5.0 \% of them), and none of the 8,646 boxes predicted to have P(s ≤ 10⁻²) = 0 has s ≤ 10⁻². On the 13,920 boxes of the eighth rule, whose square-root shape had already been seen, the same quantities are k = 1.04 and g = 1.16 [computation; numerical m\_∞ and minima; a rule fixed before counting; the uniformity of t(b₁) is assumed, not proved]. Outside the rules, the same uniform measure, applied to each of the 1,123 boxes with det(χ, γ) ≡ 0 and b₁ of depth 2 with its own discrete data (whether the denominator of x(b₁) is a power of p, the component corrections, the component), predicts 189 of them in [4, 4.5) against 192 observed, a ratio of 2.05 against 2.07: given those discrete data, the size of the excess is accounted for by the archimedean place, while how often b₁ is a point whose denominator is a power of p, and why det(χ, γ) ≡ 0 occurs with frequency 1/p above the threshold, are not [computation, outside the rules, after the fact on the same boxes]. A tenth rule fixed before counting asked whether the departure of t(b₁) from uniformity comes from the competition between b₁ and the second minimum of K₁. Since b₁ is primitive, b₁ is a shortest vector if and only if ĥ(b₁) ≤ ĥ(b₂), where b₂ completes a reduced basis; writing ĥ(b₁) = λ′\_∞(t) + N₁, where N₁ collects the logarithm of the denominator of x(b₁) and the component corrections and does not depend on t, this says λ′\_∞(t) ≤ ĥ(b₂) − N₁ [paper; an identity]. The model M takes t uniform on its component conditioned only on this inequality, with N₁ and ĥ(b₂) held fixed [an assumption]. Under M the fraction V/L is uniform on [0, 1], where V and L are the lengths of the parts of the component on which λ′\_∞ is at most λ′\_∞(t(b₁)) and at most ĥ(b₂) − N₁; and on the boxes of the ninth rule the coefficient C(E) is multiplied by 1/L [paper, under M]. Before counting we found that on those boxes b₁ is much shorter than b₂ (ĥ(b₂)/ĥ(b₁) ≥ 2.5 for 94 \% of them), so that L is within 10⁻³ of 1 for at least 90 \% of them and M changes the predicted coefficient by 0.1 \%: M cannot account for k = 1.10 and g = 0.76, and the rule therefore also tested V/L on all boxes, in the five classes of ĥ(b₂)/ĥ(b₁) used above. On 12,868 further boxes at p ∈ \{5, 7\} from curves with |A| ≤ 1000 and |B| ≤ 10⁴ whose (A, B) had not been drawn before (12,807 with β(E)/log p ∈ [0.4, 2.0], 1,890 of them with C(E) \textgreater{} 0 and no flag, 263 of these with s ≤ 10⁻²), the coefficient gives k = 1.14 (1.00 to 1.28) and g = 0.61 (0.31 to 0.92), against 0.2 \% from M, and the test of V/L rejects in the classes with ratio below 1.1 and at least 2.5; the verdict of the rule is that neither part holds. The second part of the rule was mis-specified: the classes are defined through ĥ(b₁), which depends on t, so that even under M the class with ratio near 1 selects large V/L and the class with large ratio selects small V/L. With the reference corrected for this (V uniform on the interval allowed by the class), on the new boxes the four classes with ratio below 2.5 agree with M and only the class with ratio at least 2.5 does not; cutting instead by L, which does not depend on t, the boxes with L \textless{} 0.8 agree with M on both sets of boxes (and those with 0.8 ≤ L \textless{} 0.95 on the new boxes), while those with L within 10⁻³ of 1 have t closer to the minimum of λ′\_∞ than uniform (mean of V/L 0.475 and 0.477 on the new boxes and on those of the ninth rule) [computation; numerical m\_∞; the corrected reference and the cut by L were made after the rule was fixed and after the data were seen]. The competition with b₂ thus describes t(b₁) where b₂ is close, but the departure in the coefficient comes from boxes where b₂ is far, and its origin is not known [computation]. The cause of the effect is not known [computation].

\textit{The same boxes cut by conductor, and cut by h\_min.} Cut by conductor the ratio is not monotone — 0.218, 0.316, 0.144, 0.279 on (0, 20000], (20000, 50000], (50000, 100000], (100000, 200000] — so the conductor is not the variable. The table below re-cuts three families of §5.6 and this section by u, the excess being k := v\_p(Reg\_p) − (baseline), the baseline max(r − 2, 0) in the split case and r − 2 in the non-split case of rank ≥ 2:

\begin{center}\begin{small}
\begin{tabular}{>{\raggedright\arraybackslash}p{0.199\linewidth}>{\raggedright\arraybackslash}p{0.07\linewidth}>{\raggedright\arraybackslash}p{0.05\linewidth}>{\raggedright\arraybackslash}p{0.05\linewidth}>{\raggedright\arraybackslash}p{0.05\linewidth}>{\raggedright\arraybackslash}p{0.05\linewidth}>{\raggedright\arraybackslash}p{0.05\linewidth}>{\raggedright\arraybackslash}p{0.05\linewidth}>{\raggedright\arraybackslash}p{0.175\linewidth}>{\raggedright\arraybackslash}p{0.114\linewidth}}
\hline
family & boxes & u \textless{} 2 & 2–4 & 4–6 & 6–8 & 8–12 & ≥ 12 & all & P(k ≥ 2 ∣ k ≥ 1): observed / Σ1/p \\
\hline
rank 1, split (conductor ≤ 12000) & 1,006 & — & 0.70 & 0.86 & 0.47 & 0.92 & 1.13 & \textbf{1.00} (135/135.2) & 21 / 20.9 \\
rank 2, non-split (conductor ≤ 200000) & 24,299 & \textbf{0.92} & 0.98 & 1.00 & 1.07 & 0.95 & 0.77 & \textbf{0.97} (2,489/2,556.5) & 352 / 325.1 \\
rank 2, split (conductor ≤ 200000) & 3,056 & \textbf{0.00} & \textbf{0.06} & 0.95 & 0.90 & 0.53 & 1.16 & 0.25 (100/405.4) & 25 / 13.6 \\
\hline
\end{tabular}
\end{small}\end{center}

(In this table h\_min is taken over all non-zero elements of Λ ∩ E\_1, as in Proposition 24, and not only outside pΛ as in the table above; on the split boxes of rank two the two cuts give the same ratios to two decimals. Rank one has no stratum u \textless{} 2, since there ker χ = pΛ and h\_min = p²ĥ(P).) The first row reproduces the rank-one distribution 871 / 114 / 17 / 4 of sweep E, and the second the excesses 21,810 / 2,137 / 306 / 39 / 7 for k = 0, …, 4. \textbf{An excess that asks for no denominator follows 1/p in every stratum}: the non-split excess is det T ≡ 0, a condition on a Schur complement and not on a denominator, and it sees no archimedean obstruction. The obstruction acts on the split boxes of rank two at k ≥ 1 alone. The last column of the split row (25 against 13.6) is not explained here.

\textit{Lemmas 23 and 23′ as counted, and the next layer.} Extending the sweep to conductor (200000, 500000] gives 50,707 boxes, 5,572 split: Lemma 23 is violated (D\_Q ≥ 2 with h(x(Q)) \textless{} 4 log p) in \textbf{0}, Lemma 23′ in 0 of 152,121 elements, det(χ, γ) ≡ 0 occurs in 283 against 270.4 expected above the canonical threshold u = 4 (ratio 1.047; above the naive threshold the expectation is about 493, ratio 0.57, from upper bounds; above the threshold 4 log p + β(E) of Proposition 24ᶜ, 283 against 284.4, ratio 0.995, and none below it among 3,251 boxes, table 5.12ᶜ), and D\_Q ≥ 3 holds for the shortest element in 14 boxes against 15.5 expected above 6 log p. That count depends on the element: 27 for another choice of element, 2 when three elements are required to agree, 100 when any one of three suffices — for D₁ = 2 the depth of Q + pP₁ drops to 2. The invariant of the box is min(D\_Q, D₁) ≥ 3, which holds in the two boxes below. Among the shortest candidates with s ≤ 2 log 2 (1,709 elements) all but 17 have c′ = 1, and those 17 have s = 2 log 2 exactly, the equality case of Lemma 23′. The archimedean excess max(0, log|x(Q)|) saturates near 5 as s grows, so the proportion with c′ = 1 falls from 97 \% for s \textless{} 2 to 0.9 \% for s ≥ 12. The index of the p-adic closure of Λ is p\textasciicircum{}m with m = min(D\_Q, D₁) − 1; over conductor ≤ 200000 m ≥ 2 occurs in no box (0.66 expected among the 308 boxes with u ≥ 6), and over conductor ≤ 500000 in exactly \textbf{two, 328563a1 and 477950bg1 at p = 7} (2.47 expected), both with D\_Q = D₁ = 3. At both, γ vanishes at every point computed (21 of 21 at 328563a1), γ₂ ≡ χ (mod 7), M\_Λ mod 7 is [[1, 3], [3, 2]] and [[5, 3], [3, 6]] of rank one, and v₇(Reg₇) = 1, as the paragraph on small closures in §3.15 predicts; at every point with χ(P) ≠ 0 among the 1,008 pairs of points computed there, v₇(n²h\_p(P)) = 0.

\textit{Corollaries R′\_s and K′(i) in rank two.} Over the split boxes of rank two with conductor ≤ 20000 — 111 boxes (27, 31, 19, 16, 3, 15 by p; the smallest conductors 4407, 4629, 4976, 4279, 2343, 3333) — ρ₁ was measured in 109, the two boxes with a generator in E\_1(Q\_5), 7387a1 and 8962b1 at p = 5, being left out of that sweep; in those 109, v\_p(Reg\_p) = 0 ⟺ ρ₁ = 2 holds in \textbf{109 of 109}, and ρ₁ = 2 ⟺ (χ, γ) onto in \textbf{109 of 109}, (χ, γ) being computed with no p-adic height. At 7387a1 and 8962b1, (χ, γ) is onto (D\_Q = 1) and v\_p(Reg\_p) = 0, so v\_p(Reg\_p) = 0 ⟺ (χ, γ) onto holds in \textbf{111 of 111}; the conclusion is unchanged. The three boxes where (χ, γ) is not onto are 8577a1 (p = 5) and 12208i1 (p = 13), with ρ₁ = 1 and v\_p(Reg\_p) = 1, and 16011e1 (p = 5), with ρ₁ = 0 and v\_p(Reg\_p) = 2: the three steps r − 2, r − 1, r of Corollary R′\_s all occur.

\textit{The p-adic L-function at the split boxes, and a positive control.} By §6(c), det(χ, γ) ≡ 0 in a split box of rank two satisfying (H) forces the coefficient c\_2 of T² in L\_p(E, T) to vanish modulo p, by Kato's divisibility alone and with no assumption on Ш; this is a test with a known answer of the computation of det(χ, γ) and of c\_2. A rule fixed before counting computed c\_2 modulo p as Σ\_j (j choose 2)·S\_j (§6) on two sets of boxes: the 3,056 split boxes of rank two above, and the 1,674 boxes with conductor at most 3·10⁷ among the boxes of the eighth, ninth and tenth rules (p ∈ \{5, 7\}, curves with |A| ≤ 1000 and |B| ≤ 10⁴), the conductor bound being fixed by the running time and not by det(χ, γ) or c\_2. The symbols S\_j were obtained from the p − 1 twisted values L(E, ψ, 1), ψ of order p and conductor p², by an approximate functional equation with root number w\_E·ψ(N)G(ψ)²/p²; on 14 pairs of rank two with conductor below 950 this agrees with the modular symbols of PARI to 10⁻¹⁴ and its prediction of c\_2 modulo p with \texttt{mspadicseries} in 14 of 14, and on 120 of the boxes it agrees with PARI's \texttt{lfuntwist} to 5·10⁻¹⁴, with no disagreement in c\_2 modulo p. At 446d1 and p = 5 the coefficient of T² of P\_5 is 2·5 + 2·5² + 5³ + 5⁴ modulo 5⁴, twice the value 5 + 5² + 3·5³ of [SW, §3.3] (the factor 2 is the ratio of the Néron period to E.omega[1]). (H) holds in 2,987 and 1,656 of the boxes (E[p] is irreducible in all but 15, and some q ∥ N has p ∤ ord\_q(Δ\_min) in all but 78); in every box c\_1 ≡ 0 (mod p), as rank two forces, the S\_j are integers, the root number is +1 and the given generators are saturated at p. \textbf{Among the boxes satisfying (H), det(χ, γ) ≡ 0 occurs in 95 and 205, and c\_2 ≡ 0 (mod p) in all 300}; c\_2 ≡ 0 occurs in none of the 2,892 and 1,451 boxes with det(χ, γ) ≢ 0, against about 1/p of them if c\_2 were a random residue [computation, floating point, no interval arithmetic]. The first count is the positive control, and it passes. The second is not a theorem: under (H) and the equality of Skinner and Urban it says that p ∤ \#Ш(E/Q) in these boxes, which is expected and is not proved by it. What it does prove, by §6(c) and with the computed value of c\_2 as the only numerical input, is the converse reading box by box: \textbf{for each of the 4,343 boxes with (H) and c\_2 ≢ 0 (mod p), rank E(Q) = 2, Ш(E/Q)[p\textasciicircum{}∞] = 0 and v\_p(Reg\_p) = 0} [known theorems combined, the statements checked in [Kato04], [PR92], [Sch85], [SU] and [Ski] (§6); computation for c\_2]. The value v\_p(Reg\_p) = 0 obtained in this way agrees in all of them with the value given by det(χ, γ) ≢ 0 through Corollaries K′(i) and R′\_s. A second rule asked that c\_3 vanish modulo p with frequency about 1/p in the boxes with det(χ, γ) ≡ 0, as a negative control; it vanished in all 300, because S\_j = S\_\{c−j\} with c ≡ −log\_p⟨N⟩/log\_p(1 + p) (mod p) in every box — the functional equation of L\_p seen modulo p — so that 2c\_3 ≡ (c − 2)c\_2 (mod p): the control was mis-designed and c\_3 is not independent of c\_2 [computation; the symmetry was found after the data were seen].

\subsection*{5.13 Propositions 21 and 22 on rational points}

\textit{Proposition 21 as displayed.} Five curves of rank two at p = 5 — 4407a1, 5956a1, 8577a1, 75393a1, 188032i1 — with P = u·g (|u| ≤ 2, χ(P) ≠ 0), Z = w·g (|w| ≤ 2, w ≠ 0) and a = 1, 2: F is read off the denominator of x(n(P + p\textasciicircum{}aZ)), C(P) from φ\_p′, φ\_p and 2y + a₁x + a₃ at P′, and ζ from p·nZ in E(Q\_p), the congruence being assembled from these raw parts by a separate program. The displayed congruence holds in \textbf{2,400 of 2,400} pairs at a = 1 and \textbf{2,400 of 2,400} at a = 2. Dropping the term C(P)ζ(P), or reversing its sign, leaves 560 of 2,400 — the pairs with χ(Z) = 0, where both sides vanish, the pairs with ζ(P) = 0, and coincidences at the rate 1/p. Steps (1) and (2) of the proof, Θ(P + S) − Θ(P) ≡ C(P)ζ(S) and q(P + S) − q(P) − q(S) ≡ C(P)ζ(S) for S ∈ ker χ, hold in 400 of 400. In rank one with Z = P, F((1 + p)P) ≡ F(P) + (q(P) − C(P)ζ(P)/2) (mod p) holds in \textbf{643 of 643} boxes of the filter of §5.10 and F((1 + p²)P) ≡ F(P) + p(q(P) − C(P)ζ(P)/2) (mod p²) in \textbf{574 of 574} (the boxes with coordinates too large for the computation left out), against 96 of 643 with the term C(P)ζ(P) dropped; ζ(P) = 0 in exactly the 97 boxes with e ≥ 3.

\textit{Proposition 21 (iii) and (ii) realised.} Points P with v₅(q(P)) ≥ 3 and S ∈ ker χ with p²S of depth exactly three, at 4407a1 and 5956a1 (p = 5), both of rank two, so that v\_p(n²h\_p(P)) is read pointwise in place of v\_p(Reg\_p): in \textbf{32 pairs of 32} the points P and P + p²S have F mod p² in common and are separated by min(v\_p, 3). An explicit pair, recomputed for this version by a separate script: at 4407a1, p = 5, n = 2, P = −5g₁ + 6g₂ and R = 25(g₁ − 2g₂) give F(P) ≡ 49 and F(P + R) ≡ 24 (mod 125), so F agrees modulo 25, and v₅(n²h\_p) = 7 and 2. With R = 125(g₁ − 2g₂) the two points share F mod p³ and are separated by min(v\_p, 4) (2 of 2). With R = 25(−g₁ + g₂) instead — χ(Z) ≠ 0, ζ(Z) = 0 — the two points share t₁ mod p⁴, F(P + R) ≡ 29 differs from 49 at valuation exactly one, and v₅(n²h\_p) = 7 and 2: the necessity of the second digit of F, realised by two points of one curve. At 8577a1, 75393a1 and 188032i1 there is no such S in the range searched, since γ vanishes on ker χ there (the boxes with det(χ, γ) ≡ 0 of §5.12): the hypothesis γ|\_\{ker χ\} ≢ 0 of (iii) is needed, and this is its negative control. A search that allows only translates R = p²Z with χ(Z) ≠ 0 finds 80 separating pairs at 4407a1 and 5956a1, and in all 80 the second digit of F moves as well; that is part (ii) of Proposition 21 (e(P) = 2 and p ∣ q(P) in all 80), and not an identity.

\textit{Proposition 22 as displayed.} Over the 648 boxes of §5.10 (e distributed as [551, 82, 13, 2], which reproduces the filter) the congruence F ≡ A/2 − pA/2 + pA²/8 (mod p²) holds in \textbf{648 of 648}, and agrees in 648 of 648 with an independent inversion of Theorem N* by the series of log and exp. Dropping c/(p²x(P″)) leaves 72 of the 551 boxes with e = 2 and all 97 with e ≥ 3, where that term has valuation 2e − 2 ≥ 4 and is invisible; dropping (p²/2)Θ² leaves 111 of 648, and dropping −pA/2 leaves 204.

\section*{6. What is proved in Lean and what is not}

Thirty files, each checking from \texttt{import\ \allowbreak{}Mathlib} alone: \texttt{PadicHeightValuation\{,2,…,29\}.\allowbreak{}lean} and \texttt{PadicHeightValuationCanonical.\allowbreak{}lean} (namespace \texttt{PadicCanonicalHeight}), in the namespace \texttt{PadicHeight} except file 23 (\texttt{PadicHeightStar}), file 24 (\texttt{PadicHeightLayer}), file 25 (\texttt{PadicHeightCM}), files 26–29, whose namespaces are listed below, and the canonical file; and a thirty-first, \texttt{AnomalousEulerFactor.\allowbreak{}lean}, which belongs to the companion note. Files 22 to 25 hold the arithmetic of §§3.9, 3.11–3.12 and 3.13 — Hasse's two cases for p ∣ \#E(F\_p), the coefficients of ℒ, the digit at which each term of Theorem N* is visible, and the comparison of the two rates of Proposition 20 with the divisibility step that carries it to the minimality of Theorem G⁽ᵏ⁾. File 24 holds nineteen statements, with no warning and no \texttt{sorry}; the rates themselves (v\_p(ΔΘ) = j − 1 and v\_p(Δz) = 1 + j) enter it as hypotheses, since they need the formal group. File 25 holds eighteen statements on the two places where c enters: that the lag of the ledger is uniform in m and uses neither c ≠ 0 nor v\_p(μ\_\{2m\}) ≥ 0, the condition for the m = 1 term of Proposition 20 (iii) to lead and its automatic case p ∤ c, the cancellation criterion (1 + w)\textasciicircum{}\{2m\} = 1 of the boundary row with its two verdicts at p = 7, the residue arithmetic of 4p = 1 + 3f² that admits p = 7 and p = 19 and excludes p = 13, and the emptiness of the list of local invariants before the eleventh digit at a cell with c = 0. Files 18–21 hold the arithmetic of §3.14 and §4.13: the two isogeny identities of Lemma 13 as identities in a commutative ring, the residue step of the equivalence (R′) ∧ (L), the cancellations 2n\_p + 1 = p and 4n\_p² − 1 = (2n\_p−1)(2n\_p+1), the local expansion D²H/DH = 2κ − 1 and the telescoping Σ\_\{k\textless{}i\}(q\textasciicircum{}k−1)/(1−q\textasciicircum{}k) = −i of §4.13(c), the balanced representation of every integer ≥ −n\_p as pm + j with |j| ≤ n\_p, the counting identities for m\_S and m\_K, and the exponent arithmetic behind ε\_p = −1 ⟺ p ≡ 3 (mod 8). Every statement is about integers, matrices, finite sums or p-adic congruences; Mathlib has neither the p-adic sigma function nor division polynomials nor Vélu's formulas nor the Tate curve, so the geometry and the p-adic analysis enter as hypotheses. The companion \texttt{statements-and-dependencies.\allowbreak{}md} lists all Lean names; the main line is

\begin{small}
\begin{longtable}{>{\raggedright\arraybackslash}p{0.242\linewidth}>{\raggedright\arraybackslash}p{0.574\linewidth}>{\raggedright\arraybackslash}p{0.143\linewidth}}
\hline
\textbf{Statement} & \textbf{Lean name} & \textbf{File} \\
\hline\endfirsthead
\hline
Statement & Lean name & File \\
\hline\endhead
Theorem A & \texttt{norm\_\allowbreak{}det\_\allowbreak{}div\_\allowbreak{}sq\_\allowbreak{}le} & 1–4 \\
Theorem B, step 3 & \texttt{fermatQuot\_\allowbreak{}sub\_\allowbreak{}iff\_\allowbreak{}of\_\allowbreak{}sq\_\allowbreak{}eq\_\allowbreak{}cube}, \texttt{dvd\_\allowbreak{}three\_\allowbreak{}mul\_\allowbreak{}sub\_\allowbreak{}two\_\allowbreak{}mul}, \texttt{sub\_\allowbreak{}ne\_\allowbreak{}zero\_\allowbreak{}of\_\allowbreak{}two\_\allowbreak{}mul\_\allowbreak{}eq\_\allowbreak{}three\_\allowbreak{}mul} & 1–4, 8 \\
Theorem B, the model with p ∣ a₁ & \texttt{sq\_\allowbreak{}sub\_\allowbreak{}cube\_\allowbreak{}dvd\_\allowbreak{}of\_\allowbreak{}dvd\_\allowbreak{}a1} (and \texttt{sq\_\allowbreak{}sub\_\allowbreak{}cube\_\allowbreak{}dvd\_\allowbreak{}of\_\allowbreak{}sq\_\allowbreak{}dvd\_\allowbreak{}psi} for e ≥ 2) & 8 \\
Lemmas 4′ and 5 & \texttt{eval\_\allowbreak{}sub\_\allowbreak{}eval\_\allowbreak{}dvd\_\allowbreak{}sq} & 1–4 \\
Corollary D & \texttt{sq\_\allowbreak{}dvd\_\allowbreak{}pow\_\allowbreak{}sub\_\allowbreak{}one\_\allowbreak{}of\_\allowbreak{}isPow} & 1–4 \\
Corollary E & \texttt{sq\_\allowbreak{}dvd\_\allowbreak{}alpha\_\allowbreak{}sub\_\allowbreak{}beta\_\allowbreak{}of\_\allowbreak{}isPow} & 1–4 \\
Theorem H & \texttt{fermatQuot\_\allowbreak{}scale} & 1–4 \\
Theorem J & \texttt{dvd\_\allowbreak{}det\_\allowbreak{}of\_\allowbreak{}col\_\allowbreak{}dvd}, \texttt{pow\_\allowbreak{}succ\_\allowbreak{}dvd\_\allowbreak{}det\_\allowbreak{}of\_\allowbreak{}col} & 1–4 \\
Main theorem, linear algebra & \texttt{regulator\_\allowbreak{}valuation\_\allowbreak{}of\_\allowbreak{}split\_\allowbreak{}anomalous} & 5 \\
Main theorem (ii), equivalence & \texttt{regulator\_\allowbreak{}valuation\_\allowbreak{}of\_\allowbreak{}split\_\allowbreak{}anomalous\_\allowbreak{}iff\_\allowbreak{}split} & 5 \\
§4.6 (2), the single slope and the affine zero & \texttt{slope\_\allowbreak{}of\_\allowbreak{}rho}, \texttt{fermatQuot\_\allowbreak{}deform}, \texttt{affine\_\allowbreak{}zero\_\allowbreak{}unique}, \texttt{lambda\_\allowbreak{}eq\_\allowbreak{}zero\_\allowbreak{}iff\_\allowbreak{}split}, \texttt{lambda\_\allowbreak{}dichotomy} & 5 \\
§4.6 (3), the weight count and the Newton polygon & \texttt{weight\_\allowbreak{}split}, \texttt{twelve\_\allowbreak{}dvd\_\allowbreak{}weight}, \texttt{sum\_\allowbreak{}pole\_\allowbreak{}orders}, \texttt{cusp\_\allowbreak{}sum\_\allowbreak{}eq\_\allowbreak{}m}, \texttt{mul\_\allowbreak{}compl\_\allowbreak{}le\_\allowbreak{}half\_\allowbreak{}mul\_\allowbreak{}half}, \texttt{eq\_\allowbreak{}C\_\allowbreak{}mul\_\allowbreak{}of\_\allowbreak{}dvd\_\allowbreak{}of\_\allowbreak{}natDegree\_\allowbreak{}eq} & 6 \\
§4.6 (4), the residues & \texttt{inv\_\allowbreak{}one\_\allowbreak{}sub\_\allowbreak{}add\_\allowbreak{}inv\_\allowbreak{}one\_\allowbreak{}sub\_\allowbreak{}inv}, \texttt{first\_\allowbreak{}term\_\allowbreak{}residue}, \texttt{second\_\allowbreak{}term\_\allowbreak{}leading\_\allowbreak{}ne\_\allowbreak{}zero}, \texttt{residues\_\allowbreak{}do\_\allowbreak{}not\_\allowbreak{}cancel}, \texttt{twelve\_\allowbreak{}mul\_\allowbreak{}sum\_\allowbreak{}cusp\_\allowbreak{}orders} & 7 \\
Theorem S, the ledger & \texttt{exists\_\allowbreak{}shift\_\allowbreak{}dvd\_\allowbreak{}a1}, \texttt{add\_\allowbreak{}val\_\allowbreak{}exactly\_\allowbreak{}one}, \texttt{not\_\allowbreak{}dvd\_\allowbreak{}of\_\allowbreak{}val\_\allowbreak{}one} & 8 \\
Main theorem (⟹ of (i)) & \texttt{quadratic\_\allowbreak{}value\_\allowbreak{}mem\_\allowbreak{}of\_\allowbreak{}entries\_\allowbreak{}mem} & 8 \\
Theorem R, the linear algebra (R0) & \texttt{bilin\_\allowbreak{}rank\_\allowbreak{}one\_\allowbreak{}of\_\allowbreak{}zero\_\allowbreak{}iff}, \texttt{not\_\allowbreak{}dvd\_\allowbreak{}of\_\allowbreak{}congr\_\allowbreak{}ne\_\allowbreak{}zero} & 9 \\
Corollary R′, the determinant & \texttt{det\_\allowbreak{}conj\_\allowbreak{}diagonal}, \texttt{det\_\allowbreak{}conj\_\allowbreak{}diagonal\_\allowbreak{}one\_\allowbreak{}p}, \texttt{reg\_\allowbreak{}eq\_\allowbreak{}of\_\allowbreak{}det\_\allowbreak{}eq} & 9 \\
Theorem R\_s, the linear algebra (R0s) and rank ≤ 2 & \texttt{bilin\_\allowbreak{}eq\_\allowbreak{}of\_\allowbreak{}vanishing\_\allowbreak{}on\_\allowbreak{}hyperplane}, \texttt{bilin\_\allowbreak{}zero\_\allowbreak{}of\_\allowbreak{}mem\_\allowbreak{}radical}, \texttt{gram\_\allowbreak{}det\_\allowbreak{}three\_\allowbreak{}eq\_\allowbreak{}zero} & 10 \\
Corollary R′\_s, the determinant & \texttt{det\_\allowbreak{}two\_\allowbreak{}dvd\_\allowbreak{}iff}, \texttt{pow\_\allowbreak{}dvd\_\allowbreak{}det\_\allowbreak{}of\_\allowbreak{}offdiag\_\allowbreak{}block} & 10 \\
Proposition R″\_s, the second Fermat quotient & \texttt{fermatQuot2\_\allowbreak{}add}, \texttt{dvd\_\allowbreak{}three\_\allowbreak{}mul\_\allowbreak{}sub\_\allowbreak{}two\_\allowbreak{}mul\_\allowbreak{}two}, \texttt{split\_\allowbreak{}const\_\allowbreak{}eq} & 10 \\
Proposition R⁗ (the slope) and Corollary R⁗′ & \texttt{eval\_\allowbreak{}shift\_\allowbreak{}of\_\allowbreak{}deriv\_\allowbreak{}eq}, \texttt{sq\_\allowbreak{}dvd\_\allowbreak{}pow\_\allowbreak{}add\_\allowbreak{}sub}, \texttt{fermatQuot2\_\allowbreak{}move}, \texttt{lambda2\_\allowbreak{}affine}, \texttt{lambda2\_\allowbreak{}const\_\allowbreak{}of\_\allowbreak{}slope\_\allowbreak{}zero}, \texttt{lambda2\_\allowbreak{}lift\_\allowbreak{}bijective}, \texttt{lambda2\_\allowbreak{}unique\_\allowbreak{}zero\_\allowbreak{}lift}, \texttt{b\_\allowbreak{}eq\_\allowbreak{}slope\_\allowbreak{}mul\_\allowbreak{}kappa} & 11 \\
Lemma 7, the first digit of the x-coordinate & \texttt{eval\_\allowbreak{}shift\_\allowbreak{}first\_\allowbreak{}digit}, \texttt{first\_\allowbreak{}digit\_\allowbreak{}eq}, \texttt{eval\_\allowbreak{}shift\_\allowbreak{}depth\_\allowbreak{}two}, \texttt{first\_\allowbreak{}digit\_\allowbreak{}ne\_\allowbreak{}zero\_\allowbreak{}iff} & 12 \\
Theorem K and Corollary K′, the bilinear algebra and the three layers & \texttt{master}, \texttt{master\_\allowbreak{}symm}, \texttt{master\_\allowbreak{}add\_\allowbreak{}left}, \texttt{master\_\allowbreak{}smul\_\allowbreak{}left}, \texttt{master\_\allowbreak{}vanish\_\allowbreak{}on\_\allowbreak{}ker\_\allowbreak{}chi}, \texttt{master\_\allowbreak{}mem\_\allowbreak{}radical}, \texttt{master\_\allowbreak{}eq\_\allowbreak{}of\_\allowbreak{}gam\_\allowbreak{}eq\_\allowbreak{}smul}, \texttt{master\_\allowbreak{}offdiag\_\allowbreak{}ne\_\allowbreak{}iff}, \texttt{master\_\allowbreak{}diag\_\allowbreak{}eq\_\allowbreak{}zero\_\allowbreak{}iff}, \texttt{master\_\allowbreak{}bad\_\allowbreak{}class\_\allowbreak{}unique}, \texttt{det\_\allowbreak{}three\_\allowbreak{}layer}, \texttt{det\_\allowbreak{}eq\_\allowbreak{}zero\_\allowbreak{}of\_\allowbreak{}rank\_\allowbreak{}one\_\allowbreak{}block}, \texttt{det\_\allowbreak{}sq\_\allowbreak{}dvd\_\allowbreak{}iff} & 12 \\
Proposition 9, the second digit under lifting & \texttt{X\_\allowbreak{}pow\_\allowbreak{}four\_\allowbreak{}dvd\_\allowbreak{}sub\_\allowbreak{}trunc}, \texttt{eval\_\allowbreak{}quartic\_\allowbreak{}dvd}, \texttt{eval\_\allowbreak{}shift\_\allowbreak{}second\_\allowbreak{}digit}, \texttt{fermatQuot\_\allowbreak{}second\_\allowbreak{}move} & 13 \\
Proposition 10, the bad lift and the layers & \texttt{root\_\allowbreak{}factor}, \texttt{dvd\_\allowbreak{}iff\_\allowbreak{}dvd\_\allowbreak{}sub\_\allowbreak{}root}, \texttt{unique\_\allowbreak{}bad\_\allowbreak{}class}, \texttt{exists\_\allowbreak{}bad\_\allowbreak{}class} & 13 \\
Theorem N, the ledger and the cancellation of a₁ & \texttt{second\_\allowbreak{}digit\_\allowbreak{}zero\_\allowbreak{}iff}, \texttt{second\_\allowbreak{}digit\_\allowbreak{}criterion}, \texttt{a1\_\allowbreak{}correction\_\allowbreak{}cancels} & 13 \\
Corollary N′ and the three-fold split at the second digit & \texttt{dvd\_\allowbreak{}iff\_\allowbreak{}dvd\_\allowbreak{}first}, \texttt{sq\_\allowbreak{}dvd\_\allowbreak{}sub\_\allowbreak{}crit}, \texttt{val\_\allowbreak{}one\_\allowbreak{}of\_\allowbreak{}crit\_\allowbreak{}ne}, \texttt{sq\_\allowbreak{}dvd\_\allowbreak{}of\_\allowbreak{}crit}, \texttt{trichotomy}, \texttt{min\_\allowbreak{}two\_\allowbreak{}ne\_\allowbreak{}iff}, \texttt{min\_\allowbreak{}two\_\allowbreak{}eq\_\allowbreak{}of\_\allowbreak{}two\_\allowbreak{}le}, \texttt{exists\_\allowbreak{}same\_\allowbreak{}data\_\allowbreak{}diff\_\allowbreak{}val}, \texttt{q2\_\allowbreak{}additive}, \texttt{scaling\_\allowbreak{}shift\_\allowbreak{}square} & 14 \\
Corollary N⁽³⁾ and Corollary N″, the four-fold split & \texttt{inv\_\allowbreak{}pred\_\allowbreak{}cube}, \texttt{neg\_\allowbreak{}inv\_\allowbreak{}pred\_\allowbreak{}cube}, \texttt{q2\_\allowbreak{}quotient}, \texttt{q2\_\allowbreak{}quotient\_\allowbreak{}dvd}, \texttt{sq\_\allowbreak{}cube\_\allowbreak{}step}, \texttt{denominator\_\allowbreak{}third\_\allowbreak{}digit}, \texttt{cube\_\allowbreak{}dvd\_\allowbreak{}sub\_\allowbreak{}crit}, \texttt{cube\_\allowbreak{}dvd\_\allowbreak{}of\_\allowbreak{}crit}, \texttt{not\_\allowbreak{}cube\_\allowbreak{}dvd\_\allowbreak{}of\_\allowbreak{}crit\_\allowbreak{}ne}, \texttt{cube\_\allowbreak{}dvd\_\allowbreak{}iff\_\allowbreak{}crit}, \texttt{quadrichotomy}, \texttt{exists\_\allowbreak{}same\_\allowbreak{}second\_\allowbreak{}diff\_\allowbreak{}third}, \texttt{min\_\allowbreak{}three\_\allowbreak{}ne\_\allowbreak{}iff}, \texttt{min\_\allowbreak{}three\_\allowbreak{}eq\_\allowbreak{}of\_\allowbreak{}three\_\allowbreak{}le} & 15 \\
Proposition 12, Corollary N⁽⁴⁾ and Corollary N‴, the five-fold split & \texttt{low\_\allowbreak{}terms\_\allowbreak{}vanish}, \texttt{p\_\allowbreak{}mult\_\allowbreak{}second\_\allowbreak{}digit}, \texttt{w\_\allowbreak{}local}, \texttt{inv\_\allowbreak{}x\_\allowbreak{}correction}, \texttt{w\_\allowbreak{}ne\_\allowbreak{}zero\_\allowbreak{}iff}, \texttt{inv\_\allowbreak{}pred\_\allowbreak{}quartic}, \texttt{neg\_\allowbreak{}inv\_\allowbreak{}pred\_\allowbreak{}quartic}, \texttt{denominator\_\allowbreak{}fourth\_\allowbreak{}digit}, \texttt{theta\_\allowbreak{}eta\_\allowbreak{}cancels}, \texttt{cubic\_\allowbreak{}term\_\allowbreak{}vanishes}, \texttt{fourth\_\allowbreak{}digit\_\allowbreak{}reduction}, \texttt{quartic\_\allowbreak{}dvd\_\allowbreak{}iff\_\allowbreak{}crit}, \texttt{exists\_\allowbreak{}same\_\allowbreak{}third\_\allowbreak{}diff\_\allowbreak{}fourth}, \texttt{min\_\allowbreak{}four\_\allowbreak{}ne\_\allowbreak{}iff}, \texttt{min\_\allowbreak{}four\_\allowbreak{}eq\_\allowbreak{}of\_\allowbreak{}four\_\allowbreak{}le}, \texttt{quintichotomy} & 16 \\
Theorem G′ (the sensitivity) and the sharpness in Theorem G″ & \texttt{pow\_\allowbreak{}sub\_\allowbreak{}pow\_\allowbreak{}expand}, \texttt{exists\_\allowbreak{}exact\_\allowbreak{}part}, \texttt{unit\_\allowbreak{}pow\_\allowbreak{}diff\_\allowbreak{}exact}, \texttt{shift\_\allowbreak{}depth\_\allowbreak{}exact}, \texttt{shift\_\allowbreak{}depth\_\allowbreak{}dvd}, \texttt{quotient\_\allowbreak{}by\_\allowbreak{}p\_\allowbreak{}sq}, \texttt{theta\_\allowbreak{}sensitivity\_\allowbreak{}exact}, \texttt{theta\_\allowbreak{}digit\_\allowbreak{}determined}, \texttt{three\_\allowbreak{}layers\_\allowbreak{}determine}, \texttt{layer\_\allowbreak{}three\_\allowbreak{}moves\_\allowbreak{}digit\_\allowbreak{}two}, \texttt{layer\_\allowbreak{}digit\_\allowbreak{}correspondence}, \texttt{nondegeneracy\_\allowbreak{}needed} & 17 \\
§4.12, v\_p(ψ\_p′(x(T))) = 1 exactly, and (ND) as a condition on ψ\_p mod p³ & \texttt{sq\_\allowbreak{}valuation\_\allowbreak{}one}, \texttt{psi\_\allowbreak{}deriv\_\allowbreak{}valuation\_\allowbreak{}one}, \texttt{bracket\_\allowbreak{}divisible}, \texttt{nd\_\allowbreak{}iff\_\allowbreak{}bracket}, \texttt{nd\_\allowbreak{}of\_\allowbreak{}bracket\_\allowbreak{}exact} & 18 \\
Proposition 16, the weight of Res\_x(G, ψ\_p mod p) and m\_p = m\_K + m\_S & \texttt{resultant\_\allowbreak{}weight}, \texttt{resultant\_\allowbreak{}weight\_\allowbreak{}split}, \texttt{m\_\allowbreak{}p\_\allowbreak{}integral}, \texttt{m\_\allowbreak{}g\_\allowbreak{}integral}, \texttt{m\_\allowbreak{}K\_\allowbreak{}add\_\allowbreak{}m\_\allowbreak{}S}, \texttt{resultant\_\allowbreak{}weight\_\allowbreak{}K}, \texttt{resultant\_\allowbreak{}weight\_\allowbreak{}S}, \texttt{twentyfour\_\allowbreak{}dvd\_\allowbreak{}sq\_\allowbreak{}sub\_\allowbreak{}one} & 18 \\
Theorem ND, the step "the constant is non-zero ⟹ (ND)" & \texttt{ne\_\allowbreak{}zero\_\allowbreak{}of\_\allowbreak{}const\_\allowbreak{}mul\_\allowbreak{}pow}, \texttt{eval\_\allowbreak{}ne\_\allowbreak{}zero\_\allowbreak{}of\_\allowbreak{}isCoprime}, \texttt{nondegeneracy\_\allowbreak{}summary}, \texttt{nd\_\allowbreak{}iff\_\allowbreak{}S\_\allowbreak{}ne\_\allowbreak{}zero} & 18 \\
§3.14, D kills Δ₀; the cusp count of §4.13; c\_S ≠ 0 given γ\_p ≠ 0; the chain rule for D & \texttt{serre\_\allowbreak{}deriv\_\allowbreak{}disc\_\allowbreak{}zero}, \texttt{gauss\_\allowbreak{}sum\_\allowbreak{}closed}, \texttt{min\_\allowbreak{}sum\_\allowbreak{}closed}, \texttt{cusp\_\allowbreak{}ord\_\allowbreak{}total}, \texttt{resultant\_\allowbreak{}weight\_\allowbreak{}dPsi}, \texttt{m\_\allowbreak{}S\_\allowbreak{}eq}, \texttt{weight\_\allowbreak{}two\_\allowbreak{}p\_\allowbreak{}vanishes}, \texttt{pow\_\allowbreak{}card\_\allowbreak{}self}, \texttt{cS\_\allowbreak{}ne\_\allowbreak{}zero}, \texttt{deriv\_\allowbreak{}zero}, \texttt{deriv\_\allowbreak{}one}, \texttt{deriv\_\allowbreak{}sum}, \texttt{deriv\_\allowbreak{}pow\_\allowbreak{}succ}, \texttt{deriv\_\allowbreak{}eval\_\allowbreak{}chain} & 19 \\
Lemma 13, and the reduction of Proposition 15 to (R′) ∧ (L) & \texttt{isogeny\_\allowbreak{}A1}, \texttt{isogeny\_\allowbreak{}A1\_\allowbreak{}domain}, \texttt{isogeny\_\allowbreak{}A2\_\allowbreak{}core}, \texttt{isogeny\_\allowbreak{}A2}, \texttt{deriv\_\allowbreak{}eval\_\allowbreak{}of\_\allowbreak{}factor}, \texttt{residue\_\allowbreak{}iff}, \texttt{const\_\allowbreak{}mod\_\allowbreak{}p\_\allowbreak{}deriv}, \texttt{lead\_\allowbreak{}cancel}, \texttt{lc\_\allowbreak{}difference}, \texttt{deg\_\allowbreak{}K\_\allowbreak{}eq}, \texttt{weight\_\allowbreak{}L}, \texttt{two\_\allowbreak{}n\_\allowbreak{}add\_\allowbreak{}one\_\allowbreak{}eq\_\allowbreak{}zero} & 20 \\
§4.13(b), (c): the balanced exponents, D²H/DH = 2κ − 1, the telescoped sum, i ↦ p − i & \texttt{balanced\_\allowbreak{}repr}, \texttt{D\_\allowbreak{}eval\_\allowbreak{}zero}, \texttt{D2\_\allowbreak{}eval\_\allowbreak{}zero}, \texttt{log\_\allowbreak{}deriv\_\allowbreak{}ratio}, \texttt{telescope\_\allowbreak{}term}, \texttt{telescope\_\allowbreak{}sum}, \texttt{residue\_\allowbreak{}value}, \texttt{residue\_\allowbreak{}value\_\allowbreak{}zmod}, \texttt{sign\_\allowbreak{}flip\_\allowbreak{}invariance} & 21 \\
Proposition 16, m\_K, its weight, and ε\_p = −1 ⟺ p ≡ 3 (mod 8) & \texttt{ordq\_\allowbreak{}total}, \texttt{weight\_\allowbreak{}ResK}, \texttt{mK\_\allowbreak{}even}, \texttt{exponent\_\allowbreak{}two\_\allowbreak{}nsq}, \texttt{exponent\_\allowbreak{}case\_\allowbreak{}1mod4}, \texttt{exponent\_\allowbreak{}case\_\allowbreak{}3mod4}, \texttt{sign\_\allowbreak{}exponent}, \texttt{eps\_\allowbreak{}cases} & 21 \\
Proposition R⁺, the p-part of the reduction and Hasse's two cases & \texttt{card\_\allowbreak{}cases}, \texttt{two\_\allowbreak{}p\_\allowbreak{}forces\_\allowbreak{}five}, \texttt{sq\_\allowbreak{}not\_\allowbreak{}dvd}, \texttt{dvd\_\allowbreak{}of\_\allowbreak{}unit\_\allowbreak{}mul}, \texttt{dvd\_\allowbreak{}of\_\allowbreak{}unit\_\allowbreak{}sq\_\allowbreak{}mul} & 22 \\
The coefficients of σ\_p in a general model, and the depth of a term of a logarithm & \texttt{sigma\_\allowbreak{}third\_\allowbreak{}coeff}, \texttt{sigma\_\allowbreak{}fourth\_\allowbreak{}coeff}, \texttt{padicVal\_\allowbreak{}eq\_\allowbreak{}zero\_\allowbreak{}of\_\allowbreak{}lt}, \texttt{log\_\allowbreak{}term\_\allowbreak{}depth\_\allowbreak{}small}, \texttt{log\_\allowbreak{}term\_\allowbreak{}depth\_\allowbreak{}five}, \texttt{log\_\allowbreak{}term\_\allowbreak{}depth\_\allowbreak{}ne\_\allowbreak{}five}, \texttt{log\_\allowbreak{}term\_\allowbreak{}factor}, \texttt{five\_\allowbreak{}fifth\_\allowbreak{}term} & 22 \\
Proposition 16, the constant c\_p in one line & \texttt{cp\_\allowbreak{}one\_\allowbreak{}line} & 22 \\
Proposition 17, the coefficients of ℒ and the reparametrisation in t & \texttt{mu2\_\allowbreak{}eq}, \texttt{mu4\_\allowbreak{}eq}, \texttt{mu6\_\allowbreak{}eq}, \texttt{mu6\_\allowbreak{}no\_\allowbreak{}b2\_\allowbreak{}cube}, \texttt{lambda3\_\allowbreak{}eq} & 23 \\
Proposition 18, the digit at which each term is visible & \texttt{lambda\_\allowbreak{}depth\_\allowbreak{}mono}, \texttt{lambda\_\allowbreak{}depth\_\allowbreak{}m2\_\allowbreak{}e2}, \texttt{x\_\allowbreak{}form\_\allowbreak{}suffices}, \texttt{val\_\allowbreak{}thirty}, \texttt{val\_\allowbreak{}three\_\allowbreak{}sixty}, \texttt{val\_\allowbreak{}thirty\_\allowbreak{}seven}, \texttt{mu6\_\allowbreak{}visible\_\allowbreak{}at\_\allowbreak{}ten}, \texttt{theta\_\allowbreak{}depth\_\allowbreak{}j5\_\allowbreak{}p5}, \texttt{theta\_\allowbreak{}depth\_\allowbreak{}j5\_\allowbreak{}p7}, \texttt{log\_\allowbreak{}term\_\allowbreak{}j10\_\allowbreak{}five}, \texttt{log\_\allowbreak{}term\_\allowbreak{}j10\_\allowbreak{}seven} & 23 \\
Proposition 19 and Theorem G⁽ᵏ⁾, the depth of a layer and the growth of the data & \texttt{layer\_\allowbreak{}depth}, \texttt{data\_\allowbreak{}count\_\allowbreak{}linear} & 23 \\
Proposition 20 and Theorem G⁽ᵏ⁾, the two rates and the minimality at every k & \texttt{theta\_\allowbreak{}rate\_\allowbreak{}digit}, \texttt{ledger\_\allowbreak{}rate\_\allowbreak{}digit}, \texttt{gap\_\allowbreak{}is\_\allowbreak{}e}, \texttt{ledger\_\allowbreak{}lags}, \texttt{ledger\_\allowbreak{}lags\_\allowbreak{}general}, \texttt{uniform\_\allowbreak{}in\_\allowbreak{}k}, \texttt{exact\_\allowbreak{}pred\_\allowbreak{}not\_\allowbreak{}at}, \texttt{no\_\allowbreak{}cancel}, \texttt{at\_\allowbreak{}most\_\allowbreak{}one\_\allowbreak{}vanishes}, \texttt{layer\_\allowbreak{}is\_\allowbreak{}needed}, \texttt{fermat\_\allowbreak{}digit}, \texttt{fermat\_\allowbreak{}digits\_\allowbreak{}needed}, \texttt{data\_\allowbreak{}count}, \texttt{data\_\allowbreak{}count\_\allowbreak{}step} & 24 \\
Proposition 20 (iii), the boundary row and the CM cells & \texttt{ledger\_\allowbreak{}lags\_\allowbreak{}any\_\allowbreak{}m}, \texttt{cm\_\allowbreak{}lags\_\allowbreak{}ten}, \texttt{m1\_\allowbreak{}leads\_\allowbreak{}iff}, \texttt{m1\_\allowbreak{}leads\_\allowbreak{}of\_\allowbreak{}unit\_\allowbreak{}c}, \texttt{boundary\_\allowbreak{}cancel\_\allowbreak{}m1}, \texttt{boundary\_\allowbreak{}cancel\_\allowbreak{}pow}, \texttt{boundary\_\allowbreak{}cancel\_\allowbreak{}of\_\allowbreak{}card}, \texttt{boundary\_\allowbreak{}p7\_\allowbreak{}m3}, \texttt{boundary\_\allowbreak{}p7\_\allowbreak{}exception}, \texttt{boundary\_\allowbreak{}m1\_\allowbreak{}vs\_\allowbreak{}m3\_\allowbreak{}p7}, \texttt{cm\_\allowbreak{}prime\_\allowbreak{}mod\_\allowbreak{}three}, \texttt{cm\_\allowbreak{}no\_\allowbreak{}anomalous\_\allowbreak{}of\_\allowbreak{}mod\_\allowbreak{}two}, \texttt{cm\_\allowbreak{}seven\_\allowbreak{}nineteen}, \texttt{cm\_\allowbreak{}not\_\allowbreak{}thirteen}, \texttt{cm\_\allowbreak{}ledger\_\allowbreak{}invisible\_\allowbreak{}before\_\allowbreak{}eleventh}, \texttt{cm\_\allowbreak{}no\_\allowbreak{}local\_\allowbreak{}invariant\_\allowbreak{}before\_\allowbreak{}eleventh}, \texttt{layer\_\allowbreak{}is\_\allowbreak{}needed\_\allowbreak{}cm}, \texttt{cm\_\allowbreak{}gap\_\allowbreak{}at\_\allowbreak{}k} & 25 \\
Corollary K′(i) and Corollary R′\_s at r = 2, as linear algebra & \texttt{prime\_\allowbreak{}dvd\_\allowbreak{}det\_\allowbreak{}two\_\allowbreak{}iff}, \texttt{det\_\allowbreak{}two\_\allowbreak{}split\_\allowbreak{}ne\_\allowbreak{}zero\_\allowbreak{}iff}, \texttt{exists\_\allowbreak{}pair\_\allowbreak{}eq\_\allowbreak{}one}, \texttt{onto\_\allowbreak{}iff\_\allowbreak{}det\_\allowbreak{}ne\_\allowbreak{}zero}, \texttt{onto\_\allowbreak{}iff\_\allowbreak{}reduced\_\allowbreak{}det\_\allowbreak{}ne\_\allowbreak{}zero} & 26 \\
Proposition V (§6, conditional), the integer arithmetic of its four cases & \texttt{leading\_\allowbreak{}val\_\allowbreak{}ge\_\allowbreak{}rank}, \texttt{rank\_\allowbreak{}le\_\allowbreak{}leading\_\allowbreak{}val}, \texttt{leading\_\allowbreak{}val\_\allowbreak{}eq\_\allowbreak{}rank\_\allowbreak{}of\_\allowbreak{}nonsplit\_\allowbreak{}rank\_\allowbreak{}one}, \texttt{leading\_\allowbreak{}val\_\allowbreak{}gt\_\allowbreak{}rank\_\allowbreak{}of\_\allowbreak{}split\_\allowbreak{}le\_\allowbreak{}one}, \texttt{leading\_\allowbreak{}val\_\allowbreak{}eq\_\allowbreak{}rank\_\allowbreak{}iff\_\allowbreak{}of\_\allowbreak{}normal}, \texttt{exists\_\allowbreak{}two\_\allowbreak{}ranks\_\allowbreak{}same\_\allowbreak{}leading\_\allowbreak{}val}, \texttt{admissible\_\allowbreak{}leading\_\allowbreak{}two\_\allowbreak{}of\_\allowbreak{}rank\_\allowbreak{}le\_\allowbreak{}two} & 26 \\
Lemma 23, depth and height & \texttt{log\_\allowbreak{}denom\_\allowbreak{}ge}, \texttt{height\_\allowbreak{}ge\_\allowbreak{}of\_\allowbreak{}depth}, \texttt{depth\_\allowbreak{}lt\_\allowbreak{}of\_\allowbreak{}height\_\allowbreak{}lt} & 27 \\
Proposition 24, the surjectivity of (χ, γ) from one shallow element of ker χ & \texttt{surjective\_\allowbreak{}pair\_\allowbreak{}iff}, \texttt{surjective\_\allowbreak{}of\_\allowbreak{}shallow\_\allowbreak{}kernel\_\allowbreak{}element} & 27 \\
Proposition 24ᶜ, the lower bound R ≥ β(E) & \texttt{beta\_\allowbreak{}le\_\allowbreak{}remainder\_\allowbreak{}of\_\allowbreak{}local\_\allowbreak{}bounds} (hypotheses: 0 ≤ log c′, m\_∞ ≤ λ′\_∞, m\_ℓ ≤ corr\_ℓ at each ℓ, and the decomposition R = 2 log c′ + λ′\_∞ + Σ corr\_ℓ) & Canonical \\
Proposition 24ᶜ, depth against canonical height & \texttt{canonical\_\allowbreak{}ge\_\allowbreak{}of\_\allowbreak{}depth}, \texttt{depth\_\allowbreak{}lt\_\allowbreak{}of\_\allowbreak{}canonical\_\allowbreak{}lt} (hypotheses: ĥ = 2D log p + R, β ≤ R, k ≤ D, 0 ≤ log p) & Canonical \\
Proposition 24ᶜ, the surjectivity & \texttt{surjective\_\allowbreak{}of\_\allowbreak{}canonical\_\allowbreak{}lt\_\allowbreak{}four\_\allowbreak{}log\_\allowbreak{}add\_\allowbreak{}beta} (hypotheses: χ onto; χ(y) = γ(y) = 0 ⇒ D(y) ≥ 2; the decomposition and β ≤ R at every element; one element of ker χ with ĥ \textless{} 4 log p + β), \texttt{surjective\_\allowbreak{}of\_\allowbreak{}naive\_\allowbreak{}or\_\allowbreak{}canonical} (below either threshold), \texttt{surjective\_\allowbreak{}pair\_\allowbreak{}iff'} & Canonical \\
§3.15, the constants m\_ℓ at I\_N & \texttt{four\_\allowbreak{}mul\_\allowbreak{}comp\_\allowbreak{}le\_\allowbreak{}sq}, \texttt{comp\_\allowbreak{}le\_\allowbreak{}half} (hypothesis: i ≤ N) & Canonical \\
The difference of the two slacks & \texttt{remainder\_\allowbreak{}sub\_\allowbreak{}slack} (hypotheses: the naive and the canonical decompositions) & Canonical \\
§5.12, the index of the p-adic closure and the null model for det(χ, γ) & \texttt{index\_\allowbreak{}zero\_\allowbreak{}iff}, \texttt{index\_\allowbreak{}ge\_\allowbreak{}two\_\allowbreak{}iff}, \texttt{index\_\allowbreak{}le}, \texttt{det\_\allowbreak{}zero\_\allowbreak{}iff}, \texttt{card\_\allowbreak{}det\_\allowbreak{}zero\_\allowbreak{}of\_\allowbreak{}ne}, \texttt{card\_\allowbreak{}det\_\allowbreak{}zero\_\allowbreak{}of\_\allowbreak{}fst\_\allowbreak{}eq\_\allowbreak{}zero}, \texttt{density\_\allowbreak{}det\_\allowbreak{}zero} & 27 \\
Lemma 23′, the slack & \texttt{cofactor\_\allowbreak{}log\_\allowbreak{}le}, \texttt{cofactor\_\allowbreak{}eq\_\allowbreak{}one\_\allowbreak{}of\_\allowbreak{}slack\_\allowbreak{}lt}, \texttt{max\_\allowbreak{}eq\_\allowbreak{}iff} & 28 \\
§3.15, small closures: the layers factor through χ & \texttt{factor\_\allowbreak{}through\_\allowbreak{}chi}, \texttt{two\_\allowbreak{}ne\_\allowbreak{}zero\_\allowbreak{}of\_\allowbreak{}ne\_\allowbreak{}two}, \texttt{diag\_\allowbreak{}ne\_\allowbreak{}zero} & 28 \\
Proposition 20 (v) read on a pair of points, and a count of §5.12 & \texttt{exact\_\allowbreak{}of\_\allowbreak{}deeper}, \texttt{pair\_\allowbreak{}law\_\allowbreak{}j3}, \texttt{pair\_\allowbreak{}law\_\allowbreak{}F\_\allowbreak{}leads}, \texttt{card\_\allowbreak{}two\_\allowbreak{}routes} & 28 \\
Proposition 21, steps (2)–(4) and parts (i)–(iii) & \texttt{ratio\_\allowbreak{}of\_\allowbreak{}ker\_\allowbreak{}le}, \texttt{vanish\_\allowbreak{}of\_\allowbreak{}chi\_\allowbreak{}zero}, \texttt{translate\_\allowbreak{}quadratic}, \texttt{second\_\allowbreak{}difference}, \texttt{shift\_\allowbreak{}quadratic}, \texttt{ledger\_\allowbreak{}digit}, \texttt{kappa\_\allowbreak{}ne\_\allowbreak{}zero}, \texttt{separates}, \texttt{counterexample\_\allowbreak{}digits} & 29 \\
Proposition 22, the inversion at k = 2 & \texttt{inverse\_\allowbreak{}n3} & 29 \\
The Euler factor at an anomalous prime (companion note) & \texttt{ap\_\allowbreak{}of\_\allowbreak{}anomalous}, \texttt{unit\_\allowbreak{}root\_\allowbreak{}identity}, \texttt{sub\_\allowbreak{}one\_\allowbreak{}associated\_\allowbreak{}p}, \texttt{sub\_\allowbreak{}one\_\allowbreak{}associated\_\allowbreak{}p\_\allowbreak{}of\_\allowbreak{}pred} & \texttt{AnomalousEulerFactor} \\
\hline
\end{longtable}
\end{small}

\textbf{How the main theorem appears there.} \texttt{regulator\_\allowbreak{}valuation\_\allowbreak{}of\_\allowbreak{}split\_\allowbreak{}anomalous} takes as hypotheses that the entries of the height matrix on Λ₀ are divisible by p and that det M\_\{Λ₀\} = p²·det M\_Λ, and concludes p\textasciicircum{}\{r−2\} ∣ det M\_Λ; \texttt{regulator\_\allowbreak{}valuation\_\allowbreak{}of\_\allowbreak{}split\_\allowbreak{}anomalous\_\allowbreak{}iff\_\allowbreak{}split} takes the two halves — split ⟹ Reg\_p ∈ Z\_p, and non-split ⟹ Reg\_p = z/p with p ∤ z — and turns them into the equivalence of (ii). That is, \textbf{integrality of the height matrix is a hypothesis there, not a conclusion}: what the kernel checks is the bookkeeping, and the geometry that produces the hypotheses is on paper.

\textbf{Trust base.} \texttt{\#print\ \allowbreak{}axioms} is run at the end of each of the files, and a log is kept beside each source. Every line reads \texttt{depends\ \allowbreak{}on\ \allowbreak{}axioms:} followed by a subset of [propext, Classical.choice, Quot.sound], or else "does not depend on any axioms" (\texttt{deriv\_\allowbreak{}zero} and \texttt{deriv\_\allowbreak{}one} of file 19). There is no \texttt{sorryAx} and no \texttt{Lean.\allowbreak{}ofReduceBool}, that is no \texttt{native\_\allowbreak{}decide}, and each file checks with exit code 0. Files 1–21 check with no warnings; \textbf{files 22 and 23 emit warnings and no errors} — one unused-binder linter note in file 22 and three \texttt{Try\ \allowbreak{}this} suggestions from \texttt{decide} in file 23 — so the claim of a silent build is made for the first twenty-one files, and for files 26–29 and \texttt{AnomalousEulerFactor}, which also check with no warning. Files 24 and 25 are not part of that claim: they check with exit code 0 and the axioms above, but whether they check with no warning at all was not re-checked for this version. Two things are stated exactly. First, **one \texttt{decide} is load-bearing**: \texttt{twelve\_\allowbreak{}ne\_\allowbreak{}zero} of file 7 closes with \texttt{interval\_\allowbreak{}cases\ \allowbreak{}p\ \allowbreak{}\textless{};\textgreater{}\ \allowbreak{}revert\ \allowbreak{}hp\ \allowbreak{}hd\ \allowbreak{}\textless{};\textgreater{}\ \allowbreak{}decide}, and since that is the kernel's \texttt{decide} and not \texttt{native\_\allowbreak{}decide}, no extra axiom appears in any log. Second, the \texttt{\#print\ \allowbreak{}axioms} lines of files 12 and 13 are two fewer than the theorems those files prove, the omissions being \texttt{X\_\allowbreak{}pow\_\allowbreak{}two\_\allowbreak{}dvd\_\allowbreak{}sub\_\allowbreak{}trunc}, \texttt{eval\_\allowbreak{}sq\_\allowbreak{}dvd}, \texttt{X\_\allowbreak{}pow\_\allowbreak{}four\_\allowbreak{}dvd\_\allowbreak{}sub\_\allowbreak{}trunc} and \texttt{sq\_\allowbreak{}dvd\_\allowbreak{}pow\_\allowbreak{}add\_\allowbreak{}sub}; the last of these appears in the table above, so "every theorem listed" is not literally true of them, and they are covered only by the file checking clean.

\textbf{What file 17 does and does not do.} It records the valuation arithmetic of Theorem G′: the isolation of the linear term of a binomial expansion, the exactness of the valuation of (A + D)\textasciicircum{}n − A\textasciicircum{}n for a unit A and n prime to p, the depth count p·g·Δ + p·Δ² ·B + Δ⁵·C, the division by p², and the two directions (determination and sharpness). The three geometric inputs — that x ∘ τ\_T lies in Z\_p[[t]] with unit derivative, that all divided derivatives of φ\_p below the p-th are divisible by p, and the non-degeneracy (ND) — enter as hypotheses, in the shape of the displayed decomposition and of \texttt{¬\ \allowbreak{}p\ \allowbreak{}∣\ \allowbreak{}g}. \texttt{nondegeneracy\_\allowbreak{}needed} exhibits a witness showing that the conclusion fails when \texttt{p\ \allowbreak{}∣\ \allowbreak{}g}, so the hypothesis is not decorative.

\textbf{What files 18–21 leave on paper, and why Theorem ND is not a Lean theorem.} Files 18–21 hold valuation arithmetic, commutative-ring identities, finite sums and exponent counts, and one non-vanishing argument over a field. The following steps of §3.14 and §4.13 are on paper, and each of them is needed for Theorem ND:

\begin{enumerate}[leftmargin=1.4em,itemsep=.2em,topsep=.3em]
\item the canonical subgroup C = Ê[p] and the congruence M ≡ A₀ (mod p) for its kernel polynomial — the heart of Proposition 14;
\item the injectivity of q-expansions in a fixed weight (kernel generated by A − 1), which is what lifts Proposition 15 from the Tate family to every p;
\item the uniformisation of the Tate curve and the theta product for x\_Q(W₁) − x\_Q(W₂);
\item the resultants themselves — Res\_X(DΨ, Ψ) = γ\_pΔ₀\textasciicircum{}M, Res\_x(K, ψ\_p) = c\_KΔ₀\textasciicircum{}\{m\_K\}, Res\_x(S₁, ψ\_p) = c\_SΔ₀\textasciicircum{}\{m\_S\}: Lean has the weights and the q-orders, while the step "no weight-homogeneous irreducible factor other than Δ₀" is on paper, resting on the dichotomy of §4.6(3);
\item γ\_p ≠ 0 itself (\texttt{cS\_\allowbreak{}ne\_\allowbreak{}zero} receives it as a hypothesis), which is the leading-coefficient count at the cusp;
\item the root form of the Kodaira–Spencer identity (R′), and φ₁ = −2σ, σ = −E₂/24 in (L).
\end{enumerate}

So \textbf{Theorem ND is not a Lean theorem}, and neither are Propositions 14, 15 and 16; their grade is [paper], as §3.14 states. What files 18–21 add is that the arithmetic along which those proofs run is kernel-checked, in exactly the sense in which files 1–17 do it for §§3.2–3.11.

\textbf{What is not in Lean.} The p-adic sigma function and the formula for h\_p; Mazur and Tate's differential equation for σ\_p, the Weierstrass ℘-series and the formal logarithm, so that \textbf{Theorem N*, Propositions 17, 18, 19 and 20, Lemma 21 and Theorem G⁽ᵏ⁾ are not Lean theorems} — what files 22, 23 and 24 hold is the coefficient arithmetic, the valuation arithmetic and the digit comparisons they consist of, the analysis being a hypothesis; the Iwasawa logarithm; division polynomials, their composition rule and their shape modulo p; Vélu's formulas and the degree count of Theorem C; [p] = V ∘ F and the group scheme E[p]; Serre–Tate theory; the Tate parametrisation and the valuations read off it; Deuring's lifting theorem; the structure of E(Q\_p) used in Lemma 4; the height pairing itself. So \textbf{the main theorem is not a Lean theorem}: what is machine-checked is the arithmetic skeleton along which its proof runs, and its grade is [paper].

**Files 26–29 and \texttt{AnomalousEulerFactor}.** File 26 (\texttt{PadicHeightLeading}, twelve statements) holds the linear algebra of Corollary K′(i) and Corollary R′\_s at r = 2 and the integer arithmetic of Proposition V below; file 27 (\texttt{PadicDepthHeight}, \texttt{PadicSurjectivity}, \texttt{PadicClosureIndex}, \texttt{PadicNullModel}, twelve statements) the inequality of Lemma 23 given the divisibility p\textasciicircum{}\{2k\} ∣ d², the surjectivity step of Proposition 24 and the counts behind §5.12; file 28 (\texttt{PadicCofactor}, \texttt{PadicClosureCollapse}, \texttt{PadicPairLaw}, \texttt{PadicTwoRoutes}, ten statements) Lemma 23′ and the factorisation through χ of §3.15; file 29 (\texttt{PadicTranslateLinear}, \texttt{PadicTranslateDigit}, \texttt{PadicInverseLedger}, ten statements) the linear algebra and the digit bookkeeping of Propositions 21 and 22. \texttt{AnomalousEulerFactor.\allowbreak{}lean} (\texttt{AnomalousEuler}, four statements) is the arithmetic of the companion note. All five check from \texttt{import\ \allowbreak{}Mathlib} with exit code 0, no warning, no \texttt{sorry} and no \texttt{native\_\allowbreak{}decide}, and every \texttt{\#print\ \allowbreak{}axioms} line reads a subset of [propext, Classical.choice, Quot.sound]. On paper remain: that Q ∈ E\_k(Q\_p) forces p\textasciicircum{}k ∣ d(Q) (Lemma 23); step (1) of Proposition 21, which uses that φ\_p mod p is a polynomial in x\textasciicircum{}p and the invariant differential; the isometry of log on 1 + pZ\_p; and Proposition 20 (iv). So \textbf{Propositions 21, 22, 24 and 24ᶜ are not Lean theorems}; their grade is [paper].

**\texttt{PadicHeightValuationCanonical}** (namespace \texttt{PadicCanonicalHeight}, nine statements) holds the arithmetic and the logic of Proposition 24ᶜ: the termwise lower bound, the inequality between depth and canonical height, the surjectivity, and the constants of I\_N; it checks from \texttt{import\ \allowbreak{}Mathlib} with exit code 0, no warning, no \texttt{sorry} and no \texttt{native\_\allowbreak{}decide}, and its nine \texttt{\#print\ \allowbreak{}axioms} lines read [propext, Classical.choice, Quot.sound]. The local height formulas, the decomposition ĥ = Σ\_v λ′\_v, the table of component corrections and the value of m\_∞ are not in Lean; they are hypotheses of the Lean statements. Lemma 25 is not in Lean.

\textbf{One conditional statement, kept apart.} File 26 also holds the arithmetic of the one statement in this paper that assumes a conjecture. It is stated here, and not in §3, because every theorem of §3 is unconditional and the grades must not mix; the companion note \textit{The valuation of the leading term of the p-adic L-function at an anomalous prime} (the file \texttt{note-padic-L-leading-term.\allowbreak{}md}) states it with its evidence.

\begin{namedthm}[Proposition V (the leading term bounds the rank from above) [paper; conditional on the p-adic conjecture of Birch and Swinnerton-Dyer of Mazur, Tate and Teitelbaum and on p ∤ \#Ш; the integer arithmetic in Lean]]
Let p ≥ 5 be a prime of good ordinary reduction, p ∤ ∏\_ℓ c\_ℓ, p ∤ \#E(Q)\_tors, ι, r = rank E(Q), and L\_p the p-adic L-function of [MTT] on the unit-root eigenspace. \textbf{Assuming the p-adic BSD conjecture and p ∤ \#Ш},
\dispf{\textbf{v\_p(L\_p\textasciicircum{}\{(r)\}(E, χ⁰)/r!) ≥ r}}
in each of four cases: non-anomalous, by Theorem A; split anomalous, as 2 + max(r − 2, 0) ≥ r by Corollary R′\_s; non-split anomalous of rank one, as 2 + (−1) = 1 = r exactly by Theorem S; non-split anomalous of rank r ≥ 2, as 2 + (r − 2) + v\_p(det T) ≥ r by Corollary R′. The 2 is the valuation of the Euler factor (1 − α\textasciicircum{}\{−1\})² at an anomalous prime, exactly 2 for every anomalous p ≥ 5 (Lemma E of the note).
\end{namedthm}

The bound is on the rank \textbf{from above}: a bound from below would need an upper bound for v\_p(Reg\_p), which the local data do not give (§7(3)); Proposition 24 gives one from an archimedean datum, and with it, in the split case of rank two and under the same assumptions, v\_p(L\_p\textasciicircum{}\{(2)\}(E, χ⁰)/2!) = 2 exactly whenever h\_min \textless{} 4 log p or ĥ\_min \textless{} 4 log p + β(E) (Propositions 24 and 24ᶜ). The same value 2 of the leading valuation is compatible with every rank r ≤ 2 (Lean \texttt{admissible\_\allowbreak{}leading\_\allowbreak{}two\_\allowbreak{}of\_\allowbreak{}rank\_\allowbreak{}le\_\allowbreak{}two}). This is not the exceptional zero of [MTT], which occurs at split multiplicative reduction and raises the order of vanishing; here the reduction is good ordinary, the order is unchanged, and it is the valuation of the leading coefficient that rises by 2.

\textbf{Where the assumption can be replaced by theorems [known theorems, combined; the statements and their hypotheses are checked in the original papers [Kato04], [SU], [Ski], [Sch85], [PR92], [Gre99] and [GV]; the integral form of Kato's divisibility under (H) is the argument of [Ski, §2.5], which we have checked against the proof in [Kato04] (§13.14 and the results it uses)].} Write L\_p(E, T) = Σ\_j c\_j T\textasciicircum{}j with 1 + T = (1 + p)\textasciicircum{}s, as in [SW, (3.1)]. Then L\_p\textasciicircum{}\{(r)\}(E, χ⁰)/r! = c\_r·log\_p(1 + p)\textasciicircum{}r whenever c\_0 = … = c\_\{r−1\} = 0, so that v\_p(L\_p\textasciicircum{}\{(r)\}(E, χ⁰)/r!) = v\_p(c\_r) + r [paper, from the definitions; at 446d1 and p = 5, PARI gives valuation 3 and c\_2 has valuation 1, as in [SW, §3.3]]. Let p ≥ 5 be a prime of good ordinary reduction with E[p] irreducible. (a) \textbf{The inequality of Proposition V holds without the conjecture and without p ∤ \#Ш}: L\_p lies in Z\_p[[T]] [GV, Prop. 3.7], and r independent rational points give ord\_T L\_p ≥ r through the bound ord\_T f\_E ≥ r for the characteristic series f\_E of the Selmer group over the cyclotomic Z\_p-extension (Mazur's control theorem [Gre99, Thm 1.2]) and Kato's divisibility of L\_p by f\_E away from p [Kato04, Thm 17.4(1)–(2)], which needs only p ∤ N and a\_p a unit (the chain is written out in [SW, §10]); hence c\_0 = … = c\_\{r−1\} = 0 and v\_p(L\_p\textasciicircum{}\{(r)\}/r!) = v\_p(c\_r) + r ≥ r. The content of Proposition V is therefore in its cases of equality. (b) Assume moreover \textbf{(H): some prime q ∥ N has p ∤ ord\_q(Δ\_min)}. Then the theorem of Skinner and Urban gives (f\_E) = (L\_p) in Z\_p[[T]] ⊗ Q\_p [SU, Thms 3.6.4 and 3.6.9: p odd, E[p] irreducible, the condition (dist), automatic for good ordinary p ≥ 3, and the prime q], and in Z\_p[[T]] itself once Kato's divisibility holds integrally; [Kato04, Thm 17.4(3)] asks for p ≠ 2 and an image of Galois containing SL\_2(Z\_p), and [Ski, §2.5 and Thm 2.5.2] shows that E[p] irreducible and (H) suffice. In [Kato04] the hypothesis on the image enters Theorem 17.4(3) only through Theorem 12.5(4) (§17.13), whose proof (§13.14) uses three things: that every Galois-stable lattice is a scalar multiple of a fixed one (through Remark 12.8 and the proof of Lemma 14.7, the only place where the image is used; it follows from E[p] irreducible [Ski, Lemma 2.1.1]); that H¹(T) is free [Kato04, Thm 12.4(3): p ≠ 2 and E[p] irreducible]; and the Euler system bound [Kato04, Thm 13.4(3)], which asks for some σ in Gal(/Q(ζ\_\{p\textasciicircum{}∞\})) with Coker(1 − σ) free of rank one. A generator of tame inertia at q gives such a σ: inertia at q acts on the Tate module through non-trivial unipotent matrices [DDT, Prop. 2.12(a)], non-trivial modulo p exactly when p ∤ ord\_q(Δ\_min) [DDT, Prop. 2.12(c)], and a unipotent matrix that is non-trivial modulo p has cokernel free of rank one. Section 13 of [Kato04] assumes that f has no complex multiplication, which holds here because j is not integral at q. With the formula for the leading term of f\_E up to a unit of Z\_p [Sch85, Thm 2′], [PR92, §3.4.3], a curve of rank two with a\_p = 1 and ord\_T L\_p = 2 has \textbf{v\_p(c\_2) = Σ\_ℓ v\_p(c\_ℓ) + v\_p(\#Ш(E/Q)[p\textasciicircum{}∞]) + v\_p(Reg\_p)}, and the inequality ≥ needs Kato's divisibility alone. In the boxes of Proposition 24ᶜ that also satisfy (H), v\_p(L\_p\textasciicircum{}\{(2)\}(E, χ⁰)/2!) = 2 holds if and only if Ш(E/Q)[p\textasciicircum{}∞] = 0 — the assumption p ∤ \#Ш of Proposition V becomes the statement it is equivalent to. (c) Under (H), with a\_p = 1 and two independent rational points, \textbf{if c\_2 ≢ 0 (mod p) then, by Kato's divisibility alone, rank E(Q) = 2, Ш(E/Q)[p\textasciicircum{}∞] = 0, the p-adic height pairing is non-degenerate, p ∤ ∏\_ℓ c\_ℓ and v\_p(Reg\_p) = 0} (that ord\_T f\_E = 2 forces Ш(E/Q)[p\textasciicircum{}∞] to be finite and the pairing to be non-degenerate is [PR92, Thm 3.4.2 and §3.4.3] together with [Kato04, Thm 17.4(1)]; [Sch85, Thm 2′] is the converse) — the last by a route independent of Proposition 24ᶜ — and, contrapositively, det(χ, γ) ≡ 0 in a split box forces c\_2 ≡ 0 (mod p) (Corollaries K′(i) and R′\_s, applied to Λ = E(Q), give v\_p(Reg\_p) ≥ 1). c\_2 modulo p is a finite sum of modular symbols with denominator p² [SW, Prop. 3.5]: with S\_j := Σ\_\{a=1\}\textasciicircum{}\{p−1\} [ω(a)(1 + p)\textasciicircum{}j/p²]\textasciicircum{}+ it is Σ\_\{j=0\}\textasciicircum{}\{p−1\} (j choose 2)·S\_j modulo p when a\_p ≡ 1 (mod p) and L(E, 1) = 0 [paper, from the Hecke relation; for general a\_p the sum is multiplied by α\textasciicircum{}\{−2\}, α the unit root; checked against PARI's \texttt{mspadicseries} on 14 pairs (E, p) of rank two and conductor below 950, six of them with a\_p ≡ 1]. The method is that of [SW, Algorithm 11.1], here at anomalous primes, and we claim nothing new in it; the p-adic L-functions of [SW] (period Ω\_E), of PARI (period Ω⁺ = E.omega[1], with Ω\_E/Ω⁺ ∈ \{1, 2\} [SW, §3]), of [Kato04] (a Néron differential and an integral basis γ\textasciicircum{}±, §17.5) and of [SU] (cohomological periods, §3.3.3) differ by units of Z\_p when p is odd and E[p] is irreducible: the Manin constant of an optimal curve is prime to p when p ∤ 2N [Maz78, Cor. 4.1], and the isogenies in the class then have degree prime to p [paper, from the definitions: for the curve that is optimal for X\_1(N), pulling back identifies its integral cohomology with the λ\_f-part of that of X\_1(N), so the periods of [SU, §3.3.3] are its Néron periods divided by its Manin constant, and Mazur's corollary reaches it through an isogeny of degree prime to p; [SU, §3.6.7] states the comparison of L\_E with L\_f without proof]. §5.12 applies (c) to the split boxes.

\textbf{The first digit of the conjecture of Mazur, Tate and Teitelbaum [paper for the reduction; computation, against a rule fixed before counting, for the check; the conjecture is read in [MTT, II §10] and used in the form of [SW, Conj. 5.1]].} The conjecture asserts an exact equality and not an equality up to a unit [MTT, II §10] (there L\_p is not divided by a period, the leading term is (1/r!)·d\textasciicircum{}r/ds\textasciicircum{}r in the variable s, and the right-hand side carries Ω\_E\textasciicircum{}+, the integral of the Néron differential over E(R), which is Ω\_E [MTT, II §8]); in the form of [SW, §5.1], for good ordinary p, L*\_p(E, 0) = ε\_p·∏\_v c\_v·\#Ш(E/Q)·Reg\_γ(E/Q)/(\#E(Q)\_tors)², with ε\_p = (1 − 1/α)², Reg\_γ = Reg\_p/log\_p(1 + p)\textasciicircum{}r and L\_p normalised by the Néron period Ω\_E [SW, §§3–5]. Statements (b) and (c) see it only up to a unit of Z\_p. Its first p-adic digit can be read off in rank two with a\_p = 1 and c\_2 ≢ 0 (mod p): α ≡ 1 − p (mod p²) gives ε\_p/p² ≡ 1 and log\_p(1 + p)/p ≡ 1 (mod p), the sign of the height and any constant factor in it cancel in a 2 × 2 determinant (so Reg\_p in our normalisation equals that of [SW]), and with c\_2 taken in the normalisation of PARI by Ω⁺ = E.omega[1] the conjecture says

\begin{quote}
c\_2 ≡ ω·n·\#Ш(E/Q)·Reg\_p (mod p),  ω := Ω\_E/Ω⁺ ∈ \{1, 2\},  n := ∏\_ℓ c\_ℓ,   (MTT₁)
\end{quote}

where ω = 2 exactly when Δ \textgreater{} 0 [paper, from the definitions; the trivial torsion of the boxes below is used]. In a split box Theorem K gives Reg\_p ≡ −B₀²·det(χ, γ)² (mod p) — the form 2λ\_\{T,1\}χ⊗χ + B₀(χ⊗γ + γ⊗χ) has determinant −B₀² in the coordinates (χ, γ) — and \#Ш(E/Q), when finite, is a square (Cassels). So (MTT₁) implies, with ( · /p) the Legendre symbol,

\begin{quote}
(−c\_2·n·ω / p) = +1,   (Q)
\end{quote}

a statement in which neither the regulator nor Ш appears: under the conjecture, the quadratic character of c\_2 is fixed by the Tamagawa product and the sign of the discriminant. Neither (MTT₁) nor (Q) follows from (b) or (c). The constraints that the functional equation puts on the S\_j (S\_j = S\_\{c−j\}, §5.12) and the vanishing of c\_0 and c\_1 are linear, and c\_2 is a linear form that does not vanish on the space they cut out, so the structure alone does not fix the quadratic character of c\_2; on random integer vectors satisfying these constraints it is a square in 0.49 to 0.58 of 5,000 cases for each p ≤ 19 [computation]. A rule fixed before counting tested both statements on the split boxes of rank two in Cremona's tables with 200000 \textless{} N ≤ 499000 — 5,556 boxes with p ∈ \{5, …, 19\}, p ∤ N, a\_p = 1, p ∤ ∏\_ℓ c\_ℓ, trivial torsion and ι, disjoint from those of §5.12; one of them was lost to memory in the computation of heights — with c\_2 computed as in §5.12, Reg\_p modulo p computed both from PARI's p-adic heights and from Theorem K, and \#Ш replaced by the order Ш\_an = (L″(E, 1)/2)/(Ω\_E·∏\_ℓ c\_ℓ·Reg\_∞) predicted by the complex conjecture. The two routes to Reg\_p modulo p agree in 5,273 of 5,273 boxes. Of the 5,454 boxes satisfying (H), c\_2 ≢ 0 in 5,180, which are exactly those with v\_p(Reg\_p) = 0; Ш\_an = 1 in all of them, and \textbf{(Q) and (MTT₁) both hold in 5,180 of 5,180} [computation, floating point, no interval arithmetic]. If the quadratic character of c\_2 were independent of n and ω, agreement in all of them would have probability 2\textasciicircum{}\{−5180\}; in each of the twelve strata by p and the sign of Δ the one-sided binomial test rejects that hypothesis (largest p-value 8·10\textasciicircum{}\{−28\}, Holm). The controls behave as the formula requires: replacing c\_2 by a random unit gives agreement in 0.49 of the boxes; dropping the factor −1 gives agreement in none of the boxes with p ≡ 3 (mod 4) and in all with p ≡ 1 (mod 4); dropping ω leaves exactly the boxes with Δ \textless{} 0 when p ≡ ±3 (mod 8); dropping n leaves exactly the boxes with (n/p) = +1. On the 761 non-split boxes with 200000 \textless{} N ≤ 205000, where Reg\_p has no closed form modulo p, (Q) holds in 323 of the 671 with c\_2 ≢ 0 (no prediction was made), while (MTT₁), with Reg\_p computed, holds in 671 of 671. On the 4,343 boxes of §5.12 with c\_2 ≢ 0, whose values of c\_2 were known before this rule was fixed, (Q) holds in 4,343 of 4,343 and (MTT₁) in the 2,901 of them with N ≤ 600000 where Ш\_an was computed. So the conjecture of Mazur, Tate and Teitelbaum agrees with the numbers at its first digit, unit included, in every box tested — a check that the main conjecture cannot supply, and not a proof of anything; it rests on the computed c\_2, on Ш\_an in place of \#Ш, and on the form of the conjecture as stated in [SW]. The normalisations it confirms — the factor ω and the sign — are those of [SW], with L\_p computed from twisted L-values, so that no Manin constant enters.

\textbf{The second digit, at the boxes where the first one vanishes [paper for the reduction; computation, against a rule fixed before counting, for the check; the conjecture in the form of [SW, Conj. 5.1]].} Without dividing, the conjecture says c\_2 = ω·n·\#Ш(E/Q)·U·Reg\_p with U := (ε\_p/p²)·(p/log\_p(1 + p))², an equality in Q\_p (rank two, trivial torsion). For a\_p = 1 one has α ≡ 1 − p − p² (mod p³), hence ε\_p/p² ≡ 1 + 4p, log\_p(1 + p)/p ≡ 1 − p/2 and U ≡ 1 + 5p (mod p²). In a split box with det(χ, γ) ≡ 0, Theorem K gives Reg\_p ≡ 0 (mod p) and (c) gives c\_2 ≡ 0 (mod p), so (MTT₁) reads 0 ≡ 0. One digit further, the conjecture says

\begin{quote}
c\_2/p ≡ ω·n·\#Ш(E/Q)·(Reg\_p/p) (mod p)  when v\_p(Reg\_p) = 1 and p ∤ \#Ш(E/Q),   (MTT₂)
\end{quote}

and c\_2 ≡ 0 (mod p²) when v\_p(Reg\_p) ≥ 2 or p | \#Ш(E/Q) (then p² | \#Ш, Cassels); U enters only modulo p, where it is 1. A box with v\_p(Reg\_p) = 1 and c\_2 ≡ 0 (mod p²) would, under the conjecture, have p | \#Ш(E/Q): this is where the p-part of Ш would be read from the side of the L-function. Since L\_p ≡ P\_3 modulo (1 + T)\textasciicircum{}\{p²\} − 1, whose coefficients of T\textasciicircum{}k for k \textless{} p are divisible by p², c\_2 modulo p² is the coefficient of T² in P\_3, built from μ\_α on the classes modulo p³; we computed it from the twisted L-values of the characters of conductor p³ and p² (for the latter through the Hecke relation Θ\textasciicircum{}\{(3)\} = a\_p·Θ\textasciicircum{}\{(2)\}), as in §5.12, and checked it against PARI's modular-symbol series (\texttt{mspadicseries}) on 14 pairs of conductor ≤ 944, among them 446d1 with p = 5, where it gives c\_2 ≡ 10 (mod 25), twice the value of [SW]. Theorem K gives no closed form for Reg\_p/p. We computed Reg\_p modulo p\textasciicircum{}\{v+2\} from PARI's p-adic heights and, independently of PARI's algorithm for heights, from a p-adic σ function built from the formal group together with the denominators of multiples of the points (sharing with PARI only s₂ and the formal expansions); the two agree in 674 of 674 boxes. The rule covered all 574 split boxes with det(χ, γ) ≡ 0 available to us: the 274 boxes with (H) and c\_2 ≡ 0 among those of the previous paragraph, and the 300 boxes of §5.12 (whose first digit c\_2 ≡ 0 had been seen). Ш\_an = 1 in all of them. In the 398 boxes with v\_p(Reg\_p) = 1, \textbf{(MTT₂) holds in 398 of 398} (agreement by chance has probability 1/p; rejected in each of the six classes of p, Holm); in the 176 boxes with v\_p(Reg\_p) ≥ 2, \textbf{c\_2 ≡ 0 (mod p²) in 176 of 176} [computation, floating point without interval arithmetic, Ш\_an in place of \#Ш]. Read the other way, as in (c): in each of the 398 boxes, (H), the computed c\_2 ≢ 0 (mod p²) and v\_p(Reg\_p) ≥ 1 (the split corollary, the generators being saturated at p) give by Kato's divisibility alone 1 = v\_p(c\_2) ≥ v\_p(\#Ш(E/Q)[p\textasciicircum{}∞]) + v\_p(Reg\_p) ≥ v\_p(\#Ш(E/Q)[p\textasciicircum{}∞]) + 1, hence rank E(Q) = 2, Ш(E/Q)[p\textasciicircum{}∞] = 0 and v\_p(Reg\_p) = 1 for these curves and primes [known theorems combined, as in (c); computation for c\_2]. As a positive control, in the 100 boxes of smallest conductor with c\_2 ≢ 0 of the previous paragraph, c\_2 ≡ ω·n·Ш\_an·U·Reg\_p (mod p²) in 100 of 100; replacing U by 1 leaves exactly the 34 boxes with p = 5, where U ≡ 1 (mod 25). Dropping ω or n from (MTT₂) leaves exactly the boxes in which that factor is ≡ 1 (mod p), and replacing Reg\_p/p by a random unit gives agreement at the rate 1/(p − 1). So the conjecture agrees with the numbers at the second digit as well, in every box tested. But no box had p | Ш\_an, so the case in which the p-part of Ш would be read from the L-function did not occur, and the check does not separate the conjecture from the weaker statement with \#Ш replaced by 1. It is a check and not a proof; it rests on c\_2 computed in floating point, on Ш\_an, on s₂ from PARI, and on the form of the conjecture in [SW].

\textbf{Relation with the refined conjectures of Mazur and Tate [secondary sources; [MT87] not retrieved].} As restated in [PB19, §2] and [Ota, Conj. 2.4], the leading-term formula of [MT87, Conj. 4] concerns levels in which the primes of good reduction occur to the first power, with the tamely ramified trivialisation at those primes (values in F\_p\textasciicircum{}×). The coefficient c\_2 modulo p (resp. p²) is the image of the p-part of the Mazur–Tate element of level p² (resp. p³) in I²/I³ ≅ Z/p (resp. Z/p²), on the wildly ramified quotient. As far as these restatements show, (MTT₁) and (MTT₂) lie outside the leading-term formula of [MT87, Conj. 4] and are instances of [MTT] at the first layers of the cyclotomic Z\_p-extension. The form −B₀²·det(χ, γ)², with χ the discrete logarithm on E(F\_p) ≅ Z/p, resembles the discrete-logarithm regulators of [MT87], and the boxes with det(χ, γ) ≡ 0 resemble the situation, a vanishing regulator, addressed by the derived heights of [BD94]; we have not checked either resemblance against the sources. Numerical studies of [MT87] in rank one and at split multiplicative levels are [PB19] and [LC24]; [BH25] proves large parts of the conjectures of [MT87].

\textbf{Where the p-part of Ш enters [computation, against a rule fixed before counting, for the check; paper for the reduction; the conjecture in the form of [SW, Conj. 5.1]].} The case left open in the previous paragraph does not occur in rank two within Cremona's tables: in the files \texttt{allbsd} of [Cre] (conductor \textless{} 500000; 348,672 curves of rank two) the analytic order Ш\_an of every curve of rank two is 1 or 4 (4 for 42 curves), so no curve of rank two has p | Ш\_an with p ≥ 3; on the 42 curves with Ш\_an = 4 and on all 742 curves of rank two in three windows of conductor, PARI's value of Ш\_an agrees with the tables. We therefore took the check one rank down. In rank zero it is empty: the interpolation property gives c\_0 = ε\_p·L(E, 1)/Ω⁺, and Ш\_an is defined from the same number. In rank one, with P a generator and t = \#E(Q)\_tors, the height of [SW, (4.1)] is 2·log\_p(e(P)/σ\_p(P)), the negative of h\_p(P), and the conjecture reads c\_1 = −ω·n·\#Ш(E/Q)·U₁·p·h\_p(P)/t² with U₁ := (ε\_p/p²)·(p/log\_p(1 + p)) ≡ 1 + 9p/2 (mod p²) for p ≥ 5 (for p = 3 the term p²/3 of the logarithm gives U₁ ≡ 7 (mod 9)). With e := 1 + v\_p(h\_p(P)) + v\_p(\#Ш(E/Q)), it says v\_p(c\_1) = e and

\begin{quote}
c\_1/p\textasciicircum{}e ≡ −ω·n·(\#Ш/p\textasciicircum{}\{v\_p(\#Ш)\})·(h\_p(P)/p\textasciicircum{}\{v\_p(h\_p(P))\})·t\textasciicircum{}\{−2\} (mod p),
\end{quote}

with \#Ш = \#Ш(E/Q); here, by the integrality theorem of the introduction (when P reduces onto E(F\_p)), v\_p(h\_p(P)) = −1 in the non-split case and ≥ 0 in the split case. Reading it needs c\_1 modulo p\textasciicircum{}\{e+1\}, the coefficient of T in P\_\{e+2\}. For P\_m with m ≥ 4 the Hecke relation of the previous paragraph acquires a second term, Θ\textasciicircum{}\{(m)\}\_\{pk\} = a\_p·Θ\textasciicircum{}\{(m−1)\}\_\{k\} − p·Θ\textasciicircum{}\{(m−2)\}\_\{k/p\} when p | k (the sum of χ(c)·x(c/p\textasciicircum{}\{m−2\}) over the units c modulo p\textasciicircum{}\{m−1\} is p·Θ\textasciicircum{}\{(m−2)\} when χ has conductor at most p\textasciicircum{}\{m−2\}); it does not arise for m = 3. With it, the coefficients c\_1 and c\_2 of P\_m agree with PARI's modular-symbol series on 25 pairs (p = 3, 5, 7, m ≤ 6, nine of them with a\_p = 1), and for 446d1 with p = 5 we find c\_2 ≡ 185 (mod 625) and c\_3 ≡ 295 (mod 625), twice the values of [SW]. The rule covered every box of rank one in the tables with p² | Ш\_an, p ≥ 3, p ∤ N, a\_p = 1, p ∤ ∏\_ℓ c\_ℓ, p ∤ t and (H): 41 boxes, 40 with p = 3 and Ш\_an = 9 and one with p = 5 and Ш\_an = 25 (the other 64 candidates, among them all but one with p = 5 and both with p = 7, have E[p] reducible). Ш\_an from the tables and from PARI agree in all 41. In \textbf{41 of 41 boxes v\_p(c\_1) = e} — the valuation exceeds by exactly v\_p(Ш\_an) = 2 the value predicted with Ш replaced by 1, which agrees in 0 of 41 — \textbf{and the first digit agrees in 41 of 41} (agreement by chance has probability p\textasciicircum{}\{−3\} per box; rejected for p = 3 and for p = 5, Holm) [computation, floating point without interval arithmetic, Ш\_an in place of \#Ш]. The sign and the factor U₁ were fixed before on 90 boxes of rank one with Ш\_an = 1 (p = 3, 5, 7, conductor ≤ 20000): the sign −1 gives 90 of 90 and the opposite sign none, and in the 82 non-split boxes c\_1 ≡ −ω·n·U₁·p·h\_p(P)/t² (mod p²) holds in 82 of 82 and in none with U₁ replaced by 1. The sign is that of the form of [SW], whose height 2 log\_p(e/σ) [SW, (4.1)] is −h\_p; the analytic height of [MTT, II §5], ⟨P, P⟩ = log\_p(σ²/d) with d the denominator of x(P), has the sign of h\_p, so that with the height as printed there the formula of [MTT, II §10] has the opposite sign in odd rank ([Sch85, Thm 6] relates the algebraic and the analytic height by a sign); in rank two the sign cancels in the determinant [paper, from the definitions; computation for the 90 boxes]. The implementation of P\_m reproduces c\_2 modulo p² in 50 of 50 boxes of the previous paragraph. So, at the first digit and in every box available, the p-part of Ш predicted by the complex conjecture is read correctly from the p-adic L-function: the p-adic and the complex conjectures agree on the p-part of Ш there. In rank one this agreement is largely forced by theorems. Kato's divisibility, the theorem of Skinner and Urban (whose hypotheses include (H) and allow p = 3: [SU] assumes p odd, and its condition (dist) holds for good ordinary p ≥ 3) and the leading term on the algebraic side give v\_p(c\_1) = 1 + v\_p(h\_p(P)) + v\_p(\#Ш(E/Q)[p\textasciicircum{}∞]), so that, read the other way, the computed valuations give \#Ш(E/Q)[p\textasciicircum{}∞] = p² for these 41 curves [known theorems combined; computation for c\_1]; for semistable curves [JSW] proves the p-part of the complex formula; and we expect, without having checked the sources, that the p-adic and complex Gross–Zagier formulas make the first digit a consequence as well. The check is therefore mainly one of P\_m, of the heights and of the normalisations; it is a check and not a proof, and in rank two the case p | Ш remains untested within the tables.

\textbf{A search outside the tables [computation, against a rule fixed before counting; the heuristics recalled, not checked].} To find a box of rank two with p² | Ш\_an, p ≥ 5, a\_p = 1 and (H) we searched quadratic twists E\textasciicircum{}\{(d)\}. Twisting keeps E[p] irreducible and keeps (H) at the primes q ∤ d, and a\_p(E\textasciicircum{}\{(d)\}) = χ\_d(p)·a\_p(E), so from a curve with a\_p(E) = ±1 the twists with χ\_d(p) = a\_p(E) have a\_p = 1. Nothing links Ш(E)[p] to Ш(E\textasciicircum{}\{(d)\})[p]; as bases we took the 30 curves of rank zero in [Cre] with 25 | Ш\_an (and the other conditions) of smallest Ω\_E·∏\_ℓ c\_ℓ, which raises the scale √|d|/(Ω\_E·∏\_ℓ c\_ℓ) of Ш·Reg. Among the twists with root number +1 and conductor at most 1.2·10⁹, and a part, cut by time, of those up to 3.1·10¹⁰, 82 have analytic rank two and 2-descent upper bound two; for 60 of them PARI found two points, which we saturated at all primes up to 50, and Ш\_an = 1 in all 60. The same scan with p = 3 (conductor at most 2.9·10⁹) gave 30 twists of rank two, 24 of them read, all with Ш\_an = 1. In the same scans the twists of rank zero have Ш\_an \textgreater{} 1 in 1278 of 1730 cases (25 | Ш\_an in 303 of them), while in rank two the product Ш\_an·Reg, of comparable or larger size (median 37 to 140), is carried entirely by the regulator. The same functions return Ш\_an = 4 on all 39 curves of rank two with Ш\_an = 4 in [Cre] on which we ran them. So no box of rank two with p | Ш\_an was found in this range. This is consistent with heuristics of Delaunay's type, which as we recall them (without having checked) give a probability of about p\textasciicircum{}\{−(2r+1)\} for p | \#Ш in rank r, and which the model of [PR] is stated to be compatible with. Even if such a box exists, reading c\_2 modulo p\textasciicircum{}\{e+1\}, e = v\_p(Reg\_p) + v\_p(\#Ш), by the method of the previous paragraph needs about p\textasciicircum{}\{e+2\}·√N terms of the L-series; for p = 5 and e = 2 this limits the conductor to about 5·10⁷ within our memory bound. The general form v\_p(c\_2) = v\_p(Reg\_p) + v\_p(\#Ш) of the conjecture therefore remains untested here in the case p | Ш.

\section*{7. What is open}

\begin{enumerate}[leftmargin=1.4em,itemsep=.2em,topsep=.3em]
\item \textbf{The quadratic residue class of the slope.} Over every anomalous curve over F\_p with p ≤ 157 the slope ρ of §4.6 satisfies (ρ/p) = (−3/p), so that −3ρ is a square and ρ() = −3μ()² for a homomorphism μ : E(F\_p) → F\_p determined up to sign. This is [computation] and nothing more; μ has no closed form here. At the cusp ρ\_j ≈ 9j²/Δ, and the shape ρ = 9L()²/Δ with L a homomorphism is excluded by the residue class together with the squareness of −Δ.
\item \textbf{One source under §4.6.} The step "E\textasciicircum{}\{(p)\} is the member of the family at parameter t\textasciicircum{}p" uses that the Tate parameter q, as a series in 1/j, has coefficients in the prime field. j(q) ∈ Z[[q]] is standard; the inverse series is what we did not locate in a source (§9), and [ATAEC] itself we did not retrieve.
\item \textbf{Exact values rather than bounds.} Corollaries R′ and R′\_s turn part (iii) into v\_p(Reg\_p) = (r − 2) + v\_p(det T), so what is left open is exactly whether det T is a unit — a non-degeneracy of p\textasciicircum{}\{−1\}h\_p on the orthogonal complement of e₁, of the same kind as Schneider's conjecture, and not decided by any local datum here. It fails in 36 of the 7,759 non-split boxes and in 5 of the 48 split boxes of rank 3 (§5.6). On the split side of rank one the value is decided, by λ\textasciicircum{}\{(2)\} (Proposition R″\_s), but λ\textasciicircum{}\{(2)\} is not a function of the reduction: it needs E mod p³ and the point modulo p². The behaviour at p = 5 with \#E(F\_5) = 10 remains outside everything after Lemma 4. One archimedean datum decides the split case of rank two in part: v\_p(Reg\_p) = 0 whenever h\_min \textless{} 4 log p (Proposition 24). Whether archimedean data decide it in general is open.
\item \textbf{The geometry in Lean.} §6 lists what would have to exist in Mathlib for the main theorem to become a Lean theorem.
\item \textbf{The layer one deeper.} Corollary K′(iv) reduces the exact value in rank ≥ 3 to det Ĉ ≢ 0, the non-degeneracy of p\textasciicircum{}\{−1\}h\_p on Λ₁ = Λ ∩ E\_2(Q\_p). Nothing here decides it; the local data of Theorem K do not enter it.
\item \textbf{What q\_p(d′) is.} Theorem N and Corollaries N⁽³⁾, N⁽⁴⁾ show that the digits of the height past the first are not functions of the position of the point in the filtration of E(Q\_p): the displacement is carried by the Fermat quotient of the denominator d′, to one digit at the second layer and to two at the third. Two partial answers are in §4.12: the denominator descends at rate p\textasciicircum{}k while the σ-part descends at rate p\textasciicircum{}\{ke−2\}, so the two non-local ingredients of Lemma 3 arrive at different speeds and cannot be merged into one homomorphism; and q\_p(d′) is paired with the normalisation of the model, since the right-hand sides of those theorems change by 2p·q\_p(u) under x ↦ u²x. Theorem N* sharpens the first of the two: in the closed form the two non-filtration ingredients are not symmetric, because by Proposition 19 the third term ℒ(z(P″)) \textbf{is} a function of the layers t₁, while the Fermat quotient F is not a function of anything in E(Q\_p). So the whole of the non-filtration content of the height sits in one of the two Iwasawa logarithms, and in no part of the formal group. An intrinsic description of q\_p(d′) alone we still do not have, and by the remark on the model we no longer expect one.
\item \textbf{The class dependence of λ\_T one digit deeper.} (L2) gives λ\_T(m) = λ\_\{T,1\}m² modulo p. Modulo p² neither Θ\_T(m) = m²Θ\_T(1) with the standard lift of m nor the same identity with the Teichmüller lift holds (§5.7), and no law is proposed here. Theorems N and N⁽³⁾ do not need one, since the class of nP is fixed there.
\item \textbf{What is left around (ND).} The non-degeneracy itself is settled (Theorem ND). Three smaller things around it are not. (i) The primary sources for the classical inputs of §4.13 — [Kat73] on the canonical subgroup and the theta operator, [Igu] on the irreducibility of the tower, [SwD73] on the kernel of the q-expansion map, [ATAEC] for the Tate parametrisation — we did not retrieve; we used the statements as standard and verified their consequences on q-expansions ourselves. (ii) Whether Propositions 14 and 15 are in the literature we do not know; see §8. (iii) The counting of §4.13 gives the constants c\_S, c\_K, c\_p in closed form, but the \textit{ratio} of the two sides of (R′) at a supersingular point, where h degenerates, is outside everything here.
\item \textbf{What the minimality of Theorem G⁽ᵏ⁾ does and does not say.} The necessity half is a statement about the parameter space that Theorem N* defines — t₁ mod p\textasciicircum{}\{k+1\} and F mod p\textasciicircum{}\{k−1\}, the local invariants fixed. Both kinds of necessity are realised by rational points of one curve: the layer γ\_k by P and P + p\textasciicircum{}\{k−1\}S with S ∈ ker χ, γ(S) ≠ 0 (Proposition 21 (iii); k = 3 and k = 4 at 4407a1, p = 5, read pointwise — v\_p(n²h\_p(P)) in place of v\_p(Reg\_p) — since 4407a1 has rank two), and the second digit of F by P = −5g₁ + 6g₂ and P + 25(−g₁ + g₂) at the same curve (§5.13). What is not available on rational points is independence of the two kinds of datum along p-power translates (Propositions 21 and 22); item 14 records what is open beyond them.
\item \textbf{The valuation of the coefficients of ℒ.} Proposition 18 puts the m-th coefficient of ℒ at the digit 2me − 2 + v\_p(μ\_\{2m\}), and v\_p(μ\_\{2m\}) can be negative: v\_5(μ₆) = −1 is why μ₆ is visible one digit earlier at p = 5. All we have here is v\_p(μ\_\{2m\}) ≥ −v\_p((2m + 2)!) from the shape of the double integral, which is far from sharp. Without a sharp bound the digit of a given local invariant is known for m ≤ 3 only, and the ordering of the invariants μ₄, μ₆, μ₈, … by depth is not known for all p. One special case the computation settles: in the twenty cells with j = 0 of §5.11 the series has \textbf{only the coefficients μ\_\{6k\}} — μ₂ = μ₄ = μ₈ = μ₁₀ = μ₁₄ = μ₁₆ = 0 while μ₆, μ₁₂ ≠ 0, measured to z\textasciicircum{}\{16\} — which is [computation], and is why the leading term there is m = 3.
\item \textbf{Corollary N‴ has almost no test.} In conductor ≤ 20000 there is exactly \textbf{one} box with v\_p(Reg\_p) ≥ 4, so §5.8 records the absence of a violation and not a test of the criterion. Widening to conductor 60000 would give tens of boxes with v\_p ≥ 3.
\item \textbf{Rank ≥ 2 at the third and fourth digits.} Corollaries N⁽³⁾ and N⁽⁴⁾ are proved for any rank, but all the sweeps of §5.8 are of rank one; the rank-two entries checked for Theorem N (§5.7) have no counterpart there.
\item \textbf{Above the archimedean threshold.} Where the deficit ends is explained: below the canonical threshold 4 log p + β(E) of Proposition 24ᶜ the frequency is 0 by the proposition, and that threshold carries the whole deficit in expectation (100.6 \% and 99.7 \% on two disjoint sets of split boxes, §5.12), where the naive threshold of Proposition 24 carries 68 \% and 56 \%. \textbf{What is not explained is the frequency above it}: there det(χ, γ) ≡ 0 occurs 383 times against 382.6 for 1/p on the sets from Cremona's tables (ratio 0.995 to 1.018 on the two sets; 0.91 to 1.46 by p); on two further sets of curves outside Cremona's tables the boxes with det(χ, γ) ≡ 0 crowd just above the threshold (1.95 and 1.82 times 1/p in [4, 4.5), §5.12), while the frequency above the threshold stays at 1/p on the set not built from two given points; and nothing here says why γ should be a random multiple of χ once no element of ker χ is short enough, nor why the boxes with det(χ, γ) ≡ 0 should crowd just above the threshold (choosing the shortest element without the p-adic part of its height leaves the crowding unchanged, §5.12). Lemma 25 makes 1/p a local count, so both are correlations between a local class and a global height. Whether the size of the crowding follows the conductor, the construction of the set or the position β(E)/log p is not decided by our data (§5.12). Nor is the box-by-box statement a theorem: β(E) contains the numerical minimum m\_∞ of λ′\_∞. Two further things are left. In the split row of rank two, P(k ≥ 2 ∣ k ≥ 1) is 25 against 13.6 for 1/p; counting two independent routes to k ≥ 2 (ρ₁ = 0, and ρ₁ = 1 with v\_p(h\_p(Q)) ≥ 2 for Q ∈ ker χ) gives 2/p − 1/p² and an expectation of 25.05, but that count was found after the fact and has not been tested box by box. And closures of index p² occur in two boxes up to conductor 500000 (§5.12), too few to test what the non-degeneracy of the third layer says there. [Maz72] treats the p-adic closure of the Mordell–Weil group, and a conjecture attributed to Lang concerns it; we retrieved neither, and state nothing from them.
\item \textbf{Translates that are not p-th powers.} Proposition 21 determines F along P + p\textasciicircum{}aΛ by local quantities. For R ∈ Λ ∩ E\_k(Q\_p) not in p\textasciicircum{}\{k−1\}Λ — for instance R = −13g₁ + 19g₂ ∈ E\_3(Q\_5) at 75393a1 — it says nothing, and a pair P, P + R separated by min(v\_p, k) would move F by something global. How often a split box of rank two has a point with v\_p(n²h\_p(P)) ≥ k at all, the hypothesis of Proposition 21 (iii), is a question about isotropic vectors of q modulo p\textasciicircum{}k that we have not counted.
\end{enumerate}

\section*{8. What is not claimed}

\begin{enumerate}[leftmargin=1.4em,itemsep=.2em,topsep=.3em]
\item \textbf{Nothing here bears on the conjecture of Birch and Swinnerton-Dyer}, nor on Ш, on L-functions or on the rank: every statement is about the valuation of a p-adic regulator under hypotheses on the reduction at one prime, except Proposition V of §6, which assumes the p-adic conjecture and is marked as such. Schneider's conjecture is not proved either; its CM case is [known], due to [Ber82], which we did not retrieve.
\item \textbf{The baseline is not ours.} Theorem A is elementary and known (§9); what we add is the splitting hypothesis and its mechanism.
\item \textbf{The local criterion is not ours.} The equivalence of Proposition 6 is [known] (Gross's criterion [Gro90], quoted as [LR, Thm 4.1] under the hypotheses j ≠ 0, 1728 and E[p] irreducible; [OY, §2.1.1, Lem. 4.1] without them; [DW, Lemma 3.2 without hypotheses on j, Cor. 3.5 with j ≠ 0, 1728]); we give an independent proof and a presentation through a Fermat quotient, and we do not give it a name. [Gro90] itself we did not retrieve, and cite it as quoted in [LR] and [OY].
\item \textbf{The theorem is a lower bound in rank ≥ 2.} The exactness of §5.3 in the non-split case is Theorem S; Corollaries R′ and R′\_s reduce the rest to the valuation of a Schur complement, which is not predicted here on either side, and Theorem R\_s does not say which of the three values ρ₁ ∈ \{0, 1, 2\} occurs. Deuring's lifting theorem, Vélu's formulas, Serre–Tate theory, the composition rule for division polynomials, the coprimality of φ\_p and ψ\_p², and the Tate parametrisation are used as standard facts, and for the last we cite sources that quote [ATAEC] rather than [ATAEC] itself.
\item \textbf{The hypotheses are as used, not the weakest possible.} Everything from Lemma 4 on assumes p ≥ 5 and \#E(F\_p) = p; Proposition R⁺ (§3.9) covers \#E(F\_p) = 2p, which by Hasse happens for p = 5 alone, and that is the whole of the extension — the main theorem itself is not extended, and neither is anything from §3.10 on.
\item \textbf{The computations are not kernel-checked, and the statistics are not statements.} All of §5 is PARI/GP with controls, at stated precision and over stated ranges; an earlier round used the height belonging to σ ≠ σ\_p, and every count here has been recomputed with h\_p.
\item \textbf{Two inputs are used as verified, not as proved.} Katz's congruence E₂\textasciicircum{}* ≡ E\_\{p+1\}/E\_\{p−1\} (mod p) ([Kat73]), on which the Remark of §2 rests, is [known] and \textbf{the original we did not retrieve}; what we did was to check its consequence against PARI's \texttt{ellpadics2} over 315,953 boxes and to verify its two classical ingredients (E\_\{p−1\} ≡ 1, E\_\{p+1\} ≡ E₂) on q-expansions ourselves. Likewise the definition of \texttt{ellpadics2} is taken from PARI's documentation and calibrated on the one worked example there; \textbf{the source code we did not read}. So the identification of c mod p with a function of (Ē, ) is [paper] granted a [known] congruence we have checked but not proved, and every number in this paper that involves c inherits that status.
\item \textbf{The digits past the fourth are truncations carried out by machine.} Theorem N* is an exact identity, so the selection rule of Proposition 18 is exhaustive: no term can appear at a later digit that is not one of the three series. What is [computation] in Corollaries N⁽⁵⁾, N⁽⁶⁾ and N⁽⁷⁾ is that the truncation was performed and checked in 648 boxes rather than written out and proved term by term; and since the sharp lower bound for v\_p(μ\_\{2m\}) is not known here (§7), the digit at which a given coefficient of ℒ becomes visible is proved for m ≤ 3 only. The valuation table of the Remark opening §4 remains a table: it lists the terms of the expansion we carried out by hand, and it is Proposition 18 and not the table that is a statement about every digit. Each of the three corollaries is checked \textbf{in the form in which it is displayed}, and not only in the form the generating rule produces — the two can differ, and at one place in an earlier version they did (§5.10).
\item \textbf{Propositions 14 and 15 may be known.} Their ingredients are classical — [p] = V∘F, the canonical subgroup, the theta product for the Tate curve, the injectivity of q-expansions in a fixed weight, and Lemma 13 — and the shape of Proposition 15, a Kodaira–Spencer derivative of the étale p-division points, is the kind of statement that belongs to the standard theory of the theta operator on p-adic modular forms. We did not find either statement in the sources we read (§9), and we do not call them new; what we claim is that they hold and that the proofs in §4.13 are complete.
\item \textbf{Theorem N* is an assembly of standard tools.} Its ingredients are classical: Mazur and Tate's differential equation for σ\_p, the Iwasawa logarithm of a unit written through its (p − 1)-st power, and the ℘-expansion of x. That log(x·σ\_p(t)²) is the local term of the p-adic height is close to the definition of σ\_p itself, so the closed form is the kind of statement that may well be written somewhere. We looked for it as a displayed statement in the sources of §9 and did not find one; \textbf{we do not call it new.} What we do claim is the bookkeeping: that the congruences of Theorem N and of Corollaries N⁽³⁾–N⁽⁷⁾ are truncations of that one identity, that the series ℒ is even with coefficients μ₂ = −c, μ₄ = −b₄/12, μ₆ = −b₂b₄/360 − b₆/30 in which E₂ occurs in μ₂ alone, and that a local invariant becomes visible at the digit 2me − 2 + v\_p(μ\_\{2m\}). Proposition 17 is [paper], Proposition 18 is [paper] for m ≤ 3, and the digits past the fourth are [computation].
\item \textbf{No claim of priority, and no external review.} We searched [MST], [SW], [Har], [BMS], [BD], [Ban], [Wut], [BG], [LR], [OY], [DW], Sage's documentation and arXiv (§9): the baseline we found, the local criterion we found, the equivalence with integrality of the height matrix we did not, in a narrow search — no MathSciNet or zbMATH, and several primary sources not retrieved — so the statements may well be known. Observation, proofs, programs, Lean development and draft were produced by the same agent, and Lean's kernel, which covers §6 alone, is the only check not internal to the authors.
\item \textbf{Lemmas 23 and 23′ are elementary and should be regarded as known.} That a rational point deep in the formal filtration has a large denominator, and hence a large naive height, is the starting point of the study of valuations of elliptic divisibility sequences and division polynomials ([NV] and the work it cites); we did not look for Lemma 23′ in that literature and do not call either lemma new. Proposition 24 is their composition with Corollary K′(i), and Proposition 24ᶜ is the same composition with the decomposition of the canonical height into local heights, which is known [SilH]; we do not call it new. The statistics of §5.12 are [computation]: they say that the threshold of Proposition 24 accounts for about 68 \% of the observed deficit of det(χ, γ) ≡ 0 and that of Proposition 24ᶜ for all of it in expectation, on the boxes measured, and not that nothing else acts; the frequency 1/p above the threshold is not a theorem.
\item \textbf{Propositions 21 and 22 are within a known framework.} Proposition 21 follows from the quadraticity of the height, Theorem N* and the fact that φ\_p mod p is a polynomial in x\textasciicircum{}p; that the denominator term of the p-adic height behaves p-adically along p-power multiples is the kind of statement that belongs to the theory of [MST]. We did not find it displayed, and we do not call it new. Proposition 22 is Theorem N* solved for F.
\item \textbf{Proposition V assumes a conjecture.} It is the only statement of the paper that does, it is stated in §6 and in the companion note and not in §3, and nothing in §3 depends on it. The extra valuation 2 of the Euler factor at an anomalous prime is an elementary consequence of the interpolation factor of [MTT] and is probably known; in [MTT, II §10] it is the factor (1 − 1/α)², and the conjecture is used in the form of [SW, Conj. 5.1] and of PARI's \texttt{ellpadicbsd}.
\item \textbf{The canonical threshold is applied with a numerical constant.} β(E) contains m\_∞ = min λ′\_∞ over E(R), computed on a grid with refinement and checked by a second route, with no interval arithmetic; at I\_n* with c\_ℓ = 2 the rational component was not determined and the smaller correction used. The statement "no box with det(χ, γ) ≡ 0 below the threshold" is protected against both — a sampled minimum is at least the true one, which only makes the check stricter, and the smaller correction at I\_n* is still a valid lower bound — but the shares of §5.12 are not. The samples A and B are the curves of Cremona's tables satisfying the box conditions, C is a family of curves through two small points, and D a family of random short Weierstrass equations on which two small independent points were found, C and D with Mordell–Weil groups assumed saturated beyond 200; none is a random sample, and the effect just above the threshold is stated for these sets only.
\item \textbf{The effect just above the threshold is a verdict of a rule, not an explanation.} "Effect" in §5.12 is the name of an outcome of a rule fixed before counting (Poisson z \textgreater{} 2 on a set alone); it names no cause. We have not found the cause, and the tests of whether its size follows the conductor, β(E)/log p or the construction of the set came out undecided with little power; we claim no dependence and no independence. Lemma 25 is elementary and [paper], and its check on the classes modulo p² is [computation].
\end{enumerate}

\section*{9. Related work}

\textbf{Where the baseline already is.} In the tables of [MST, §§4.2–4.6], for curves of rank 1 to 5, three entries have negative valuation — 37A at p = 53, 5077A at p = 5, and the rank-four curve y² + xy = x³ − x² − 79x + 289 of conductor 234446 at p = 5, all three equal to −2 — and after the shift by r of §2 only the first of them stays negative, at −1. The third of those entries we recomputed in full, including the hypotheses [MST] does not state: it has \#E(F\_5) = 10, so it is a box of Proposition R⁺ and not of the main theorem (§5.3). At the first, the text notes that \#E(F\_p) = p, "i.e., p is anomalous"; at the second it notes a\_p ≡ 1 (mod 5) for the same reason. So the phenomenon is there, in three places, with no bound stated and no index in sight. [Har] has the mechanism implicitly rather than as a statement: the height of a general point is h\_p(Q) = h\_p(n\_\{lcm\}Q)/n\_\{lcm\}² with n\_\{lcm\} = lcm(\#E(F\_p), lcm\_ℓ c\_ℓ) (§2), so an anomalous p costs 2 v\_p(n\_\{lcm\}) digits of working precision — which is the "cost 2" of Theorem A read as a loss of precision. That reading is ours; what is in [Har] is the normalisation itself, the n² growth of the coordinates and the general warning about loss of precision, and we did not find the quantity 2 v\_p(n\_\{lcm\}) written there. [Ban, §2.2] uses the contrapositive as a hypothesis, non-anomalousness being what "extends the regulator estimate of Proposition 4.2 from E°(Q) to all of E(Q)", and that Proposition 4.2 is the positive half of Theorem A in rank two; Sage's \texttt{padic\_\allowbreak{}regulator} documents the same example. So the baseline is known to practitioners; we did not find it written as a proposition, and we do not need it to be new.

\textbf{Where the local identity (L1) sits: the anomalous discrete logarithm.} (L1) says that on a split anomalous curve the discrete logarithm χ of the reduction is given by an algebraic expression in the coordinates, χ = s·c₁/B₀ — that is, by the second p-adic digit of φ\_p′ evaluated at the point itself. Since \#E(F\_p) = p, the group Hom(E(F\_p), F\_p\textasciicircum{}+) is one-dimensional over F\_p, so \textbf{any} additive map E(F\_p) → F\_p is a scalar multiple of any other, and (L1) is necessarily a multiple of the map of [Sem98]. What the sources do and do not contain is worth stating precisely, because the three papers usually cited together are of two different kinds.

[Sem98] stays in characteristic p: for Q of order p he takes f\_Q with divisor p(Q) − p(O) and sets φ(Q) = (f′\_Q/f\_Q)(R) at an auxiliary point R, proving additivity from f\_\{Q₁+Q₂\} = c·f\_\{Q₁\}f\_\{Q₂\}; his Lemma 1 gives (f′\_Q) = (f\_Q) − (y), so f′\_Q/f\_Q = c\_Q/y and the map is a logarithmic derivative along the invariant differential ω = dx/(2y + a₁x + a₃). He evaluates it not in closed form but through O(log p) linear functions along a doubling chain (his Lemma 3 and equation (2)), and he notes that the construction also follows from a general result of Serre. Neither division polynomials nor Fermat quotients occur there. [Sma99] and [SA98] are of the other kind: they lift P and Q to E(Q\_p) (or to Z/p²) by Hensel, multiply by p, and read the ratio of p-adic elliptic logarithms modulo p; nothing in them is a rational function of the reduced coordinates, and [SA98, Thm 3.2] is the statement that the resulting map is independent of the lift. Our s·c₁ is a logarithmic derivative of the \textit{universal} polynomial φ\_p, divided by p, evaluated at the point itself; it is of Semaev's kind and, by the dimension count, a scalar multiple of his map. We claim nothing for it beyond the proof given in §4.11, which goes through Theorem H and not through any of these sources; and we note that a formula of degree p² in x is not an algorithm, both cited methods being exponentially cheaper.

One point of contact is worth recording. [Vol2] compares the two approaches and shows that the p-adic one degenerates exactly when the lift is the canonical lift, i.e. when the Serre–Tate parameter satisfies q = 1 in Z/p², which by [DW, Lemma 3.2] is exactly our splitting condition E(Q\_p)[p] ≠ 0. \textbf{Corollary D — λ ≡ 0 in the split case — is therefore the same statement as the failure of the anomalous-curve attack on the canonical lift}, read as the vanishing of a Fermat quotient; and λ\textasciicircum{}\{(2)\}, Theorem K and Theorem N are what remains when that first digit is gone. [Rüc] generalises [Sem98] to divisor class groups; we did not retrieve it. Of the three 1998–99 papers we retrieved [Sem98] as text and [Sma99] in the author's technical-report form; [SA98] we did not retrieve, and quote it as quoted in [Yas, §2] and [Vol2].

\textbf{Where the local criterion already is.} Proposition 6 is the case n = 1 of Gross's tameness criterion. [LR, Thm 4.1] states it as "Gross; see [9], p. 514; see also §§14–15", [9] being [Gro90]: for E/Q with ordinary good reduction at p, j ≠ 0, 1728 and E[p] an irreducible Gal(/Q)-module, a decomposition group D\_n at p is diagonalisable if and only if j\_E ≡ j₀ (mod p\textasciicircum{}\{n+1\}) for odd p, with j₀ the j-invariant of the canonical lift, determined by the modular equation T\_p(j₀, j₀) = 0 ([LR, Thm 4.2, Cor. 4.3], after Deuring). Those two hypotheses are absent from Proposition 6, and the passage from diagonalisability to a rational point of order p is [OY, §2.1.1 and Lem. 4.1], which for good ordinary reduction and odd p proves E(F)[p] ≠ 0 ⟺ E[p] ≅ F\_p(1) ⊕ F\_p ⟺ E[p] tame with F the unramified extension of degree the order of a\_p mod p — for us a\_p ≡ 1, so F = Q\_p — and quotes [Gro90, Prop. 13.7 and §17] for the equivalence with j\_E ≡ j₀ (mod p²). It is also [DW, Lemma 3.2 and Cor. 3.5], where a prime with E(Q\_p)[p] ≠ 0 is a \textit{local torsion prime} and the criterion is that the lift to Z/p² be the canonical one — of the two, only Lemma 3.2 is free of hypotheses on the j-invariant, Cor. 3.5 assuming j ≠ 0, 1728 as [LR, Thm 4.1] does; and the word "anomalous" does not occur in [DW], the content of "a local torsion prime is anomalous" sitting inside the proof of their Prop. 2.1(1) — of the p\textasciicircum{}d lifts exactly one has the larger p-rank — together with [DW, Prop. 2.1] for the fact that a local torsion prime is anomalous. The same statement circulates in cryptography, where the attacks on anomalous curves fail exactly for the lifts congruent to the canonical one modulo p². What we did not find is the criterion in the form q\_p(φ\_p(x₀)) = 0 — a Fermat quotient of a division polynomial at one point, which needs neither the j-invariant nor root-finding; the proof of §4.6 is independent of the sources above. Searches returning nothing on point (2026-09-24): \texttt{elliptic\ \allowbreak{}Wieferich\ \allowbreak{}prime}, \texttt{Fermat\ \allowbreak{}quotient\ \allowbreak{}elliptic\ \allowbreak{}curve}, \texttt{anomalous\ \allowbreak{}prime\ \allowbreak{}elliptic\ \allowbreak{}curve\ \allowbreak{}p-torsion}, \texttt{canonical\ \allowbreak{}lift\ \allowbreak{}p-torsion\ \allowbreak{}criterion\ \allowbreak{}mod\ \allowbreak{}p\textasciicircum{}2}. Near but different: the elliptic Wieferich primes of [Sil88] and [Vol] concern N\_p·P modulo p² for a rational point P, not rationality of the p-torsion, and the division-polynomial criteria of [Deb] are modulo p.

\textbf{Where the closed form of the ledger sits.} The three ingredients of Theorem N* are in [MST] and in the standard theory: the differential equation (d/dz)²log σ\_p = −(x + c) characterising the p-adic sigma function together with its normalisation, the Iwasawa logarithm of a p-adic unit, and the Weierstrass expansion of x in the formal parameter. Their combination — that the local term of the height is log(x(P″)·σ\_p(t(P″))²), an even series in the formal logarithm whose coefficients are polynomials in b₂, b₄, b₆ with E₂ entering the first one only — is the kind of statement that belongs to the same theory, and we searched for it as a displayed statement in [MST], [SW], [Har], [BMS], [BD], [Wut], [BG], [Ban], Sage's and PARI's documentation, and on arXiv (searches on 2026-09-25: \texttt{p-adic\ \allowbreak{}sigma\ \allowbreak{}function\ \allowbreak{}differential\ \allowbreak{}equation\ \allowbreak{}coefficients}, \texttt{p-adic\ \allowbreak{}height\ \allowbreak{}formal\ \allowbreak{}logarithm\ \allowbreak{}expansion}, \texttt{Mazur-Tate\ \allowbreak{}sigma\ \allowbreak{}even\ \allowbreak{}series}). We did not find it, in a search of the same narrow kind as the others recorded here, and we claim nothing for it beyond the proof in §3.11. What we did not find either is the rule 2me − 2 + v\_p(μ\_\{2m\}) for the digit at which a local invariant of E becomes visible in the p-adic height, or the observation that the digit can depend on p through the denominator of a coefficient. One caveat belongs here. A further search (\texttt{p-adic\ \allowbreak{}sigma\ \allowbreak{}function\ \allowbreak{}coefficients}, 2026-09-26) returned exactly one paper, [BG], which builds the universal p-adic σ out of congruences between the Laurent coefficients of the elliptic functions — the nearest published work to Proposition 17. \textbf{We have not retrieved it}, so we cannot say whether the coefficients μ₂, μ₄, μ₆ are in it; the search \texttt{p-adic\ \allowbreak{}height\ \allowbreak{}sigma\ \allowbreak{}series} returned nothing at all.

\textbf{Where depth and height meet.} Lemma 23 is the elementary end of a known subject: the p-adic valuation of the denominators of the multiples of a rational point — elliptic divisibility sequences — and the common valuations of division polynomials are studied in [NV], whose abstract we retrieved (arXiv:2203.02015); we did not read the body, and we cite it only to place Lemma 23 in a field and not for any statement. The index of the p-adic closure of the Mordell–Weil group, which §5.12 computes as p\textasciicircum{}\{min(D\_Q, D₁) − 1\}, is a central quantity of Chabauty's method and is treated in [Maz72]; we did not retrieve that paper (a search on arXiv for \texttt{p-adic\ \allowbreak{}closure\ \allowbreak{}Mordell-Weil} returned two papers, neither on this question), and we state nothing from it.

\textbf{What we did not find} is the splitting hypothesis E(Q\_p)[p] ≠ 0 in connection with the valuation of Reg\_p, or any statement that the height pairing is p-integral exactly under it. [SW] is our source for Schneider's conjecture; their §3.3 anomalous example is 446d1 at p = 5 with \#E(F\_5) = 10, the p = 5 case of §3.1, and its ord₅(Reg₅) is +1, not negative. The word "anomalous" does not occur in [BMS]. [Wut] studies the quantity e of §3.3 from the global side and [BG] the universal p-adic sigma function; [Kat] on Serre–Tate local moduli, the natural language for §4.6, we could not obtain. The nearest computation in the literature is [Ban] (2026), five rank-two CM curves at every good ordinary prime below 30,000 with a non-unit at 3 of 8,050 primes — the slice orthogonal to ours — which also formalises part of its argument in Lean 4.

\begin{thebibliography}{MRR15}

\bibitem[ATAEC]{ATAEC} J. H. Silverman, \textit{Advanced Topics in the Arithmetic of Elliptic Curves}, GTM 151. [not retrieved; Thm V.5.3, Thm V.3.1, Prop. V.6.1 and Exercise V.5.13 are cited as quoted in [Ked] and [Sym]; the minimality of the Tate model and the integrality of the coefficients of q as a series in 1/j we could not attach to numbered statements]
\bibitem[Ban]{Ban} B. S. Banwait, \textit{Second derivatives of p-adic L-functions and the Shafarevich–Tate group of rank-two CM elliptic curves}, arXiv:2609.08431 (2026).
\bibitem[BMS]{BMS} J. S. Balakrishnan, J. S. Müller, W. A. Stein, \textit{A p-adic analogue of the conjecture of Birch and Swinnerton-Dyer for modular abelian varieties}, arXiv:1210.2739.
\bibitem[Ber82]{Ber82} D. Bertrand, \textit{Valeurs de fonctions thêta et hauteurs p-adiques}, Progr. Math. \textbf{22}, Birkhäuser, 1982, 1–11. [not retrieved]
\bibitem[BG]{BG} C. Blakestad, D. Grant, \textit{On the universal p-adic sigma and Weierstrass zeta functions}, arXiv:1903.02480. [not retrieved; for p \textgreater{} 3 it reconstructs the Mazur–Tate universal p-adic σ from congruences between the Laurent coefficients of the elliptic functions, which makes it the nearest published work to Proposition 17 — see §9]
\bibitem[BD]{BD} A. Burungale, D. Disegni, arXiv:1803.09268.
\bibitem[Deb]{Deb} C. Debry, \textit{Beyond two criteria for supersingularity}, J. Théor. Nombres Bordeaux \textbf{26} (2014) 595–605.
\bibitem[Dok]{Dok} T. Dokchitser, V. Dokchitser, \textit{Local invariants of isogenous elliptic curves}, arXiv:1208.5519, Thm A.1.
\bibitem[DW]{DW} C. David, T. Weston, \textit{Local torsion on elliptic curves and the deformation theory of Galois representations}, arXiv:math/0701882.
\bibitem[Gro90]{Gro90} B. H. Gross, \textit{A tameness criterion for Galois representations associated to modular forms (mod p)}, Duke Math. J. \textbf{61} (1990) 445–517. [not retrieved; cited as quoted in [LR, Thm 4.1] (p. 514, §§14–15) and in [OY] (Prop. 13.7, §17)]
\bibitem[Har]{Har} D. Harvey, \textit{Efficient computation of p-adic heights}, arXiv:0708.3404.
\bibitem[Kat]{Kat} N. Katz, \textit{Serre–Tate local moduli}. [not retrieved]
\bibitem[Igu]{Igu} J.-I. Igusa, \textit{On the algebraic theory of elliptic modular functions}, J. Math. Soc. Japan \textbf{20} (1968) 96–106. [not retrieved; cited for the irreducibility of the tower, used as standard in §4.13]
\bibitem[Kat73]{Kat73} N. Katz, \textit{p-adic properties of modular schemes and modular forms}, in: Modular Functions of One Variable III, Lecture Notes in Math. \textbf{350}, Springer, 1973, 69–190. [not retrieved; cited for p-adic modular forms, the canonical subgroup, the theta operator and the congruence E₂\textasciicircum{}* ≡ E\_\{p+1\}/E\_\{p−1\} (mod p) used in §2, §4.13 and §7(8)]
\bibitem[Ked]{Ked} K. S. Kedlaya, \textit{The Tate curve}, lecture notes, MIT 18.727. [used only as a conduit for the statements of [ATAEC]]
\bibitem[LR]{LR} Á. Lozano-Robledo, \textit{Division fields of elliptic curves with minimal ramification}, Rev. Mat. Iberoam. \textbf{31} (2015) 1311–1332; Thm 4.1 quotes [Gro90], p. 514 and §§14–15, under the hypotheses j ≠ 0, 1728 and E[p] irreducible. (The Gross criterion is \textit{not} in the same author's Math. Ann. \textbf{357} (2013) paper, where we first looked.)
\bibitem[Maz72]{Maz72} B. Mazur, \textit{Rational points of abelian varieties with values in towers of number fields}, Invent. Math. \textbf{18} (1972), 183–266. [not retrieved; bibliographic data checked against Crossref, DOI 10.1007/BF01389815; the source of "anomalous", and a treatment of the p-adic closure of the Mordell–Weil group (§7(13), §9)]
\bibitem[MST]{MST} B. Mazur, W. Stein, J. Tate, \textit{Computation of p-adic heights and log convergence}, Documenta Math. Extra Vol.: John H. Coates' Sixtieth Birthday (2006) 585–622.
\bibitem[MTT]{MTT} B. Mazur, J. Tate, J. Teitelbaum, \textit{On p-adic analogues of the conjectures of Birch and Swinnerton-Dyer}, Invent. Math. \textbf{84} (1986), no. 1, 1–48. [read in the scan of the Göttingen digitisation centre (GDZ); Conj. BSD(p) is II §10, p. 38, the periods II §8, p. 36, the analytic height II §5, p. 33; PARI 2.17.3's \texttt{ellpadicbsd}, whose documentation names Mazur, Tate, Teitelbaum, Bernardi and Perrin-Riou, states it in the form of [SW]]
\bibitem[MT91]{MT91} B. Mazur, J. Tate, \textit{The p-adic sigma function}, Duke Math. J. \textbf{62} (1991) 663–688. [not retrieved]
\bibitem[MT87]{MT87} B. Mazur, J. Tate, \textit{Refined conjectures of the "Birch and Swinnerton-Dyer type"}, Duke Math. J. \textbf{54} (1987) 711–750. [not retrieved; read in the restatements of [PB19], [Ota], [LC24] and [BH25]]
\bibitem[BD94]{BD94} M. Bertolini, H. Darmon, \textit{Derived heights and generalized Mazur–Tate regulators}, Duke Math. J. \textbf{76} (1994) 75–111. [not retrieved; cited as in [Ota, Remark 2.5]]
\bibitem[PB19]{PB19} F. X. Portillo-Bobadilla, \textit{Experimental evidence on a refined conjecture of the BSD type}, Bol. Soc. Mat. Mex. (3) \textbf{25} (2019) 529–541; arXiv:1709.01204.
\bibitem[Ota]{Ota} K. Ota, \textit{Kato's Euler system and the Mazur–Tate refined conjecture of BSD type}, arXiv:1509.00682.
\bibitem[LC24]{LC24} J.-P. Llerena-Córdova, \textit{Numerical study of refined conjectures of the BSD type}, arXiv:2412.17703 (2024).
\bibitem[BH25]{BH25} D. Bullach, M. H. L. Honnor, \textit{On the refined 'Birch–Swinnerton-Dyer type' conjectures of Mazur and Tate}, arXiv:2511.07203 (2025).
\bibitem[JSW]{JSW} D. Jetchev, C. Skinner, X. Wan, \textit{The Birch and Swinnerton-Dyer formula for elliptic curves of analytic rank one}, arXiv:1512.06894 (2015). [abstract read]
\bibitem[PR]{PR} B. Poonen, E. Rains, \textit{Random maximal isotropic subspaces and Selmer groups}, arXiv:1009.0287 (2010). [abstract read]
\bibitem[Cre]{Cre} J. Cremona, \textit{ecdata}, files \texttt{allbsd/allbsd.\allowbreak{}*} (conductor \textless{} 500000), retrieved from the repository JohnCremona/ecdata on GitHub.
\bibitem[NV]{NV} B. Naskręcki, M. Verzobio, \textit{Common valuations of division polynomials}, arXiv:2203.02015 (2022); published version DOI 10.1017/prm.2024.7 (2024). [abstract page retrieved; cited to place Lemma 23 in a field]
\bibitem[OY]{OY} Y. Ozeki, M. Yoshida, \textit{Fields of definition of p-torsion points of elliptic curves and their ramification}, arXiv:2609.06554 (2026); §2.1.1 and Lemma 4.1.
\bibitem[Sch85]{Sch85} P. Schneider, \textit{p-adic height pairings II}, Invent. Math. \textbf{79} (1985) 329–374. [read in the GDZ scan; Thm 1 on p. 336, Prop. 1 and Thms 2 and 2′ on pp. 341–342; Thm 6 is quoted from the introduction, p. 330] Part I, Invent. Math. \textbf{69} (1982) 401–409, was not retrieved.
\bibitem[Ski]{Ski} C. Skinner, \textit{Multiplicative reduction and the cyclotomic main conjecture for GL2}, arXiv:1407.1093 (2014); §2.5 restates Kato's divisibility and the theorem of [SU] under the hypotheses used in §6, and replaces the hypothesis of [Kato04, Thm 17.4(3)] on the image of Galois by E[p] irreducible and (H); its Lemma 2.1.1 is the lattice lemma used there. (The theorem of [Kato04] proved in §13.14 is numbered 12.5(4) in the published version.)
\bibitem[DDT]{DDT} H. Darmon, F. Diamond, R. Taylor, \textit{Fermat's Last Theorem}, version of 9 September 2007 on the first author's page. [Prop. 1.5 (the Tate curve) and Prop. 2.12 (multiplicative reduction), pp. 19 and 57–58 of that version]
\bibitem[SU]{SU} C. Skinner, E. Urban, \textit{The Iwasawa main conjectures for GL2}, Invent. Math. \textbf{195} (2014) 1–277. [read in the version on the second author's page, which may differ slightly from the published one; the numbers 3.3.3, 3.6.4, 3.6.9 and 3.6.11 are those of that version]
\bibitem[Gre99]{Gre99} R. Greenberg, \textit{Iwasawa theory for elliptic curves}, in: Arithmetic Theory of Elliptic Curves (Cetraro, 1997), Lecture Notes in Math. \textbf{1716}, Springer, 1999, 51–144. [read in the author's version on his page; Thm 1.2, Prop. 2.2 and Thm 4.1]
\bibitem[Maz78]{Maz78} B. Mazur, \textit{Rational isogenies of prime degree}, Invent. Math. \textbf{44} (1978) 129–162. [read in the GDZ scan; Cor. 4.1, p. 144]
\bibitem[Kato04]{Kato04} K. Kato, \textit{p-adic Hodge theory and values of zeta functions of modular forms}, Astérisque \textbf{295} (2004) 117–290. [read on Numdam; Theorems 12.4–12.6 and condition 12.5.2 are on pp. 221–222, Theorem 13.4 on p. 226, §13.14 on p. 234, Lemma 14.7 on p. 238, Theorem 17.4 on p. 273 and its proof §17.13 on pp. 279–280]
\bibitem[GV]{GV} R. Greenberg, V. Vatsal, \textit{On the Iwasawa invariants of elliptic curves}, Invent. Math. \textbf{142} (2000) 17–63. [read in the authors' version on the first author's page; Prop. 3.7]
\bibitem[PR92]{PR92} B. Perrin-Riou, \textit{Théorie d'Iwasawa et hauteurs p-adiques}, Invent. Math. \textbf{109} (1992) 137–185. [read in the GDZ scan; Théorèmes 3.3.4, 3.3.6 and 3.4.2 and §3.4.3, pp. 165–167]
\bibitem[Wu17]{Wu17} C. Wuthrich, \textit{Numerical modular symbols for elliptic curves}, arXiv:1608.06423 (2017).
\bibitem[Ser73]{Ser73} J.-P. Serre, \textit{Formes modulaires et fonctions zêta p-adiques}, in: Modular Functions of One Variable III, Lecture Notes in Math. \textbf{350}, Springer, 1973, 191–268. [not retrieved; cited for E\_\{p−1\} ≡ 1 (mod p) and, with [SwD73], for the kernel of the q-expansion map]
\bibitem[SwD73]{SwD73} H. P. F. Swinnerton-Dyer, \textit{On ℓ-adic representations and congruences for coefficients of modular forms}, in: Modular Functions of One Variable III, Lecture Notes in Math. \textbf{350}, Springer, 1973, 1–55. [not retrieved; cited for the kernel of the q-expansion map on F\_p[E₄, E₆] being generated by A − 1]
\bibitem[Rüc]{Rüc} H.-G. Rück, \textit{On the discrete logarithm in the divisor class group of curves}, Math. Comp. \textbf{68} (1999) 805–806. [not retrieved]
\bibitem[SA98]{SA98} T. Satoh, K. Araki, \textit{Fermat quotients and the polynomial time discrete log algorithm for anomalous elliptic curves}, Comment. Math. Univ. St. Pauli \textbf{47} (1998) 81–92. [not retrieved; volume 47 is not in the publisher's repository. Cited as quoted in [Yas, §2] and [Vol2]]
\bibitem[Sem98]{Sem98} I. A. Semaev, \textit{Evaluation of discrete logarithms in a group of p-torsion points of an elliptic curve in characteristic p}, Math. Comp. \textbf{67} (1998) 353–356. [retrieved as text; the PDF at the publisher's site is behind a challenge]
\bibitem[Sil88]{Sil88} J. H. Silverman, \textit{Wieferich's criterion and the abc-conjecture}, J. Number Theory \textbf{30} (1988) 226–237.
\bibitem[SilH]{SilH} J. H. Silverman, \textit{Computing heights on elliptic curves}, Math. Comp. \textbf{51} (1988) 339–358, doi:10.1090/S0025-5718-1988-0942161-4. [Bibliographic data checked through Crossref; the text was not retrieved here.]
\bibitem[SW]{SW} W. Stein, C. Wuthrich, \textit{Algorithms for the arithmetic of elliptic curves using Iwasawa theory}, Math. Comp. \textbf{82} (2013) 1757–1792.
\bibitem[Sma99]{Sma99} N. P. Smart, \textit{The discrete logarithm problem on elliptic curves of trace one}, J. Cryptology \textbf{12} (1999) 193–196; quoted from the author's technical report HPL-97-128 (1997), same title.
\bibitem[Sym]{Sym} \textit{Symplectic criteria for elliptic curves, revisited}, arXiv:2509.19938, Prop. 2.2. [used only as a conduit for [ATAEC, Prop. V.6.1]]
\bibitem[Vél]{Vél} J. Vélu, \textit{Isogénies entre courbes elliptiques} (1971). [not retrieved]
\bibitem[Vol]{Vol} J. F. Voloch, J. Number Theory \textbf{81} (2000) 205–209.
\bibitem[Vol2]{Vol2} J. F. Voloch, \textit{Relating the Smart–Satoh–Araki and Semaev approaches to the discrete logarithm problem on anomalous elliptic curves}, note on the author's page.
\bibitem[Yas]{Yas} M. Yasuda, \textit{A p-adic approach to the discrete logarithm problem on anomalous elliptic curves}, Int. J. Pure Appl. Math. \textbf{77} (2012) 1–9; quotes [SA98, Thm 3.2].
\bibitem[Wut]{Wut} C. Wuthrich, \textit{On p-adic elliptic logarithms and p-adic approximation lattices} (2006).
\begin{itemize}[leftmargin=1.4em,itemsep=.2em,topsep=.3em]
\item Software: PARI/GP 2.17.3 with \texttt{elldata} (Cremona's tables, conductor \textless{} 500000); Sage, \texttt{src/sage/schemes/elliptic\_\allowbreak{}curves/padics.\allowbreak{}py}, docstrings of \texttt{padic\_\allowbreak{}regulator} and \texttt{padic\_\allowbreak{}height}.
\end{itemize}

\end{thebibliography}
\end{document}
